% === DEBUT DU CHUNK preamble+ch1-2 ===
\documentclass[12pt,a4paper,openany]{book}

% ============================================================
% ENCODAGE ET LANGUE
% ============================================================
\usepackage[utf8]{inputenc}
\usepackage[T1]{fontenc}
\usepackage[french]{babel}
\usepackage[margin=2.5cm]{geometry}

% ============================================================
% MATHEMATIQUES
% ============================================================
\usepackage{amsmath,amssymb,amsthm,mathrsfs,mathtools}

% ============================================================
% GRAPHIQUES
% ============================================================
\usepackage{tikz}
\usetikzlibrary{arrows.meta,calc,positioning,patterns,decorations.markings,decorations.pathmorphing,cd}
\usepackage{pgfplots}
\pgfplotsset{compat=1.18}
\usepackage{tikz-cd}

% ============================================================
% BOITES COLOREES
% ============================================================
\usepackage[most]{tcolorbox}
\tcbuselibrary{theorems,skins,breakable}

% ============================================================
% LIENS ET REFERENCES
% ============================================================
\usepackage[colorlinks=true,linkcolor=blue!60!black,citecolor=green!50!black]{hyperref}

% ============================================================
% INDEX
% ============================================================
\usepackage{makeidx}
\makeindex

% ============================================================
% MISE EN PAGE
% ============================================================
\usepackage{enumitem}
\usepackage{fancyhdr}
\usepackage{caption}
\usepackage{booktabs}
\usepackage{array}
\usepackage{longtable}
\usepackage{minitoc}
\usepackage{microtype}
\usepackage{xcolor}

% ============================================================
% EN-TETES ET PIEDS DE PAGE
% ============================================================
\pagestyle{fancy}
\fancyhf{}
\fancyhead[LE]{\leftmark}
\fancyhead[RO]{\rightmark}
\fancyfoot[C]{\thepage}
\renewcommand{\headrulewidth}{0.4pt}
\renewcommand{\footrulewidth}{0pt}

% ============================================================
% ENVIRONNEMENTS DE THEOREMES (amsthm)
% ============================================================
\theoremstyle{definition}
\newtheorem{definition}{Définition}[chapter]
\newtheorem{exemple}[definition]{Exemple}
\newtheorem{exercice}{Exercice}[chapter]

\theoremstyle{plain}
\newtheorem{theoreme}[definition]{Théorème}
\newtheorem{proposition}[definition]{Proposition}
\newtheorem{lemme}[definition]{Lemme}
\newtheorem{corollaire}[definition]{Corollaire}

\theoremstyle{remark}
\newtheorem{remarque}[definition]{Remarque}
\newtheorem{notation}[definition]{Notation}

% ============================================================
% HABILLAGE TCOLORBOX
% ============================================================
\tcolorboxenvironment{definition}{enhanced,breakable,colback=blue!5,colframe=blue!70!black,fonttitle=\bfseries,before skip=10pt,after skip=10pt,left=8pt,right=8pt}
\tcolorboxenvironment{theoreme}{enhanced,breakable,colback=red!5,colframe=red!70!black,fonttitle=\bfseries,before skip=10pt,after skip=10pt,left=8pt,right=8pt}
\tcolorboxenvironment{proposition}{enhanced,breakable,colback=orange!5,colframe=orange!70!black,fonttitle=\bfseries,before skip=10pt,after skip=10pt,left=8pt,right=8pt}
\tcolorboxenvironment{lemme}{enhanced,breakable,colback=yellow!5,colframe=yellow!50!black,fonttitle=\bfseries,before skip=10pt,after skip=10pt,left=8pt,right=8pt}
\tcolorboxenvironment{corollaire}{enhanced,breakable,colback=green!5,colframe=green!50!black,fonttitle=\bfseries,before skip=10pt,after skip=10pt,left=8pt,right=8pt}
\tcolorboxenvironment{remarque}{enhanced,breakable,colback=gray!5,colframe=gray!50!black,before skip=8pt,after skip=8pt,left=6pt,right=6pt}
\tcolorboxenvironment{exemple}{enhanced,breakable,colback=purple!5,colframe=purple!50!black,before skip=8pt,after skip=8pt,left=6pt,right=6pt}

% ============================================================
% COMMANDES PERSONNALISEES
% ============================================================
\newcommand{\N}{\mathbb{N}}
\newcommand{\Z}{\mathbb{Z}}
\newcommand{\Q}{\mathbb{Q}}
\newcommand{\R}{\mathbb{R}}
\newcommand{\C}{\mathbb{C}}
\newcommand{\RP}{\mathbb{R}\mathrm{P}}
\newcommand{\CP}{\mathbb{C}\mathrm{P}}
\DeclareMathOperator{\Id}{Id}
\DeclareMathOperator{\im}{im}
\DeclareMathOperator{\rg}{rg}
\DeclareMathOperator{\GL}{GL}
\DeclareMathOperator{\SL}{SL}
\DeclareMathOperator{\SO}{SO}
\DeclareMathOperator{\SU}{SU}
\DeclareMathOperator{\Diff}{Diff}
\DeclareMathOperator{\Hom}{Hom}
\newcommand{\dd}{\mathrm{d}}
\newcommand{\abs}[1]{\left|#1\right|}
\newcommand{\set}[1]{\left\{#1\right\}}
\newcommand{\rel}{\operatorname{rel}}

% ============================================================
% ============================================================
\begin{document}
\dominitoc

\begin{titlepage}
\begin{tikzpicture}[remember picture, overlay]
  \fill[left color=blue!15!white, right color=blue!5!white]
    (current page.south west) rectangle (current page.north east);
  \fill[color=orange!70!yellow]
    ([yshift=-2cm]current page.north west) rectangle ([yshift=-2.3cm]current page.north east);
  \fill[color=blue!80!black, rounded corners=8pt]
    ([xshift=3cm, yshift=-4cm]current page.north west) rectangle
    ([xshift=-3cm, yshift=-10cm]current page.north east);
  \node[anchor=center, text=white, font=\Huge\bfseries]
    at ([yshift=-6cm]current page.north) {Topologie Diff\'erentielle};
  \node[anchor=center, text=orange!80!yellow, font=\Large]
    at ([yshift=-7.5cm]current page.north) {Notes de Cours / Lecture Notes};
  \node[anchor=center, text=white!80, font=\large]
    at ([yshift=-8.8cm]current page.north) {Master / Doctorat --- 2025--2026};
  \node[anchor=center, text=blue!80!black, font=\Large\itshape]
    at ([yshift=-12cm]current page.north) {Yaé Ulrich Gaba};
  \draw[orange!70!yellow, line width=2pt]
    ([xshift=5cm, yshift=-13.5cm]current page.north west) --
    ([xshift=-5cm, yshift=-13.5cm]current page.north east);
  \node[anchor=center, text width=12cm, align=center, text=blue!60!black, font=\itshape]
    at ([yshift=-16cm]current page.north) {Vari\'et\'es, plongements et th\'eorie de Morse~: les structures lisses au c\oe ur de la g\'eom\'etrie.};
  \node[anchor=south, text=gray, font=\small]
    at ([yshift=1.5cm]current page.south) {\today};
\end{tikzpicture}
\end{titlepage}
\tableofcontents

% ============================================================
% TABLE DES NOTATIONS
% ============================================================
\chapter*{Table des notations}
\addcontentsline{toc}{chapter}{Table des notations}

\begin{longtable}{>{\raggedright}p{3.5cm} p{10cm}}
\toprule
\textbf{Notation} & \textbf{Signification} \\
\midrule
\endfirsthead
\toprule
\textbf{Notation} & \textbf{Signification} \\
\midrule
\endhead
\bottomrule
\endfoot
$M$, $N$, $P$ & Variétés différentielles lisses \\
$n = \dim M$ & Dimension de la variété $M$ \\
$(U, \varphi)$ & Carte locale : $U$ ouvert de $M$, $\varphi \colon U \to \R^n$ homéomorphisme sur son image \\
$\mathcal{A} = \set{(U_\alpha, \varphi_\alpha)}$ & Atlas lisse sur $M$ \\
$\varphi_\beta \circ \varphi_\alpha^{-1}$ & Changement de cartes (application de transition) \\
$C^\infty(M)$ & Algèbre des fonctions lisses $M \to \R$ \\
$C^\infty(M, N)$ & Ensemble des applications lisses de $M$ dans $N$ \\
$T_pM$ & Espace tangent à $M$ au point $p$ \\
$T_p^*M$ & Espace cotangent à $M$ au point $p$ \\
$TM$ & Fibré tangent de $M$ \\
$T^*M$ & Fibré cotangent de $M$ \\
$\dd f_p$ ou $(\dd f)_p$ & Différentielle de $f$ au point $p$ \\
$\dfrac{\partial}{\partial x^i}\Big|_p$ & Vecteur tangent associé à la $i$-ème coordonnée au point $p$ \\
$\dd x^i$ & $1$-forme différentielle coordonnée \\
$\Omega^k(M)$ & Espace des $k$-formes différentielles sur $M$ \\
$\Diff(M)$ & Groupe des difféomorphismes de $M$ \\
$\GL(n, \R)$ & Groupe linéaire général \\
$\SL(n, \R)$ & Groupe spécial linéaire \\
$\SO(n)$ & Groupe spécial orthogonal \\
$\SU(n)$ & Groupe spécial unitaire \\
$S^n$ & Sphère de dimension $n$ \\
$\RP^n$ & Espace projectif réel de dimension $n$ \\
$\CP^n$ & Espace projectif complexe de dimension $n$ \\
$\mathbb{T}^n$ & Tore de dimension $n$ \\
$\operatorname{Gr}(k,n)$ & Grassmannienne des $k$-plans dans $\R^n$ \\
$\Id_M$ & Application identité sur $M$ \\
$\rg(f)$ ou $\rg_p(f)$ & Rang de $f$ (au point $p$) \\
$\partial M$ & Bord de la variété $M$ \\
$\mathbb{H}^n$ & Demi-espace supérieur $\set{x \in \R^n : x^n \geq 0}$ \\
\end{longtable}

% ============================================================
% PREFACE
% ============================================================
\chapter*{Préface}
\addcontentsline{toc}{chapter}{Préface}

La topologie différentielle est l'étude des variétés lisses et des applications
différentiables entre elles. Elle occupe une position centrale dans les
mathématiques modernes, au carrefour de l'analyse, de la géométrie et de la
topologie algébrique.

\medskip

Pourquoi la structure lisse est-elle si importante, au-delà de la seule structure
topologique ? La réponse réside dans le fait qu'une variété topologique peut
admettre plusieurs structures différentiables non difféomorphes, voire aucune.
Le résultat le plus spectaculaire dans cette direction est dû à John Milnor
(1956) : la sphère $S^7$ admet exactement $28$ structures différentiables
distinctes, les \emph{sphères exotiques}\index{sphères exotiques}. Plus
surprenant encore, Simon Donaldson (1983) et Michael Freedman (1982) ont montré
que $\R^4$ admet une infinité non dénombrable de structures lisses exotiques,
un phénomène propre à la dimension~$4$.

\medskip

L'histoire de la topologie différentielle est jalonnée de contributions
fondatrices. Hassler Whitney\index{Whitney, Hassler} a établi dans les années
1930 les bases de la théorie des variétés et démontré le célèbre théorème de
plongement qui porte son nom : toute variété lisse de dimension $n$ se plonge
dans $\R^{2n+1}$. René Thom\index{Thom, René} a introduit la théorie du
cobordisme dans les années 1950, récompensée par la médaille Fields en 1958,
établissant des liens profonds entre la topologie différentielle et la topologie
algébrique. Stephen Smale\index{Smale, Stephen} a démontré la conjecture de
Poincaré généralisée en dimension $n \geq 5$ (1961) et développé la théorie de
Morse en dimension infinie. John Milnor\index{Milnor, John}, outre sa
découverte des sphères exotiques, a apporté des contributions décisives à la
théorie de Morse, à la $K$-théorie et à la chirurgie des variétés.

\medskip

Ce cours s'adresse aux étudiants de Master et de Doctorat. Nous supposons
acquises les bases de la topologie générale (espaces topologiques, compacité,
connexité), de l'algèbre linéaire et de l'analyse dans $\R^n$ (théorème
d'inversion locale, théorème des fonctions implicites). Le cours se veut
rigoureux : toutes les démonstrations essentielles sont présentées en détail.
Chaque chapitre comporte des exemples détaillés et des exercices de difficulté
graduée.

\medskip

\noindent\textit{Bonne lecture et bon travail.}


% ############################################################
%
%              CHAPITRE 1 : VARIETES DIFFERENTIELLES
%
% ############################################################
\chapter{Variétés différentielles --- Définitions et exemples}
\label{chap:varietes}
\minitoc

\section{Rappels topologiques}
\label{sec:rappels-topo}

Nous rappelons ici les notions topologiques indispensables à la théorie des
variétés. Le lecteur familier avec ces concepts pourra passer directement à la

ef{sec:cartes-atlas}.

\begin{definition}[Espace topologique séparé]\label{def:separe}
\index{espace séparé}\index{Hausdorff|see{espace séparé}}
Un espace topologique $(X, \mathcal{T})$ est dit \emph{séparé} (ou de
\emph{Hausdorff}) si pour tous points distincts $x, y \in X$, il existe des
ouverts $U, V \in \mathcal{T}$ tels que $x \in U$, $y \in V$ et $U \cap V =
\varnothing$.
\end{definition}

\begin{definition}[Base dénombrable]\label{def:base-denombrable}
\index{base dénombrable}\index{second dénombrable|see{base dénombrable}}
Un espace topologique $(X, \mathcal{T})$ est dit à \emph{base dénombrable}
(ou \emph{second dénombrable}) s'il existe une famille dénombrable
$\mathcal{B} = \set{B_n}_{n \in \N}$ d'ouverts de $X$ telle que tout ouvert
de $X$ s'écrive comme réunion d'éléments de $\mathcal{B}$.
\end{definition}

\begin{definition}[Paracompacité]\label{def:paracompact}
\index{paracompact}
Un espace topologique $X$ est dit \emph{paracompact} si tout recouvrement
ouvert de $X$ admet un raffinement localement fini. Plus précisément, pour
tout recouvrement ouvert $(U_\alpha)_{\alpha \in A}$ de $X$, il existe un
recouvrement ouvert $(V_\beta)_{\beta \in B}$ de $X$ tel que :
\begin{enumerate}[label=(\roman*)]
    \item pour tout $\beta \in B$, il existe $\alpha \in A$ avec $V_\beta
    \subset U_\alpha$ ;
    \item pour tout $x \in X$, il existe un voisinage ouvert $W$ de $x$ tel que
    $\set{\beta \in B : W \cap V_\beta \neq \varnothing}$ soit fini.
\end{enumerate}
\end{definition}

\begin{proposition}[Les espaces à base dénombrable séparés sont paracompacts]\label{prop:denombrable-paracompact}
\index{paracompact!et base dénombrable}
Tout espace topologique séparé à base dénombrable est paracompact.
\end{proposition}

\begin{proof}
Soit $X$ un espace séparé à base dénombrable $\mathcal{B} = \set{B_n}_{n \in \N}$.
Soit $(U_\alpha)_{\alpha \in A}$ un recouvrement ouvert de $X$. Pour chaque
$x \in X$, choisissons $\alpha(x)$ tel que $x \in U_{\alpha(x)}$, puis
$n(x) \in \N$ tel que $x \in B_{n(x)} \subset U_{\alpha(x)}$. Posons $A_k =
\set{x \in X : n(x) = k}$ pour $k \in \N$. Alors $X = \bigcup_{k \in \N} A_k$.

Pour construire un raffinement localement fini, on utilise une exhaustion
compacte. Puisque $X$ est séparé et à base dénombrable, on peut construire une
suite $(K_j)_{j \geq 1}$ de compacts tels que $K_j \subset \mathring{K}_{j+1}$
et $X = \bigcup_{j} K_j$. On pose $C_j = K_j \setminus \mathring{K}_{j-1}$
(avec $K_0 = \varnothing$). Chaque $C_j$ est compact et contenu dans l'ouvert
$W_j = \mathring{K}_{j+1} \setminus K_{j-2}$ (avec $K_{-1} = \varnothing$).

Pour chaque $j$, le recouvrement $(U_\alpha \cap W_j)_\alpha$ de $C_j$ admet
un sous-recouvrement fini $(V_{j,1}, \ldots, V_{j,m_j})$ avec $V_{j,i} \subset
W_j$. La famille $(V_{j,i})_{j,i}$ est un raffinement ouvert localement fini
de $(U_\alpha)_\alpha$, car tout point de $X$ possède un voisinage ne
rencontrant qu'un nombre fini des $W_j$.
\end{proof}

\begin{remarque}\label{rem:importance-paracompacite}
La paracompacité est essentielle pour l'existence de partitions de l'unité,
outil fondamental de la topologie différentielle (voir la

ef{sec:partitions-unite}).
\end{remarque}


\section{Cartes locales et atlas}
\label{sec:cartes-atlas}

\begin{definition}[Variété topologique]\label{def:variete-topo}
\index{variété topologique}
Une \emph{variété topologique de dimension $n$} est un espace topologique $M$
vérifiant :
\begin{enumerate}[label=(\roman*)]
    \item $M$ est séparé ;
    \item $M$ est à base dénombrable ;
    \item $M$ est \emph{localement euclidien de dimension $n$} : pour tout $p
    \in M$, il existe un ouvert $U$ contenant $p$, un ouvert $\widehat{U}
    \subset \R^n$ et un homéomorphisme $\varphi \colon U \to \widehat{U}$.
\end{enumerate}
\end{definition}

\begin{definition}[Carte locale]\label{def:carte}
\index{carte locale}
Une \emph{carte locale} (ou \emph{carte}) sur une variété topologique $M$ de
dimension $n$ est un couple $(U, \varphi)$ où $U$ est un ouvert de $M$ et
$\varphi \colon U \to \widehat{U} \subset \R^n$ est un homéomorphisme sur un
ouvert $\widehat{U}$ de $\R^n$. On dit que $U$ est le \emph{domaine} de la
carte et $\varphi$ la \emph{paramétrisation locale} (ou \emph{application
coordonnée}).

Si $\varphi(p) = (x^1(p), \ldots, x^n(p))$, les fonctions $x^1, \ldots, x^n
\colon U \to \R$ sont appelées les \emph{coordonnées locales}\index{coordonnées
locales} associées à la carte $(U, \varphi)$.
\end{definition}

\begin{definition}[Atlas]\label{def:atlas}
\index{atlas}
Un \emph{atlas} sur une variété topologique $M$ est une famille de cartes
$\mathcal{A} = \set{(U_\alpha, \varphi_\alpha)}_{\alpha \in A}$ telle que
$M = \bigcup_{\alpha \in A} U_\alpha$.
\end{definition}

\begin{definition}[Changement de cartes et atlas lisse]\label{def:atlas-lisse}
\index{changement de cartes}\index{application de transition}\index{atlas!lisse}
Soient $(U_\alpha, \varphi_\alpha)$ et $(U_\beta, \varphi_\beta)$ deux cartes
d'un atlas $\mathcal{A}$ sur $M$ avec $U_\alpha \cap U_\beta \neq \varnothing$.
L'\emph{application de transition} (ou \emph{changement de cartes}) est
l'homéomorphisme
\[
    \varphi_\beta \circ \varphi_\alpha^{-1} \colon
    \varphi_\alpha(U_\alpha \cap U_\beta) \longrightarrow
    \varphi_\beta(U_\alpha \cap U_\beta).
\]
L'atlas $\mathcal{A}$ est dit \emph{$C^\infty$-lisse} (ou simplement
\emph{lisse}) si toutes les applications de transition sont des
difféomorphismes $C^\infty$ entre ouverts de $\R^n$.
\end{definition}

\begin{center}
\begin{tikzpicture}[scale=1.1, >=Stealth]
    % Manifold
    \draw[thick, fill=blue!5, rounded corners=25pt]
        (-1.5, 1.5) -- (3.5, 2) -- (5, 0.5) -- (3, -1.5) -- (-1, -1) -- cycle;
    \node at (2, 2.3) {$M$};

    % U_alpha
    \draw[thick, blue!70!black, fill=blue!15, opacity=0.6, rounded corners=10pt]
        (-0.5, 0.8) -- (1.5, 1.2) -- (2.2, 0) -- (1, -0.8) -- (-0.5, -0.3) -- cycle;
    \node[blue!70!black] at (0.2, 1.1) {$U_\alpha$};

    % U_beta
    \draw[thick, red!70!black, fill=red!15, opacity=0.6, rounded corners=10pt]
        (1, 1.3) -- (3.2, 1.5) -- (4, 0.2) -- (2.8, -0.8) -- (1.2, -0.5) -- cycle;
    \node[red!70!black] at (3.5, 1.3) {$U_\beta$};

    % Intersection label
    \node at (1.8, 0.4) {\small $U_\alpha \cap U_\beta$};

    % phi_alpha
    \draw[->, thick, blue!70!black] (-0.3, -0.8) -- (-2, -3);
    \node[blue!70!black] at (-1.8, -1.7) {$\varphi_\alpha$};

    % phi_beta
    \draw[->, thick, red!70!black] (3.5, -0.5) -- (5.2, -3);
    \node[red!70!black] at (5, -1.5) {$\varphi_\beta$};

    % R^n left
    \draw[thick, fill=blue!8] (-4, -3.2) rectangle (-0.5, -5.5);
    \node at (-2.25, -5.8) {$\R^n$};
    \draw[thick, blue!50, fill=blue!15, rounded corners=5pt]
        (-3.5, -3.5) -- (-1.5, -3.5) -- (-1, -4.5) -- (-2, -5.2) -- (-3.5, -4.8) -- cycle;
    \node[blue!70!black, font=\small] at (-2.5, -4.3) {$\varphi_\alpha(U_\alpha)$};

    % R^n right
    \draw[thick, fill=red!8] (3.5, -3.2) rectangle (7, -5.5);
    \node at (5.25, -5.8) {$\R^n$};
    \draw[thick, red!50, fill=red!15, rounded corners=5pt]
        (4, -3.5) -- (6, -3.6) -- (6.5, -4.5) -- (5.5, -5.2) -- (4, -4.7) -- cycle;
    \node[red!70!black, font=\small] at (5.2, -4.3) {$\varphi_\beta(U_\beta)$};

    % Transition map
    \draw[->, thick, green!50!black, decorate, decoration={snake, amplitude=1.5pt, segment length=8pt}]
        (-1.2, -4.3) -- (4.2, -4.3);
    \node[green!50!black, above] at (1.5, -4.1) {$\varphi_\beta \circ \varphi_\alpha^{-1}$};
\end{tikzpicture}
\captionof{figure}{Atlas et changement de cartes sur une variété $M$.}
\end{center}

\begin{definition}[Compatibilité et atlas maximal]\label{def:atlas-maximal}
\index{atlas!maximal}\index{structure différentiable}
Deux atlas lisses $\mathcal{A}$ et $\mathcal{A}'$ sur $M$ sont dits
\emph{compatibles} si $\mathcal{A} \cup \mathcal{A}'$ est encore un atlas
lisse. La relation de compatibilité est une relation d'équivalence.

Un atlas lisse $\mathcal{A}$ est dit \emph{maximal} s'il n'est contenu dans
aucun atlas lisse strictement plus grand. Tout atlas lisse $\mathcal{A}$ est
contenu dans un unique atlas maximal $\overline{\mathcal{A}}$, défini par :
\[
    \overline{\mathcal{A}} = \set{(U, \varphi) \text{ carte sur } M :
    (U, \varphi) \text{ est compatible avec toute carte de } \mathcal{A}}.
\]
\end{definition}


\section{Structure différentiable}
\label{sec:structure-diff}

\begin{definition}[Variété différentielle lisse]\label{def:variete-lisse}
\index{variété différentielle}\index{variété lisse}
Une \emph{variété différentielle lisse} (ou simplement \emph{variété lisse})
de dimension $n$ est un couple $(M, \overline{\mathcal{A}})$ où $M$ est une
variété topologique de dimension $n$ et $\overline{\mathcal{A}}$ est un atlas
maximal lisse sur $M$. L'atlas maximal $\overline{\mathcal{A}}$ est appelé
la \emph{structure lisse} (ou \emph{structure différentiable}) de $M$.

En pratique, on spécifie une variété lisse en donnant un atlas lisse
$\mathcal{A}$ quelconque ; la structure lisse sous-jacente est alors
$\overline{\mathcal{A}}$.
\end{definition}

\begin{remarque}\label{rem:variete-notation}
Par abus de notation, on écrira souvent $M$ au lieu de $(M,
\overline{\mathcal{A}})$ lorsque la structure lisse est claire du contexte.
On parlera de variété lisse de dimension $n$, ou de $n$-variété lisse.
\end{remarque}

\begin{remarque}\label{rem:structures-exotiques}
\index{structures exotiques}
Une même variété topologique peut admettre des structures lisses non
difféomorphes. Le résultat célèbre de Milnor (1956) affirme que $S^7$ admet
exactement $28$ classes de structures lisses. En dimension~$4$, la situation
est encore plus remarquable : $\R^4$ admet une infinité non dénombrable de
structures lisses exotiques (Donaldson, Freedman).
\end{remarque}


\section{Exemples fondamentaux}
\label{sec:exemples}

\subsection{L'espace euclidien \texorpdfstring{$\R^n$}{Rn} et ses ouverts}

\begin{exemple}[$\R^n$ comme variété lisse]\label{ex:Rn}
\index{variété lisse!Rn@$\R^n$}
L'espace $\R^n$ muni de l'atlas à une seule carte $\mathcal{A} = \set{(\R^n,
\Id_{\R^n})}$ est une variété lisse de dimension $n$. Il est séparé, à base
dénombrable (les boules ouvertes à centre rationnel et rayon rationnel forment
une base), et trivialement localement euclidien.
\end{exemple}

\begin{exemple}[Ouverts de $\R^n$]\label{ex:ouverts-Rn}
\index{variété lisse!ouverts de Rn@ouverts de $\R^n$}
Plus généralement, tout ouvert $U \subset \R^n$ est une variété lisse de
dimension $n$, muni de l'atlas $\set{(U, \iota)}$ où $\iota \colon U
\hookrightarrow \R^n$ est l'inclusion. Plus généralement encore, tout ouvert
d'une variété lisse hérite naturellement d'une structure de variété lisse.
\end{exemple}

\subsection{Les sphères \texorpdfstring{$S^n$}{Sn}}

\begin{exemple}[La sphère $S^n$]\label{ex:sphere}
\index{sphère}\index{projection stéréographique}
La sphère unité de $\R^{n+1}$ est définie par
\[
    S^n = \set{x = (x^0, x^1, \ldots, x^n) \in \R^{n+1} : \sum_{i=0}^{n}
    (x^i)^2 = 1}.
\]
Elle est munie d'un atlas à deux cartes, les \emph{projections
stéréographiques} depuis le pôle nord $N = (1, 0, \ldots, 0)$ et le pôle
sud $S = (-1, 0, \ldots, 0)$:
\begin{align*}
    \varphi_N &\colon S^n \setminus \set{N} \longrightarrow \R^n, &
    \varphi_N(x^0, x^1, \ldots, x^n) &= \frac{1}{1 - x^0}(x^1, \ldots, x^n),
    \\
    \varphi_S &\colon S^n \setminus \set{S} \longrightarrow \R^n, &
    \varphi_S(x^0, x^1, \ldots, x^n) &= \frac{1}{1 + x^0}(x^1, \ldots, x^n).
\end{align*}
\end{exemple}

\begin{center}
\begin{tikzpicture}[scale=1.8, >=Stealth]
    % Sphere S^2
    \shade[ball color=blue!15, opacity=0.7] (0,0) circle (1.2);
    \draw[thick] (0,0) circle (1.2);
    % Equator
    \draw[thick, dashed] (1.2,0) arc[start angle=0, end angle=180, x radius=1.2, y radius=0.35];
    \draw[thick] (-1.2,0) arc[start angle=180, end angle=360, x radius=1.2, y radius=0.35];

    % Poles
    \fill (0, 1.2) circle (1.5pt) node[above] {$N$};
    \fill (0, -1.2) circle (1.5pt) node[below] {$S$};

    % Point P on sphere
    \fill (0.7, 0.6) circle (1.5pt) node[above right] {$p$};

    % Projection line from N through p to plane
    \draw[thick, red!70!black, ->] (0, 1.2) -- (1.82, -1.2);
    \fill[red!70!black] (1.4, -0.2) circle (1.5pt);

    % Plane (tangent at south pole level, for illustration)
    \draw[thick, gray] (-2.2, -1.2) -- (2.5, -1.2);
    \node[gray] at (2.7, -1.2) {$\R^n$};

    % Projected point
    \fill[red!70!black] (1.82, -1.2) circle (1.5pt)
        node[below, red!70!black] {$\varphi_N(p)$};

    \node at (-1.8, 0) {$S^n$};
\end{tikzpicture}
\captionof{figure}{Projection stéréographique de $S^n$ depuis le pôle nord $N$.}
\end{center}

\begin{proposition}[Lissité de l'atlas stéréographique]\label{prop:atlas-stereo}
L'atlas $\mathcal{A} = \set{(S^n \setminus \set{N}, \varphi_N), (S^n \setminus
\set{S}, \varphi_S)}$ est un atlas $C^\infty$-lisse sur $S^n$.
\end{proposition}

\begin{proof}
L'intersection des domaines est $S^n \setminus \set{N, S}$, et les images par
$\varphi_N$ et $\varphi_S$ de cet ensemble sont toutes deux $\R^n \setminus
\set{0}$. Calculons le changement de cartes. Soit $y = \varphi_N(x) =
\frac{1}{1 - x^0}(x^1, \ldots, x^n) \in \R^n \setminus \set{0}$. Alors
\[
    \|y\|^2 = \frac{(x^1)^2 + \cdots + (x^n)^2}{(1 - x^0)^2}
    = \frac{1 - (x^0)^2}{(1 - x^0)^2} = \frac{1 + x^0}{1 - x^0}.
\]
D'autre part,
\[
    \varphi_S(x) = \frac{1}{1 + x^0}(x^1, \ldots, x^n)
    = \frac{1 - x^0}{(1 + x^0)(1 - x^0)} \cdot (1 - x^0) \cdot
    \frac{(x^1, \ldots, x^n)}{1 - x^0}.
\]
Plus directement, on a
\[
    \varphi_S(x) = \frac{1}{1 + x^0}(x^1, \ldots, x^n)
    = \frac{1 - x^0}{1 - (x^0)^2}(x^1, \ldots, x^n)
    = \frac{1}{\|y\|^2}\, y.
\]
Ainsi, l'application de transition est
\[
    (\varphi_S \circ \varphi_N^{-1})(y) = \frac{y}{\|y\|^2},
    \qquad y \in \R^n \setminus \set{0}.
\]
C'est l'\emph{inversion}\index{inversion} par rapport à la sphère unité, qui
est un difféomorphisme $C^\infty$ de $\R^n \setminus \set{0}$ dans lui-même
(chaque composante est une fraction rationnelle de dénominateur non nul).
\end{proof}

\subsection{Espaces projectifs}

\begin{exemple}[L'espace projectif réel $\RP^n$]\label{ex:RPn}
\index{espace projectif réel}\index{RPn@$\RP^n$}
L'espace projectif réel $\RP^n$ est l'ensemble des droites vectorielles de
$\R^{n+1}$, c'est-à-dire le quotient
\[
    \RP^n = (\R^{n+1} \setminus \set{0}) / {\sim},
\]
où $x \sim y$ si et seulement si $x = \lambda y$ pour un certain $\lambda \in
\R^*$. On note $[x^0 : x^1 : \cdots : x^n]$ la classe de $(x^0, x^1, \ldots,
x^n)$.

Pour $i = 0, 1, \ldots, n$, posons $U_i = \set{[x^0 : \cdots : x^n] : x^i
\neq 0}$ et définissons
\[
    \varphi_i \colon U_i \longrightarrow \R^n, \qquad
    [x^0 : \cdots : x^n] \longmapsto
    \left(\frac{x^0}{x^i}, \ldots, \widehat{\frac{x^i}{x^i}}, \ldots,
    \frac{x^n}{x^i}\right),
\]
où le chapeau indique l'omission du terme $x^i/x^i = 1$. L'atlas $\set{(U_i,
\varphi_i)}_{i=0}^n$ est un atlas lisse sur $\RP^n$, ce qui en fait une variété
lisse compacte de dimension $n$.

Vérifions la lissité des changements de cartes. Pour $i \neq j$, l'image
$\varphi_i(U_i \cap U_j)$ est l'ensemble des $(t^1, \ldots, t^n) \in \R^n$
avec $t^j \neq 0$ (si $j < i$) ou $t^{j-1} \neq 0$ (si $j > i$, car la
$j$-ème coordonnée est décalée). Par exemple, pour $i = 0$ et $j = 1$, on a
$\varphi_0^{-1}(t^1, \ldots, t^n) = [1 : t^1 : t^2 : \cdots : t^n]$, donc
\[
    (\varphi_1 \circ \varphi_0^{-1})(t^1, \ldots, t^n)
    = \varphi_1([1 : t^1 : \cdots : t^n])
    = \left(\frac{1}{t^1}, \frac{t^2}{t^1}, \ldots, \frac{t^n}{t^1}\right),
\]
définie sur $\set{t \in \R^n : t^1 \neq 0}$. C'est une application $C^\infty$
(fractions rationnelles de dénominateur non nul). Le cas général est analogue.
\end{exemple}

\begin{exemple}[L'espace projectif complexe $\CP^n$]\label{ex:CPn}
\index{espace projectif complexe}\index{CPn@$\CP^n$}
De manière analogue, l'espace projectif complexe est
\[
    \CP^n = (\C^{n+1} \setminus \set{0}) / {\sim},
\]
où $z \sim w$ si $z = \lambda w$ pour $\lambda \in \C^*$. En identifiant $\C^n
\cong \R^{2n}$, les cartes $\varphi_i \colon U_i \to \C^n \cong \R^{2n}$
définissent un atlas lisse, faisant de $\CP^n$ une variété lisse compacte de
dimension réelle $2n$.
\end{exemple}

\subsection{Le tore}

\begin{exemple}[Le tore $\mathbb{T}^n$]\label{ex:tore}
\index{tore}\index{T@$\mathbb{T}^n$}
Le tore de dimension $n$ est le quotient
\[
    \mathbb{T}^n = \R^n / \Z^n,
\]
où $\Z^n$ agit par translations entières. La projection canonique $\pi \colon
\R^n \to \mathbb{T}^n$ est un homéomorphisme local. Si $p = \pi(x)$ et $U$ est
un ouvert de $\R^n$ contenant $x$ et suffisamment petit pour que $\pi|_U$ soit
injective (par exemple une boule ouverte de rayon $< 1/2$), alors $(
\pi(U), (\pi|_U)^{-1})$ est une carte locale. L'atlas ainsi obtenu est lisse
car les changements de cartes sont des translations par des vecteurs de $\Z^n$,
qui sont des difféomorphismes $C^\infty$ de $\R^n$.

Le tore $\mathbb{T}^n$ est compact (image continue du compact $[0,1]^n$) et
connexe.
\end{exemple}

\begin{remarque}\label{rem:tore-produit}
On a un difféomorphisme canonique $\mathbb{T}^n \cong (S^1)^n$, obtenu via
l'identification $\R/\Z \cong S^1$ donnée par $t \mapsto e^{2\pi i t}$.
En particulier, $\mathbb{T}^2 = S^1 \times S^1$ est le tore usuel, plongeable
dans $\R^3$ par la paramétrisation
\[
    (\theta, \phi) \longmapsto \big((R + r\cos\theta)\cos\phi,\
    (R + r\cos\theta)\sin\phi,\ r\sin\theta\big),
\]
avec $0 < r < R$.
\end{remarque}

\subsection{Groupes de Lie classiques}
\label{subsec:groupes-Lie}

\begin{definition}[Groupe de Lie]\label{def:groupe-Lie}
\index{groupe de Lie}
Un \emph{groupe de Lie} est un groupe $G$ muni d'une structure de variété lisse
telle que les applications
\begin{align*}
    \mu &\colon G \times G \to G, & (g, h) &\mapsto gh, \\
    \iota &\colon G \to G, & g &\mapsto g^{-1},
\end{align*}
soient lisses.
\end{definition}

\begin{exemple}[$\GL(n, \R)$]\label{ex:GLn}
\index{groupe de Lie!GL(n,R)@$\GL(n,\R)$}\index{GL(n,R)@$\GL(n,\R)$}
Le \emph{groupe linéaire général} $\GL(n, \R)$ est l'ensemble des matrices $n
\times n$ inversibles à coefficients réels :
\[
    \GL(n, \R) = \set{A \in M_n(\R) : \det A \neq 0}.
\]
C'est un ouvert de $M_n(\R) \cong \R^{n^2}$ (car $\det$ est continue), donc
une variété lisse de dimension $n^2$. La multiplication matricielle est
polynomiale en les coefficients, et l'inversion est donnée par la formule des
cofacteurs, donc ces opérations sont lisses. Ainsi $\GL(n, \R)$ est un groupe
de Lie de dimension $n^2$.
\end{exemple}

\begin{exemple}[$\SL(n, \R)$]\label{ex:SLn}
\index{groupe de Lie!SL(n,R)@$\SL(n,\R)$}\index{SL(n,R)@$\SL(n,\R)$}
Le \emph{groupe spécial linéaire} est
\[
    \SL(n, \R) = \set{A \in M_n(\R) : \det A = 1} = \det^{-1}(\set{1}).
\]
L'application $\det \colon M_n(\R) \to \R$ est lisse, et $1$ est une valeur
régulière de $\det$ restreinte à $\GL(n, \R)$ (on le montrera rigoureusement au

ef{chap:applications-lisses} via le théorème de la valeur régulière). Par
conséquent, $\SL(n, \R)$ est une sous-variété lisse de $M_n(\R)$ de dimension
$n^2 - 1$, et c'est un groupe de Lie.
\end{exemple}

\begin{exemple}[$\SO(n)$]\label{ex:SOn}
\index{groupe de Lie!SO(n)@$\SO(n)$}\index{SO(n)@$\SO(n)$}\index{groupe orthogonal}
Le \emph{groupe orthogonal} est $O(n) = \set{A \in M_n(\R) : A^T A = I_n}$.
Considérons l'application $\Phi \colon M_n(\R) \to \operatorname{Sym}_n(\R)$
définie par $\Phi(A) = A^T A$, où $\operatorname{Sym}_n(\R)$ est l'espace des
matrices symétriques. L'espace $\operatorname{Sym}_n(\R)$ est de dimension
$n(n+1)/2$. On montre que $I_n$ est une valeur régulière de $\Phi$, ce qui
implique que $O(n) = \Phi^{-1}(I_n)$ est une sous-variété lisse de $M_n(\R)$
de dimension
\[
    \dim O(n) = n^2 - \frac{n(n+1)}{2} = \frac{n(n-1)}{2}.
\]
Le \emph{groupe spécial orthogonal} est la composante connexe de l'identité :
\[
    \SO(n) = \set{A \in O(n) : \det A = 1}.
\]
C'est un groupe de Lie connexe compact de dimension $n(n-1)/2$.
\end{exemple}

\begin{exemple}[$\SU(n)$]\label{ex:SUn}
\index{groupe de Lie!SU(n)@$\SU(n)$}\index{SU(n)@$\SU(n)$}\index{groupe unitaire}
Le \emph{groupe unitaire} $U(n) = \set{A \in M_n(\C) : A^* A = I_n}$ (où
$A^* = \overline{A}^T$) est une sous-variété lisse de $M_n(\C) \cong
\R^{2n^2}$ de dimension $n^2$. Le \emph{groupe spécial unitaire}
\[
    \SU(n) = \set{A \in U(n) : \det A = 1}
\]
est un groupe de Lie compact connexe et simplement connexe de dimension $n^2 -
1$. En particulier, $\SU(2)$ est difféomorphe à $S^3$.
\end{exemple}

\begin{remarque}\label{rem:SU2-S3}
Le difféomorphisme $\SU(2) \cong S^3$ s'obtient en écrivant tout élément de
$\SU(2)$ sous la forme
$\begin{psmallmatrix} \alpha & -\bar\beta \\ \beta & \bar\alpha
\end{psmallmatrix}$ avec $|\alpha|^2 + |\beta|^2 = 1$, et en identifiant
$(\alpha, \beta) \in \C^2 \cong \R^4$ avec un point de $S^3$.
\end{remarque}

\subsection{Grassmanniennes}

\begin{exemple}[Grassmannienne $\operatorname{Gr}(k,n)$]\label{ex:grassmannienne}
\index{grassmannienne}\index{Gr(k,n)@$\operatorname{Gr}(k,n)$}
La \emph{grassmannienne} $\operatorname{Gr}(k,n)$ est l'ensemble des
sous-espaces vectoriels de dimension $k$ de $\R^n$. C'est une variété lisse
compacte de dimension $k(n-k)$. On a $\operatorname{Gr}(1, n+1) \cong \RP^n$.

La structure de variété est définie comme suit. Soit $W \in
\operatorname{Gr}(k,n)$ et soit $W^\perp$ son complément orthogonal. Pour
tout sous-espace $V$ suffisamment proche de $W$ (au sens où $V \cap W^\perp =
\set{0}$), $V$ est le graphe d'une unique application linéaire $A \colon W \to
W^\perp$. Ceci définit une carte locale $\varphi_W \colon U_W \to
\operatorname{Hom}(W, W^\perp) \cong \R^{k(n-k)}$. Les changements de cartes
sont lisses, ce qui munit $\operatorname{Gr}(k,n)$ d'une structure de variété
lisse.
\end{exemple}


\section{Variétés avec bord}
\label{sec:varietes-bord}

\begin{definition}[Demi-espace supérieur]\label{def:demi-espace}
\index{demi-espace}
Le \emph{demi-espace supérieur} de $\R^n$ est
\[
    \mathbb{H}^n = \set{(x^1, \ldots, x^n) \in \R^n : x^n \geq 0}.
\]
Son \emph{bord} est $\partial \mathbb{H}^n = \set{x \in \R^n : x^n = 0}
\cong \R^{n-1}$ et son \emph{intérieur} est $\operatorname{Int}(\mathbb{H}^n) =
\set{x \in \R^n : x^n > 0}$.
\end{definition}

\begin{definition}[Variété à bord]\label{def:variete-bord}
\index{variété à bord}\index{bord}
Une \emph{variété lisse à bord de dimension $n$} est un espace topologique
séparé à base dénombrable $M$ tel que tout point possède un voisinage
homéomorphe à un ouvert de $\mathbb{H}^n$, muni d'un atlas lisse (les
changements de cartes étant des difféomorphismes entre ouverts de
$\mathbb{H}^n$).

Le \emph{bord} de $M$, noté $\partial M$\index{bord!d'une variété}, est
l'ensemble des points de $M$ envoyés sur $\partial \mathbb{H}^n$ par une (et
donc toute) carte locale. L'\emph{intérieur} de $M$ est $\operatorname{Int}(M)
= M \setminus \partial M$.
\end{definition}

\begin{proposition}\label{prop:bord-variete}
Soit $M$ une variété lisse à bord de dimension $n$. Alors :
\begin{enumerate}[label=(\roman*)]
    \item $\partial M$ est une variété lisse (sans bord) de dimension $n - 1$ ;
    \item $\operatorname{Int}(M)$ est une variété lisse (sans bord) de
    dimension $n$ ;
    \item $\partial(\partial M) = \varnothing$.
\end{enumerate}
\end{proposition}

\begin{proof}
(i) Si $(U, \varphi)$ est une carte de $M$ avec $U \cap \partial M \neq
\varnothing$, alors $\varphi(U \cap \partial M) = \varphi(U) \cap \partial
\mathbb{H}^n$ est un ouvert de $\R^{n-1} \cong \partial \mathbb{H}^n$. Les
restrictions de ces cartes à $\partial M$ forment un atlas lisse. (ii) est
analogue, les cartes intérieures ayant leur image dans
$\operatorname{Int}(\mathbb{H}^n) \cong \R^n$. (iii) est immédiat car les
cartes de $\partial M$ ont leur image dans $\R^{n-1}$ (pas dans
$\mathbb{H}^{n-1}$).
\end{proof}

\begin{exemple}\label{ex:varietes-bord}
Voici des exemples classiques de variétés à bord :
\begin{enumerate}[label=(\alph*)]
    \item Le disque fermé $\overline{D}^n = \set{x \in \R^n : \|x\| \leq 1}$
    est une variété à bord de dimension $n$, avec $\partial \overline{D}^n =
    S^{n-1}$.
    \item Le cylindre $S^1 \times [0,1]$ est une variété à bord de dimension $2$,
    avec $\partial(S^1 \times [0,1]) = S^1 \times \set{0} \sqcup S^1 \times
    \set{1}$.
    \item Le ruban de Möbius est une variété à bord de dimension $2$, dont le
    bord est homéomorphe à $S^1$.
\end{enumerate}
\end{exemple}


\section{Partitions de l'unité}
\label{sec:partitions-unite}

Les partitions de l'unité sont un outil technique fondamental en topologie
différentielle. Elles permettent de « recoller » des constructions locales en
des objets globaux.

\begin{definition}[Support]\label{def:support}
\index{support}
Soit $f \colon M \to \R$ une fonction continue. Le \emph{support} de $f$ est
\[
    \operatorname{supp}(f) = \overline{\set{p \in M : f(p) \neq 0}}.
\]
\end{definition}

\begin{definition}[Partition de l'unité]\label{def:partition-unite}
\index{partition de l'unité}
Soit $(U_\alpha)_{\alpha \in A}$ un recouvrement ouvert d'une variété lisse $M$.
Une \emph{partition de l'unité subordonnée} à ce recouvrement est une famille
$(\psi_\alpha)_{\alpha \in A}$ de fonctions lisses $\psi_\alpha \colon M \to
[0, 1]$ telle que :
\begin{enumerate}[label=(\roman*)]
    \item $\operatorname{supp}(\psi_\alpha) \subset U_\alpha$ pour tout
    $\alpha \in A$ ;
    \item la famille $(\operatorname{supp}(\psi_\alpha))_{\alpha \in A}$ est
    localement finie ;
    \item $\sum_{\alpha \in A} \psi_\alpha(p) = 1$ pour tout $p \in M$
    (la somme est bien définie grâce à (ii)).
\end{enumerate}
\end{definition}

\begin{lemme}[Fonction bosse lisse]\label{lem:fonction-bosse}
\index{fonction bosse}
Il existe une fonction lisse $\rho \colon \R^n \to [0, +\infty)$ à support
compact, telle que $\rho(x) > 0$ pour $\|x\| < 1$ et $\rho(x) = 0$ pour
$\|x\| \geq 1$.
\end{lemme}

\begin{proof}
On considère la fonction lisse $h \colon \R \to \R$ définie par
\[
    h(t) = \begin{cases} e^{-1/t} & \text{si } t > 0, \\ 0 & \text{si }
    t \leq 0. \end{cases}
\]
Cette fonction est $C^\infty$ sur $\R$ entier (toutes les dérivées en $0$ sont
nulles). La fonction $g(t) = h(t) / (h(t) + h(1 - t))$ est lisse, vérifie
$g(t) = 0$ pour $t \leq 0$, $g(t) = 1$ pour $t \geq 1$, et $0 \leq g(t) \leq
1$. On pose alors
\[
    \rho(x) = g(1 - \|x\|^2).
\]
C'est une fonction lisse (car $\|x\|^2$ est lisse et $g$ est lisse),
strictement positive pour $\|x\| < 1$ et nulle pour $\|x\| \geq 1$.
\end{proof}

\begin{theoreme}[Existence de partitions de l'unité]\label{thm:partition-unite}
\index{partition de l'unité!existence}
Soit $M$ une variété lisse (sans bord ou à bord) et $(U_\alpha)_{\alpha \in A}$
un recouvrement ouvert de $M$. Alors il existe une partition de l'unité lisse
subordonnée à $(U_\alpha)_{\alpha \in A}$.
\end{theoreme}

\begin{proof}
La preuve se décompose en plusieurs étapes.

\textbf{Étape 1 : Recouvrement localement fini par des boules coordonnées.}
Puisque $M$ est à base dénombrable, elle possède une exhaustion par des
compacts : il existe une suite $(K_j)_{j \geq 1}$ de compacts avec $K_j
\subset \mathring{K}_{j+1}$ et $M = \bigcup_j K_j$. On pose $A_j =
K_j \setminus \mathring{K}_{j-1}$ (compact) et $W_j = \mathring{K}_{j+1}
\setminus K_{j-2}$ (ouvert contenant $A_j$), avec la convention $K_0 =
K_{-1} = \varnothing$.

Pour chaque point $p \in A_j$, choisissons $\alpha(p) \in A$ avec $p \in
U_{\alpha(p)}$, et une carte $(V_p, \varphi_p)$ centrée en $p$ telle que
$V_p \subset U_{\alpha(p)} \cap W_j$ et $\varphi_p(V_p)$ contienne la boule
unité fermée $\overline{B}(0, 1)$. Posons $B_p = \varphi_p^{-1}(B(0, 1))$.
Les ouverts $B_p$, pour $p \in A_j$, recouvrent le compact $A_j$, donc on
peut en extraire un sous-recouvrement fini $B_{p_1^j}, \ldots, B_{p_{m_j}^j}$.
La collection $(B_{p_i^j})_{j, i}$ est un recouvrement localement fini de $M$
(car chaque $W_j$ ne rencontre qu'un nombre fini de bandes $A_k$).

\textbf{Étape 2 : Fonctions bosses.}
Pour chaque $(j, i)$, posons $f_{j,i} = \rho \circ \varphi_{p_i^j}$ (étendue
par $0$ en dehors de $V_{p_i^j}$), où $\rho$ est la fonction bosse du

ef{lem:fonction-bosse}. Alors $f_{j,i}$ est lisse, à support dans $V_{p_i^j}
\subset U_{\alpha(p_i^j)}$, et strictement positive sur $B_{p_i^j}$.

\textbf{Étape 3 : Normalisation.}
La somme $\sigma(p) = \sum_{j,i} f_{j,i}(p)$ est bien définie et lisse (somme
localement finie de fonctions lisses), et $\sigma(p) > 0$ pour tout $p \in M$
car les $B_{p_i^j}$ recouvrent $M$. On pose $\widetilde{\psi}_{j,i} =
f_{j,i} / \sigma$.

\textbf{Étape 4 : Subordination.}
Pour chaque $\alpha \in A$, on pose
\[
    \psi_\alpha = \sum_{\substack{(j,i) \\ \alpha(p_i^j) = \alpha}}
    \widetilde{\psi}_{j,i}.
\]
Alors $(\psi_\alpha)_{\alpha \in A}$ est une partition de l'unité lisse
subordonnée à $(U_\alpha)_\alpha$.
\end{proof}

\begin{corollaire}\label{cor:partition-applications}
Soit $M$ une variété lisse. Alors :
\begin{enumerate}[label=(\roman*)]
    \item Pour tout fermé $F \subset M$ et tout ouvert $U \supset F$, il
    existe une fonction lisse $f \colon M \to [0,1]$ avec $f|_F = 1$ et
    $\operatorname{supp}(f) \subset U$.
    \item Toute fonction continue sur $M$ peut être uniformément approchée par
    des fonctions lisses.
\end{enumerate}
\end{corollaire}

\begin{proof}
(i) Le recouvrement $\set{U, M \setminus F}$ de $M$ admet une partition de
l'unité subordonnée $\set{\psi_1, \psi_2}$ avec $\operatorname{supp}(\psi_1)
\subset U$ et $\operatorname{supp}(\psi_2) \subset M \setminus F$. Alors
$f = \psi_1$ satisfait $\operatorname{supp}(f) \subset U$ et, pour $p \in F$,
$\psi_2(p) = 0$ donc $f(p) = \psi_1(p) = 1$.

(ii) découle de (i) par un argument standard de régularisation à l'aide de
partitions de l'unité.
\end{proof}


\section{Exercices}
\label{sec:exercices-ch1}

\begin{exercice}[$\star\star$]\label{exo:atlas-S1}
\index{exercice!atlas de $S^1$}
Montrer que l'atlas stéréographique de $S^1$ (deux cartes) est compatible avec
l'atlas formé des quatre cartes obtenues en projetant les demi-cercles ouverts
sur les axes de coordonnées. En déduire qu'ils définissent la même structure
lisse.
\end{exercice}

\begin{exercice}[$\star\star$]\label{exo:RPn-quotient}
\index{exercice!espace projectif}
Montrer que la projection canonique $\pi \colon S^n \to \RP^n$ (qui identifie
les points antipodaux) est un homéomorphisme local. En déduire que $\RP^n$ est
compact et que l'atlas défini dans l'
ef{ex:RPn} est bien un atlas lisse.
\end{exercice}

\begin{exercice}[$\star\star$]\label{exo:produit-varietes}
\index{exercice!produit de variétés}
Soient $M$ et $N$ deux variétés lisses de dimensions $m$ et $n$ respectivement.
Montrer que $M \times N$ admet une structure naturelle de variété lisse de
dimension $m + n$, avec l'atlas $\set{(U_\alpha \times V_\beta, \varphi_\alpha
\times \psi_\beta)}$ où $\set{(U_\alpha, \varphi_\alpha)}$ et $\set{(V_\beta,
\psi_\beta)}$ sont des atlas de $M$ et $N$.
\end{exercice}

\begin{exercice}[$\star\star$]\label{exo:CP1-S2}
\index{exercice!CP1@$\CP^1 \cong S^2$}
Montrer que $\CP^1$ est difféomorphe à $S^2$.
\emph{Indication} : Utiliser les deux cartes standard de $\CP^1$ et l'atlas
stéréographique de $S^2$, et construire explicitement un difféomorphisme.
\end{exercice}

\begin{exercice}[$\star\star\star$]\label{exo:SOn-connexe}
\index{exercice!connexité de $\SO(n)$}
Montrer que $\SO(n)$ est connexe pour tout $n \geq 1$.
\emph{Indication} : Considérer la fibration $\SO(n) \to S^{n-1}$ donnée par
la dernière colonne, et procéder par récurrence sur $n$.
\end{exercice}

\begin{exercice}[$\star\star\star$]\label{exo:grassmann-compact}
\index{exercice!grassmannienne compacte}
Montrer que la grassmannienne $\operatorname{Gr}(k,n)$ est compacte.
\emph{Indication} : Montrer que $\operatorname{Gr}(k,n)$ est un quotient du
groupe compact $O(n)$.
\end{exercice}

\begin{exercice}[$\star\star$]\label{exo:partition-metrique}
\index{exercice!partition de l'unité et métrique}
Soit $M$ une variété lisse. En utilisant les partitions de l'unité, montrer que
$M$ admet une métrique riemannienne, c'est-à-dire un produit scalaire lisse
sur chaque espace tangent, variant de manière lisse avec le point.
\emph{Remarque} : On utilisera la notion d'espace tangent du

ef{chap:applications-lisses}.
\end{exercice}

\begin{exercice}[$\star\star\star$]\label{exo:variete-bord-double}
\index{exercice!double d'une variété à bord}
Soit $M$ une variété lisse à bord. Le \emph{double} de $M$ est l'espace
topologique $D(M) = (M \sqcup M) / \sim$, où l'on identifie les deux copies
de $\partial M$. Montrer que $D(M)$ admet une structure de variété lisse (sans
bord) de même dimension que $M$.
\emph{Indication} : Utiliser un collier de $\partial M$ dans $M$, c'est-à-dire
un difféomorphisme $\partial M \times [0, 1) \to U$ où $U$ est un voisinage
ouvert de $\partial M$ dans $M$.
\end{exercice}

\begin{exercice}[$\star\star$]\label{exo:sous-groupe-fermé}
\index{exercice!sous-groupe fermé}
Soit $G$ un groupe de Lie et $H$ un sous-groupe fermé de $G$. Admettre que $H$
est une sous-variété lisse de $G$ (théorème de Cartan). En déduire que les
groupes suivants sont des groupes de Lie, et calculer leur dimension :
\begin{enumerate}[label=(\alph*)]
    \item le groupe des matrices triangulaires supérieures inversibles ;
    \item le groupe symplectique $\operatorname{Sp}(2n, \R) = \set{A \in
    \GL(2n, \R) : A^T J A = J}$ où $J = \begin{psmallmatrix} 0 & I_n \\ -I_n
    & 0 \end{psmallmatrix}$.
\end{enumerate}
\end{exercice}


% ############################################################
%
%    CHAPITRE 2 : APPLICATIONS LISSES ET DIFFEOMORPHISMES
%
% ############################################################
\chapter{Applications lisses et difféomorphismes}
\label{chap:applications-lisses}
\minitoc

\section{Applications lisses entre variétés}
\label{sec:applications-lisses}

\begin{definition}[Application lisse]\label{def:application-lisse}
\index{application lisse}
Soient $M$ et $N$ des variétés lisses de dimensions $m$ et $n$ respectivement.
Une application continue $f \colon M \to N$ est dite \emph{lisse} (ou
\emph{$C^\infty$}) si, pour toute carte $(U, \varphi)$ de $M$ et toute carte
$(V, \psi)$ de $N$ avec $f(U) \cap V \neq \varnothing$, l'application
\[
    \psi \circ f \circ \varphi^{-1} \colon
    \varphi(U \cap f^{-1}(V)) \longrightarrow \R^n
\]
est lisse (au sens usuel de l'analyse dans $\R^n$). L'application $\psi \circ f
\circ \varphi^{-1}$ est appelée la \emph{représentation locale}\index{représentation locale} (ou \emph{expression en coordonnées}) de $f$.
\end{definition}

\begin{center}
\begin{tikzpicture}[scale=1, >=Stealth]
    % M manifold
    \draw[thick, fill=blue!5, rounded corners=20pt]
        (-4, 1.5) -- (-1, 2) -- (0, 0.5) -- (-1.5, -1) -- (-4, -0.5) -- cycle;
    \node at (-3.5, 1.8) {$M$};

    % U domain
    \draw[thick, blue!70!black, fill=blue!15, opacity=0.5, rounded corners=8pt]
        (-3.2, 1) -- (-1.5, 1.2) -- (-1, 0) -- (-2, -0.5) -- (-3.2, 0) -- cycle;
    \node[blue!70!black] at (-3, 1.2) {$U$};

    % N manifold
    \draw[thick, fill=red!5, rounded corners=20pt]
        (3, 1.5) -- (6, 2) -- (7, 0.5) -- (5.5, -1) -- (3, -0.5) -- cycle;
    \node at (3.5, 1.8) {$N$};

    % V domain
    \draw[thick, red!70!black, fill=red!15, opacity=0.5, rounded corners=8pt]
        (3.8, 1) -- (5.5, 1.2) -- (6, 0) -- (5, -0.5) -- (3.8, 0) -- cycle;
    \node[red!70!black] at (5.8, 1.2) {$V$};

    % f arrow
    \draw[->, very thick, green!50!black] (-1.2, 0.6) -- (3.9, 0.6);
    \node[green!50!black, above] at (1.3, 0.7) {$f$};

    % R^m
    \draw[thick, fill=blue!8] (-4.5, -2.5) rectangle (-1, -4.5);
    \node at (-2.75, -4.8) {$\R^m$};
    \draw[->, thick, blue!70!black] (-2.5, -0.5) -- (-2.8, -2.5);
    \node[blue!70!black, left] at (-2.9, -1.5) {$\varphi$};

    % R^n
    \draw[thick, fill=red!8] (3, -2.5) rectangle (6.5, -4.5);
    \node at (4.75, -4.8) {$\R^n$};
    \draw[->, thick, red!70!black] (5, -0.5) -- (4.8, -2.5);
    \node[red!70!black, right] at (5.1, -1.5) {$\psi$};

    % Local expression
    \draw[->, thick, purple!70!black, decorate,
        decoration={snake, amplitude=1.5pt, segment length=8pt}]
        (-1.3, -3.5) -- (3.2, -3.5);
    \node[purple!70!black, above] at (0.95, -3.3) {$\psi \circ f \circ \varphi^{-1}$};
\end{tikzpicture}
\captionof{figure}{Application lisse entre variétés et sa représentation locale.}
\end{center}

\begin{exemple}\label{ex:application-lisse-concrete}
Considérons l'application $f \colon S^1 \to S^1$ définie par $f(e^{i\theta}) =
e^{2i\theta}$ (application de degré~$2$). Pour vérifier que $f$ est lisse, on
utilise les cartes stéréographiques. Si $\varphi_N \colon S^1 \setminus
\set{N} \to \R$ est la projection stéréographique, alors la représentation
locale est
\[
    (\varphi_N \circ f \circ \varphi_N^{-1})(t) = \frac{2t}{1 - t^2}
    \cdot \frac{1}{1 + \left(\frac{2t}{1-t^2}\right)^2 + \cdots},
\]
qui est une fonction rationnelle lisse sur son domaine de définition.
\end{exemple}

\begin{remarque}\label{rem:lisse-indep-cartes}
La condition de lissité ne dépend pas du choix des cartes. En effet, si
$(U', \varphi')$ et $(V', \psi')$ sont d'autres cartes, alors
\[
    \psi' \circ f \circ (\varphi')^{-1}
    = (\psi' \circ \psi^{-1}) \circ (\psi \circ f \circ \varphi^{-1})
    \circ (\varphi \circ (\varphi')^{-1}),
\]
qui est une composée d'applications lisses (les changements de cartes étant
des difféomorphismes lisses par définition de la structure lisse).
\end{remarque}

\begin{proposition}[Propriétés élémentaires]\label{prop:proprietes-lisses}
\begin{enumerate}[label=(\roman*)]
    \item L'application identité $\Id_M \colon M \to M$ est lisse.
    \item Si $f \colon M \to N$ et $g \colon N \to P$ sont lisses, alors
    $g \circ f \colon M \to P$ est lisse.
    \item Les inclusions d'ouverts $U \hookrightarrow M$ sont lisses.
    \item Les projections $M \times N \to M$ et $M \times N \to N$ sont lisses.
    \item Une application $f \colon M \to N$ est lisse si et seulement si, pour
    tout $p \in M$, il existe des cartes $(U, \varphi)$ autour de $p$ et
    $(V, \psi)$ autour de $f(p)$ avec $f(U) \subset V$ telles que
    $\psi \circ f \circ \varphi^{-1}$ soit lisse.
\end{enumerate}
\end{proposition}

\begin{proof}
Les points (i)--(iv) sont des vérifications directes à partir de la définition
et des propriétés de lissité dans $\R^n$.

(v) Le sens direct est clair. Pour la réciproque, soient $(U', \varphi')$ et
$(V', \psi')$ deux cartes quelconques. La lissité de $\psi' \circ f \circ
(\varphi')^{-1}$ au voisinage d'un point $\varphi'(p)$ résulte de l'écriture
\[
    \psi' \circ f \circ (\varphi')^{-1}
    = (\psi' \circ \psi^{-1}) \circ (\psi \circ f \circ \varphi^{-1})
    \circ (\varphi \circ (\varphi')^{-1})
\]
valable au voisinage de $\varphi'(p)$, chaque facteur étant lisse.
\end{proof}

\begin{notation}\label{not:Cinfty}
On note $C^\infty(M, N)$\index{C-infini@$C^\infty(M,N)$} l'ensemble des
applications lisses de $M$ dans $N$. Lorsque $N = \R$, on note simplement
$C^\infty(M) = C^\infty(M, \R)$, qui est une $\R$-algèbre commutative pour
les opérations ponctuelles.
\end{notation}

\begin{remarque}[Catégorie $C^\infty$]\label{rem:categorie-lisse}
\index{catégorie lisse}
Les variétés lisses et les applications lisses forment une catégorie, notée
$\mathbf{Man}$\index{Man@$\mathbf{Man}$} ou $\mathbf{Diff}$. Les objets sont
les variétés lisses, les morphismes sont les applications lisses, la composition
est la composition usuelle (lisse par la proposition précédente), et l'identité
est $\Id_M$.
\end{remarque}


\section{Difféomorphismes}
\label{sec:diffeomorphismes}

\begin{definition}[Difféomorphisme]\label{def:diffeomorphisme}
\index{difféomorphisme}
Soient $M$ et $N$ deux variétés lisses. Un \emph{difféomorphisme} de $M$ sur
$N$ est une bijection $f \colon M \to N$ telle que $f$ et $f^{-1}$ soient
lisses. On dit que $M$ et $N$ sont \emph{difféomorphes} et on note $M \cong N$
ou $M \approx N$.
\end{definition}

\begin{remarque}\label{rem:diffeo-invariance}
\index{invariant différentiel}
La notion de difféomorphisme est la notion d'isomorphisme dans la catégorie des
variétés lisses. Toute propriété définie en termes de la structure lisse est
invariante par difféomorphisme : dimension, orientabilité, compacité,
connexité, etc.
\end{remarque}

\begin{exemple}\label{ex:diffeomorphismes}
\begin{enumerate}[label=(\alph*)]
    \item L'application $f \colon \R \to \R$, $x \mapsto x^3$ est une bijection
    lisse, mais ce n'est \emph{pas} un difféomorphisme car $f^{-1}(y) = y^{1/3}$
    n'est pas dérivable en $0$.
    \item Toute boule ouverte $B(0, r) \subset \R^n$ est difféomorphe à $\R^n$
    via $x \mapsto x / \sqrt{r^2 - \|x\|^2}$.
    \item $\SU(2) \cong S^3$ (cf.\ 
ef{rem:SU2-S3}).
    \item $\CP^1 \cong S^2$ (cf.\ 
ef{exo:CP1-S2}).
\end{enumerate}
\end{exemple}

\begin{definition}[Groupe des difféomorphismes]\label{def:Diff-M}
\index{Diff(M)@$\Diff(M)$}
Pour une variété lisse $M$, le \emph{groupe des difféomorphismes} de $M$ est
\[
    \Diff(M) = \set{f \colon M \to M : f \text{ est un difféomorphisme}},
\]
muni de la loi de composition. C'est un groupe (en général de dimension infinie).
\end{definition}


\section{Théorème d'inversion locale}
\label{sec:inversion-locale}

\begin{definition}[Différentielle d'une application lisse]\label{def:differentielle}
\index{différentielle}\index{application tangente}
Soit $f \colon M \to N$ une application lisse et $p \in M$. La
\emph{différentielle} (ou \emph{application tangente}) de $f$ en $p$ est
l'application linéaire
\[
    \dd f_p \colon T_pM \longrightarrow T_{f(p)}N
\]
définie comme suit. En coordonnées locales $(U, \varphi)$ autour de $p$ et
$(V, \psi)$ autour de $f(p)$, si $\hat{f} = \psi \circ f \circ \varphi^{-1}$
est la représentation locale, alors
\[
    \dd f_p = (\dd \psi_{f(p)})^{-1} \circ D\hat{f}_{\varphi(p)} \circ
    \dd \varphi_p,
\]
où $D\hat{f}_{\varphi(p)}$ est la matrice jacobienne usuelle. En coordonnées,
si $\varphi = (x^1, \ldots, x^m)$ et $\psi = (y^1, \ldots, y^n)$, alors
\[
    \dd f_p\left(\frac{\partial}{\partial x^i}\bigg|_p\right)
    = \sum_{j=1}^{n} \frac{\partial \hat{f}^j}{\partial x^i}(\varphi(p))\,
    \frac{\partial}{\partial y^j}\bigg|_{f(p)}.
\]
\end{definition}

\begin{remarque}\label{rem:rang}
\index{rang}
Le \emph{rang de $f$ en $p$}\index{rang!d'une application lisse} est $\rg_p(f)
= \rg(\dd f_p) = \dim(\im(\dd f_p))$. C'est le rang de la matrice jacobienne
de toute représentation locale de $f$ en $p$.
\end{remarque}

\begin{theoreme}[Inversion locale --- version euclidienne]\label{thm:inversion-locale-Rn}
\index{théorème d'inversion locale!version euclidienne}
Soient $U \subset \R^n$ un ouvert et $f \colon U \to \R^n$ une application
$C^\infty$. Si la matrice jacobienne $Df_a$ est inversible en un point $a \in
U$, alors il existe des voisinages ouverts $V$ de $a$ et $W$ de $f(a)$ tels que
$f|_V \colon V \to W$ soit un difféomorphisme $C^\infty$.
\end{theoreme}

\begin{theoreme}[Inversion locale --- version variétés]\label{thm:inversion-locale}
\index{théorème d'inversion locale}
Soient $M$ et $N$ des variétés lisses de même dimension $n$, et $f \colon M \to
N$ une application lisse. Si $\dd f_p \colon T_pM \to T_{f(p)}N$ est un
isomorphisme en un point $p \in M$, alors il existe des voisinages ouverts $U$
de $p$ et $V$ de $f(p)$ tels que $f|_U \colon U \to V$ soit un difféomorphisme.
\end{theoreme}

\begin{proof}
Choisissons des cartes $(U_0, \varphi)$ autour de $p$ et $(V_0, \psi)$ autour
de $f(p)$ avec $f(U_0) \subset V_0$. La représentation locale $\hat{f} =
\psi \circ f \circ \varphi^{-1}$ vérifie $D\hat{f}_{\varphi(p)}$ inversible
(car c'est la matrice de $\dd f_p$ dans les bases coordonnées). Par le

ef{thm:inversion-locale-Rn}, il existe des voisinages ouverts $\hat{U}$ de
$\varphi(p)$ et $\hat{V}$ de $\psi(f(p))$ tels que $\hat{f}|_{\hat{U}} \colon
\hat{U} \to \hat{V}$ soit un difféomorphisme. Posons $U = \varphi^{-1}(\hat{U})$
et $V = \psi^{-1}(\hat{V})$. Alors $f|_U = \psi^{-1} \circ \hat{f} \circ
\varphi|_U \colon U \to V$ est un difféomorphisme (composée de
difféomorphismes).
\end{proof}


\section{Théorème des fonctions implicites}
\label{sec:fonctions-implicites}

\begin{theoreme}[Fonctions implicites --- version euclidienne]\label{thm:fonctions-implicites-Rn}
\index{théorème des fonctions implicites!version euclidienne}
Soient $U \subset \R^m \times \R^n$ un ouvert et $F \colon U \to \R^n$ une
application $C^\infty$. Soit $(a, b) \in U$ avec $F(a, b) = 0$. Si la matrice
$\frac{\partial F}{\partial y}(a, b) \in M_n(\R)$ (dérivée partielle par rapport
aux $n$ dernières variables) est inversible, alors il existe des voisinages
ouverts $V$ de $a$ dans $\R^m$ et $W$ de $b$ dans $\R^n$ et une unique
application lisse $g \colon V \to W$ tels que
\[
    \set{(x, y) \in V \times W : F(x, y) = 0} = \set{(x, g(x)) : x \in V}.
\]
\end{theoreme}

\begin{theoreme}[Fonctions implicites --- version variétés]\label{thm:fonctions-implicites}
\index{théorème des fonctions implicites}
Soient $M$ et $N$ des variétés lisses de dimensions $m$ et $n$ respectivement,
avec $m \geq n$. Soit $f \colon M \to N$ une application lisse et $q \in N$
une \emph{valeur régulière}\index{valeur régulière} de $f$ (c'est-à-dire que
$\dd f_p$ est surjective pour tout $p \in f^{-1}(q)$). Alors $f^{-1}(q)$, s'il
est non vide, est une sous-variété lisse de $M$ de dimension $m - n$.
\end{theoreme}

\begin{proof}
Soit $p \in f^{-1}(q)$. Choisissons des cartes $(\tilde{U}, \varphi)$ autour
de $p$ et $(\tilde{V}, \psi)$ autour de $q$ avec $\psi(q) = 0$ et
$f(\tilde{U}) \subset \tilde{V}$. La représentation locale $\hat{f} = \psi
\circ f \circ \varphi^{-1}$ vérifie $\hat{f}(\varphi(p)) = 0$ et $D\hat{f}_{
\varphi(p)}$ est surjective (de rang $n$).

Quitte à réordonner les coordonnées de $\R^m$, on peut supposer que les $n$
dernières colonnes de $D\hat{f}_{\varphi(p)}$ forment une matrice inversible.
On écrit $\varphi(p) = (a, b)$ avec $a \in \R^{m-n}$ et $b \in \R^n$. Par le

ef{thm:fonctions-implicites-Rn}, il existe un voisinage $V_0$ de $a$ et une
application lisse $g \colon V_0 \to \R^n$ tels que $\hat{f}(x', x'') = 0$
localement si et seulement si $x'' = g(x')$.

L'application $\Phi \colon V_0 \to \R^m$, $x' \mapsto (x', g(x'))$ est un
plongement lisse dont l'image est $\hat{f}^{-1}(0) \cap (V_0 \times \R^n)$.
En revenant sur la variété, l'application $\varphi^{-1} \circ \Phi$ fournit
une carte de $f^{-1}(q)$ au voisinage de $p$, montrant que $f^{-1}(q)$ est
une sous-variété de dimension $m - n$.
\end{proof}

\begin{corollaire}[Théorème de la valeur régulière]\label{cor:valeur-reguliere}
\index{théorème de la valeur régulière}
Soit $f \colon M \to N$ une application lisse entre variétés lisses, avec
$\dim M \geq \dim N$. Si $q \in N$ est une valeur régulière de $f$, alors
$f^{-1}(q)$ est soit vide, soit une sous-variété fermée de $M$ de codimension
$\dim N$.
\end{corollaire}

\begin{exemple}\label{ex:sphere-valeur-reguliere}
La sphère $S^{n-1} = f^{-1}(1)$ où $f \colon \R^n \to \R$, $f(x) = \|x\|^2$.
On a $\dd f_x(h) = 2\langle x, h \rangle$. Pour $x \neq 0$, $\dd f_x$ est
surjective, donc $1$ est une valeur régulière. Par le

ef{cor:valeur-reguliere}, $S^{n-1}$ est une sous-variété lisse de $\R^n$
de dimension $n - 1$.
\end{exemple}


\section{Immersions et submersions}
\label{sec:immersions-submersions}

\begin{definition}[Immersion, submersion]\label{def:immersion-submersion}
\index{immersion}\index{submersion}
Soit $f \colon M \to N$ une application lisse.
\begin{enumerate}[label=(\roman*)]
    \item $f$ est une \emph{immersion} si $\dd f_p$ est injective pour tout
    $p \in M$ (rang $= \dim M$ partout).
    \item $f$ est une \emph{submersion} si $\dd f_p$ est surjective pour tout
    $p \in M$ (rang $= \dim N$ partout).
    \item $f$ est un \emph{difféomorphisme local}\index{difféomorphisme local}
    si $\dd f_p$ est un isomorphisme pour tout $p \in M$ (nécessairement
    $\dim M = \dim N$).
\end{enumerate}
\end{definition}

\begin{remarque}\label{rem:immersion-dim}
Pour qu'une immersion $f \colon M \to N$ existe, il faut $\dim M \leq \dim N$.
Pour qu'une submersion existe, il faut $\dim M \geq \dim N$. Un
difféomorphisme local nécessite $\dim M = \dim N$.
\end{remarque}

\begin{definition}[Plongement]\label{def:plongement}
\index{plongement}
Une \emph{immersion injective} $f \colon M \to N$ est appelée un
\emph{plongement}\footnote{On dit aussi \emph{plongement lisse} ou
\emph{embedding}.} si $f$ est un homéomorphisme de $M$ sur $f(M)$ (muni de
la topologie induite). L'image $f(M)$ est alors une \emph{sous-variété} (au
sens plongé) de $N$.
\end{definition}

\begin{center}
\begin{tikzpicture}[scale=0.9, >=Stealth]
    % Immersion non injective (figure eight)
    \begin{scope}[xshift=-5cm]
        \node[above] at (0, 2) {\textbf{Immersion non injective}};
        \draw[thick, blue!70!black, ->]
            plot[smooth, domain=-1.4:1.4, samples=80]
            ({sin(2*\x r)}, {sin(\x r)*1.5});
        \fill[red] (0, 0) circle (2pt);
        \node[below] at (0, -1.8) {\small (figure en huit)};
    \end{scope}

    % Immersion injective non plongement
    \begin{scope}[xshift=0cm]
        \node[above] at (0, 2) {\textbf{Immersion injective}};
        \node[above] at (0, 1.6) {\textbf{(non plongement)}};
        \draw[thick, green!50!black, ->]
            plot[smooth, domain=0.15:6, samples=100]
            ({1.2*cos(\x r)/\x}, {1.2*sin(\x r)/\x});
        \node[below] at (0, -1.8) {\small (spirale s'accumulant)};
    \end{scope}

    % Plongement
    \begin{scope}[xshift=5cm]
        \node[above] at (0, 2) {\textbf{Plongement}};
        \draw[thick, red!70!black, ->]
            plot[smooth, domain=-1.3:1.3, samples=50]
            ({\x}, {0.5*\x*\x});
        \node[below] at (0, -1.8) {\small (parabole)};
    \end{scope}
\end{tikzpicture}
\captionof{figure}{Immersion, immersion injective (non plongement) et plongement.}
\end{center}

\begin{theoreme}[Forme locale des immersions]\label{thm:forme-locale-immersion}
\index{immersion!forme locale}\index{forme locale!immersion}
Soit $f \colon M^m \to N^n$ une immersion (avec $m \leq n$) et $p \in M$.
Alors il existe des cartes $(U, \varphi)$ autour de $p$ et $(V, \psi)$ autour
de $f(p)$ telles que
\[
    (\psi \circ f \circ \varphi^{-1})(x^1, \ldots, x^m)
    = (x^1, \ldots, x^m, 0, \ldots, 0).
\]
Autrement dit, toute immersion est localement l'inclusion canonique $\R^m
\hookrightarrow \R^n$.
\end{theoreme}

\begin{proof}
Choisissons des cartes $(\tilde{U}, \tilde{\varphi})$ autour de $p$ et
$(\tilde{V}, \tilde{\psi})$ autour de $f(p)$. Posons $\hat{f} = \tilde{\psi}
\circ f \circ \tilde{\varphi}^{-1}$ et $a = \tilde{\varphi}(p)$. Puisque $f$
est une immersion, $D\hat{f}_a$ est injective, donc de rang $m$. Quitte à
réordonner les coordonnées de $\R^n$, on peut supposer que les $m$ premières
lignes de $D\hat{f}_a$ forment une matrice inversible.

Écrivons $\hat{f}(x) = (\hat{f}_1(x), \hat{f}_2(x))$ avec $\hat{f}_1 \colon
\R^m \to \R^m$ et $\hat{f}_2 \colon \R^m \to \R^{n-m}$. La matrice
$D(\hat{f}_1)_a$ est inversible. Définissons $\Psi \colon \R^n \to \R^n$ par
\[
    \Psi(y_1, y_2) = (\hat{f}_1^{-1}(y_1),\ y_2 - \hat{f}_2(\hat{f}_1^{-1}(y_1)))
\]
localement autour de $\hat{f}(a)$ (en utilisant l'inversion locale pour
$\hat{f}_1$). Alors
\[
    (\Psi \circ \hat{f})(x) = (\hat{f}_1^{-1}(\hat{f}_1(x)),\ \hat{f}_2(x) -
    \hat{f}_2(x)) = (x, 0).
\]
On pose $\varphi = \tilde{\varphi}|$ et $\psi = \Psi \circ \tilde{\psi}$
(sur des domaines restreints appropriés).
\end{proof}

\begin{theoreme}[Forme locale des submersions]\label{thm:forme-locale-submersion}
\index{submersion!forme locale}\index{forme locale!submersion}
Soit $f \colon M^m \to N^n$ une submersion (avec $m \geq n$) et $p \in M$.
Alors il existe des cartes $(U, \varphi)$ autour de $p$ et $(V, \psi)$ autour
de $f(p)$ telles que
\[
    (\psi \circ f \circ \varphi^{-1})(x^1, \ldots, x^m)
    = (x^1, \ldots, x^n).
\]
Autrement dit, toute submersion est localement une projection canonique $\R^m
\twoheadrightarrow \R^n$.
\end{theoreme}

\begin{proof}
Choisissons des cartes $(\tilde{U}, \tilde{\varphi})$ autour de $p$ et
$(\tilde{V}, \tilde{\psi})$ autour de $f(p)$, et posons $\hat{f} = \tilde{\psi}
\circ f \circ \tilde{\varphi}^{-1}$, $a = \tilde{\varphi}(p)$. Puisque $f$
est une submersion, $D\hat{f}_a$ est surjective de rang $n$. Quitte à
réordonner les coordonnées de $\R^m$, les $n$ premières colonnes de
$D\hat{f}_a$ forment une matrice inversible.

Définissons $\Phi \colon \R^m \to \R^m$ par
\[
    \Phi(x) = (\hat{f}(x),\ x^{n+1}, \ldots, x^m).
\]
Alors $D\Phi_a$ est inversible (sa matrice est triangulaire par blocs avec des
blocs diagonaux inversibles). Par le théorème d'inversion locale, $\Phi$ est
un difféomorphisme local autour de $a$. On a
\[
    (\hat{f} \circ \Phi^{-1})(y^1, \ldots, y^m) = (y^1, \ldots, y^n),
\]
car la première composante de $\Phi$ est exactement $\hat{f}$. On pose
$\varphi = \Phi \circ \tilde{\varphi}$ et $\psi = \tilde{\psi}$.
\end{proof}

\begin{corollaire}\label{cor:submersion-ouverte}
\index{submersion!application ouverte}
Toute submersion est une application ouverte.
\end{corollaire}

\begin{proof}
La projection canonique $\R^m \to \R^n$ est ouverte. Par la forme locale des
submersions, une submersion est localement ouverte, donc ouverte.
\end{proof}

\begin{corollaire}\label{cor:submersion-fibre}
Soit $f \colon M \to N$ une submersion et $q \in N$. Alors $f^{-1}(q)$ est une
sous-variété lisse fermée de $M$ de dimension $\dim M - \dim N$.
\end{corollaire}

\begin{proof}
Comme $f$ est une submersion, $q$ est automatiquement une valeur régulière. Le
résultat découle du 
ef{thm:fonctions-implicites}. La fermeture provient de
la continuité de $f$.
\end{proof}


\section{Théorème de rang constant}
\label{sec:rang-constant}

\begin{theoreme}[Rang constant]\label{thm:rang-constant}
\index{théorème de rang constant}
Soit $f \colon M^m \to N^n$ une application lisse de rang constant $r$ (c'est-à-dire
$\rg_p(f) = r$ pour tout $p \in M$). Alors pour tout $p \in M$, il existe des
cartes $(U, \varphi)$ autour de $p$ et $(V, \psi)$ autour de $f(p)$ telles que
\[
    (\psi \circ f \circ \varphi^{-1})(x^1, \ldots, x^m)
    = (x^1, \ldots, x^r, 0, \ldots, 0).
\]
\end{theoreme}

\begin{proof}
Choisissons des cartes et posons $\hat{f} = \psi_0 \circ f \circ \varphi_0^{-1}$
avec $\hat{f}(0) = 0$ (quitte à translater). Puisque $\rg(D\hat{f}_0) = r$,
quitte à permuter les coordonnées, les $r$ premières lignes et $r$ premières
colonnes de $D\hat{f}_0$ forment une matrice inversible. Écrivons $x = (x',
x'') \in \R^r \times \R^{m-r}$ et $\hat{f}(x) = (P(x), Q(x))$ avec $P \colon
\R^m \to \R^r$ et $Q \colon \R^m \to \R^{n-r}$.

\textbf{Premier changement de coordonnées (dans la source).}
Posons $\Phi(x', x'') = (P(x', x''), x'')$. Alors $D\Phi_0$ est inversible
(bloc triangulaire), donc $\Phi$ est un difféomorphisme local. Dans les
nouvelles coordonnées $u = \Phi(x)$, on a
\[
    (\hat{f} \circ \Phi^{-1})(u', u'') = (u', \widetilde{Q}(u', u''))
\]
pour une certaine fonction lisse $\widetilde{Q}$. La condition de rang constant
$r$ implique que $D_{u''}\widetilde{Q} = 0$ identiquement, donc
$\widetilde{Q}(u', u'') = \widetilde{Q}(u', 0) =: R(u')$ ne dépend que de
$u'$.

\textbf{Second changement de coordonnées (dans le but).}
Posons $\Psi(y', y'') = (y',\ y'' - R(y'))$. Alors $D\Psi_0$ est inversible
et
\[
    (\Psi \circ \hat{f} \circ \Phi^{-1})(u', u'') = (u', 0).
\]
On obtient les cartes cherchées par composition.
\end{proof}

\begin{remarque}\label{rem:rang-constant-generalise}
Les théorèmes de forme locale des immersions et submersions
(
ef{thm:forme-locale-immersion,thm:forme-locale-submersion}) sont des cas
particuliers du théorème de rang constant, avec $r = m$ et $r = n$
respectivement.
\end{remarque}


\section{Exemples d'applications lisses}
\label{sec:exemples-applications}

\subsection{Projections et inclusions}

\begin{exemple}[Projection canonique]\label{ex:projection}
\index{projection}
La projection $\pi \colon M \times N \to M$, $(p, q) \mapsto p$ est une
submersion surjective. Sa différentielle en $(p, q)$ est la projection
$T_{(p,q)}(M \times N) \cong T_pM \oplus T_qN \to T_pM$.
\end{exemple}

\begin{exemple}[Inclusion de sous-variété]\label{ex:inclusion}
\index{inclusion}
Si $S \subset M$ est une sous-variété, l'inclusion $\iota \colon S
\hookrightarrow M$ est un plongement. Sa différentielle $\dd\iota_p \colon T_pS
\to T_pM$ est l'injection canonique de $T_pS$ dans $T_pM$ (identifiant $T_pS$
à un sous-espace de $T_pM$).
\end{exemple}

\subsection{L'application de Hopf}

\begin{exemple}[Fibration de Hopf]\label{ex:hopf}
\index{fibration de Hopf}\index{Hopf}
L'\emph{application de Hopf} est l'application $h \colon S^3 \to S^2$ définie
comme suit. On identifie $S^3 \subset \C^2$ par $S^3 = \set{(z_1, z_2) \in
\C^2 : |z_1|^2 + |z_2|^2 = 1}$, et $S^2 \cong \CP^1$ via la projection
stéréographique. On pose
\[
    h(z_1, z_2) = [z_1 : z_2] \in \CP^1 \cong S^2.
\]
En coordonnées réelles, si $z_1 = x_0 + ix_1$ et $z_2 = x_2 + ix_3$, on a
\[
    h(x_0, x_1, x_2, x_3) = (2(x_0 x_2 + x_1 x_3),\ 2(x_1 x_2 - x_0 x_3),
    \ x_0^2 + x_1^2 - x_2^2 - x_3^2).
\]
\end{exemple}

\begin{proposition}\label{prop:hopf-submersion}
L'application de Hopf $h \colon S^3 \to S^2$ est une submersion surjective.
Pour tout $p \in S^2$, la fibre $h^{-1}(p)$ est difféomorphe à $S^1$.
\end{proposition}

\begin{proof}
La surjectivité est claire : pour $[w_1 : w_2] \in \CP^1$, le point $(w_1,
w_2)/\|(w_1, w_2)\| \in S^3$ est dans la préimage.

Pour la submersion, il suffit de vérifier que le rang est $2$ en tout point. En
utilisant les cartes de $S^3$ et $S^2$, on calcule la matrice jacobienne et on
vérifie que son rang est $2$ (utilisant la contrainte $|z_1|^2 + |z_2|^2 = 1$).

La fibre au-dessus de $[w_1 : w_2]$ est $\set{e^{i\theta}(w_1, w_2) /
\|(w_1, w_2)\| : \theta \in [0, 2\pi)} \cong S^1$.
\end{proof}

\subsection{Difféomorphismes locaux}

\begin{exemple}\label{ex:diffeo-local}
\index{difféomorphisme local!exemples}
\begin{enumerate}[label=(\alph*)]
    \item L'application $\exp \colon \R \to \R^*_+$, $x \mapsto e^x$ est un
    difféomorphisme (global).
    \item L'application $p \colon \R \to S^1$, $t \mapsto e^{2\pi i t}$ est un
    difféomorphisme local et une submersion, mais pas un difféomorphisme global
    (elle n'est pas injective).
    \item La projection $\pi \colon S^n \to \RP^n$ (quotient par
    l'antipodie) est un difféomorphisme local (revêtement double).
\end{enumerate}
\end{exemple}

\begin{proposition}\label{prop:diffeo-local-caract}
Soit $f \colon M \to N$ une application lisse entre variétés de même
dimension. Alors $f$ est un difféomorphisme local si et seulement si $\dd f_p$
est un isomorphisme pour tout $p \in M$.
\end{proposition}

\begin{proof}
Le sens direct est clair. La réciproque est le théorème d'inversion locale
(
ef{thm:inversion-locale}).
\end{proof}

\begin{proposition}\label{prop:diffeo-global}
Soit $f \colon M \to N$ une application lisse bijective entre variétés de même
dimension. Si $f$ est une immersion (ou, de manière équivalente, une
submersion), alors $f$ est un difféomorphisme.
\end{proposition}

\begin{proof}
Si $\dim M = \dim N = n$ et $f$ est une immersion, alors $\dd f_p$ est
injective de $T_pM$ dans $T_{f(p)}N$, deux espaces de dimension $n$, donc
$\dd f_p$ est un isomorphisme. Par le théorème d'inversion locale, $f$ est un
difféomorphisme local. Comme $f$ est bijective et continue, $f^{-1}$ est
continue (toute bijection continue d'un espace compact vers un espace séparé
est un homéomorphisme, et dans le cas non compact, l'inversion locale fournit
directement la lissité de $f^{-1}$ point par point). Donc $f^{-1}$ est lisse.
\end{proof}


\section{Diagrammes commutatifs et catégorie \texorpdfstring{$C^\infty$}{C-infini}}
\label{sec:diagrammes}

La structure catégorielle de $\mathbf{Man}$ se manifeste naturellement par des
diagrammes commutatifs\index{diagramme commutatif}.

\begin{exemple}\label{ex:diagramme-quotient}
Soit $G$ un groupe qui agit librement et proprement sur une variété $M$.
Alors l'espace quotient $M/G$ admet une unique structure de variété lisse
telle que la projection $\pi \colon M \to M/G$ soit une submersion. Si $f
\colon M \to N$ est une application lisse $G$-invariante, alors il existe une
unique application lisse $\bar{f} \colon M/G \to N$ telle que le diagramme
suivant commute :
\[
\begin{tikzcd}
    M \arrow[r, "f"] \arrow[d, "\pi"'] & N \\
    M/G \arrow[ur, "\bar{f}"', dashed]
\end{tikzcd}
\]
\end{exemple}

\begin{exemple}\label{ex:diagramme-produit}
Le produit de variétés satisfait la propriété universelle suivante : pour
toutes applications lisses $f \colon P \to M$ et $g \colon P \to N$, il existe
une unique application lisse $(f, g) \colon P \to M \times N$ telle que
\[
\begin{tikzcd}
    & P \arrow[dl, "f"'] \arrow[dr, "g"] \arrow[d, dashed, "{(f,g)}"] & \\
    M & M \times N \arrow[l, "\pi_M"] \arrow[r, "\pi_N"'] & N
\end{tikzcd}
\]
commute.
\end{exemple}


\section{Exercices}
\label{sec:exercices-ch2}

\begin{exercice}[$\star\star$]\label{exo:composition-immersions}
\index{exercice!composition d'immersions}
Soient $f \colon M \to N$ et $g \colon N \to P$ des applications lisses.
\begin{enumerate}[label=(\alph*)]
    \item Montrer que si $f$ et $g$ sont des immersions, alors $g \circ f$ est
    une immersion.
    \item Montrer que si $f$ et $g$ sont des submersions, alors $g \circ f$ est
    une submersion.
    \item Montrer que si $g \circ f$ est une immersion, alors $f$ est une
    immersion.
    \item Montrer que si $g \circ f$ est une submersion, alors $g$ est une
    submersion.
\end{enumerate}
\end{exercice}

\begin{exercice}[$\star\star$]\label{exo:graphe-sous-variete}
\index{exercice!graphe comme sous-variété}
Soit $f \colon M \to N$ une application lisse. Montrer que le graphe
\[
    \Gamma_f = \set{(p, f(p)) : p \in M} \subset M \times N
\]
est une sous-variété lisse de $M \times N$, difféomorphe à $M$.
\end{exercice}

\begin{exercice}[$\star\star$]\label{exo:det-submersion}
\index{exercice!déterminant comme submersion}
Montrer que l'application $\det \colon \GL(n, \R) \to \R^*$ est une submersion.
En déduire que $\SL(n, \R) = \det^{-1}(1)$ est une sous-variété lisse de
$\GL(n, \R)$ de codimension~$1$.
\emph{Indication} : calculer $\dd(\det)_A(H) = \det(A) \cdot
\operatorname{tr}(A^{-1}H)$.
\end{exercice}

\begin{exercice}[$\star\star\star$]\label{exo:hopf-fibration-detail}
\index{exercice!fibration de Hopf}
Soit $h \colon S^3 \to S^2$ l'application de Hopf.
\begin{enumerate}[label=(\alph*)]
    \item Montrer en détail que $h$ est lisse.
    \item Calculer $\dd h_{(1,0)}$ dans les cartes stéréographiques et vérifier
    qu'il est surjectif.
    \item Montrer que deux fibres distinctes $h^{-1}(p)$ et $h^{-1}(q)$ sont
    des cercles enlacés dans $S^3$.
\end{enumerate}
\end{exercice}

\begin{exercice}[$\star\star$]\label{exo:inversion-locale-application}
\index{exercice!théorème d'inversion locale}
Soit $f \colon \R^2 \to \R^2$ définie par $f(x, y) = (e^x \cos y, e^x \sin y)$.
\begin{enumerate}[label=(\alph*)]
    \item Montrer que $f$ est un difféomorphisme local en tout point.
    \item Déterminer $f(\R^2)$ et montrer que $f$ n'est pas un difféomorphisme
    global.
    \item Trouver un ouvert maximal $U \subset \R^2$ tel que $f|_U$ soit un
    difféomorphisme sur son image.
\end{enumerate}
\end{exercice}

\begin{exercice}[$\star\star$]\label{exo:submersion-section}
\index{exercice!section d'une submersion}
Soit $f \colon M \to N$ une submersion surjective.
\begin{enumerate}[label=(\alph*)]
    \item Montrer que pour tout $q \in N$, il existe un voisinage ouvert $V$ de
    $q$ et une application lisse $s \colon V \to M$ (appelée \emph{section
    locale}\index{section locale}) telle que $f \circ s = \Id_V$.
    \item Donner un exemple montrant qu'il n'existe pas toujours de section
    globale lisse $s \colon N \to M$.
\end{enumerate}
\end{exercice}

\begin{exercice}[$\star\star\star$]\label{exo:submersion-propre}
\index{exercice!submersion propre}
Soit $f \colon M \to N$ une submersion propre surjective (propre signifie que
la préimage de tout compact est compacte). Montrer que $f$ est une
\emph{fibration localement triviale}\index{fibration localement triviale} :
pour tout $q \in N$, il existe un voisinage ouvert $V$ de $q$ et un
difféomorphisme $\Phi \colon f^{-1}(V) \to V \times f^{-1}(q)$ tel que
$\pi_1 \circ \Phi = f|_{f^{-1}(V)}$, où $\pi_1$ est la projection sur le
premier facteur.
\emph{Indication} : C'est le théorème d'Ehresmann\index{théorème d'Ehresmann}.
Utiliser un champ de vecteurs horizontal et son flot.
\end{exercice}

\begin{exercice}[$\star\star$]\label{exo:image-inverse-sous-variete}
\index{exercice!image inverse d'une sous-variété}
Soit $f \colon M \to N$ une application lisse et $S \subset N$ une sous-variété.
On dit que $f$ est \emph{transverse} à $S$, et on note $f \pitchfork S$, si
pour tout $p \in f^{-1}(S)$,
\[
    \im(\dd f_p) + T_{f(p)}S = T_{f(p)}N.
\]
Montrer que si $f \pitchfork S$, alors $f^{-1}(S)$ est une sous-variété de $M$
de codimension $\operatorname{codim}(S)$.
\emph{Indication} : Se ramener localement au cas où $S = f_0^{-1}(0)$ pour une
submersion $f_0$, puis appliquer le théorème de la valeur régulière à
$f_0 \circ f$.
\end{exercice}

\begin{exercice}[$\star\star\star$]\label{exo:rang-semi-continu}
\index{exercice!semi-continuité du rang}
Soit $f \colon M \to N$ une application lisse.
\begin{enumerate}[label=(\alph*)]
    \item Montrer que la fonction $p \mapsto \rg_p(f)$ est semi-continue
    inférieurement : pour tout $p_0 \in M$, il existe un voisinage $U$ de
    $p_0$ tel que $\rg_p(f) \geq \rg_{p_0}(f)$ pour tout $p \in U$.
    \item En déduire que l'ensemble $\set{p \in M : \rg_p(f) \geq r}$ est
    ouvert dans $M$ pour tout $r \in \N$.
    \item Montrer que si $M$ est connexe et $\rg_p(f) = r$ pour tout $p$ dans
    un ouvert dense de $M$, alors l'ensemble des points où le rang est
    strictement inférieur à $r$ est contenu dans un fermé d'intérieur vide.
\end{enumerate}
\end{exercice}

% === FIN DU CHUNK preamble+ch1-2 ===
% === DEBUT DU CHUNK ch3-4 ===

% ============================================================
%  CHAPITRE 3 : FIBRÉS TANGENTS ET COTANGENTS
% ============================================================
\chapter{Fibrés tangents et cotangents}\label{ch:fibres-tangents}
\minitoc\bigskip

\section{L'espace tangent en un point}\label{sec:espace-tangent}

Soit $M$ une variété lisse\index{variété lisse} de dimension $n$ et $p \in M$.
L'espace tangent\index{espace tangent} en $p$ admet plusieurs définitions équivalentes,
chacune éclairant un aspect différent de la géométrie différentielle.

% ----------------------------------------------------------
%  Définition via les dérivations
% ----------------------------------------------------------
\subsection{Définition via les dérivations}\label{subsec:derivations}

Notons $C^\infty(M)$\index{fonction lisse} l'algèbre des fonctions lisses $M \to \R$.

\begin{definition}[Dérivation en un point]\label{def:derivation}
Une \textbf{dérivation}\index{dérivation} en $p \in M$ est une application
$\R$-linéaire $v : C^\infty(M) \to \R$ satisfaisant la \emph{règle de Leibniz} :
\[
  v(fg) = f(p)\,v(g) + g(p)\,v(f), \qquad \forall f,g \in C^\infty(M).
\]
L'ensemble de toutes les dérivations en $p$ est noté $T_pM$\index{TpM@$T_pM$}
et s'appelle l'\textbf{espace tangent à $M$ en $p$}.
\end{definition}

\begin{lemme}[Annulation sur les constantes]\label{lem:derivation-constante}
Si $v \in T_pM$ et $c \in \R$ est une fonction constante, alors $v(c) = 0$.
\end{lemme}

\begin{proof}
Posons $f = g = 1$ dans la règle de Leibniz :
$v(1) = 1 \cdot v(1) + 1 \cdot v(1) = 2\,v(1)$, d'où $v(1) = 0$.
Par linéarité, $v(c) = c \cdot v(1) = 0$.
\end{proof}

\begin{proposition}[Base de l'espace tangent]\label{prop:base-tangent}
Soit $(U, \varphi = (x^1, \ldots, x^n))$ une carte\index{carte locale} autour de $p$.
Les dérivations partielles\index{dérivée partielle}
\[
  \left.\frac{\partial}{\partial x^i}\right|_p : f \longmapsto
  \frac{\partial (f \circ \varphi^{-1})}{\partial r^i}\bigg|_{\varphi(p)},
  \qquad i = 1, \ldots, n,
\]
forment une base de $T_pM$. En particulier, $\dim T_pM = n$.
\end{proposition}

\begin{proof}
\emph{Indépendance linéaire.}
Si $\sum_i a^i \left.\frac{\partial}{\partial x^i}\right|_p = 0$, on évalue
cette dérivation sur la fonction coordonnée $x^j$ :
\[
  0 = \sum_i a^i \left.\frac{\partial x^j}{\partial x^i}\right|_p
    = \sum_i a^i \delta^j_i = a^j.
\]
Donc $a^j = 0$ pour tout $j$.

\emph{Engendrement.}
Soit $v \in T_pM$ et posons $v^i = v(x^i)$. Montrons que
$v = \sum_i v^i \left.\frac{\partial}{\partial x^i}\right|_p$.
Pour toute $f \in C^\infty(M)$, on écrit $\hat{f} = f \circ \varphi^{-1}$ au voisinage
de $\hat{p} = \varphi(p)$. Par le lemme de Hadamard\index{lemme de Hadamard},
il existe des fonctions lisses $g_i$ au voisinage de $\hat{p}$ telles que
\[
  \hat{f}(x) = \hat{f}(\hat{p}) + \sum_{i=1}^n (x^i - \hat{p}^i)\,g_i(x),
  \qquad g_i(\hat{p}) = \frac{\partial \hat{f}}{\partial r^i}(\hat{p}).
\]
En appliquant $v$ et en utilisant le 
ef{lem:derivation-constante} :
\[
  v(f) = \sum_{i=1}^n v(x^i - p^i)\,g_i(\hat{p})
       = \sum_{i=1}^n v^i \frac{\partial \hat{f}}{\partial r^i}(\hat{p})
       = \sum_{i=1}^n v^i \left.\frac{\partial f}{\partial x^i}\right|_p. \qedhere
\]
\end{proof}

% ----------------------------------------------------------
%  Définition via les courbes
% ----------------------------------------------------------
\subsection{Définition via les classes d'équivalence de courbes}\label{subsec:tangent-courbes}

\begin{definition}[Vecteur tangent via les courbes]\label{def:tangent-courbe}
Soit $p \in M$. Deux courbes lisses\index{courbe lisse} $\gamma_1, \gamma_2 : (-\varepsilon, \varepsilon) \to M$
avec $\gamma_1(0) = \gamma_2(0) = p$ sont dites \textbf{équivalentes}\index{équivalence de courbes}
si, pour une (et donc toute) carte $(U, \varphi)$ autour de $p$,
\[
  (\varphi \circ \gamma_1)'(0) = (\varphi \circ \gamma_2)'(0).
\]
L'espace tangent $T_pM$ est l'ensemble des classes d'équivalence, noté $[\gamma]_p$.
\end{definition}

\begin{remarque}\label{rem:independance-carte-courbe}
La relation d'équivalence ne dépend pas du choix de carte : si $(V, \psi)$
est une autre carte autour de $p$, alors
\[
  (\psi \circ \gamma)'(0) = D(\psi \circ \varphi^{-1})_{\varphi(p)} \cdot (\varphi \circ \gamma)'(0),
\]
et la matrice jacobienne $D(\psi \circ \varphi^{-1})_{\varphi(p)}$ est inversible
car $\psi \circ \varphi^{-1}$ est un difféomorphisme local.
\end{remarque}

% ----------------------------------------------------------
%  Équivalence des définitions
% ----------------------------------------------------------
\subsection{Équivalence des définitions}\label{subsec:equivalence-tangent}

\begin{theoreme}[Équivalence des définitions de l'espace tangent]\label{thm:equiv-tangent}
Soit $M$ une variété lisse de dimension $n$ et $p \in M$.
L'application
\[
  \Phi : \set{[\gamma]_p} \longrightarrow T_pM, \qquad
  [\gamma]_p \longmapsto v_\gamma,
  \quad \text{où } v_\gamma(f) = (f \circ \gamma)'(0),
\]
est un isomorphisme d'espaces vectoriels\index{isomorphisme!espace tangent}.
\end{theoreme}

\begin{proof}
\emph{Bonne définition.}
Si $\gamma_1 \sim \gamma_2$, alors pour toute $f \in C^\infty(M)$ et
toute carte $(U, \varphi)$ autour de $p$ :
\[
  (f \circ \gamma_i)'(0) = D(f \circ \varphi^{-1})_{\varphi(p)} \cdot (\varphi \circ \gamma_i)'(0).
\]
Comme $(\varphi \circ \gamma_1)'(0) = (\varphi \circ \gamma_2)'(0)$,
on obtient $v_{\gamma_1} = v_{\gamma_2}$.

\emph{$\Phi$ est une dérivation.}
Pour $f, g \in C^\infty(M)$ :
\[
  v_\gamma(fg) = ((fg) \circ \gamma)'(0) = f(p)\,(g \circ \gamma)'(0) + g(p)\,(f \circ \gamma)'(0)
  = f(p)\,v_\gamma(g) + g(p)\,v_\gamma(f).
\]

\emph{Injectivité.}
Si $v_\gamma = 0$, alors pour chaque coordonnée $x^i$ :
$0 = v_\gamma(x^i) = (x^i \circ \gamma)'(0)$, ce qui signifie $(\varphi \circ \gamma)'(0) = 0$,
d'où $[\gamma]_p$ est la classe du vecteur nul.

\emph{Surjectivité.}
Soit $v = \sum_i a^i \left.\frac{\partial}{\partial x^i}\right|_p \in T_pM$.
Définissons $\gamma(t) = \varphi^{-1}(\varphi(p) + t\,a)$ où $a = (a^1, \ldots, a^n)$.
Alors $\gamma(0) = p$ et $(\varphi \circ \gamma)'(0) = a$, de sorte que
$v_\gamma(x^i) = a^i = v(x^i)$ pour tout $i$. Par la 
ef{prop:base-tangent},
$v_\gamma = v$.
\end{proof}

\begin{figure}[ht]
\centering
\begin{tikzpicture}[scale=1.2]
  % Variété (ellipse)
  \draw[thick, fill=blue!5] (0,0) ellipse (2.5 and 1.5);
  \node at (1.8, -1.2) {$M$};
  % Point
  \fill (0,0.2) circle (1.5pt) node[below left] {$p$};
  % Plan tangent
  \draw[blue!50, fill=blue!10, opacity=0.4]
    (-1.8, 1.8) -- (1.8, 1.8) -- (2.2, -0.2) -- (-1.4, -0.2) -- cycle;
  \node[blue!70!black] at (2.0, 1.6) {$T_pM$};
  % Courbe gamma
  \draw[thick, red!70!black, ->]
    (-1.2, -0.5) .. controls (-0.3, -0.1) and (0.3, 0.5) .. (1.2, 0.8);
  \node[red!70!black] at (1.4, 0.6) {$\gamma$};
  % Vecteur tangent
  \draw[thick, ->, blue!70!black] (0, 0.2) -- (1.0, 1.0);
  \node[blue!70!black] at (0.6, 1.1) {$\gamma'(0)$};
\end{tikzpicture}
\caption{Un vecteur tangent comme vecteur vitesse d'une courbe sur la variété.}\label{fig:vecteur-tangent}
\end{figure}

% ----------------------------------------------------------
%  Différentielle
% ----------------------------------------------------------
\section{Différentielle d'une application lisse}\label{sec:differentielle}

\begin{definition}[Différentielle]\label{def:differentielle}
Soient $M, N$ des variétés lisses et $f : M \to N$ une application lisse\index{application lisse}.
La \textbf{différentielle}\index{différentielle} (ou \emph{application tangente})
de $f$ en $p \in M$ est l'application linéaire
\[
  \dd f_p : T_pM \longrightarrow T_{f(p)}N
\]
définie par
\[
  (\dd f_p(v))(g) = v(g \circ f), \qquad \forall\, v \in T_pM,\; g \in C^\infty(N).
\]
\end{definition}

\begin{remarque}\label{rem:diff-courbe}
En termes de courbes : si $v = [\gamma]_p$, alors
$\dd f_p(v) = [f \circ \gamma]_{f(p)}$.
\end{remarque}

\begin{proposition}[Expression en coordonnées]\label{prop:diff-coordonnees}
Soient $(U, \varphi = (x^1, \ldots, x^m))$ et $(V, \psi = (y^1, \ldots, y^n))$
des cartes autour de $p$ et $f(p)$ respectivement. Si $\hat{f} = \psi \circ f \circ \varphi^{-1}$,
alors
\[
  \dd f_p\!\left(\left.\frac{\partial}{\partial x^i}\right|_p\right)
  = \sum_{j=1}^n \frac{\partial \hat{f}^j}{\partial r^i}(\varphi(p))
    \left.\frac{\partial}{\partial y^j}\right|_{f(p)}.
\]
La matrice de $\dd f_p$ dans ces bases est la \textbf{matrice jacobienne}\index{matrice jacobienne}
$J_{\hat{f}}(\varphi(p))$.
\end{proposition}

\begin{proof}
Pour toute $g \in C^\infty(N)$, posons $\hat{g} = g \circ \psi^{-1}$. Alors :
\begin{align*}
  \dd f_p\!\left(\left.\frac{\partial}{\partial x^i}\right|_p\right)(g)
  &= \left.\frac{\partial}{\partial x^i}\right|_p(g \circ f)
   = \frac{\partial(\hat{g} \circ \hat{f})}{\partial r^i}(\varphi(p)) \\
  &= \sum_{j=1}^n \frac{\partial \hat{g}}{\partial r^j}(\hat{f}(\varphi(p)))
     \cdot \frac{\partial \hat{f}^j}{\partial r^i}(\varphi(p))
   = \sum_{j=1}^n \frac{\partial \hat{f}^j}{\partial r^i}(\varphi(p))
     \left.\frac{\partial g}{\partial y^j}\right|_{f(p)}. \qedhere
\end{align*}
\end{proof}

\begin{theoreme}[Règle de chaîne]\label{thm:chaine}
\index{règle de chaîne}
Soient $f : M \to N$ et $g : N \to P$ des applications lisses. Pour tout $p \in M$,
\[
  \dd(g \circ f)_p = \dd g_{f(p)} \circ \dd f_p.
\]
\end{theoreme}

\begin{proof}
Pour $v \in T_pM$ et $h \in C^\infty(P)$ :
\[
  \dd(g \circ f)_p(v)(h) = v(h \circ g \circ f)
  = \dd f_p(v)(h \circ g)
  = \dd g_{f(p)}(\dd f_p(v))(h). \qedhere
\]
\end{proof}

\begin{corollaire}[Fonctorialité]\label{cor:fonctorialite}
\index{fonctorialité}
L'assignation $M \mapsto TM$, $f \mapsto \dd f$ définit un foncteur
de la catégorie des variétés lisses vers la catégorie des espaces vectoriels.
En particulier :
\begin{enumerate}[label=(\roman*)]
  \item $\dd(\Id_M)_p = \Id_{T_pM}$ pour tout $p \in M$.
  \item Si $f : M \to N$ est un difféomorphisme\index{difféomorphisme}, alors
    $\dd f_p$ est un isomorphisme et $(\dd f_p)^{-1} = \dd(f^{-1})_{f(p)}$.
\end{enumerate}
\end{corollaire}

\begin{proof}
Le point (i) est immédiat. Pour (ii), la règle de chaîne donne
$\dd(f^{-1} \circ f)_p = \dd(f^{-1})_{f(p)} \circ \dd f_p = \Id_{T_pM}$
et de même dans l'autre sens.
\end{proof}

\begin{exemple}[Différentielle de la projection stéréographique]\label{ex:diff-stereo}
\index{projection stéréographique!différentielle}
Soit $\sigma_N : S^n \setminus \set{N} \to \R^n$ la projection stéréographique depuis
le pôle nord $N = (0, \ldots, 0, 1)$, définie par
\[
  \sigma_N(x^1, \ldots, x^{n+1}) = \frac{1}{1 - x^{n+1}}(x^1, \ldots, x^n).
\]
C'est un difféomorphisme, et sa différentielle en un point $p = (p^1, \ldots, p^{n+1}) \in S^n$
est l'isomorphisme linéaire
\[
  \dd(\sigma_N)_p : T_pS^n \longrightarrow T_{\sigma_N(p)}\R^n \cong \R^n.
\]
Pour $v \in T_pS^n \subset \R^{n+1}$ (avec $\langle p, v \rangle = 0$), un calcul
direct donne
\[
  \dd(\sigma_N)_p(v) = \frac{1}{1 - p^{n+1}}\left(v^1, \ldots, v^n\right)
  + \frac{v^{n+1}}{(1 - p^{n+1})^2}\left(p^1, \ldots, p^n\right).
\]
\end{exemple}

\begin{exemple}[Différentielle de l'exponentielle matricielle]\label{ex:diff-exp}
\index{exponentielle matricielle!différentielle}
L'application exponentielle $\exp : M_n(\R) \to \GL(n, \R)$ définie par
$\exp(A) = \sum_{k=0}^\infty \frac{A^k}{k!}$ est lisse. En identifiant
$T_0 M_n(\R) \cong M_n(\R)$ et $T_{I_n}\GL(n,\R) \cong M_n(\R)$, sa
différentielle en $0$ est
\[
  \dd(\exp)_0(B) = \left.\frac{\dd}{\dd t}\right|_{t=0} \exp(tB) = B.
\]
Autrement dit, $\dd(\exp)_0 = \Id_{M_n(\R)}$, ce qui montre que $\exp$ est un
difféomorphisme local au voisinage de $0$.
\end{exemple}

% ----------------------------------------------------------
%  Fibré tangent TM
% ----------------------------------------------------------
\section{Le fibré tangent}\label{sec:fibre-tangent}

\begin{definition}[Fibré tangent]\label{def:fibre-tangent}
Soit $M$ une variété lisse de dimension $n$. Le \textbf{fibré tangent}\index{fibré tangent}
de $M$ est l'union disjointe
\[
  TM = \bigsqcup_{p \in M} T_pM = \set{(p, v) : p \in M,\; v \in T_pM},
\]
muni de la projection canonique $\pi : TM \to M$, $\pi(p, v) = p$.
\end{definition}

\begin{theoreme}[Structure de variété lisse sur $TM$]\label{thm:TM-variete}
\index{fibré tangent!structure de variété}
Soit $M$ une variété lisse de dimension $n$. Alors $TM$ admet une unique structure
de variété lisse de dimension $2n$ telle que, pour toute carte $(U, \varphi)$ de $M$,
l'application
\[
  \widetilde{\varphi} : \pi^{-1}(U) \longrightarrow \varphi(U) \times \R^n, \qquad
  \left(p,\, \sum_{i=1}^n v^i \left.\frac{\partial}{\partial x^i}\right|_p\right)
  \longmapsto (\varphi(p),\, v^1, \ldots, v^n),
\]
soit un difféomorphisme. De plus, $\pi : TM \to M$ est une submersion\index{submersion} lisse.
\end{theoreme}

\begin{proof}
\emph{Topologie.}
On munit $TM$ de la topologie la plus fine rendant toutes les applications
$\widetilde{\varphi}$ des homéomorphismes. Les ouverts de cette topologie sont
les unions d'ensembles de la forme $\widetilde{\varphi}^{-1}(W)$ avec $W$ ouvert
de $\varphi(U) \times \R^n$. Cette topologie est de Hausdorff et à base dénombrable
car $M$ possède ces propriétés.

\emph{Atlas lisse.}
Les changements de cartes sont de la forme
\[
  \widetilde{\psi} \circ \widetilde{\varphi}^{-1} : \varphi(U \cap V) \times \R^n
  \longrightarrow \psi(U \cap V) \times \R^n,
\]
\[
  (x, v) \longmapsto \left(\psi \circ \varphi^{-1}(x),\;
  D(\psi \circ \varphi^{-1})_x \cdot v\right).
\]
Comme $\psi \circ \varphi^{-1}$ est un difféomorphisme, cette application est lisse,
et la collection $\set{(\pi^{-1}(U), \widetilde{\varphi})}$ forme un atlas lisse sur $TM$.

\emph{Dimension.}
Chaque carte $\widetilde{\varphi}$ prend ses valeurs dans un ouvert de $\R^n \times \R^n = \R^{2n}$,
donc $\dim TM = 2n$.

\emph{$\pi$ est une submersion.}
En coordonnées locales, $\pi$ se lit $\varphi \circ \pi \circ \widetilde{\varphi}^{-1}(x, v) = x$,
qui est la projection sur les $n$ premières composantes, de rang $n$ partout.
\end{proof}

\begin{exemple}[Fibré tangent de $\R^n$]\label{ex:TRn}
On a $T\R^n \cong \R^n \times \R^n = \R^{2n}$. L'identification est donnée
par la carte globale $(\R^n, \Id)$.
\end{exemple}

\begin{exemple}[Fibré tangent de $S^1$]\label{ex:TS1}
\index{cercle!fibré tangent}
Le fibré tangent $TS^1$ est difféomorphe à $S^1 \times \R$, car $S^1$
est un groupe de Lie et tout groupe de Lie a un fibré tangent trivialisable.
\end{exemple}

% Diagramme commutatif
Le diagramme suivant résume la situation pour une application lisse $f : M \to N$ :
\[
\begin{tikzcd}
  TM \arrow[r, "\dd f"] \arrow[d, "\pi_M"'] & TN \arrow[d, "\pi_N"] \\
  M \arrow[r, "f"'] & N
\end{tikzcd}
\]
\index{diagramme commutatif!fibré tangent}

% ----------------------------------------------------------
%  Fibré cotangent
% ----------------------------------------------------------
\section{Le fibré cotangent}\label{sec:fibre-cotangent}

\begin{definition}[Espace cotangent]\label{def:cotangent}
L'\textbf{espace cotangent}\index{espace cotangent} en $p \in M$ est le dual de
l'espace tangent :
\[
  T_p^*M = (T_pM)^* = \Hom(T_pM, \R).
\]
Un élément de $T_p^*M$ est appelé \textbf{covecteur}\index{covecteur} en $p$.
\end{definition}

\begin{definition}[Différentielle d'une fonction]\label{def:diff-fonction}
Pour $f \in C^\infty(M)$, la \textbf{différentielle}\index{différentielle!d'une fonction}
de $f$ en $p$ est le covecteur
\[
  \dd f_p \in T_p^*M, \qquad \dd f_p(v) = v(f), \quad \forall\, v \in T_pM.
\]
\end{definition}

\begin{proposition}[Base duale]\label{prop:base-duale}
Si $(x^1, \ldots, x^n)$ sont des coordonnées locales autour de $p$, les covecteurs
$\dd x^1|_p, \ldots, \dd x^n|_p$ forment la base duale\index{base duale} de
$\left(\left.\frac{\partial}{\partial x^1}\right|_p, \ldots,
\left.\frac{\partial}{\partial x^n}\right|_p\right)$, c'est-à-dire
\[
  \dd x^i|_p \!\left(\left.\frac{\partial}{\partial x^j}\right|_p\right) = \delta^i_j.
\]
Tout covecteur $\omega \in T_p^*M$ s'écrit $\omega = \sum_{i=1}^n \omega_i\, \dd x^i|_p$
avec $\omega_i = \omega\!\left(\left.\frac{\partial}{\partial x^i}\right|_p\right)$.
\end{proposition}

\begin{definition}[Fibré cotangent]\label{def:fibre-cotangent}
Le \textbf{fibré cotangent}\index{fibré cotangent} de $M$ est
\[
  T^*M = \bigsqcup_{p \in M} T_p^*M,
\]
muni de la projection $\pi^* : T^*M \to M$. C'est une variété lisse de dimension $2n$.
\end{definition}

Le diagramme des projections s'écrit :
\[
\begin{tikzcd}
  T^*M \arrow[dr, "\pi^*"'] & & TM \arrow[dl, "\pi"] \\
  & M \arrow[ul, "\text{section}"] \arrow[ur, "\text{section}"'] &
\end{tikzcd}
\]

\begin{proposition}[Application cotangente]\label{prop:application-cotangente}
\index{application cotangente}
Soit $f : M \to N$ une application lisse. L'\textbf{application cotangente}
(ou \emph{pull-back}\index{pull-back}) est l'application
\[
  \dd f_p^* : T_{f(p)}^*N \longrightarrow T_p^*M, \qquad
  (\dd f_p^*\,\omega)(v) = \omega(\dd f_p(v)),
\]
pour $\omega \in T_{f(p)}^*N$ et $v \in T_pM$.
Contrairement à $\dd f_p$, l'application cotangente va en sens
\emph{inverse} : elle est \textbf{contravariante}\index{contravariante}.
\end{proposition}

\begin{proof}
La linéarité de $\dd f_p^*$ est claire. Pour vérifier la bonne définition,
il suffit de noter que $\dd f_p^*\,\omega$ est la composée de l'application
linéaire $\dd f_p : T_pM \to T_{f(p)}N$ et de la forme linéaire
$\omega : T_{f(p)}N \to \R$, ce qui donne bien un élément de $T_p^*M$.
\end{proof}

\begin{remarque}[Notation duale]\label{rem:notation-duale}
Le diagramme commutatif complet des fibrés tangent et cotangent est :
\[
\begin{tikzcd}
  T_p^*M & T_{f(p)}^*N \arrow[l, "\dd f_p^*"'] \\
  T_pM \arrow[r, "\dd f_p"'] & T_{f(p)}N
\end{tikzcd}
\]
L'application tangente va dans le même sens que $f$ (covariante) tandis que
l'application cotangente va en sens inverse (contravariante).
\end{remarque}

\begin{exemple}[Pull-back de la différentielle]\label{ex:pullback-diff}
Si $h \in C^\infty(N)$ et $f : M \to N$ est lisse, alors
\[
  \dd f_p^*(\dd h_{f(p)}) = \dd(h \circ f)_p.
\]
En effet, pour tout $v \in T_pM$ :
$\dd f_p^*(\dd h_{f(p)})(v) = \dd h_{f(p)}(\dd f_p(v)) = (\dd f_p(v))(h) = v(h \circ f) = \dd(h \circ f)_p(v)$.
\end{exemple}

% ----------------------------------------------------------
%  Champs de vecteurs
% ----------------------------------------------------------
\section{Champs de vecteurs}\label{sec:champs-vecteurs}

\begin{definition}[Champ de vecteurs]\label{def:champ-vecteurs}
Un \textbf{champ de vecteurs}\index{champ de vecteurs} sur $M$ est une section lisse
du fibré tangent, c'est-à-dire une application lisse $X : M \to TM$ telle que
$\pi \circ X = \Id_M$. L'ensemble des champs de vecteurs sur $M$ est noté
$\mathfrak{X}(M)$\index{X(M)@$\mathfrak{X}(M)$}.
\end{definition}

\begin{remarque}\label{rem:champ-derivation}
Un champ de vecteurs $X \in \mathfrak{X}(M)$ agit comme une dérivation sur
$C^\infty(M)$ :
\[
  X : C^\infty(M) \longrightarrow C^\infty(M), \qquad (Xf)(p) = X_p(f).
\]
En coordonnées locales $(x^1, \ldots, x^n)$, on écrit
$X = \sum_{i=1}^n X^i \frac{\partial}{\partial x^i}$ où $X^i \in C^\infty(U)$.
\end{remarque}

\begin{exemple}[Champs de vecteurs sur $S^1$]\label{ex:champ-S1}
\index{cercle!champ de vecteurs}
Le cercle $S^1 = \set{(\cos\theta, \sin\theta) : \theta \in \R}$ admet le champ de
vecteurs $X = \frac{\partial}{\partial \theta}$, qui ne s'annule jamais.
C'est une conséquence de la trivialité de $TS^1$.
\end{exemple}

\begin{figure}[ht]
\centering
\begin{tikzpicture}[scale=1.6]
  % Cercle
  \draw[thick] (0,0) circle (1);
  % Champ de vecteurs
  \foreach \a in {0,30,...,330} {
    \pgfmathsetmacro{\ca}{cos(\a)}
    \pgfmathsetmacro{\sa}{sin(\a)}
    \draw[->, blue!70!black, thick]
      (\ca, \sa) -- +({-0.35*\sa}, {0.35*\ca});
  }
  \node at (0, -1.5) {Champ $\frac{\partial}{\partial\theta}$ sur $S^1$};
\end{tikzpicture}
\caption{Le champ de vecteurs tangent unitaire sur le cercle $S^1$.}\label{fig:champ-S1}
\end{figure}

% ----------------------------------------------------------
%  Crochet de Lie
% ----------------------------------------------------------
\section{Le crochet de Lie}\label{sec:crochet-lie}

\begin{definition}[Crochet de Lie]\label{def:crochet-lie}
Le \textbf{crochet de Lie}\index{crochet de Lie} de deux champs de vecteurs
$X, Y \in \mathfrak{X}(M)$ est le champ de vecteurs $[X, Y]$ défini par
\[
  [X, Y](f) = X(Y(f)) - Y(X(f)), \qquad \forall\, f \in C^\infty(M).
\]
\end{definition}

\begin{proposition}[Le crochet de Lie est bien un champ de vecteurs]\label{prop:crochet-derivation}
Pour tous $X, Y \in \mathfrak{X}(M)$, l'opérateur $[X, Y]$ est une dérivation
de $C^\infty(M)$, et donc définit un champ de vecteurs lisse sur $M$.
\end{proposition}

\begin{proof}
La linéarité est claire. Il suffit de vérifier la règle de Leibniz :
\begin{align*}
  [X,Y](fg) &= X(Y(fg)) - Y(X(fg)) \\
  &= X(fY(g) + gY(f)) - Y(fX(g) + gX(f)) \\
  &= X(f)Y(g) + fX(Y(g)) + X(g)Y(f) + gX(Y(f)) \\
  &\quad\; - Y(f)X(g) - fY(X(g)) - Y(g)X(f) - gY(X(f)) \\
  &= f[X,Y](g) + g[X,Y](f). \qedhere
\end{align*}
\end{proof}

\begin{proposition}[Expression en coordonnées locales]\label{prop:crochet-coordonnees}
\index{crochet de Lie!en coordonnées}
Soient $X = \sum_i X^i \frac{\partial}{\partial x^i}$ et
$Y = \sum_j Y^j \frac{\partial}{\partial x^j}$. Alors
\[
  [X, Y] = \sum_{k=1}^n \left(\sum_{i=1}^n X^i \frac{\partial Y^k}{\partial x^i}
  - Y^i \frac{\partial X^k}{\partial x^i}\right) \frac{\partial}{\partial x^k}.
\]
\end{proposition}

\begin{proof}
Appliquons $[X,Y]$ à une fonction coordonnée $x^k$. On a :
\begin{align*}
  [X,Y](x^k) &= X(Y^k) - Y(X^k)
  = \sum_i X^i \frac{\partial Y^k}{\partial x^i}
  - \sum_i Y^i \frac{\partial X^k}{\partial x^i}. \qedhere
\end{align*}
\end{proof}

\begin{theoreme}[Propriétés du crochet de Lie]\label{thm:proprietes-crochet}
\index{crochet de Lie!propriétés}
L'espace $(\mathfrak{X}(M), [\cdot, \cdot])$ est une algèbre de Lie\index{algèbre de Lie}
de dimension infinie. Plus précisément, pour tous $X, Y, Z \in \mathfrak{X}(M)$
et $f, g \in C^\infty(M)$ :
\begin{enumerate}[label=(\roman*)]
  \item \textbf{Bilinéarité sur $\R$ :} $[aX + bY, Z] = a[X,Z] + b[Y,Z]$ pour $a, b \in \R$.
  \item \textbf{Antisymétrie :} $[X, Y] = -[Y, X]$.
  \item \textbf{Identité de Jacobi\index{identité de Jacobi} :}
    $[X, [Y, Z]] + [Y, [Z, X]] + [Z, [X, Y]] = 0$.
  \item \textbf{$C^\infty(M)$-comportement :} $[fX, gY] = fg[X,Y] + f(Xg)Y - g(Yf)X$.
\end{enumerate}
\end{theoreme}

\begin{proof}
Les points (i) et (ii) résultent directement de la définition. Pour (iii), il s'agit
d'un calcul direct :
\begin{align*}
  [X,[Y,Z]](f) &= X(Y(Z(f)) - Z(Y(f))) - (Y(Z(X(f))) - Z(Y(X(f)))) \\
                &= X(Y(Z(f))) - X(Z(Y(f))) - Y(Z(X(f))) + Z(Y(X(f))).
\end{align*}
En sommant cycliquement sur $(X,Y,Z)$, chaque terme apparaît une fois avec un signe
positif et une fois avec un signe négatif, d'où l'annulation.

Pour (iv), calculons $[fX, gY](h)$ :
\begin{align*}
  [fX, gY](h) &= fX(gY(h)) - gY(fX(h)) \\
  &= fg\,X(Y(h)) + f(Xg)\,Y(h) - gf\,Y(X(h)) - g(Yf)\,X(h) \\
  &= fg[X,Y](h) + f(Xg)Y(h) - g(Yf)X(h). \qedhere
\end{align*}
\end{proof}

% ----------------------------------------------------------
%  Flot d'un champ de vecteurs
% ----------------------------------------------------------
\section{Flot d'un champ de vecteurs}\label{sec:flot}

\begin{definition}[Courbe intégrale]\label{def:courbe-integrale}
Soit $X \in \mathfrak{X}(M)$. Une \textbf{courbe intégrale}\index{courbe intégrale}
de $X$ issue de $p \in M$ est une courbe lisse $\gamma : I \to M$, où $I$ est un
intervalle ouvert contenant $0$, telle que
\[
  \gamma(0) = p \quad \text{et} \quad \gamma'(t) = X_{\gamma(t)}, \quad \forall\, t \in I.
\]
\end{definition}

\begin{theoreme}[Existence et unicité des courbes intégrales]\label{thm:cauchy-lipschitz}
\index{théorème!de Cauchy--Lipschitz}\index{Cauchy--Lipschitz}
Soit $X \in \mathfrak{X}(M)$ et $p \in M$. Il existe un intervalle ouvert maximal
$I_p \ni 0$ et une unique courbe intégrale maximale $\gamma_p : I_p \to M$ de $X$
issue de $p$.
\end{theoreme}

\begin{proof}
\emph{Existence et unicité locales.}
En coordonnées locales $(x^1, \ldots, x^n)$ autour de $p$, l'équation $\gamma'(t) = X_{\gamma(t)}$
se traduit par le système d'équations différentielles ordinaires
\[
  \frac{\dd x^i}{\dd t}(t) = X^i(x^1(t), \ldots, x^n(t)), \qquad i = 1, \ldots, n,
\]
avec condition initiale $(x^1(0), \ldots, x^n(0)) = \varphi(p)$.
Comme les $X^i$ sont lisses (donc localement lipschitziennes), le théorème de
Cauchy--Lipschitz classique donne l'existence et l'unicité d'une solution sur un
intervalle $(-\varepsilon, \varepsilon)$.

\emph{Maximalité.}
Si $\gamma_1 : I_1 \to M$ et $\gamma_2 : I_2 \to M$ sont deux courbes intégrales
issues de $p$, alors par unicité locale elles coïncident sur $I_1 \cap I_2$.
On définit $I_p = \bigcup \set{I : \text{il existe une courbe intégrale sur } I}$
et $\gamma_p$ en recollant les solutions.
\end{proof}

\begin{definition}[Flot]\label{def:flot}
Le \textbf{flot}\index{flot} du champ $X$ est l'application
\[
  \theta : \mathcal{D} \longrightarrow M, \qquad \theta(t, p) = \theta_t(p) = \gamma_p(t),
\]
où $\mathcal{D} = \set{(t, p) \in \R \times M : t \in I_p}$ est le \emph{domaine du flot}.
\end{definition}

\begin{proposition}[Propriétés du flot]\label{prop:proprietes-flot}
\index{flot!propriétés}
Soit $\theta$ le flot de $X \in \mathfrak{X}(M)$.
\begin{enumerate}[label=(\roman*)]
  \item $\theta_0 = \Id_M$.
  \item $\theta_{s+t} = \theta_s \circ \theta_t$ là où les deux membres sont définis
    (\emph{propriété de groupe à un paramètre}).
  \item Si $X$ est \emph{complet}\index{champ complet} (i.e., $I_p = \R$ pour tout $p$),
    alors $t \mapsto \theta_t$ est un homomorphisme de groupes $(\R, +) \to \Diff(M)$.
\end{enumerate}
\end{proposition}

\begin{proof}
Le point (i) est immédiat par définition. Pour (ii), fixons $p \in M$ et considérons
les courbes $t \mapsto \theta_{s+t}(p)$ et $t \mapsto \theta_s(\theta_t(p))$. Les deux sont
des courbes intégrales de $X$ passant par $\theta_s(p)$ en $t = 0$. Par unicité, elles
coïncident. Le point (iii) en découle directement.
\end{proof}

\begin{theoreme}[Compacité implique complétude]\label{thm:compact-complet}
\index{champ complet}
Si $M$ est compacte, alors tout champ de vecteurs sur $M$ est complet.
\end{theoreme}

\begin{proof}[Esquisse de preuve]
Par compacité, on peut recouvrir $M$ par un nombre fini d'ouverts de cartes.
Sur chacun, le théorème de Cauchy--Lipschitz fournit un temps d'existence
$\varepsilon_i > 0$ uniforme. On pose $\varepsilon = \min_i \varepsilon_i > 0$.
Par la propriété de groupe, la solution se prolonge de $\varepsilon$ en $\varepsilon$
sur tout $\R$.
\end{proof}

\begin{exemple}[Flot d'un champ de rotation]\label{ex:flot-rotation}
\index{flot!rotation}
Sur $\R^2$, considérons le champ de vecteurs $X = -y\,\frac{\partial}{\partial x}
+ x\,\frac{\partial}{\partial y}$. Le système d'équations différentielles associé est
\[
  \dot{x} = -y, \qquad \dot{y} = x,
\]
dont la solution avec condition initiale $(x_0, y_0)$ est
\[
  \theta_t(x_0, y_0) = (x_0\cos t - y_0\sin t,\; x_0\sin t + y_0\cos t).
\]
C'est la rotation d'angle $t$ autour de l'origine. Le champ $X$ est complet
et le flot est un groupe à un paramètre de rotations :
$\theta_t = R_t \in \SO(2)$ pour tout $t \in \R$.
Les courbes intégrales sont les cercles centrés à l'origine.
\end{exemple}

\begin{proposition}[Lien entre différentielle et flot]\label{prop:diff-flot}
\index{différentielle!et flot}
Soient $X \in \mathfrak{X}(M)$ de flot $\theta_t$ et $f : M \to N$ une application
lisse. Pour tout $p \in M$ :
\[
  \dd f_p(X_p) = \left.\frac{\dd}{\dd t}\right|_{t=0} f(\theta_t(p)).
\]
\end{proposition}

\begin{proof}
Soit $\gamma(t) = \theta_t(p)$ la courbe intégrale de $X$ issue de $p$.
Alors $\gamma'(0) = X_p$ et, par définition de la différentielle via les courbes :
\[
  \dd f_p(X_p) = \dd f_p(\gamma'(0)) = (f \circ \gamma)'(0)
  = \left.\frac{\dd}{\dd t}\right|_{t=0} f(\theta_t(p)). \qedhere
\]
\end{proof}

\begin{proposition}[Commutativité des flots et crochet nul]\label{prop:flots-commutent}
\index{flot!commutativité}
Soient $X, Y \in \mathfrak{X}(M)$ des champs complets de flots $\theta_t$ et $\psi_s$.
Alors $[X, Y] = 0$ si et seulement si $\theta_t \circ \psi_s = \psi_s \circ \theta_t$
pour tous $s, t \in \R$.
\end{proposition}

\begin{proof}[Esquisse de preuve]
Le sens direct utilise la dérivée de Lie : $[X, Y] = 0$ signifie que $Y$ est invariant
par le flot de $X$, c'est-à-dire $(\theta_t)_* Y = Y$ pour tout $t$. Ceci implique que
le flot de $Y$ commute avec $\theta_t$. La réciproque s'obtient en dérivant l'égalité
$\theta_t \circ \psi_s = \psi_s \circ \theta_t$ par rapport à $t$ en $t = 0$.
\end{proof}

\begin{remarque}[Crochet de Lie et flot]\label{rem:crochet-flot}
Le crochet de Lie possède une interprétation dynamique fondamentale : si $\theta_t$
et $\psi_s$ désignent les flots de $X$ et $Y$ respectivement, alors
\[
  [X, Y]_p = \lim_{t \to 0} \frac{(\theta_{-t})_* Y_{\theta_t(p)} - Y_p}{t}
  = \left.\frac{\dd}{\dd t}\right|_{t=0} (\theta_t^* Y)_p,
\]
c'est-à-dire que $[X,Y]$ est la \emph{dérivée de Lie}\index{dérivée de Lie} de $Y$
le long du flot de $X$.
\end{remarque}

\begin{figure}[ht]
\centering
\begin{tikzpicture}[scale=1.0]
  % Plan tangent à S²
  % Sphère
  \shade[ball color=blue!15, opacity=0.5] (0,0) circle (2);
  \draw[thick] (0,0) circle (2);
  % Equateur
  \draw[thick, dashed] (2,0) arc (0:180:2 and 0.6);
  \draw[thick] (-2,0) arc (180:360:2 and 0.6);
  % Point
  \fill (0.3, 1.95) circle (1.5pt) node[above right] {$p$};
  % Plan tangent (parallelogramme)
  \draw[blue!40, fill=blue!10, opacity=0.5]
    (-1.5, 2.8) -- (2.1, 2.8) -- (2.1, 1.1) -- (-1.5, 1.1) -- cycle;
  \node[blue!70!black] at (2.4, 2.5) {$T_pS^2$};
  % Vecteurs tangents
  \draw[thick, ->, red!70!black] (0.3, 1.95) -- (1.3, 2.45);
  \draw[thick, ->, green!50!black] (0.3, 1.95) -- (-0.5, 1.55);
  \node[red!70!black] at (1.5, 2.6) {$v_1$};
  \node[green!50!black] at (-0.7, 1.35) {$v_2$};
\end{tikzpicture}
\caption{Le plan tangent $T_pS^2$ à la sphère en un point $p$.}\label{fig:plan-tangent-S2}
\end{figure}

% ----------------------------------------------------------
%  Exercices du chapitre 3
% ----------------------------------------------------------
\section{Exercices}\label{sec:exo-ch3}

\begin{exercice}\label{exo:tangent-Rn}
Montrer que l'espace tangent $T_p\R^n$ est canoniquement isomorphe à $\R^n$
via l'identification $v \mapsto (v(x^1), \ldots, v(x^n))$ où $(x^1, \ldots, x^n)$
sont les coordonnées standard.
\end{exercice}

\begin{exercice}\label{exo:tangent-S1}
Soit $S^1 = \set{(x,y) \in \R^2 : x^2 + y^2 = 1}$.
En utilisant la paramétrisation $\gamma(\theta) = (\cos\theta, \sin\theta)$,
montrer que $T_{(1,0)}S^1$ est engendré par le vecteur $(0,1)$.
\end{exercice}

\begin{exercice}\label{exo:diff-lineaire}
Soit $L : \R^m \to \R^n$ une application linéaire. Montrer que pour tout
$p \in \R^m$, sous l'identification $T_p\R^m \cong \R^m$, on a $\dd L_p = L$.
\end{exercice}

\begin{exercice}\label{exo:diff-inclusion}
Soit $U \subset M$ un ouvert et $\iota : U \hookrightarrow M$ l'inclusion.
Montrer que $\dd\iota_p : T_pU \to T_pM$ est un isomorphisme pour tout $p \in U$.
\end{exercice}

\begin{exercice}\label{exo:crochet-R3}
\index{crochet de Lie!calcul}
Soient $X = y\,\frac{\partial}{\partial x} - x\,\frac{\partial}{\partial y}$
et $Y = \frac{\partial}{\partial z}$ sur $\R^3$. Calculer $[X, Y]$ et interpréter
géométriquement.
\end{exercice}

\begin{exercice}\label{exo:flot-lineaire}
\index{flot!système linéaire}
Soit $A \in M_n(\R)$ et $X_A$ le champ de vecteurs sur $\R^n$ défini par
$X_A(x) = Ax$. Montrer que le flot de $X_A$ est $\theta_t(x) = e^{tA} x$.
\end{exercice}

\begin{exercice}\label{exo:crochet-f}
Soient $X, Y \in \mathfrak{X}(M)$ et $f : M \to N$ un difféomorphisme.
Montrer que $f_*[X,Y] = [f_*X, f_*Y]$, où $f_*X$ est le champ image directe
défini par $(f_*X)_{f(p)} = \dd f_p(X_p)$.
\end{exercice}

\begin{exercice}\label{exo:euler}
\index{champ d'Euler}
Le \emph{champ d'Euler} sur $\R^n$ est $E = \sum_{i=1}^n x^i \frac{\partial}{\partial x^i}$.
\begin{enumerate}[label=(\alph*)]
  \item Calculer le flot de $E$.
  \item Montrer que $[E, X] = -X$ pour tout champ de vecteurs $X$ homogène de degré $0$.
  \item Montrer que $Ef = kf$ si et seulement si $f$ est homogène de degré $k$.
\end{enumerate}
\end{exercice}

\begin{exercice}\label{exo:cotangent-changement}
\index{fibré cotangent!changement de cartes}
Soient $(U, \varphi = (x^i))$ et $(V, \psi = (y^j))$ deux cartes sur $M$ avec
$U \cap V \neq \emptyset$. Montrer que la loi de transformation d'un covecteur
$\omega = \sum_i \omega_i\,\dd x^i$ dans la carte $(V, \psi)$ est
\[
  \omega = \sum_j \left(\sum_i \omega_i \frac{\partial x^i}{\partial y^j}\right) \dd y^j.
\]
Comparer avec la loi de transformation des composantes d'un vecteur tangent.
\end{exercice}

\begin{exercice}\label{exo:lie-matrice}
\index{algèbre de Lie!des matrices}
Soit $G = \GL(n, \R)$\index{GL(n,R)@$\GL(n,\R)$}. Identifier $T_I G$ avec $M_n(\R)$
et montrer que le crochet de Lie de deux champs de vecteurs invariants à gauche correspond
au commutateur $[A, B] = AB - BA$.
\end{exercice}


% ============================================================
%  CHAPITRE 4 : SOUS-VARIÉTÉS ET THÉORÈME DE LA VALEUR RÉGULIÈRE
% ============================================================
\chapter{Sous-variétés et théorème de la valeur régulière}\label{ch:sous-varietes}
\minitoc\bigskip

\section{Immersions, submersions, plongements}\label{sec:immersions}

\begin{definition}[Immersion, submersion]\label{def:immersion-submersion}
Soit $f : M \to N$ une application lisse entre variétés.
\begin{enumerate}[label=(\roman*)]
  \item $f$ est une \textbf{immersion}\index{immersion} si $\dd f_p$ est injective
    pour tout $p \in M$.
  \item $f$ est une \textbf{submersion}\index{submersion} si $\dd f_p$ est surjective
    pour tout $p \in M$.
  \item Un point $p \in M$ est un \textbf{point régulier}\index{point régulier}
    de $f$ si $\dd f_p$ est surjective. Sinon, $p$ est un \textbf{point critique}\index{point critique}.
  \item Une valeur $c \in N$ est une \textbf{valeur régulière}\index{valeur régulière}
    de $f$ si tout point de $f^{-1}(c)$ est régulier.
    Une valeur qui n'est pas régulière est une \textbf{valeur critique}\index{valeur critique}.
\end{enumerate}
\end{definition}

\begin{remarque}\label{rem:valeur-reguliere-vide}
Par convention, si $f^{-1}(c) = \emptyset$, alors $c$ est une valeur régulière
(la condition est vide).
\end{remarque}

\begin{definition}[Plongement]\label{def:plongement}
Une application lisse $f : M \to N$ est un \textbf{plongement}\index{plongement}
(ou \emph{plongement lisse}) si elle est une immersion injective qui est un
homéomorphisme sur son image $f(M) \subset N$ (muni de la topologie induite).
\end{definition}

\begin{proposition}[Immersion injective propre]\label{prop:immersion-propre}
\index{immersion!injective propre}
Toute immersion injective propre\index{application propre} est un plongement.
En particulier, si $M$ est compacte, toute immersion injective $f : M \to N$
est un plongement.
\end{proposition}

\begin{proof}
Soit $f : M \to N$ une immersion injective propre. Il suffit de montrer que
$f$ est un homéomorphisme sur son image. Comme $f$ est propre et $N$ est localement
compact, $f$ est une application fermée. Une bijection continue et fermée
est un homéomorphisme. Pour la seconde assertion, si $M$ est compacte, toute
application continue $M \to N$ est propre (image réciproque d'un compact est un
fermé de $M$ compact, donc compact).
\end{proof}

\begin{theoreme}[Forme locale d'une immersion]\label{thm:forme-locale-immersion}
\index{immersion!forme locale}
Soit $f : M^m \to N^n$ une immersion ($m \leqslant n$). Pour tout $p \in M$,
il existe des cartes $(U, \varphi)$ autour de $p$ et $(V, \psi)$ autour de $f(p)$
telles que
\[
  \psi \circ f \circ \varphi^{-1}(x^1, \ldots, x^m) = (x^1, \ldots, x^m, 0, \ldots, 0).
\]
\end{theoreme}

\begin{proof}
Comme $\dd f_p$ est injective, la matrice jacobienne de $\hat{f} = \psi \circ f \circ \varphi^{-1}$
en $\varphi(p)$ est de rang $m$. Quitte à permuter les coordonnées de $N$, on peut
supposer que les $m$ premières lignes forment une sous-matrice inversible.
Définissons $F : \varphi(U) \times \R^{n-m} \to \R^n$ par
\[
  F(x, y) = \hat{f}(x) + (0, \ldots, 0, y^1, \ldots, y^{n-m}).
\]
Alors $DF_{(\varphi(p), 0)}$ est inversible, donc $F$ est un difféomorphisme local.
En posant $\psi' = F^{-1} \circ \psi$ au voisinage de $f(p)$, on obtient
$\psi' \circ f \circ \varphi^{-1}(x) = (x, 0)$.
\end{proof}

\begin{theoreme}[Forme locale d'une submersion]\label{thm:forme-locale-submersion}
\index{submersion!forme locale}
Soit $f : M^m \to N^n$ une submersion ($m \geqslant n$). Pour tout $p \in M$,
il existe des cartes $(U, \varphi)$ autour de $p$ et $(V, \psi)$ autour de $f(p)$
telles que
\[
  \psi \circ f \circ \varphi^{-1}(x^1, \ldots, x^m) = (x^1, \ldots, x^n).
\]
\end{theoreme}

\begin{proof}
La surjectivité de $\dd f_p$ implique que les $n$ premières colonnes de la jacobienne
$D\hat{f}_{\varphi(p)}$ (quitte à permuter) forment une matrice inversible.
Définissons $G : \varphi(U) \to \R^m$ par
\[
  G(x) = (\hat{f}(x),\, x^{n+1}, \ldots, x^m).
\]
Alors $DG_{\varphi(p)}$ est inversible et $\hat{f} \circ G^{-1}(y^1, \ldots, y^m) = (y^1, \ldots, y^n)$.
En posant $\varphi' = G \circ \varphi$, on obtient la forme canonique souhaitée.
\end{proof}

\begin{exemple}[Immersion et submersion]\label{ex:immersion-submersion}
\begin{enumerate}[label=(\roman*)]
  \item L'application $f : \R \to \R^2$, $f(t) = (t^2, t^3)$, n'est \emph{pas}
    une immersion car $\dd f_0 = 0$. En revanche, $f|_{\R \setminus \set{0}}$ est une immersion.
  \item La projection $\pi : \R^n \to \R^k$, $\pi(x^1, \ldots, x^n) = (x^1, \ldots, x^k)$
    avec $k \leqslant n$, est une submersion (la jacobienne est $(I_k \mid 0)$, de rang $k$).
  \item Le déterminant $\det : \GL(n, \R) \to \R^*$ est une submersion.
\end{enumerate}
\end{exemple}

% ----------------------------------------------------------
%  Sous-variétés
% ----------------------------------------------------------
\section{Sous-variétés}\label{sec:sous-varietes}

\begin{definition}[Sous-variété plongée]\label{def:sous-variete-plongee}
\index{sous-variété!plongée}
Un sous-ensemble $S \subset M$ est une \textbf{sous-variété plongée} de dimension $k$
si, pour tout $p \in S$, il existe une carte $(U, \varphi)$ de $M$ autour de $p$ telle que
\[
  \varphi(U \cap S) = \varphi(U) \cap (\R^k \times \set{0}),
\]
où l'on identifie $\R^k \times \set{0} \subset \R^n$. Une telle carte est appelée
\textbf{carte de sous-variété}\index{carte de sous-variété} (ou \emph{carte adaptée}).
La \textbf{codimension}\index{codimension} de $S$ dans $M$ est $n - k$.
\end{definition}

\begin{definition}[Sous-variété immergée]\label{def:sous-variete-immergee}
\index{sous-variété!immergée}
Un sous-ensemble $S \subset M$ est une \textbf{sous-variété immergée} de dimension $k$
si $S$ admet une structure de variété lisse de dimension $k$ pour laquelle l'inclusion
$\iota : S \hookrightarrow M$ est une immersion lisse.
\end{definition}

\begin{remarque}\label{rem:plongee-vs-immergee}
Toute sous-variété plongée est immergée, mais la réciproque est fausse.
La topologie d'une sous-variété immergée peut différer de la topologie induite.
\end{remarque}

\begin{exemple}[Sous-variété plongée : $S^n \subset \R^{n+1}$]\label{ex:sphere-plongee}
\index{sphère!comme sous-variété}
La sphère $S^n = \set{x \in \R^{n+1} : \abs{x}^2 = 1}$ est une sous-variété
plongée de $\R^{n+1}$ de codimension $1$. Nous le démontrerons rigoureusement
au 
ef{thm:valeur-reguliere} comme préimage d'une valeur régulière.
\end{exemple}

\begin{exemple}[Sous-variété immergée non plongée]\label{ex:huit-immerse}
\index{sous-variété!immergée non plongée}
L'application $\gamma : \R \to \R^2$ définie par
$\gamma(t) = (\sin(2t), \sin(t))$ est une immersion dont l'image est une
figure en huit. L'image, munie de la topologie de $\R$ via $\gamma$, est une
sous-variété immergée de $\R^2$ qui n'est pas plongée (au point double $(0,0)$).
\end{exemple}

\begin{exemple}[Droite irrationnelle sur le tore]\label{ex:droite-irrationnelle}
\index{tore!droite irrationnelle}
Soit $\alpha \in \R \setminus \Q$ et considérons l'application
$\gamma : \R \to T^2 = \R^2/\Z^2$ définie par $\gamma(t) = [(t, \alpha t)]$.
C'est une immersion injective dont l'image est dense dans $T^2$.
L'image est une sous-variété immergée de dimension $1$ qui n'est pas plongée.
\end{exemple}

% ----------------------------------------------------------
%  Théorème de la valeur régulière
% ----------------------------------------------------------
\section{Le théorème de la valeur régulière}\label{sec:valeur-reguliere}

Le théorème de la valeur régulière est un des résultats fondamentaux de la
topologie différentielle. Il fournit un procédé systématique pour construire
des sous-variétés.

\begin{theoreme}[Théorème de la valeur régulière]\label{thm:valeur-reguliere}
\index{théorème!de la valeur régulière}
Soit $f : M^m \to N^n$ une application lisse entre variétés lisses, avec $m \geqslant n$.
Si $c \in N$ est une \textbf{valeur régulière} de $f$, alors :
\begin{enumerate}[label=(\roman*)]
  \item $S = f^{-1}(c)$ est une sous-variété plongée de $M$ de
    codimension $n$, c'est-à-dire $\dim S = m - n$.
  \item Pour tout $p \in S$, l'espace tangent à $S$ en $p$ est
    \[
      T_pS = \ker \dd f_p = \set{v \in T_pM : \dd f_p(v) = 0}.
    \]
\end{enumerate}
\end{theoreme}

\begin{proof}
Soit $p \in S = f^{-1}(c)$. Par hypothèse, $\dd f_p : T_pM \to T_cN$ est surjective,
donc de rang $n$. On procède en plusieurs étapes.

\emph{Étape 1 : Réduction en coordonnées.}
Choisissons des cartes $(U, \varphi)$ autour de $p$ et $(V, \psi)$ autour de $c$
avec $f(U) \subset V$, $\varphi(p) = 0$ et $\psi(c) = 0$.
L'application $\hat{f} = \psi \circ f \circ \varphi^{-1} : \varphi(U) \to \psi(V)$
est lisse et $D\hat{f}_0$ est surjective.

\emph{Étape 2 : Application du théorème des fonctions implicites.}
Quitte à réordonner les coordonnées, on peut supposer que les $n$ dernières
colonnes de $D\hat{f}_0$ forment une matrice $n \times n$ inversible.
Écrivons $\varphi(U) \ni x = (x', x'')$ avec $x' \in \R^{m-n}$ et $x'' \in \R^n$.
Le théorème des fonctions implicites\index{théorème!des fonctions implicites}
donne l'existence d'un voisinage ouvert $W$ de $0$ dans $\R^{m-n}$ et d'une
application lisse $g : W \to \R^n$ telle que $g(0) = 0$ et
\[
  \hat{f}(x', x'') = 0 \quad \Longleftrightarrow \quad x'' = g(x').
\]

\emph{Étape 3 : Carte de sous-variété.}
Définissons $\Phi : W \times \R^n \supset U' \to \R^m$ par
$\Phi(x', y) = (x', g(x') + y)$. Alors $D\Phi_0 = \begin{psmallmatrix} I_{m-n} & 0 \\ Dg_0 & I_n \end{psmallmatrix}$
est inversible, donc $\Phi$ est un difféomorphisme local. L'application
$\widetilde{\varphi} = \Phi^{-1} \circ \varphi$ est une carte autour de $p$ telle que
\[
  \widetilde{\varphi}(U \cap S) = \widetilde{\varphi}(U) \cap (\R^{m-n} \times \set{0}).
\]
Ceci montre que $S$ est une sous-variété plongée de codimension $n$.

\emph{Étape 4 : Espace tangent.}
Soit $v \in T_pS$. On peut écrire $v = [\gamma]_p$ pour une courbe $\gamma : (-\varepsilon, \varepsilon) \to S$
avec $\gamma(0) = p$. Comme $f \circ \gamma(t) = c$ pour tout $t$, on a
$\dd f_p(v) = (f \circ \gamma)'(0) = 0$, d'où $T_pS \subset \ker \dd f_p$.
Réciproquement, $\dim T_pS = m - n = m - \rg(\dd f_p) = \dim \ker \dd f_p$
par le théorème du rang, d'où l'égalité.
\end{proof}

\begin{figure}[ht]
\centering
\begin{tikzpicture}[scale=1.0]
  % Domaine M
  \draw[thick, fill=orange!8] (-3, -2) rectangle (3, 3);
  \node at (2.5, 2.6) {$M$};

  % Lignes de niveau
  \draw[blue!30, thick] (-2.5, -1.5) .. controls (-1, -0.5) and (1, 0.5) .. (2.5, -0.2);
  \draw[blue!30, thick] (-2.5, -0.5) .. controls (-1, 0.5) and (1, 1.5) .. (2.5, 0.8);
  \draw[red!70!black, very thick] (-2.5, 0.5) .. controls (-1, 1.5) and (1, 2.5) .. (2.5, 1.8);
  \draw[blue!30, thick] (-2.5, 1.5) .. controls (-1, 2.5) and (1, 2.8) .. (2.2, 2.8);
  \node[red!70!black] at (-2.8, 0.5) {$f^{-1}(c)$};

  % Flèche f
  \draw[thick, ->] (3.5, 0.5) -- (5.0, 0.5) node[midway, above] {$f$};

  % Codomain N
  \draw[thick] (5.5, -1.5) -- (5.5, 3);
  \node at (5.8, 2.6) {$N$};
  % Points
  \fill[blue!30] (5.5, -0.8) circle (2pt);
  \fill[blue!30] (5.5, -0.1) circle (2pt);
  \fill[red!70!black] (5.5, 0.7) circle (2.5pt) node[right] {$c$};
  \fill[blue!30] (5.5, 1.5) circle (2pt);
\end{tikzpicture}
\caption{Le théorème de la valeur régulière : la préimage $f^{-1}(c)$
d'une valeur régulière $c$ est une sous-variété.}\label{fig:valeur-reguliere}
\end{figure}

% ----------------------------------------------------------
%  Applications du théorème
% ----------------------------------------------------------
\section{Applications fondamentales}\label{sec:applications-VR}

\subsection{La sphère $S^{n-1}$}\label{subsec:sphere}

\begin{exemple}[$S^{n-1}$ comme préimage d'une valeur régulière]\label{ex:sphere-preimage}
\index{sphère!valeur régulière}
Considérons l'application $f : \R^n \to \R$ définie par $f(x) = \abs{x}^2 = \sum_{i=1}^n (x^i)^2$.
La différentielle en $x$ est
\[
  \dd f_x(v) = 2\sum_{i=1}^n x^i v^i = 2\langle x, v\rangle.
\]
Pour $x \neq 0$, l'application $\dd f_x$ est surjective (il suffit de prendre $v = x$).
En particulier, $1$ est une valeur régulière de $f$, et
\[
  S^{n-1} = f^{-1}(1)
\]
est une sous-variété plongée de $\R^n$ de dimension $n-1$. L'espace tangent est
\[
  T_xS^{n-1} = \ker \dd f_x = \set{v \in \R^n : \langle x, v\rangle = 0} = x^\perp.
\]
\end{exemple}

\subsection{Le groupe orthogonal $O(n)$}\label{subsec:On}

\begin{exemple}[$O(n)$ comme sous-variété de $\GL(n, \R)$]\label{ex:On-preimage}
\index{O(n)@$O(n)$!comme sous-variété}\index{groupe orthogonal}
Considérons l'espace $M_n(\R)$ des matrices $n \times n$ (identifié à $\R^{n^2}$)
et le sous-espace $\operatorname{Sym}_n(\R)$\index{matrices symétriques} des matrices
symétriques (de dimension $n(n+1)/2$). Définissons
\[
  f : M_n(\R) \longrightarrow \operatorname{Sym}_n(\R), \qquad f(A) = A^\top A.
\]
Calculons la différentielle. Pour $A \in M_n(\R)$ et $B \in T_A M_n(\R) \cong M_n(\R)$ :
\[
  \dd f_A(B) = \left.\frac{\dd}{\dd t}\right|_{t=0} (A + tB)^\top(A + tB)
  = B^\top A + A^\top B.
\]
Montrons que $\dd f_A$ est surjective pour $A \in O(n) = f^{-1}(I_n)$.
Soit $S \in \operatorname{Sym}_n(\R)$. Posons $B = \frac{1}{2}AS$. Alors
\[
  \dd f_A(B) = \frac{1}{2}(S A^\top A + A^\top A S) = \frac{1}{2}(S + S) = S,
\]
car $A^\top A = I_n$. Donc $I_n$ est une valeur régulière et
\[
  O(n) = f^{-1}(I_n)
\]
est une sous-variété plongée de $M_n(\R)$ de dimension
\[
  \dim O(n) = n^2 - \frac{n(n+1)}{2} = \frac{n(n-1)}{2}.
\]
L'espace tangent en $I_n$ est
\[
  T_{I_n}O(n) = \ker \dd f_{I_n} = \set{B \in M_n(\R) : B^\top + B = 0}
  = \mathfrak{o}(n),
\]
l'algèbre de Lie\index{algèbre de Lie!o(n)@$\mathfrak{o}(n)$} des matrices antisymétriques.
Plus généralement, pour $A \in O(n)$ :
\[
  T_A O(n) = \set{B \in M_n(\R) : B^\top A + A^\top B = 0} = A \cdot \mathfrak{o}(n).
\]
\end{exemple}

\subsection{Le groupe spécial linéaire $\SL(n, \R)$}\label{subsec:SLn}

\begin{exemple}[$\SL(n, \R)$ comme sous-variété de $\GL(n, \R)$]\label{ex:SLn-preimage}
\index{SL(n,R)@$\SL(n,\R)$!comme sous-variété}\index{groupe spécial linéaire}
Considérons le déterminant $\det : M_n(\R) \to \R$. Sa différentielle en
$A \in \GL(n, \R)$ est donnée par
\[
  \dd(\det)_A(B) = \det(A) \cdot \operatorname{tr}(A^{-1}B).
\]
Pour $A \in \SL(n, \R) = \det^{-1}(1)$, on a $\det(A) = 1$, donc
$\dd(\det)_A(B) = \operatorname{tr}(A^{-1}B)$, qui est surjective
(prendre par exemple $B = \frac{1}{n}A$, ce qui donne $\operatorname{tr}(I_n/n) \cdot n = 1$).

Ainsi, $1$ est une valeur régulière du déterminant et
\[
  \SL(n, \R) = \det^{-1}(1)
\]
est une sous-variété plongée de $\GL(n, \R)$ de dimension $n^2 - 1$.
L'espace tangent en $I_n$ est
\[
  T_{I_n}\SL(n, \R) = \set{B \in M_n(\R) : \operatorname{tr}(B) = 0} = \mathfrak{sl}(n, \R).
\]
\end{exemple}

\subsection{Préimages dans $\R^n \to \R^m$}\label{subsec:preimages-Rn-Rm}

\begin{exemple}[Tore comme préimage]\label{ex:tore-preimage}
\index{tore!comme préimage}
Considérons $f : \R^3 \to \R$ définie par
\[
  f(x, y, z) = \left(\sqrt{x^2 + y^2} - R\right)^2 + z^2,
\]
où $R > r > 0$. La valeur $r^2$ est une valeur régulière de $f$
(on le vérifie par un calcul direct de $\dd f$), et
\[
  T^2 = f^{-1}(r^2)
\]
est un tore de révolution, sous-variété plongée de $\R^3$ de dimension $2$.
\end{exemple}

\begin{exemple}[Hyperboloïde à une nappe]\label{ex:hyperboloide}
\index{hyperboloïde}
L'application $f : \R^3 \to \R$, $f(x,y,z) = x^2 + y^2 - z^2$, a pour différentielle
$\dd f_{(x,y,z)} = (2x, 2y, -2z)$, qui s'annule uniquement à l'origine. Donc toute
valeur $c \neq 0$ est régulière et $f^{-1}(c)$ est une surface lisse.
Pour $c = 1$, on obtient l'hyperboloïde à une nappe.
\end{exemple}

\begin{exemple}[Préimage dans $\R^4 \to \R^2$]\label{ex:preimage-R4-R2}
Considérons $f : \R^4 \to \R^2$ définie par
\[
  f(x_1, x_2, x_3, x_4) = (x_1^2 + x_2^2 - 1,\; x_3^2 + x_4^2 - 1).
\]
La matrice jacobienne est
\[
  Df_{(x_1,x_2,x_3,x_4)} = \begin{pmatrix} 2x_1 & 2x_2 & 0 & 0 \\ 0 & 0 & 2x_3 & 2x_4 \end{pmatrix}.
\]
Cette matrice est de rang $2$ en tout point de $f^{-1}(0,0)$ (car $(x_1, x_2) \neq 0$
et $(x_3, x_4) \neq 0$ sur cet ensemble). Donc $(0,0)$ est une valeur régulière et
\[
  f^{-1}(0,0) = S^1 \times S^1 = T^2
\]
est une sous-variété plongée de $\R^4$ de dimension $2$, difféomorphe au tore.
\end{exemple}

\subsection{Le théorème de Sard}\label{subsec:sard}

Le théorème de la valeur régulière est d'autant plus puissant que la plupart
des valeurs sont régulières, au sens suivant.

\begin{theoreme}[Théorème de Sard]\label{thm:sard}
\index{théorème!de Sard}\index{Sard (théorème de)}
Soit $f : M \to N$ une application lisse. L'ensemble des valeurs critiques
de $f$ est de mesure de Lebesgue nulle\index{mesure nulle} dans $N$.
En particulier, l'ensemble des valeurs régulières est dense dans $N$.
\end{theoreme}

\begin{remarque}\label{rem:sard-utilite}
La preuve du théorème de Sard, bien que technique, repose essentiellement sur
le théorème de Fubini et des estimations de volume. Nous l'admettons ici et
renvoyons à \cite{milnor1965topology} pour une démonstration complète.
Ce théorème, combiné au théorème de la valeur régulière, montre que pour
une application lisse \og{}générique\fg{}, presque toutes les fibres sont des
sous-variétés lisses.
\end{remarque}

% ----------------------------------------------------------
%  Espace tangent d'une sous-variété
% ----------------------------------------------------------
\section{Espace tangent d'une sous-variété}\label{sec:tangent-sous-variete}

\begin{proposition}[Espace tangent d'une sous-variété plongée]\label{prop:tangent-sous-variete}
\index{espace tangent!sous-variété}
Soit $S \subset M$ une sous-variété plongée et $\iota : S \hookrightarrow M$ l'inclusion.
Pour tout $p \in S$, l'application $\dd\iota_p : T_pS \to T_pM$ est injective et
identifie $T_pS$ au sous-espace
\[
  \dd\iota_p(T_pS) = \set{v \in T_pM : v(f) = 0 \;\text{pour toute}\; f \in C^\infty(M)
  \;\text{avec}\; f|_S = 0}.
\]
\end{proposition}

\begin{proof}
L'inclusion $\iota$ est un plongement, donc $\dd\iota_p$ est injective.
Soit $v = [\gamma]_p \in T_pS$ et $f \in C^\infty(M)$ avec $f|_S = 0$. Alors
$\dd\iota_p(v)(f) = v(f \circ \iota) = v(f|_S) = v(0) = 0$.
Réciproquement, soit $v \in T_pM$ annulant toute $f$ qui s'annule sur $S$.
En carte de sous-variété $(U, \varphi = (x^1, \ldots, x^k, x^{k+1}, \ldots, x^n))$
avec $U \cap S = \set{x^{k+1} = \cdots = x^n = 0}$, les fonctions
$x^{k+1}, \ldots, x^n$ s'annulent sur $S$, donc $v(x^j) = 0$ pour $j > k$.
Ainsi $v = \sum_{i=1}^k v(x^i) \frac{\partial}{\partial x^i}\big|_p \in \dd\iota_p(T_pS)$.
\end{proof}

\begin{corollaire}[Espace tangent d'une préimage]\label{cor:tangent-preimage}
\index{espace tangent!préimage}
Si $c$ est une valeur régulière de $f : M \to N$ et $S = f^{-1}(c)$, alors
pour tout $p \in S$ :
\[
  T_pS = \ker \dd f_p.
\]
\end{corollaire}

% ----------------------------------------------------------
%  Théorème du rang constant
% ----------------------------------------------------------
\section{Théorème du rang constant}\label{sec:rang-constant}

\begin{theoreme}[Théorème du rang constant]\label{thm:rang-constant}
\index{théorème!du rang constant}
Soit $f : M^m \to N^n$ une application lisse de rang constant $r$ sur un
voisinage de $p \in M$. Alors il existe des cartes $(U, \varphi)$ centrées en $p$
et $(V, \psi)$ centrées en $f(p)$ telles que
\[
  \psi \circ f \circ \varphi^{-1}(x^1, \ldots, x^m) = (x^1, \ldots, x^r, 0, \ldots, 0).
\]
\end{theoreme}

\begin{proof}
Par un changement de coordonnées préliminaire, on peut supposer que le mineur
$r \times r$ en haut à gauche de la matrice jacobienne $D\hat{f}_0$ est inversible.
Définissons $\Phi : \R^m \to \R^m$ par
\[
  \Phi(x) = (\hat{f}^1(x), \ldots, \hat{f}^r(x), x^{r+1}, \ldots, x^m).
\]
Comme $D\Phi_0$ est inversible (le mineur $r \times r$ et le bloc identité
se combinent), $\Phi$ est un difféomorphisme local. En composant, on obtient
\[
  (\hat{f} \circ \Phi^{-1})(u) = (u^1, \ldots, u^r, h^{r+1}(u), \ldots, h^n(u))
\]
pour certaines fonctions lisses $h^j$. L'hypothèse de rang constant $r$ implique
que $\frac{\partial h^j}{\partial u^i} = 0$ pour tout $i > r$ et $j > r$.
Donc $h^j(u) = h^j(u^1, \ldots, u^r, 0, \ldots, 0) =: k^j(u^1, \ldots, u^r)$.
On définit alors $\Psi : \R^n \to \R^n$ par
\[
  \Psi(y^1, \ldots, y^n) = (y^1, \ldots, y^r, y^{r+1} - k^{r+1}(y^1, \ldots, y^r),
  \ldots, y^n - k^n(y^1, \ldots, y^r)).
\]
C'est un difféomorphisme local et
$(\Psi \circ \hat{f} \circ \Phi^{-1})(u) = (u^1, \ldots, u^r, 0, \ldots, 0)$.
\end{proof}

\begin{corollaire}[Version sous-variété du théorème du rang]\label{cor:rang-sous-variete}
\index{théorème!du rang constant!version sous-variété}
Sous les hypothèses du 
ef{thm:rang-constant} :
\begin{enumerate}[label=(\roman*)]
  \item Si $f$ est une immersion ($r = m$), alors localement $f$ s'écrit comme
    une inclusion linéaire $\R^m \hookrightarrow \R^n$.
  \item Si $f$ est une submersion ($r = n$), alors localement $f$ s'écrit comme
    une projection linéaire $\R^m \twoheadrightarrow \R^n$.
  \item Si $f$ est de rang constant $r$, alors pour tout $c \in f(M)$, $f^{-1}(c)$
    est une sous-variété plongée de dimension $m - r$.
\end{enumerate}
\end{corollaire}

\begin{proof}
Les points (i) et (ii) sont des cas particuliers de la forme normale. Pour (iii),
en coordonnées normales, $f^{-1}(c) \cap U$ est défini par les $r$ premières
coordonnées constantes, ce qui donne la structure de sous-variété.
\end{proof}

% ----------------------------------------------------------
%  Sous-variétés à bord
% ----------------------------------------------------------
\section{Sous-variétés à bord}\label{sec:sous-varietes-bord}

Rappelons le demi-espace\index{demi-espace} $\mathbb{H}^n = \set{x \in \R^n : x^n \geqslant 0}$.

\begin{definition}[Variété à bord]\label{def:variete-bord}
\index{variété à bord}
Une \textbf{variété à bord} de dimension $n$ est un espace topologique $M$ de
Hausdorff, à base dénombrable, muni d'un atlas maximal de cartes à valeurs dans
des ouverts de $\mathbb{H}^n$ dont les changements de cartes sont lisses.
Le \textbf{bord}\index{bord} de $M$ est
\[
  \partial M = \set{p \in M : \varphi(p) \in \partial \mathbb{H}^n \text{ pour une (toute) carte } (U, \varphi)},
\]
et l'\textbf{intérieur} est $\operatorname{Int}(M) = M \setminus \partial M$.
\end{definition}

\begin{proposition}[Le bord est une variété sans bord]\label{prop:bord-sans-bord}
\index{bord!variété sans bord}
Si $M$ est une variété à bord de dimension $n$, alors $\partial M$ est une variété
lisse (sans bord) de dimension $n - 1$.
\end{proposition}

\begin{definition}[Sous-variété à bord propre]\label{def:sous-variete-bord}
\index{sous-variété!à bord}
Soit $M$ une variété sans bord. Un sous-ensemble $S \subset M$ est une
\textbf{sous-variété à bord propre} de $M$ si, pour tout $p \in S$, il existe une
carte $(U, \varphi)$ de $M$ telle que
\[
  \varphi(U \cap S) = \varphi(U) \cap \mathbb{H}^k
\]
pour un certain $k \leqslant n$.
\end{definition}

\begin{exemple}[Disque comme sous-variété à bord]\label{ex:disque-bord}
\index{disque!sous-variété à bord}
Le disque fermé $\overline{D}^2 = \set{x \in \R^2 : \abs{x}^2 \leqslant 1}$ est une
sous-variété à bord de $\R^2$, de bord $\partial\overline{D}^2 = S^1$.
\end{exemple}

% ----------------------------------------------------------
%  Voisinage tubulaire
% ----------------------------------------------------------
\section{Voisinage tubulaire}\label{sec:voisinage-tubulaire}

Le théorème du voisinage tubulaire affirme qu'une sous-variété plongée admet
un voisinage qui est difféomorphe à son fibré normal. C'est un outil fondamental
pour la théorie de l'intersection et la transversalité.

\begin{definition}[Fibré normal]\label{def:fibre-normal}
\index{fibré normal}
Soit $S \subset M$ une sous-variété plongée et supposons $M$ munie d'une métrique
riemannienne $g$. Le \textbf{fibré normal} de $S$ dans $M$ est le fibré vectoriel
\[
  NS = \bigsqcup_{p \in S} N_pS, \qquad N_pS = (T_pS)^\perp \subset T_pM,
\]
où l'orthogonalité est prise par rapport à $g_p$.
\end{definition}

\begin{theoreme}[Voisinage tubulaire]\label{thm:voisinage-tubulaire}
\index{théorème!du voisinage tubulaire}\index{voisinage tubulaire}
Soit $S$ une sous-variété plongée compacte d'une variété lisse $M$.
Il existe un voisinage ouvert $U$ de $S$ dans $M$ et un difféomorphisme
\[
  \Phi : NS \supset V \xrightarrow{\;\sim\;} U
\]
d'un voisinage $V$ de la section nulle dans $NS$ sur $U$, tel que $\Phi$ restreint
à la section nulle soit l'identité $S \to S$.
\end{theoreme}

\begin{remarque}\label{rem:voisinage-tubulaire-construction}
La preuve (que nous omettons) repose sur l'application exponentielle riemannienne :
on définit $\Phi(p, v) = \exp_p(v)$ pour $v \in N_pS$ suffisamment petit.
La compacité de $S$ assure l'existence d'un $\varepsilon > 0$ uniforme tel que
$\Phi$ soit défini et soit un difféomorphisme sur le $\varepsilon$-voisinage tubulaire.
\end{remarque}

\begin{figure}[ht]
\centering
\begin{tikzpicture}[scale=1.2]
  % Sous-variété S (courbe)
  \draw[very thick, red!70!black]
    (-3, 0) .. controls (-1.5, 1) and (1.5, -1) .. (3, 0);
  \node[red!70!black] at (3.3, 0) {$S$};

  % Voisinage tubulaire
  \draw[thick, blue!40, fill=blue!5, opacity=0.6]
    (-3, 0.8) .. controls (-1.5, 1.8) and (1.5, -0.2) .. (3, 0.8)
    -- (3, -0.8) .. controls (1.5, -1.8) and (-1.5, 0.2) .. (-3, -0.8) -- cycle;
  \node[blue!60!black] at (0, 1.6) {$U \cong NS$};

  % Fibres normales
  \foreach \x in {-2, -1, 0, 1, 2} {
    \pgfmathsetmacro{\ybase}{0.22*\x*\x - 0.11*\x*\x*\x/3 + 0.05*\x}
    \draw[gray, thin, ->] (\x, \ybase - 0.6) -- (\x, \ybase + 0.6);
  }
\end{tikzpicture}
\caption{Un voisinage tubulaire d'une sous-variété $S$ : chaque fibre
normale est représentée comme un segment transverse.}\label{fig:voisinage-tubulaire}
\end{figure}

% ----------------------------------------------------------
%  Exercices du chapitre 4
% ----------------------------------------------------------
\section{Exercices}\label{sec:exo-ch4}

\begin{exercice}\label{exo:sphere-variete}
\index{sphère!sous-variété}
Soit $f : \R^{n+1} \to \R$ définie par $f(x) = \sum_{i=1}^{n+1} (x^i)^2$.
\begin{enumerate}[label=(\alph*)]
  \item Montrer que $1$ est une valeur régulière de $f$.
  \item En déduire que $S^n$ est une sous-variété plongée de $\R^{n+1}$ de dimension $n$.
  \item Calculer $T_xS^n$ pour tout $x \in S^n$.
\end{enumerate}
\end{exercice}

\begin{exercice}\label{exo:SOn-variete}
\index{SO(n)@$\SO(n)$!comme sous-variété}
Montrer que $\SO(n) = \set{A \in O(n) : \det A = 1}$ est une sous-variété
plongée de $M_n(\R)$ de même dimension que $O(n)$.
\textit{Indication :} montrer que $\SO(n)$ est une composante connexe de $O(n)$.
\end{exercice}

\begin{exercice}\label{exo:stiefel}
\index{variété de Stiefel}
La \textbf{variété de Stiefel} $V_k(\R^n)$ est l'ensemble des $k$-repères
orthonormaux dans $\R^n$, c'est-à-dire
\[
  V_k(\R^n) = \set{A \in M_{n \times k}(\R) : A^\top A = I_k}.
\]
Montrer que $V_k(\R^n)$ est une sous-variété plongée de $M_{n \times k}(\R)$
et calculer sa dimension.
\end{exercice}

\begin{exercice}\label{exo:grassmannienne-dim}
\index{grassmannienne}
En utilisant le fait que $\GL(k, \R)$ agit librement sur $V_k(\R^n)$,
montrer que la grassmannienne $\operatorname{Gr}_k(\R^n) = V_k(\R^n)/\GL(k, \R)$
est une variété lisse de dimension $k(n - k)$.
\end{exercice}

\begin{exercice}\label{exo:rang-constant-surjectif}
Soit $f : M \to N$ une application lisse de rang constant $r$ avec $M$ compacte et connexe.
Montrer que $f(M)$ est une sous-variété plongée compacte de $N$ de dimension $r$.
\end{exercice}

\begin{exercice}\label{exo:produit-sous-variete}
\index{produit de sous-variétés}
Soient $S_1 \subset M_1$ et $S_2 \subset M_2$ des sous-variétés plongées.
Montrer que $S_1 \times S_2$ est une sous-variété plongée de $M_1 \times M_2$
et que $\operatorname{codim}(S_1 \times S_2) = \operatorname{codim}(S_1) + \operatorname{codim}(S_2)$.
\end{exercice}

\begin{exercice}\label{exo:submersion-sections}
Soit $\pi : M \to N$ une submersion surjective.
\begin{enumerate}[label=(\alph*)]
  \item Montrer que $\pi$ est une application ouverte.
  \item Montrer que pour tout $q \in N$, la fibre $\pi^{-1}(q)$ est une
    sous-variété plongée de $M$ de codimension $\dim N$.
  \item Si de plus $\pi$ est propre, montrer que $\pi$ est un fibré lisse
    localement trivial (théorème d'Ehresmann\index{théorème!d'Ehresmann}).
\end{enumerate}
\end{exercice}

\begin{exercice}\label{exo:preimage-transverse}
\index{transversalité}
Soient $f : M \to N$ une application lisse et $S \subset N$ une sous-variété plongée.
On dit que $f$ est \textbf{transverse} à $S$ (noté $f \pitchfork S$) si pour tout
$p \in f^{-1}(S)$ :
\[
  \dd f_p(T_pM) + T_{f(p)}S = T_{f(p)}N.
\]
Montrer que si $f \pitchfork S$, alors $f^{-1}(S)$ est une sous-variété plongée
de $M$ de codimension $\operatorname{codim}(S)$.
\textit{Indication :} se ramener au théorème de la valeur régulière en composant
$f$ avec une submersion locale définissant $S$.
\end{exercice}

\begin{exercice}\label{exo:SU2-S3}
\index{SU(2)@$\SU(2)$}
Montrer que le groupe unitaire spécial $\SU(2) = \set{A \in M_2(\C) : A\bar{A}^\top = I_2,\; \det A = 1}$
est difféomorphe à $S^3$. \textit{Indication :} écrire la forme générale d'un
élément de $\SU(2)$ et identifier avec $S^3 \subset \R^4 \cong \C^2$.
\end{exercice}

\begin{exercice}\label{exo:niveau-morse}
\index{fonction de Morse}
Soit $f : \R^3 \to \R$ définie par $f(x,y,z) = x^2 + y^2 - z^2$.
\begin{enumerate}[label=(\alph*)]
  \item Déterminer les points critiques et les valeurs critiques de $f$.
  \item Pour chaque valeur régulière $c$, décrire la topologie de $f^{-1}(c)$.
  \item Dessiner les lignes de niveau pour $c = -1, 0, 1$.
  \item Que se passe-t-il au passage par la valeur critique $c = 0$ ?
\end{enumerate}
\end{exercice}

\begin{figure}[ht]
\centering
\begin{tikzpicture}[scale=0.9]
  % Lignes de niveau de f(x,y,z) = x² + y² - z², dans le plan (x,z)
  \begin{scope}[shift={(-4.5,0)}]
    \draw[->] (-2.3, 0) -- (2.3, 0) node[right] {$x$};
    \draw[->] (0, -2.3) -- (0, 2.3) node[above] {$z$};
    % c = 1 : hyperbole x² - z² = 1
    \draw[thick, blue] plot[domain=-1.5:1.5, samples=50] ({cosh(\x)}, {sinh(\x)});
    \draw[thick, blue] plot[domain=-1.5:1.5, samples=50] ({-cosh(\x)}, {sinh(\x)});
    \node[blue] at (2.0, -1.8) {$c = 1$};
    \node at (0, -2.8) {(section $y = 0$)};
  \end{scope}
  \begin{scope}
    \draw[->] (-2.3, 0) -- (2.3, 0) node[right] {$x$};
    \draw[->] (0, -2.3) -- (0, 2.3) node[above] {$z$};
    % c = 0 : cône x² - z² = 0
    \draw[thick, red!70!black] (-2, -2) -- (0,0) -- (2, 2);
    \draw[thick, red!70!black] (-2, 2) -- (0,0) -- (2, -2);
    \node[red!70!black] at (2.0, -1.8) {$c = 0$};
  \end{scope}
  \begin{scope}[shift={(4.5,0)}]
    \draw[->] (-2.3, 0) -- (2.3, 0) node[right] {$x$};
    \draw[->] (0, -2.3) -- (0, 2.3) node[above] {$z$};
    % c = -1 : hyperbole z² - x² = 1
    \draw[thick, green!50!black] plot[domain=-1.5:1.5, samples=50] ({sinh(\x)}, {cosh(\x)});
    \draw[thick, green!50!black] plot[domain=-1.5:1.5, samples=50] ({sinh(\x)}, {-cosh(\x)});
    \node[green!50!black] at (2.0, -1.8) {$c = -1$};
  \end{scope}
\end{tikzpicture}
\caption{Sections des ensembles de niveau de $f(x,y,z) = x^2 + y^2 - z^2$
dans le plan $y = 0$, pour les valeurs $c = 1, 0, -1$.}\label{fig:niveaux-morse}
\end{figure}

% === FIN DU CHUNK ch3-4 ===
% === DEBUT DU CHUNK ch5-6 ===

% ============================================================
%  CHAPITRE 5 : TRANSVERSALITÉ
% ============================================================
\chapter{Transversalité}\label{ch:transversalite}
\minitoc\vspace{1cm}

La transversalité est l'un des concepts les plus puissants de la topologie
différentielle. Intuitivement, deux sous-variétés se coupent transversalement
lorsqu'elles se croisent de la manière la plus \og{}générale\fg{} possible, sans
tangence. Ce chapitre développe cette notion, établit le théorème fondamental
reliant transversalité et structure de sous-variété, puis montre que la
transversalité est une condition \emph{générique} au sens de Baire. Les
applications aux nombres d'intersection et au théorème de Lefschetz
illustrent la richesse de cette théorie.

% ============================================================
\section{Intersection transverse}\label{sec:intersection-transverse}
% ============================================================

\subsection{Définition pour une application}\label{subsec:transverse-application}

\begin{definition}[Transversalité d'une application à une sous-variété]%
\label{def:transverse-application}\index{transversalité!d'une application}
Soient $f \colon M \to N$ une application lisse entre variétés lisses et
$W \subset N$ une sous-variété (plongée). On dit que $f$ est
\textbf{transverse} à $W$, et on note $f \pitchfork W$, si pour tout
$x \in f^{-1}(W)$ on a
\[
  \dd f_x(T_x M) + T_{f(x)} W = T_{f(x)} N.
\]
Autrement dit, l'image de la différentielle de $f$ en $x$ et l'espace
tangent à $W$ en $f(x)$ engendrent ensemble tout l'espace tangent à $N$.
\end{definition}

\begin{definition}[Transversalité de deux sous-variétés]%
\label{def:transverse-sous-varietes}\index{transversalité!de sous-variétés}
Soient $M$ et $W$ deux sous-variétés d'une variété ambiante $N$. On dit
que $M$ et $W$ \textbf{se coupent transversalement}, et on note
$M \pitchfork W$, si pour tout $p \in M \cap W$ on a
\[
  T_p M + T_p W = T_p N.
\]
\end{definition}

\begin{remarque}\label{rem:transverse-inclusion}
La définition~\ref{def:transverse-sous-varietes} est un cas particulier de
la définition~\ref{def:transverse-application} : il suffit de prendre
$f \colon M \hookrightarrow N$ l'inclusion. Notons aussi que si
$f^{-1}(W) = \varnothing$, la condition de transversalité est
automatiquement satisfaite (il n'y a aucun point à vérifier).
\end{remarque}

\begin{remarque}\label{rem:transverse-codimension}
Si $\dim M + \dim W < \dim N$, alors la transversalité $M \pitchfork W$
implique $M \cap W = \varnothing$ (pour des raisons de dimension, la somme
$T_p M + T_p W$ ne peut jamais être $T_p N$ si les deux sous-espaces sont
en position générale).
\end{remarque}

% -------- TikZ : transverse vs non-transverse --------
\begin{figure}[ht]
\centering
\begin{tikzpicture}[scale=1.1]
  % Intersection transverse
  \begin{scope}[shift={(-3.5,0)}]
    \draw[thick, blue!70!black] (-1.8,0) -- (1.8,0)
      node[right]{\small $M$};
    \draw[thick, red!70!black] (0,-1.5) -- (0,1.5)
      node[above]{\small $W$};
    \fill (0,0) circle (2pt);
    \node[below right] at (0.1,-0.2) {\small $p$};
    \node at (0,-2.2) {$M \pitchfork W$ : transverse};
  \end{scope}
  % Intersection non transverse
  \begin{scope}[shift={(3.5,0)}]
    \draw[thick, blue!70!black] (-1.8,0) .. controls (-0.5,0)
      and (0.5,0.8) .. (1.8,0.8)
      node[right]{\small $M$};
    \draw[thick, red!70!black] (-1.8,0) .. controls (-0.5,0)
      and (0.5,-0.8) .. (1.8,-0.8)
      node[right]{\small $W$};
    \fill (-1.2,0) circle (2pt);
    \node[above] at (-1.2,0.15) {\small $p$};
    \node at (0,-2.2) {Non transverse : tangence en $p$};
  \end{scope}
\end{tikzpicture}
\caption{À gauche, deux courbes dans $\R^2$ se coupant transversalement ;
  à droite, une intersection non transverse (tangente).}
\label{fig:transverse-vs-non}
\end{figure}

\subsection{Exemples et non-exemples}\label{subsec:exemples-transversalite}

\begin{exemple}[Droites dans le plan]\label{ex:droites-plan}%
\index{transversalité!exemples}
Dans $\R^2$, deux droites distinctes et non parallèles se coupent
transversalement. Deux droites parallèles distinctes vérifient trivialement
la transversalité (intersection vide). En revanche, une droite considérée
comme sous-variété de $\R^2$ se coupe elle-même de manière non transverse
(la somme $T_p L + T_p L = T_p L \neq T_p \R^2$).
\end{exemple}

\begin{exemple}[Sphère et plan]\label{ex:sphere-plan}
Dans $\R^3$, la sphère $S^2$ et le plan $\set{z = c}$ se coupent
transversalement si et seulement si $c \neq \pm 1$. Aux pôles $c = \pm 1$,
le plan est tangent à la sphère.
\end{exemple}

\begin{exemple}[Courbes sur une surface]\label{ex:courbes-surface}
Sur le tore $T^2 = S^1 \times S^1$, les cercles
$\gamma_1 = S^1 \times \set{1}$ et $\gamma_2 = \set{1} \times S^1$ se
coupent transversalement en un unique point.
\end{exemple}

\begin{exemple}[Application non transverse]\label{ex:application-non-transverse}
L'application $f \colon \R \to \R^2$, $t \mapsto (t, t^2)$ n'est pas
transverse à l'axe des abscisses $W = \R \times \set{0}$ en $t = 0$, car
$\dd f_0(\R) = \R \times \set{0} = T_{(0,0)} W$.
\end{exemple}

% -------- TikZ : courbes sur le tore --------
\begin{figure}[ht]
\centering
\begin{tikzpicture}[scale=0.9]
  % Tore simplifié : représentation rectangle
  \draw[thick] (0,0) rectangle (4,3);
  \draw[thick, blue!70!black, ->] (0,1.5) -- (4,1.5)
    node[midway, above]{\small $\gamma_1$};
  \draw[thick, red!70!black, ->] (2,0) -- (2,3)
    node[midway, right]{\small $\gamma_2$};
  \fill (2,1.5) circle (2.5pt);
  \node[above right] at (2.1,1.6) {\small $p$};
  % Identifications
  \draw[thick, ->, green!50!black] (4.3,0.5) -- (4.3,2.5);
  \draw[thick, ->, green!50!black] (0.5,3.3) -- (3.5,3.3);
  \node[right, green!50!black] at (4.4,1.5) {\small $\sim$};
  \node[above, green!50!black] at (2,3.4) {\small $\sim$};
  \node at (2,-0.6) {Domaine fondamental du tore $T^2$};
\end{tikzpicture}
\caption{Deux cercles générateurs sur le tore, se coupant
  transversalement en un point.}
\label{fig:courbes-tore}
\end{figure}

% ============================================================
\section{Théorème de transversalité}\label{sec:thm-transversalite}
% ============================================================

Le résultat fondamental suivant montre que l'image réciproque d'une
sous-variété par une application transverse est encore une sous-variété,
de la codimension attendue.

\begin{theoreme}[Préimage transverse]\label{thm:preimage-transverse}%
\index{théorème!de la préimage transverse}\index{transversalité!théorème de la préimage}
Soient $f \colon M \to N$ une application lisse et $W \subset N$ une
sous-variété de codimension $k$. Si $f \pitchfork W$, alors :
\begin{enumerate}[label=\textup{(\roman*)}]
  \item $f^{-1}(W)$ est une sous-variété lisse de $M$ ;
  \item $\operatorname{codim}_M f^{-1}(W) = \operatorname{codim}_N W = k$ ;
  \item pour tout $x \in f^{-1}(W)$, on a
    $T_x(f^{-1}(W)) = (\dd f_x)^{-1}(T_{f(x)} W)$.
\end{enumerate}
\end{theoreme}

\begin{proof}
Soit $x_0 \in f^{-1}(W)$ et posons $y_0 = f(x_0) \in W$.

\emph{Étape 1 : Linéarisation locale de $W$.}
Comme $W$ est une sous-variété de codimension $k$ de $N$, il existe un
voisinage ouvert $V$ de $y_0$ dans $N$ et une submersion
$g \colon V \to \R^k$ telle que $W \cap V = g^{-1}(0)$.

\emph{Étape 2 : La composée $g \circ f$ est une submersion.}
Considérons l'application $h = g \circ f \colon f^{-1}(V) \to \R^k$.
Pour $x \in f^{-1}(W \cap V)$, on a $f(x) \in W$, donc
$\ker \dd g_{f(x)} = T_{f(x)} W$. La condition de transversalité
$\dd f_x(T_x M) + T_{f(x)} W = T_{f(x)} N$ et la surjectivité de
$\dd g_{f(x)} \colon T_{f(x)} N \to \R^k$ donnent :
\[
  \dd h_x(T_x M) = \dd g_{f(x)}(\dd f_x(T_x M))
    = \dd g_{f(x)}\bigl(\dd f_x(T_x M) + T_{f(x)} W\bigr)
    = \dd g_{f(x)}(T_{f(x)} N) = \R^k.
\]
Donc $h = g \circ f$ est une submersion en tout point de $h^{-1}(0)$.

\emph{Étape 3 : Application du théorème de la valeur régulière.}
Puisque $h$ est une submersion en tout point de $h^{-1}(0)$, le théorème
de la valeur régulière assure que $h^{-1}(0) = f^{-1}(W) \cap f^{-1}(V)$
est une sous-variété de $f^{-1}(V)$ de codimension $k$. Comme ceci vaut
au voisinage de chaque point $x_0 \in f^{-1}(W)$, on obtient que
$f^{-1}(W)$ est une sous-variété de $M$ de codimension $k$.

\emph{Étape 4 : Espace tangent.}
Pour $x \in f^{-1}(W)$, on a
\[
  T_x(f^{-1}(W)) = T_x(h^{-1}(0)) = \ker \dd h_x
    = \ker(\dd g_{f(x)} \circ \dd f_x)
    = (\dd f_x)^{-1}(\ker \dd g_{f(x)})
    = (\dd f_x)^{-1}(T_{f(x)} W). \qedhere
\]
\end{proof}

\begin{corollaire}\label{cor:transverse-sous-varietes}%
\index{transversalité!intersection de sous-variétés}
Si $M$ et $W$ sont deux sous-variétés de $N$ avec $M \pitchfork W$,
alors $M \cap W$ est une sous-variété de $N$ et
\[
  \operatorname{codim}(M \cap W) = \operatorname{codim}(M) +
    \operatorname{codim}(W).
\]
De manière équivalente,
$\dim(M \cap W) = \dim M + \dim W - \dim N$.
\end{corollaire}

\begin{proof}
On applique le théorème~\ref{thm:preimage-transverse} à l'inclusion
$\iota \colon M \hookrightarrow N$, qui est transverse à $W$
précisément lorsque $M \pitchfork W$. On a
$\iota^{-1}(W) = M \cap W$.
\end{proof}

\begin{exemple}\label{ex:intersection-spheres}
Dans $\R^3$, considérons les sphères
$S_1 = \set{x^2 + y^2 + z^2 = 1}$ et
$S_2 = \set{(x-1)^2 + y^2 + z^2 = 1}$. Elles se coupent
transversalement car en tout point d'intersection, les normales sont
distinctes. Par le corollaire, $S_1 \cap S_2$ est une sous-variété de
dimension $2 + 2 - 3 = 1$ : c'est un cercle.
\end{exemple}

% ============================================================
\section{Transversalité et stabilité}\label{sec:stabilite-transversalite}
% ============================================================

La transversalité est une condition \og{}ouverte\fg{} : elle persiste
sous de petites perturbations.

\begin{proposition}[Stabilité de la transversalité]%
\label{prop:stabilite-transversalite}\index{transversalité!stabilité}
Soient $M$ compacte, $N$ une variété lisse et $W \subset N$ une
sous-variété fermée. L'ensemble
\[
  \pitchfork(M, N; W) = \set{f \in C^\infty(M, N) \mid f \pitchfork W}
\]
est un ouvert de $C^\infty(M, N)$ pour la topologie $C^1$ (et a fortiori
pour la topologie $C^\infty$).
\end{proposition}

\begin{proof}[Idée de preuve]
La condition de transversalité en un point est une condition de rang
maximal sur une application linéaire, donc ouverte. Comme $M$ est
compacte, on peut recouvrir $f^{-1}(W)$ par un nombre fini de cartes et
utiliser la continuité de la différentielle en topologie $C^1$ pour
conclure.
\end{proof}

\begin{lemme}[Transversalité et rang]\label{lem:transversalite-rang}
Soit $f \colon M \to N$ lisse et $W \subset N$ une sous-variété de
codimension $k$ définie localement comme $g^{-1}(0)$ pour une submersion
$g \colon V \to \R^k$. Alors $f \pitchfork W$ sur $f^{-1}(V)$ si et
seulement si $g \circ f$ est de rang $k$ en tout point de
$(g \circ f)^{-1}(0)$.
\end{lemme}

\begin{proof}
C'est exactement le calcul effectué dans l'étape~2 de la preuve du
théorème~\ref{thm:preimage-transverse} : la surjectivité de
$\dd(g \circ f)_x$ est équivalente à la condition de transversalité.
\end{proof}

\begin{exemple}[Perturbation transverse]\label{ex:perturbation}
Considérons dans $\R^2$ la parabole $C = \set{y = x^2}$ et l'axe
$W = \set{y = 0}$. L'intersection $C \cap W = \set{(0,0)}$ n'est pas
transverse car $T_{(0,0)} C = T_{(0,0)} W = \R \times \set{0}$. Cependant,
pour tout $\varepsilon > 0$, la translation $C_\varepsilon = \set{y = x^2 - \varepsilon}$
coupe $W$ transversalement en deux points $(\pm\sqrt{\varepsilon}, 0)$.
\end{exemple}

\begin{remarque}\label{rem:perturbation-transverse}
En pratique, si une application $f$ n'est pas transverse à $W$, on peut
souvent la \og{}perturber\fg{} légèrement pour obtenir une application
$\tilde{f}$, homotope à $f$, telle que $\tilde{f} \pitchfork W$.
C'est le contenu du théorème de transversalité de Thom.
\end{remarque}

% ============================================================
\section{Théorème de transversalité de Thom}\label{sec:thom-transversalite}
% ============================================================

Le théorème suivant affirme que la transversalité est une condition
\emph{générique}\index{généricité}, c'est-à-dire vérifiée par un
ensemble résiduel d'applications.

\begin{theoreme}[Transversalité de Thom]\label{thm:thom-transversalite}%
\index{théorème!de transversalité de Thom}\index{Thom!théorème de transversalité}
Soit $F \colon M \times S \to N$ une application lisse, où $M$, $S$, $N$
sont des variétés lisses avec $M$ de dimension finie. Soit
$W \subset N$ une sous-variété. Si $F \pitchfork W$, alors l'ensemble
\[
  \set{s \in S \mid F_s \pitchfork W},
  \quad \text{où } F_s = F(\cdot, s) \colon M \to N,
\]
est résiduel dans $S$ (et donc dense si $S$ est un espace de Baire).
De plus, si $M$ est compacte, cet ensemble est un ouvert dense.
\end{theoreme}

\begin{proof}[Esquisse de preuve]
\emph{Étape 1.}
Par hypothèse $F \pitchfork W$, donc le
théorème~\ref{thm:preimage-transverse} assure que
$\mathcal{W} = F^{-1}(W)$ est une sous-variété de $M \times S$
de codimension $\operatorname{codim}(W)$.

\emph{Étape 2.}
Considérons la projection $\pi \colon \mathcal{W} \to S$,
$\pi(x, s) = s$. On montre que $s \in S$ est une valeur régulière
de $\pi$ si et seulement si $F_s \pitchfork W$.

En effet, fixons $(x_0, s_0) \in \mathcal{W}$. L'espace tangent
$T_{(x_0,s_0)}\mathcal{W} = (\dd F_{(x_0,s_0)})^{-1}(T_{F(x_0,s_0)}W)$.
La condition que $\dd\pi$ soit surjective en $(x_0,s_0)$ signifie que la
composante en $S$ de cet espace tangent est tout $T_{s_0}S$, ce qui
se reformule exactement comme $F_{s_0} \pitchfork W$ en $x_0$.

\emph{Étape 3.}
Le théorème de Sard\index{Sard!théorème de} appliqué à $\pi$ donne que
l'ensemble des valeurs régulières de $\pi$ est résiduel (dense et de
mesure pleine). Donc $\set{s \in S \mid F_s \pitchfork W}$ est résiduel.

La partie sur l'ouverture lorsque $M$ est compacte découle de la
proposition~\ref{prop:stabilite-transversalite}.
\end{proof}

\begin{remarque}\label{rem:thom-interpretation}
Le théorème de Thom admet l'interprétation suivante : si l'on dispose
d'une famille d'applications suffisamment riche (paramétrée par $S$)
et si la \og{}famille totale\fg{} $F$ est transverse à $W$, alors
\og{}presque tout\fg{} membre de la famille est transverse à $W$. Le
paramètre $s \in S$ joue le rôle d'une perturbation, et l'espace $S$
doit être assez grand pour \og{}résoudre\fg{} tous les défauts de
transversalité.
\end{remarque}

\begin{exemple}[Application du théorème de Thom]\label{ex:thom-plongements}
Considérons $M$ une variété compacte plongée dans $\R^N$ et
$W \subset \R^N$ une sous-variété. Pour $s \in \R^N$, posons
$M_s = M + s$ (translaté). La famille $F \colon M \times \R^N \to \R^N$,
$F(x, s) = x + s$ est une submersion, donc $F \pitchfork W$. Par le
théorème de Thom, pour presque tout $s \in \R^N$ (au sens de la mesure
et de Baire), $M_s \pitchfork W$.
\end{exemple}

\begin{corollaire}[Densité de la transversalité]\label{cor:densite-transversalite}%
\index{transversalité!densité}
Soient $M$ compacte, $N$ une variété et $W \subset N$ une sous-variété
fermée. Alors $\set{f \in C^\infty(M, N) \mid f \pitchfork W}$ est un
ouvert dense de $C^\infty(M, N)$ pour la topologie $C^\infty$.
\end{corollaire}

\begin{proof}
La densité s'obtient en appliquant le théorème~\ref{thm:thom-transversalite}
avec un choix judicieux de famille paramétrée. Par exemple, si $N$ est
plongée dans $\R^K$ (par le théorème de Whitney), on pose
$F \colon M \times B_\varepsilon(0) \to N$, $(x, s) \mapsto \pi_N(f(x) + s)$
où $\pi_N$ est une projection tubulaire. On vérifie que $F \pitchfork W$
pour $\varepsilon$ assez petit. L'ouverture découle de la
proposition~\ref{prop:stabilite-transversalite}.
\end{proof}

% ============================================================
\section{Applications de la transversalité}\label{sec:applications-transversalite}
% ============================================================

\subsection{Nombre d'intersection}\label{subsec:nombre-intersection}

\begin{definition}[Nombre d'intersection orienté]%
\label{def:nombre-intersection}\index{nombre d'intersection}
Soient $M^m$ et $N^n$ deux variétés lisses compactes orientées sans bord,
et soit $f \colon M \to N$ une application lisse avec $f \pitchfork W$,
où $W \subset N$ est une sous-variété fermée orientée de codimension $m$.
Alors $f^{-1}(W)$ est un ensemble fini de points. En chaque point
$x \in f^{-1}(W)$, on définit un \textbf{signe}\index{signe d'intersection}
$\varepsilon(x) = \pm 1$ selon que l'isomorphisme
\[
  \dd f_x \colon T_x M / T_x(f^{-1}(W)) \xrightarrow{\;\sim\;} T_{f(x)} N / T_{f(x)} W
\]
préserve ou renverse l'orientation. Le \textbf{nombre d'intersection} de
$f$ et $W$ est
\[
  I(f, W) = \sum_{x \in f^{-1}(W)} \varepsilon(x) \in \Z.
\]
\end{definition}

\begin{theoreme}[Invariance homotopique du nombre d'intersection]%
\label{thm:invariance-intersection}\index{nombre d'intersection!invariance homotopique}
Le nombre d'intersection $I(f, W)$ ne dépend que de la classe d'homotopie
de $f$. Plus précisément, si $f_0 \simeq f_1$ et les deux sont
transverses à $W$, alors $I(f_0, W) = I(f_1, W)$.
\end{theoreme}

\begin{proof}[Idée de preuve]
Soit $H \colon M \times [0,1] \to N$ une homotopie lisse entre $f_0$
et $f_1$. Par le théorème de Thom, on peut supposer $H \pitchfork W$
(quitte à perturber). Alors $H^{-1}(W)$ est une $1$-variété compacte à
bord dans $M \times [0,1]$, dont le bord est
$f_0^{-1}(W) \sqcup f_1^{-1}(W)$. Un argument de comptage des signes sur
le bord d'une $1$-variété orientée montre l'égalité.
\end{proof}

\begin{definition}[Nombre d'intersection de deux sous-variétés]%
\label{def:intersection-sous-varietes}\index{nombre d'intersection!de sous-variétés}
Soient $M^n$ une variété orientée compacte sans bord, et soient
$A^a, B^b \subset M$ deux sous-variétés compactes orientées avec
$a + b = n$ et $A \pitchfork B$. Le nombre d'intersection est
\[
  I(A, B) = \sum_{p \in A \cap B} \varepsilon(p),
\]
où $\varepsilon(p) = +1$ si l'orientation de
$T_p A \oplus T_p B$ coïncide avec celle de $T_p M$, et $-1$ sinon.
\end{definition}

\begin{proposition}\label{prop:proprietes-intersection}
Le nombre d'intersection satisfait :
\begin{enumerate}[label=\textup{(\roman*)}]
  \item $I(A, B) = (-1)^{ab}\, I(B, A)$ où $a = \operatorname{codim} A$
    et $b = \operatorname{codim} B$ ;
  \item $I(A, B)$ ne dépend que des classes d'homologie de $A$ et $B$ ;
  \item si $A$ et $B$ sont homologiquement triviaux, alors $I(A, B) = 0$.
\end{enumerate}
\end{proposition}

\begin{proof}[Preuve de (i)]
Il suffit d'observer que le changement de l'ordre des facteurs dans
$T_p A \oplus T_p B \cong T_p M$ introduit un signe
$(-1)^{\dim A \cdot \dim B} = (-1)^{(n-a)(n-b)}$. Comme
$a + b = n$, on vérifie que $(n-a)(n-b) = ab \pmod{2}$, d'où le résultat.
\end{proof}

\begin{exemple}[Intersection sur le tore]\label{ex:intersection-tore-calcul}
Sur le tore $T^2 = S^1 \times S^1$, orienté par la forme d'aire standard,
les cercles $\alpha = S^1 \times \set{1}$ et $\beta = \set{1} \times S^1$
se coupent transversalement en $(1,1)$. L'orientation de
$T_{(1,1)}\alpha \oplus T_{(1,1)}\beta$ coïncide avec l'orientation de
$T^2$, donc $I(\alpha, \beta) = +1$. De même, $I(\beta, \alpha) = +1$
car $1 \cdot 1 = 1$ est impair, mais $(-1)^{1 \cdot 1} \cdot (+1) = -1$...
En fait, avec les conventions de codimension, $a = b = 1$ et
$I(\alpha, \beta) = (-1)^{1 \cdot 1} I(\beta, \alpha)$, donc
$I(\alpha, \beta) = -I(\beta, \alpha)$ si l'on échange l'ordre.
Ceci dépend de la convention d'orientation choisie ; avec la convention
usuelle, $I(\alpha, \beta) = +1$.
\end{exemple}

% -------- TikZ : signe d'intersection --------
\begin{figure}[ht]
\centering
\begin{tikzpicture}[scale=1.2]
  % Intersection positive
  \begin{scope}[shift={(-3,0)}]
    \draw[thick, blue!70!black, ->] (-1.3,0) -- (1.3,0)
      node[right]{\small $A$};
    \draw[thick, red!70!black, ->] (0,-1.3) -- (0,1.3)
      node[above]{\small $B$};
    \fill (0,0) circle (1.5pt);
    \node at (0,-1.8) {\small $\varepsilon = +1$};
  \end{scope}
  % Intersection négative
  \begin{scope}[shift={(3,0)}]
    \draw[thick, blue!70!black, ->] (-1.3,0) -- (1.3,0)
      node[right]{\small $A$};
    \draw[thick, red!70!black, ->] (0,1.3) -- (0,-1.3)
      node[below]{\small $B$};
    \fill (0,0) circle (1.5pt);
    \node at (0,-1.8) {\small $\varepsilon = -1$};
  \end{scope}
\end{tikzpicture}
\caption{Signe de l'intersection : $\varepsilon = +1$ lorsque les
  orientations de $A$ et $B$ donnent l'orientation ambiante,
  $\varepsilon = -1$ sinon.}
\label{fig:signe-intersection}
\end{figure}

\subsection{Nombre d'intersection modulo 2}\label{subsec:intersection-mod2}

\begin{definition}[Nombre d'intersection mod 2]%
\label{def:intersection-mod2}\index{nombre d'intersection!mod 2}
Si $M$, $A$, $B$ ne sont pas nécessairement orientables, on définit le
\textbf{nombre d'intersection modulo $2$} :
\[
  I_2(A, B) = \# (A \cap B) \mod 2 \in \Z/2\Z,
\]
à condition que $A \pitchfork B$ et que les dimensions soient
complémentaires. Ce nombre est un invariant d'homotopie.
\end{definition}

\begin{exemple}\label{ex:intersection-RP2}
Dans $\RP^2$, la droite projective standard $\RP^1 \subset \RP^2$ a un
nombre d'auto-intersection $I_2(\RP^1, \RP^1) = 1 \in \Z/2\Z$ (on
perturbe une copie et on compte les intersections modulo $2$).
\end{exemple}

\subsection{Nombre d'auto-intersection}\label{subsec:auto-intersection}

\begin{definition}[Nombre d'auto-intersection]%
\label{def:auto-intersection}\index{auto-intersection}
Soit $A$ une sous-variété compacte orientée de codimension $k$ dans une
variété orientée $M^n$. Le \textbf{nombre d'auto-intersection} de $A$
est
\[
  I(A, A) = I(A, A'),
\]
où $A'$ est une petite perturbation de $A$ telle que $A \pitchfork A'$.
Ce nombre est bien défini par l'invariance homotopique.
\end{definition}

\begin{exemple}\label{ex:auto-intersection-S2}
La section nulle de $S^2$ dans son fibré tangent $TS^2$ a un nombre
d'auto-intersection égal à $\chi(S^2) = 2$, la caractéristique d'Euler.
Plus généralement, le nombre d'auto-intersection de la diagonale
$\Delta \subset M \times M$ est $\chi(M)$.
\end{exemple}

\subsection{Lien avec le théorème de Lefschetz}\label{subsec:lefschetz}

\begin{theoreme}[Théorème du point fixe de Lefschetz --- énoncé]%
\label{thm:lefschetz}\index{Lefschetz!théorème du point fixe}%
\index{théorème!du point fixe de Lefschetz}
Soit $f \colon M \to M$ une application lisse d'une variété compacte
orientée sans bord. Le \textbf{nombre de Lefschetz} est
\[
  L(f) = \sum_{k=0}^{\dim M} (-1)^k \operatorname{tr}\bigl(f_* \colon
    H_k(M; \Q) \to H_k(M; \Q)\bigr).
\]
Si $L(f) \neq 0$, alors $f$ admet un point fixe.
\end{theoreme}

\begin{remarque}\label{rem:lefschetz-transversalite}
Le lien avec la transversalité est le suivant. Le graphe
$\Gamma_f = \set{(x, f(x)) \mid x \in M} \subset M \times M$ et la
diagonale $\Delta = \set{(x, x) \mid x \in M}$ sont des sous-variétés de
$M \times M$. Les points fixes de $f$ correspondent exactement aux points
de $\Gamma_f \cap \Delta$. Si $\Gamma_f \pitchfork \Delta$, alors le
nombre de Lefschetz est le nombre d'intersection :
$L(f) = I(\Gamma_f, \Delta)$.
\end{remarque}

% ============================================================
\section{Exercices}\label{sec:exo-transversalite}
% ============================================================

\begin{exercice}\label{exo:transverse-courbes-R2}
Dans $\R^2$, déterminer pour quelles valeurs de $a \in \R$ les courbes
$\set{y = x^2}$ et $\set{y = a}$ se coupent transversalement.
\end{exercice}

\begin{exercice}\label{exo:transverse-spheres}
Montrer que dans $\R^{n+1}$, les sphères
$S_r = \set{\abs{x} = r}$ et $S'_R = \set{\abs{x - a} = R}$
(avec $a \in \R^{n+1}$) se coupent transversalement si et seulement si
elles ne sont pas tangentes intérieurement ou extérieurement.
\end{exercice}

\begin{exercice}\label{exo:preimage-transverse-cercle}
Soit $f \colon \R^3 \to \R$, $(x, y, z) \mapsto x^2 + y^2 - z^2$.
\begin{enumerate}[label=\textup{(\alph*)}]
  \item Montrer que $f \pitchfork \set{1}$.
  \item En déduire que $f^{-1}(1)$ est une sous-variété de $\R^3$.
    Quelle est sa dimension ?
  \item Que se passe-t-il pour $f^{-1}(0)$ ?
\end{enumerate}
\end{exercice}

\begin{exercice}\label{exo:codimension-intersection}
Soient $V$ et $W$ deux sous-espaces vectoriels de $\R^n$ avec
$V \pitchfork W$ (au sens de sous-variétés de $\R^n$). Montrer que
$V + W = \R^n$ et en déduire
$\dim(V \cap W) = \dim V + \dim W - n$.
\end{exercice}

\begin{exercice}[Transversalité et submersions]\label{exo:submersion-transverse}
Montrer que si $f \colon M \to N$ est une submersion, alors $f$ est
transverse à toute sous-variété $W \subset N$.
\end{exercice}

\begin{exercice}\label{exo:intersection-tore}
Sur le tore $T^2$, considérer les courbes
$\alpha(t) = (e^{2\pi i t}, 1)$ et
$\beta(t) = (e^{2\pi i p t}, e^{2\pi i q t})$ avec $p, q \in \Z$
et $\gcd(p, q) = 1$. Montrer que $\alpha \pitchfork \beta$ et
calculer $I(\alpha, \beta)$.
\end{exercice}

\begin{exercice}[Densité de la transversalité]\label{exo:densite-transversalite}
Soit $f \colon S^1 \to \R^2$ un plongement lisse. Montrer qu'il existe
une droite $L \subset \R^2$ telle que $f(S^1) \pitchfork L$.
\textit{Indication : utiliser le théorème de Sard.}
\end{exercice}

\begin{exercice}\label{exo:nombre-intersection-degre}
Soit $f \colon M^n \to N^n$ une application lisse entre variétés
orientées compactes sans bord, et soit $q \in N$ une valeur régulière.
Montrer que $I(f, \set{q}) = \deg(f)$, le degré topologique de $f$.
\end{exercice}

\begin{exercice}[Auto-intersection de la diagonale]\label{exo:auto-intersection-diag}
Soit $M^n$ une variété compacte orientée. Considérons la diagonale
$\Delta \subset M \times M$.
\begin{enumerate}[label=\textup{(\alph*)}]
  \item Montrer que le fibré normal de $\Delta$ dans $M \times M$ est
    isomorphe au fibré tangent $TM$.
  \item En déduire que $I(\Delta, \Delta) = \chi(M)$, la
    caractéristique d'Euler de $M$.
\end{enumerate}
\end{exercice}

\begin{exercice}\label{exo:hairy-ball}
En utilisant le résultat de l'exercice~\ref{exo:auto-intersection-diag},
retrouver le fait que $S^{2n}$ ne possède pas de champ de vecteurs ne
s'annulant jamais. \textit{(Théorème de la boule chevelue.)\index{boule chevelue (théorème)}}
\end{exercice}

\begin{exercice}[Transversalité dans les espaces projectifs]%
\label{exo:transverse-projectif}
Dans $\CP^n$, montrer que deux sous-espaces projectifs
$\CP^k$ et $\CP^\ell$ (en position générale) se coupent transversalement.
Quelle est la dimension de l'intersection lorsque $k + \ell \geq n$ ?
\end{exercice}

\begin{exercice}[Application du théorème de Thom]\label{exo:application-thom}
Soit $M^m$ une variété compacte et $f \colon M \to \R^n$ une application
lisse avec $m < n$.
\begin{enumerate}[label=\textup{(\alph*)}]
  \item Montrer que pour presque tout $y \in \R^n$, $f^{-1}(y) = \varnothing$.
  \item En déduire que l'image $f(M)$ est de mesure nulle dans $\R^n$.
    \textit{(Retrouver ainsi une conséquence du théorème de Sard.)}
\end{enumerate}
\end{exercice}

% ============================================================
% ============================================================
%  CHAPITRE 6 : FORMES DIFFÉRENTIELLES ET INTÉGRATION
% ============================================================
% ============================================================
\chapter{Formes différentielles et intégration sur les variétés}%
\label{ch:formes-differentielles}
\minitoc\vspace{1cm}

Les formes différentielles fournissent le cadre naturel pour
l'intégration sur les variétés. Ce chapitre rappelle d'abord l'algèbre
extérieure, puis construit les formes différentielles et leur
différentielle extérieure. Nous introduisons ensuite le tiré en arrière,
le produit intérieur et la dérivée de Lie, avant de traiter l'orientation
et l'intégration.

% ============================================================
\section{Algèbre extérieure}\label{sec:algebre-exterieure}
% ============================================================

\begin{definition}[Puissance extérieure]\label{def:puissance-exterieure}%
\index{algèbre extérieure}\index{puissance extérieure}
Soit $V$ un espace vectoriel réel de dimension $n$. La $k$-ième
\textbf{puissance extérieure} $\Lambda^k V^*$ est l'espace vectoriel
des formes $k$-linéaires alternées sur $V$ :
\[
  \Lambda^k V^* = \set{\omega \colon \underbrace{V \times \cdots
    \times V}_{k} \to \R \;\Big|\;
    \omega \text{ est $k$-linéaire et alternée}}.
\]
Par convention, $\Lambda^0 V^* = \R$ et $\Lambda^k V^* = 0$ pour $k > n$.
\end{definition}

\begin{proposition}\label{prop:dim-exterieure}
$\dim \Lambda^k V^* = \binom{n}{k}$. Si $(e_1, \ldots, e_n)$ est une
base de $V$ et $(e^1, \ldots, e^n)$ la base duale, alors les éléments
\[
  e^{i_1} \wedge \cdots \wedge e^{i_k}, \quad
  1 \leq i_1 < \cdots < i_k \leq n,
\]
forment une base de $\Lambda^k V^*$.
\end{proposition}

\begin{definition}[Produit extérieur]\label{def:produit-exterieur}%
\index{produit extérieur}\index{$\wedge$}
Le \textbf{produit extérieur}
$\wedge \colon \Lambda^k V^* \times \Lambda^\ell V^* \to
  \Lambda^{k+\ell} V^*$
est défini par
\[
  (\alpha \wedge \beta)(v_1, \ldots, v_{k+\ell}) =
    \frac{1}{k!\,\ell!} \sum_{\sigma \in \mathfrak{S}_{k+\ell}}
    \operatorname{sgn}(\sigma)\, \alpha(v_{\sigma(1)}, \ldots,
    v_{\sigma(k)})\, \beta(v_{\sigma(k+1)}, \ldots, v_{\sigma(k+\ell)}).
\]
\end{definition}

\begin{proposition}[Propriétés du produit extérieur]%
\label{prop:proprietes-wedge}
Pour $\alpha \in \Lambda^k V^*$, $\beta \in \Lambda^\ell V^*$,
$\gamma \in \Lambda^m V^*$ :
\begin{enumerate}[label=\textup{(\roman*)}]
  \item \emph{Associativité :}
    $(\alpha \wedge \beta) \wedge \gamma =
      \alpha \wedge (\beta \wedge \gamma)$ ;
  \item \emph{Anti-commutativité graduée :}
    $\alpha \wedge \beta = (-1)^{k\ell}\, \beta \wedge \alpha$ ;
  \item \emph{Bilinéarité :} $\wedge$ est bilinéaire en chaque argument.
\end{enumerate}
\end{proposition}

\begin{definition}[Algèbre extérieure]\label{def:algebre-exterieure-totale}%
\index{algèbre extérieure!totale}
L'\textbf{algèbre extérieure} de $V^*$ est la somme directe graduée
\[
  \Lambda^* V^* = \bigoplus_{k=0}^{n} \Lambda^k V^*,
\]
munie du produit extérieur. C'est une algèbre associative graduée de
dimension $2^n$.
\end{definition}

\begin{exemple}\label{ex:algebre-ext-R3}
Pour $V = \R^3$ avec base duale $(e^1, e^2, e^3)$ :
\begin{itemize}
  \item $\Lambda^0(\R^3)^* \cong \R$ (dimension $1$) ;
  \item $\Lambda^1(\R^3)^*$ a pour base $\set{e^1, e^2, e^3}$
    (dimension $3$) ;
  \item $\Lambda^2(\R^3)^*$ a pour base
    $\set{e^1 \wedge e^2,\; e^1 \wedge e^3,\; e^2 \wedge e^3}$
    (dimension $3$) ;
  \item $\Lambda^3(\R^3)^*$ a pour base
    $\set{e^1 \wedge e^2 \wedge e^3}$ (dimension $1$).
\end{itemize}
On retrouve la dimension totale $1 + 3 + 3 + 1 = 8 = 2^3$.
L'isomorphisme $\Lambda^1 \cong \Lambda^2$ en dimension $3$ est à
l'origine du produit vectoriel et de la dualité de Hodge.
\end{exemple}

\begin{proposition}[Déterminant comme forme extérieure]%
\label{prop:determinant-exterieur}\index{déterminant}
Si $\dim V = n$, alors $\Lambda^n V^*$ est de dimension $1$. Pour
$\varphi \colon V \to V$ linéaire, l'application induite
$\Lambda^n \varphi^* \colon \Lambda^n V^* \to \Lambda^n V^*$ est la
multiplication par $\det(\varphi)$.
\end{proposition}

% ============================================================
\section{Formes différentielles}\label{sec:formes-differentielles}
% ============================================================

\begin{definition}[Forme différentielle de degré $k$]%
\label{def:forme-differentielle}\index{forme différentielle}
Soit $M$ une variété lisse de dimension $n$. Une \textbf{forme
différentielle de degré $k$} (ou $k$-forme) sur $M$ est une section
lisse du fibré $\Lambda^k T^*M$ :
\[
  \omega \in \Omega^k(M) = \Gamma^\infty(\Lambda^k T^*M).
\]
Autrement dit, $\omega$ associe à chaque point $p \in M$ un élément
$\omega_p \in \Lambda^k T_p^* M$, de manière lisse.
\end{definition}

\begin{remarque}\label{rem:0-formes}
$\Omega^0(M) = C^\infty(M)$ : les $0$-formes sont les fonctions lisses.
$\Omega^k(M) = 0$ pour $k > \dim M$.
\end{remarque}

\begin{proposition}[Expression locale]\label{prop:expression-locale-formes}%
\index{forme différentielle!expression locale}
Si $(U, \varphi = (x^1, \ldots, x^n))$ est une carte locale de $M$,
toute $k$-forme $\omega \in \Omega^k(M)$ s'écrit sur $U$ :
\[
  \omega\big|_U = \sum_{1 \leq i_1 < \cdots < i_k \leq n}
    \omega_{i_1 \cdots i_k}\, \dd x^{i_1} \wedge \cdots \wedge
    \dd x^{i_k},
\]
où les $\omega_{i_1 \cdots i_k} \in C^\infty(U)$ sont les
\textbf{fonctions composantes} de $\omega$.
\end{proposition}

\begin{exemple}\label{ex:formes-R3}
Sur $\R^3$ avec coordonnées $(x, y, z)$ :
\begin{itemize}
  \item $0$-forme : $f(x, y, z)$ ;
  \item $1$-forme : $\omega = P\,\dd x + Q\,\dd y + R\,\dd z$ ;
  \item $2$-forme :
    $\eta = A\,\dd y \wedge \dd z + B\,\dd z \wedge \dd x +
      C\,\dd x \wedge \dd y$ ;
  \item $3$-forme : $\mu = \rho\,\dd x \wedge \dd y \wedge \dd z$.
\end{itemize}
\end{exemple}

% ============================================================
\section{Différentielle extérieure}\label{sec:differentielle-exterieure}
% ============================================================

\begin{theoreme}[Différentielle extérieure]\label{thm:differentielle-exterieure}%
\index{différentielle extérieure}\index{$\dd$}
Il existe un unique opérateur linéaire
$\dd \colon \Omega^k(M) \to \Omega^{k+1}(M)$, pour chaque $k \geq 0$,
vérifiant les propriétés suivantes :
\begin{enumerate}[label=\textup{(\roman*)}]
  \item pour $f \in \Omega^0(M) = C^\infty(M)$, $\dd f$ est la
    différentielle usuelle ;
  \item \emph{Règle de Leibniz graduée :} pour
    $\alpha \in \Omega^k(M)$, $\beta \in \Omega^\ell(M)$,
    \[
      \dd(\alpha \wedge \beta) = \dd\alpha \wedge \beta +
        (-1)^k \alpha \wedge \dd\beta ;
    \]
  \item \emph{Nilpotence :} $\dd \circ \dd = 0$ ;
  \item $\dd$ est local : si $\omega\big|_U = \eta\big|_U$, alors
    $\dd\omega\big|_U = \dd\eta\big|_U$.
\end{enumerate}
\end{theoreme}

\begin{proof}
\emph{Unicité.}
Sur un ouvert de carte $(U, x^1, \ldots, x^n)$, les propriétés (i) et
(ii) imposent :
\[
  \dd\Bigl(\sum_I \omega_I\, \dd x^I\Bigr)
    = \sum_I \dd\omega_I \wedge \dd x^I
    = \sum_I \sum_{j=1}^n \frac{\partial \omega_I}{\partial x^j}
      \dd x^j \wedge \dd x^I,
\]
où nous utilisons la multi-notation $\dd x^I = \dd x^{i_1} \wedge \cdots
\wedge \dd x^{i_k}$. Ceci détermine $\dd$ de manière unique.

\emph{Existence.}
On définit $\dd$ par la formule ci-dessus en coordonnées locales. Il faut
vérifier que la définition est indépendante du choix de carte (ce qui
découle de la règle de changement de variables pour les dérivées
partielles) et que les propriétés (i)--(iv) sont satisfaites.

Vérifions (iii) : $\dd^2 = 0$. Il suffit de le vérifier sur les
générateurs. Pour $f \in C^\infty(M)$ :
\[
  \dd^2 f = \dd\Bigl(\sum_j \frac{\partial f}{\partial x^j}
    \dd x^j\Bigr) = \sum_{i,j} \frac{\partial^2 f}{\partial x^i
    \partial x^j} \dd x^i \wedge \dd x^j = 0,
\]
car $\frac{\partial^2 f}{\partial x^i \partial x^j}$ est symétrique en
$(i,j)$ tandis que $\dd x^i \wedge \dd x^j$ est antisymétrique. La
propriété s'étend aux $k$-formes par la règle de Leibniz.
\end{proof}

\begin{proposition}[Formule intrinsèque]\label{prop:formule-intrinseque-d}%
\index{différentielle extérieure!formule intrinsèque}
Pour $\omega \in \Omega^k(M)$ et des champs de vecteurs
$X_0, \ldots, X_k \in \mathfrak{X}(M)$ :
\begin{align*}
  \dd\omega(X_0, \ldots, X_k) &= \sum_{i=0}^k (-1)^i\, X_i\bigl(
    \omega(X_0, \ldots, \widehat{X_i}, \ldots, X_k)\bigr) \\
    &\quad + \sum_{i < j} (-1)^{i+j}\, \omega\bigl([X_i, X_j],
    X_0, \ldots, \widehat{X_i}, \ldots, \widehat{X_j}, \ldots, X_k\bigr),
\end{align*}
où $\widehat{X_i}$ signifie que le terme $X_i$ est omis.
\end{proposition}

\begin{exemple}\label{ex:d-1-forme}
Pour une $1$-forme $\omega$ et des champs $X$, $Y$ :
\[
  \dd\omega(X, Y) = X(\omega(Y)) - Y(\omega(X)) - \omega([X, Y]).
\]
\end{exemple}

\begin{exemple}[Lien avec l'analyse vectorielle]\label{ex:analyse-vectorielle}%
\index{analyse vectorielle}
Sur $\R^3$ :
\begin{itemize}
  \item $\dd f = \nabla f \cdot \dd\mathbf{x}$ (gradient) ;
  \item $\dd(P\,\dd x + Q\,\dd y + R\,\dd z)$ correspond au rotationnel ;
  \item $\dd(A\,\dd y \wedge \dd z + B\,\dd z \wedge \dd x +
    C\,\dd x \wedge \dd y) = \operatorname{div}(A,B,C)\,
    \dd x \wedge \dd y \wedge \dd z$.
\end{itemize}
L'identité $\dd^2 = 0$ encode simultanément
$\operatorname{rot}(\nabla f) = 0$ et
$\operatorname{div}(\operatorname{rot} \mathbf{F}) = 0$.
\end{exemple}

\subsection{Formes fermées et formes exactes}\label{subsec:fermees-exactes}

\begin{definition}[Formes fermées et exactes]\label{def:fermee-exacte}%
\index{forme fermée}\index{forme exacte}
Une $k$-forme $\omega \in \Omega^k(M)$ est dite :
\begin{itemize}
  \item \textbf{fermée} si $\dd\omega = 0$ ;
  \item \textbf{exacte} s'il existe $\eta \in \Omega^{k-1}(M)$ telle que
    $\omega = \dd\eta$.
\end{itemize}
Puisque $\dd^2 = 0$, toute forme exacte est fermée. La réciproque est
en général fausse.
\end{definition}

\begin{definition}[Cohomologie de de Rham]\label{def:cohomologie-de-rham}%
\index{cohomologie!de de Rham}\index{de Rham!cohomologie}
Le $k$-ième \textbf{groupe de cohomologie de de Rham} de $M$ est
\[
  H^k_{\mathrm{dR}}(M) =
    \frac{\ker(\dd \colon \Omega^k(M) \to \Omega^{k+1}(M))}
         {\im(\dd \colon \Omega^{k-1}(M) \to \Omega^k(M))}
    = \frac{Z^k(M)}{B^k(M)},
\]
où $Z^k(M)$ est l'espace des $k$-formes fermées et $B^k(M)$ celui des
$k$-formes exactes. C'est un espace vectoriel réel qui mesure
l'obstruction à ce que toute forme fermée soit exacte.
\end{definition}

\begin{theoreme}[Lemme de Poincaré]\label{thm:lemme-poincare}%
\index{Poincaré!lemme de}\index{lemme!de Poincaré}
Si $U \subset \R^n$ est un ouvert étoilé (ou plus généralement
contractile), alors $H^k_{\mathrm{dR}}(U) = 0$ pour tout $k \geq 1$.
Autrement dit, toute forme fermée sur $U$ est exacte.
\end{theoreme}

\begin{proof}[Idée de preuve]
On construit un \emph{opérateur d'homotopie}
$K \colon \Omega^k(U) \to \Omega^{k-1}(U)$ tel que
$\dd K + K \dd = \Id$ sur $\Omega^k(U)$ pour $k \geq 1$. Si $U$ est
étoilé par rapport à l'origine, on pose
\[
  (K\omega)_x(v_1, \ldots, v_{k-1}) = \int_0^1 t^{k-1}\,
    \omega_{tx}(x, v_1, \ldots, v_{k-1})\, \dd t.
\]
Alors pour $\omega$ fermée : $\omega = \dd(K\omega) + K(\dd\omega)
= \dd(K\omega)$.
\end{proof}

% ============================================================
\section{Tiré en arrière}\label{sec:tire-en-arriere}
% ============================================================

\begin{definition}[Tiré en arrière]\label{def:tire-en-arriere}%
\index{tiré en arrière}\index{$f^*$}
Soit $f \colon M \to N$ une application lisse. Le \textbf{tiré en arrière}
(ou \textit{pullback}) est l'application linéaire
$f^* \colon \Omega^k(N) \to \Omega^k(M)$ définie par
\[
  (f^*\omega)_p(v_1, \ldots, v_k) = \omega_{f(p)}(\dd f_p(v_1),
    \ldots, \dd f_p(v_k)),
\]
pour $p \in M$ et $v_1, \ldots, v_k \in T_p M$.
\end{definition}

\begin{theoreme}[Propriétés du tiré en arrière]%
\label{thm:proprietes-pullback}\index{tiré en arrière!propriétés}
Soient $f \colon M \to N$ et $g \colon N \to P$ des applications lisses.
\begin{enumerate}[label=\textup{(\roman*)}]
  \item \emph{Fonctorialité :}\index{fonctorialité}
    $(g \circ f)^* = f^* \circ g^*$ ;
  \item $(\Id_M)^* = \Id_{\Omega^*(M)}$ ;
  \item $f^*(\alpha \wedge \beta) = f^*\alpha \wedge f^*\beta$ ;
  \item \emph{Naturalité de $\dd$ :}\index{naturalité}
    $f^* \circ \dd = \dd \circ f^*$, c'est-à-dire le diagramme
    suivant commute pour tout $k$ :
    \[
    \begin{tikzcd}
      \Omega^k(N) \arrow[r, "\dd"] \arrow[d, "f^*"'] &
        \Omega^{k+1}(N) \arrow[d, "f^*"] \\
      \Omega^k(M) \arrow[r, "\dd"'] & \Omega^{k+1}(M)
    \end{tikzcd}
    \]
  \item pour $h \in C^\infty(N)$ : $f^*(h) = h \circ f$.
\end{enumerate}
\end{theoreme}

\begin{proof}
(i) Pour $\omega \in \Omega^k(P)$ et $v_1, \ldots, v_k \in T_p M$ :
\begin{align*}
  ((g \circ f)^*\omega)_p(v_1, \ldots, v_k)
    &= \omega_{g(f(p))}(\dd(g \circ f)_p(v_1), \ldots,
      \dd(g \circ f)_p(v_k)) \\
    &= \omega_{g(f(p))}(\dd g_{f(p)} \circ \dd f_p(v_1), \ldots,
      \dd g_{f(p)} \circ \dd f_p(v_k)) \\
    &= (g^*\omega)_{f(p)}(\dd f_p(v_1), \ldots, \dd f_p(v_k)) \\
    &= (f^*(g^*\omega))_p(v_1, \ldots, v_k).
\end{align*}

(iv) Il suffit de vérifier en coordonnées locales. Pour
$\omega = \sum_I \omega_I\, \dd y^I$ sur $N$ :
\[
  f^*(\dd\omega) = f^*\Bigl(\sum_I \dd\omega_I \wedge \dd y^I\Bigr)
    = \sum_I \dd(\omega_I \circ f) \wedge f^*(\dd y^I)
    = \dd\Bigl(\sum_I (\omega_I \circ f)\, f^*(\dd y^I)\Bigr)
    = \dd(f^*\omega),
\]
en utilisant $\dd^2 = 0$ et la règle de Leibniz.
\end{proof}

\begin{remarque}\label{rem:pullback-pas-pushforward}
Le tiré en arrière $f^*$ est défini pour \emph{toute} application lisse
$f$, pas seulement pour les difféomorphismes. Il va dans le sens
\og{}inverse\fg{}, ce qui en fait un foncteur contravariant. Il n'existe
en général pas d'analogue covariant (\og{}pushforward\fg{}) pour les
formes différentielles, sauf si $f$ est un difféomorphisme.
\end{remarque}

% ============================================================
\section{Produit intérieur et dérivée de Lie}\label{sec:produit-interieur-lie}
% ============================================================

\begin{definition}[Produit intérieur]\label{def:produit-interieur}%
\index{produit intérieur}\index{$\iota_X$}
Soit $X \in \mathfrak{X}(M)$ un champ de vecteurs. Le \textbf{produit
intérieur} (ou contraction) par $X$ est l'opérateur linéaire
$\iota_X \colon \Omega^k(M) \to \Omega^{k-1}(M)$ défini par
\[
  (\iota_X \omega)(Y_1, \ldots, Y_{k-1}) = \omega(X, Y_1, \ldots, Y_{k-1}).
\]
Pour $k = 0$, on pose $\iota_X f = 0$.
\end{definition}

\begin{proposition}[Propriétés du produit intérieur]%
\label{prop:proprietes-contraction}
\begin{enumerate}[label=\textup{(\roman*)}]
  \item $\iota_X$ est une \emph{anti-dérivation} de degré $-1$ :
    \[
      \iota_X(\alpha \wedge \beta) = (\iota_X \alpha) \wedge \beta +
        (-1)^k \alpha \wedge (\iota_X \beta),
      \quad \alpha \in \Omega^k(M) ;
    \]
  \item $\iota_X \circ \iota_X = 0$ ;
  \item $\iota_X \circ \iota_Y + \iota_Y \circ \iota_X = 0$ ;
  \item $\iota_{fX} \omega = f \cdot \iota_X \omega$ pour
    $f \in C^\infty(M)$.
\end{enumerate}
\end{proposition}

\begin{definition}[Dérivée de Lie des formes]\label{def:derivee-lie-formes}%
\index{dérivée de Lie!des formes différentielles}\index{$\mathcal{L}_X$}
Soit $X \in \mathfrak{X}(M)$ un champ de vecteurs et $\varphi_t$ son
flot. La \textbf{dérivée de Lie} d'une $k$-forme $\omega$ par rapport à
$X$ est
\[
  \mathcal{L}_X \omega = \frac{\dd}{\dd t}\bigg|_{t=0} \varphi_t^* \omega.
\]
\end{definition}

\begin{theoreme}[Formule de Cartan]\label{thm:formule-cartan}%
\index{formule de Cartan}\index{Cartan!formule magique}
Pour tout champ de vecteurs $X \in \mathfrak{X}(M)$ et toute forme
$\omega \in \Omega^k(M)$ :
\[
  \boxed{\mathcal{L}_X \omega = \dd(\iota_X \omega) +
    \iota_X(\dd\omega) = (\dd \circ \iota_X + \iota_X \circ \dd)\omega.}
\]
\end{theoreme}

\begin{proof}
On procède par récurrence sur le degré $k$.

\emph{Cas $k = 0$.} Pour $f \in C^\infty(M)$ :
$\mathcal{L}_X f = X(f) = \dd f(X) = \iota_X(\dd f)$,
et $\dd(\iota_X f) = \dd(0) = 0$. Donc
$\mathcal{L}_X f = \dd(\iota_X f) + \iota_X(\dd f)$.

\emph{Cas $k = 1$ pour les formes exactes.}
Pour $\omega = \dd f$ :
$\mathcal{L}_X(\dd f) = \dd(\mathcal{L}_X f) = \dd(X(f))$
(car $\mathcal{L}_X$ et $\dd$ commutent). D'autre part,
$\dd(\iota_X \dd f) + \iota_X(\dd^2 f) = \dd(X(f)) + 0$. D'accord.

\emph{Cas général.}
On utilise le fait que $\mathcal{L}_X$, $\dd \circ \iota_X + \iota_X
\circ \dd$, et les formes localement engendrées par des produits de
fonctions et de leurs différentielles. Les deux opérateurs sont des
dérivations de l'algèbre $\Omega^*(M)$ qui coïncident sur $\Omega^0(M)$
et sur $\dd(\Omega^0(M))$, donc ils coïncident partout puisque
$\Omega^*(M)$ est localement engendrée par ces éléments.
\end{proof}

\begin{corollaire}\label{cor:lie-proprietes}
\begin{enumerate}[label=\textup{(\roman*)}]
  \item $\mathcal{L}_X \circ \dd = \dd \circ \mathcal{L}_X$ ;
  \item $\mathcal{L}_X(\alpha \wedge \beta) =
    (\mathcal{L}_X \alpha) \wedge \beta +
    \alpha \wedge (\mathcal{L}_X \beta)$ ;
  \item $\mathcal{L}_{[X,Y]} = \mathcal{L}_X \circ \mathcal{L}_Y -
    \mathcal{L}_Y \circ \mathcal{L}_X$.
\end{enumerate}
\end{corollaire}

\begin{remarque}\label{rem:algebre-derivations}
Les opérateurs $\dd$, $\iota_X$, $\mathcal{L}_X$ satisfont les relations
d'une \emph{algèbre de Weil}\index{algèbre de Weil} :
\[
  [\dd, \dd] = 0, \quad [\iota_X, \iota_Y] = 0, \quad
  [\dd, \iota_X] = \mathcal{L}_X, \quad
  [\mathcal{L}_X, \dd] = 0, \quad
  [\mathcal{L}_X, \iota_Y] = \iota_{[X,Y]},
\]
où $[A, B] = A \circ B - (-1)^{|A||B|} B \circ A$ est le crochet gradué.
\end{remarque}

\begin{exemple}[Dérivée de Lie d'une $1$-forme]\label{ex:lie-1-forme}
Soit $X = \partial_x$ et $\omega = y\,\dd x + x^2\,\dd y$ sur $\R^2$.
Calculons $\mathcal{L}_X \omega$ via la formule de Cartan :
\begin{align*}
  \iota_X \omega &= y, \\
  \dd(\iota_X \omega) &= \dd y, \\
  \dd\omega &= \dd y \wedge \dd x + 2x\,\dd x \wedge \dd y
    = (2x - 1)\,\dd x \wedge \dd y, \\
  \iota_X(\dd\omega) &= (2x-1)\,\dd y.
\end{align*}
Donc $\mathcal{L}_X \omega = \dd y + (2x-1)\,\dd y = 2x\,\dd y$.
On peut vérifier : le flot de $X$ est $\varphi_t(x,y) = (x+t, y)$, donc
$\varphi_t^*\omega = y\,\dd x + (x+t)^2\,\dd y$ et
$\frac{\dd}{\dd t}\big|_{t=0} \varphi_t^*\omega = 2x\,\dd y$. Conforme.
\end{exemple}

\begin{proposition}[Formule de Cartan pour les fonctions]%
\label{prop:cartan-fonctions}
Pour $f \in C^\infty(M)$ et $X \in \mathfrak{X}(M)$ :
\[
  \mathcal{L}_X f = \iota_X(\dd f) = \dd f(X) = X(f).
\]
La dérivée de Lie d'une fonction est simplement la dérivée directionnelle.
\end{proposition}

\begin{proposition}[Formes fermées et dérivée de Lie]%
\label{prop:fermee-lie}
Soit $\omega$ une forme fermée ($\dd\omega = 0$). Alors
$\mathcal{L}_X \omega = \dd(\iota_X \omega)$ est exacte. Autrement dit,
la dérivée de Lie d'une forme fermée est exacte.
\end{proposition}

\begin{proof}
Par la formule de Cartan :
$\mathcal{L}_X \omega = \dd(\iota_X \omega) + \iota_X(\dd\omega)
= \dd(\iota_X \omega) + 0 = \dd(\iota_X \omega)$.
\end{proof}

% ============================================================
\section{Orientation}\label{sec:orientation}
% ============================================================

\begin{definition}[Orientation d'une variété]\label{def:orientation}%
\index{orientation}\index{variété!orientée}
Soit $M$ une variété lisse de dimension $n$. Une \textbf{orientation} de
$M$ est le choix d'une classe d'équivalence de $n$-formes ne s'annulant
en aucun point, où $\omega_1 \sim \omega_2$ si
$\omega_1 = f \cdot \omega_2$ avec $f > 0$ partout. Une variété est dite
\textbf{orientable} si une telle forme existe, et \textbf{orientée} si
un tel choix a été fait.
\end{definition}

\begin{proposition}[Critères d'orientabilité]\label{prop:criteres-orientabilite}%
\index{orientabilité!critères}
Pour une variété $M^n$ connexe, les conditions suivantes sont
équivalentes :
\begin{enumerate}[label=\textup{(\roman*)}]
  \item $M$ est orientable ;
  \item il existe un atlas $\set{(U_\alpha, \varphi_\alpha)}$ dont
    toutes les fonctions de transition ont un jacobien positif ;
  \item le fibré $\Lambda^n T^*M$ est trivialisable (isomorphe à
    $M \times \R$) ;
  \item il existe une $n$-forme $\omega \in \Omega^n(M)$ qui ne
    s'annule nulle part.
\end{enumerate}
\end{proposition}

\begin{exemple}\label{ex:orientabilite}
\begin{enumerate}[label=\textup{(\alph*)}]
  \item Toute variété simplement connexe est orientable (faux en
    général, mais vrai pour les surfaces).
  \item $\R^n$, $S^n$, $T^n$, $\GL^+(n, \R)$ sont orientables.
  \item $\RP^n$ est orientable si et seulement si $n$ est impair.
  \item La bouteille de Klein\index{bouteille de Klein} et le ruban de
    Möbius\index{Möbius!ruban de} ne sont pas orientables.
\end{enumerate}
\end{exemple}

% -------- TikZ : orientation --------
\begin{figure}[ht]
\centering
\begin{tikzpicture}[scale=1.0]
  % Surface orientée avec une normale
  \begin{scope}[shift={(-3.5,0)}]
    \draw[thick, fill=blue!10] (0,0) ellipse (1.5 and 0.8);
    \draw[->, thick, red!70!black] (0,0) -- (0,1.5)
      node[above]{\small $\mathbf{n}$};
    \draw[->, blue!60!black] (0,0) -- (1,0.2)
      node[right]{\small $e_1$};
    \draw[->, blue!60!black] (0,0) -- (-0.2,0.5)
      node[above left]{\small $e_2$};
    \node at (0,-1.5) {Surface orientée};
  \end{scope}
  % Ruban de Möbius (schéma)
  \begin{scope}[shift={(3.5,0)}]
    \draw[thick] (-1.5,-0.6) -- (1.5,0.6);
    \draw[thick] (-1.5,0.6) -- (1.5,-0.6);
    \draw[thick] (-1.5,-0.6) -- (-1.5,0.6);
    \draw[thick] (1.5,-0.6) -- (1.5,0.6);
    \draw[thick, ->] (-1.5,-0.3) -- (-1.5,0.3);
    \draw[thick, ->] (1.5,0.3) -- (1.5,-0.3);
    \node at (0,-1.5) {Ruban de Möbius (non orientable)};
  \end{scope}
\end{tikzpicture}
\caption{À gauche, une surface orientée avec un champ normal ;
  à droite, le domaine fondamental du ruban de Möbius
  (les flèches indiquent l'identification avec renversement).}
\label{fig:orientation}
\end{figure}

\begin{definition}[Orientation induite sur le bord]\label{def:orientation-bord}%
\index{orientation!du bord}
Soit $M$ une variété orientée à bord. L'\textbf{orientation induite} sur
$\partial M$ est définie par : une base
$(e_1, \ldots, e_{n-1})$ de $T_p(\partial M)$ est positivement orientée
si $(-\nu, e_1, \ldots, e_{n-1})$ est positivement orientée dans
$T_p M$, où $\nu$ est le vecteur normal unitaire sortant.
\end{definition}

\begin{exemple}[Orientation du bord d'un disque]\label{ex:orientation-bord-disque}
Le disque $D^2 = \set{(x,y) \in \R^2 \mid x^2 + y^2 \leq 1}$, muni de
l'orientation standard $\dd x \wedge \dd y$, a pour bord $\partial D^2 = S^1$.
Le vecteur normal sortant en $(x,y) \in S^1$ est $\nu = (x,y)$. L'orientation
induite sur $S^1$ est le sens trigonométrique (antihoraire).
\end{exemple}

\begin{exemple}[Orientation du bord d'un intervalle]\label{ex:orientation-bord-intervalle}
Pour $M = [a,b]$ orienté dans le sens croissant, $\partial M = \set{a, b}$.
Le vecteur sortant en $b$ est $+\partial_x$ (sortant vers la droite), et
en $a$ c'est $-\partial_x$ (sortant vers la gauche). L'orientation
induite donne $\varepsilon(b) = +1$ et $\varepsilon(a) = -1$. Ceci est
cohérent avec le théorème fondamental de l'analyse :
$\int_a^b f'(t)\,\dd t = f(b) - f(a)$.
\end{exemple}

\begin{lemme}\label{lem:forme-volume-orientation}%
\index{forme volume}
Soit $M^n$ une variété orientée et compacte sans bord. Pour toute
$n$-forme $\omega$ compatible avec l'orientation (i.e.\
$\omega_p$ est positive en chaque point), on a $\int_M \omega > 0$.
En particulier, $\omega$ n'est pas exacte.
\end{lemme}

\begin{proof}
La positivité de l'intégrale résulte du fait que dans chaque carte orientée,
$\omega$ s'écrit $f\,\dd x^1 \wedge \cdots \wedge \dd x^n$ avec $f > 0$,
et la partition de l'unité $\rho_\alpha \geq 0$. Si $\omega = \dd\eta$,
le théorème de Stokes donnerait $\int_M \omega = \int_{\partial M} \eta = 0$
(car $\partial M = \varnothing$), contradiction.
\end{proof}

% ============================================================
\section{Intégration des formes différentielles}%
\label{sec:integration-formes}
% ============================================================

\subsection{Intégration sur un ouvert de $\R^n$}%
\label{subsec:integration-Rn}

\begin{definition}\label{def:integration-Rn}%
\index{intégration!sur $\R^n$}
Soit $\omega = f(x)\, \dd x^1 \wedge \cdots \wedge \dd x^n$ une
$n$-forme à support compact sur un ouvert $U \subset \R^n$. On définit
\[
  \int_U \omega = \int_U f(x)\, \dd x^1 \cdots \dd x^n,
\]
où le membre de droite est l'intégrale de Lebesgue usuelle.
\end{definition}

\begin{lemme}[Changement de variables]\label{lem:changement-variables}%
\index{changement de variables}
Soit $\varphi \colon U \to V$ un difféomorphisme entre ouverts de $\R^n$
préservant l'orientation (i.e.\ $\det \dd\varphi > 0$). Alors pour toute
$n$-forme $\omega$ à support compact sur $V$ :
\[
  \int_V \omega = \int_U \varphi^*\omega.
\]
Si $\varphi$ renverse l'orientation, on a
$\int_V \omega = -\int_U \varphi^*\omega$.
\end{lemme}

\begin{proof}
Écrivons $\omega = g(y)\, \dd y^1 \wedge \cdots \wedge \dd y^n$.
Alors $\varphi^*\omega = g(\varphi(x))\, \det(\dd\varphi_x)\,
\dd x^1 \wedge \cdots \wedge \dd x^n$. Par la formule de changement de
variables de l'intégrale de Lebesgue :
\[
  \int_U \varphi^*\omega = \int_U g(\varphi(x))\, \det(\dd\varphi_x)\,
    \dd x^1 \cdots \dd x^n
    = \int_V g(y)\, \dd y^1 \cdots \dd y^n = \int_V \omega,
\]
à condition que $\det \dd\varphi > 0$. Si $\det \dd\varphi < 0$, un signe
$-$ apparaît.
\end{proof}

\subsection{Intégration sur une variété orientée}%
\label{subsec:integration-variete}

\begin{definition}[Intégrale d'une $n$-forme]\label{def:integrale-variete}%
\index{intégration!sur une variété}
Soit $M^n$ une variété orientée et $\omega \in \Omega^n(M)$ une
$n$-forme à support compact. Soit $\set{(U_\alpha, \varphi_\alpha)}$
un atlas orienté et $\set{\rho_\alpha}$ une partition de l'unité
subordonnée\index{partition de l'unité}. On définit
\[
  \int_M \omega = \sum_\alpha \int_{\varphi_\alpha(U_\alpha)}
    (\varphi_\alpha^{-1})^*(\rho_\alpha \omega).
\]
\end{definition}

\begin{theoreme}\label{thm:integrale-bien-definie}
L'intégrale $\int_M \omega$ est bien définie, c'est-à-dire indépendante
du choix d'atlas orienté et de partition de l'unité.
\end{theoreme}

\begin{proof}
Soient $\set{(U_\alpha, \varphi_\alpha), \rho_\alpha}$ et
$\set{(V_\beta, \psi_\beta), \sigma_\beta}$ deux choix. Alors
\begin{align*}
  \sum_\alpha \int_{\varphi_\alpha(U_\alpha)}
    (\varphi_\alpha^{-1})^*(\rho_\alpha \omega)
  &= \sum_\alpha \sum_\beta \int_{\varphi_\alpha(U_\alpha)}
    (\varphi_\alpha^{-1})^*(\rho_\alpha \sigma_\beta \omega) \\
  &= \sum_\alpha \sum_\beta \int_{\psi_\beta(V_\beta)}
    (\psi_\beta^{-1})^*(\rho_\alpha \sigma_\beta \omega) \\
  &= \sum_\beta \int_{\psi_\beta(V_\beta)}
    (\psi_\beta^{-1})^*(\sigma_\beta \omega),
\end{align*}
où la deuxième égalité utilise le lemme~\ref{lem:changement-variables}
avec le changement de cartes
$\psi_\beta \circ \varphi_\alpha^{-1}$ (qui préserve l'orientation
par hypothèse).
\end{proof}

\begin{proposition}[Propriétés de l'intégrale]\label{prop:proprietes-integrale}%
\index{intégration!propriétés}
\begin{enumerate}[label=\textup{(\roman*)}]
  \item \emph{Linéarité :}
    $\int_M (a\omega + b\eta) = a \int_M \omega + b \int_M \eta$ ;
  \item \emph{Changement d'orientation :} si $\overline{M}$ désigne $M$
    avec l'orientation opposée,
    $\int_{\overline{M}} \omega = -\int_M \omega$ ;
  \item \emph{Additivité :} si $M = M_1 \sqcup M_2$ (union disjointe),
    $\int_M \omega = \int_{M_1} \omega + \int_{M_2} \omega$ ;
  \item \emph{Fonctorialité :} si $f \colon M \to N$ est un
    difféomorphisme préservant l'orientation,
    $\int_N \omega = \int_M f^*\omega$.
\end{enumerate}
\end{proposition}

% -------- TikZ : forme différentielle --------
\begin{figure}[ht]
\centering
\begin{tikzpicture}[scale=1.0]
  % Variété (surface)
  \draw[thick, fill=blue!8] plot[smooth cycle, tension=0.7]
    coordinates {(-2.5,-0.5) (-1.5,1.2) (0.5,1.5) (2.2,0.8)
    (2.5,-0.5) (1.0,-1.2) (-1.2,-1.0)};
  \node at (2.8,0.0) {\small $M$};
  % Forme : petits parallélogrammes orientés
  \foreach \x/\y in {-1.0/0.3, 0.2/0.5, 1.2/0.1, -0.5/-0.4, 0.8/-0.5} {
    \draw[red!70!black, fill=red!20, opacity=0.6]
      (\x,\y) -- ++(0.3,0.1) -- ++(0.05,0.25) -- ++(-0.3,-0.1) -- cycle;
    \draw[->, red!70!black, thin] (\x+0.15,\y+0.15)
      arc[start angle=180, end angle=520, radius=0.08];
  }
  \node[red!70!black] at (0,-2.0) {$\omega \in \Omega^2(M)$ :
    mesure infinitésimale orientée};
\end{tikzpicture}
\caption{Visualisation d'une $2$-forme sur une surface : en chaque point,
  la forme associe une aire orientée à chaque paire de vecteurs tangents.}
\label{fig:forme-differentielle}
\end{figure}

% ============================================================
\section{Formes à support compact}\label{sec:support-compact}
% ============================================================

\begin{definition}[Support d'une forme]\label{def:support-forme}%
\index{support!d'une forme différentielle}
Le \textbf{support} d'une forme $\omega \in \Omega^k(M)$ est
\[
  \operatorname{supp}(\omega) = \overline{\set{p \in M \mid
    \omega_p \neq 0}}.
\]
On note $\Omega^k_c(M)$ l'espace des $k$-formes à support compact.
\end{definition}

\begin{remarque}\label{rem:support-compact-important}
L'hypothèse de support compact est essentielle pour que l'intégrale
converge sur les variétés non compactes. Si $M$ est compacte, toute
forme est automatiquement à support compact.
\end{remarque}

\begin{proposition}\label{prop:d-support-compact}
Si $\omega \in \Omega^k_c(M)$, alors $\dd\omega \in \Omega^{k+1}_c(M)$.
Autrement dit, $\dd$ restreint à $\Omega^*_c(M)$ donne un complexe de
chaînes
\[
  0 \to \Omega^0_c(M) \xrightarrow{\dd} \Omega^1_c(M) \xrightarrow{\dd}
    \cdots \xrightarrow{\dd} \Omega^n_c(M) \to 0.
\]
La cohomologie de ce complexe est la \textbf{cohomologie de de Rham à
support compact}\index{cohomologie!de de Rham à support compact}
$H^k_c(M)$.
\end{proposition}

\begin{proposition}[Tiré en arrière propre]\label{prop:pullback-propre}%
\index{application propre}
Si $f \colon M \to N$ est \emph{propre} (i.e.\ la préimage de tout
compact est compacte), alors
$f^* \colon \Omega^k_c(N) \to \Omega^k_c(M)$.
\end{proposition}

% ============================================================
\section{Vers le théorème de Stokes}\label{sec:vers-stokes}
% ============================================================

Nous énonçons ici le théorème de Stokes, qui sera démontré au chapitre
suivant.

\begin{theoreme}[Stokes]\label{thm:stokes-enonce}%
\index{théorème!de Stokes}\index{Stokes!théorème de}
Soit $M$ une variété orientée de dimension $n$ à bord $\partial M$ (muni
de l'orientation induite). Pour toute $(n-1)$-forme
$\omega \in \Omega^{n-1}_c(M)$ :
\[
  \boxed{\int_M \dd\omega = \int_{\partial M} \omega.}
\]
Si $\partial M = \varnothing$, alors $\int_M \dd\omega = 0$.
\end{theoreme}

\begin{remarque}\label{rem:stokes-cas-particuliers}
Le théorème de Stokes unifie :
\begin{itemize}
  \item le théorème fondamental de l'analyse : $\int_a^b f'(t)\, \dd t =
    f(b) - f(a)$ ;
  \item le théorème de Green dans le plan ;
  \item le théorème de la divergence (Gauss--Ostrogradski) dans $\R^3$ ;
  \item le théorème de Stokes classique pour les surfaces dans $\R^3$.
\end{itemize}
\end{remarque}

% ============================================================
\section{Exercices}\label{sec:exo-formes}
% ============================================================

\begin{exercice}[Calculs en coordonnées]\label{exo:calculs-coordonnees}
Sur $\R^3$, soient $\omega = xz\,\dd x + y^2\,\dd y - z\,\dd z$ et
$\eta = y\,\dd x \wedge \dd z + x\,\dd y \wedge \dd z$.
\begin{enumerate}[label=\textup{(\alph*)}]
  \item Calculer $\dd\omega$.
  \item Calculer $\dd\eta$.
  \item Vérifier que $\dd^2\omega = 0$.
  \item Calculer $\omega \wedge \dd\omega$.
\end{enumerate}
\end{exercice}

\begin{exercice}[Formes sur $S^2$]\label{exo:formes-S2}
Considérons $S^2 \subset \R^3$ munie des coordonnées sphériques
$(\theta, \phi)$. Soit $\omega = \sin\theta\, \dd\theta \wedge \dd\phi$.
\begin{enumerate}[label=\textup{(\alph*)}]
  \item Montrer que $\omega$ est une forme d'aire sur $S^2$.
  \item Calculer $\int_{S^2} \omega$.
  \item Montrer que $\omega$ n'est pas exacte.
\end{enumerate}
\end{exercice}

\begin{exercice}[Tiré en arrière]\label{exo:tire-en-arriere}
Soit $f \colon \R^2 \to \R^3$ définie par
$f(u,v) = (\cos u, \sin u, v)$.
\begin{enumerate}[label=\textup{(\alph*)}]
  \item Calculer $f^*(\dd x)$, $f^*(\dd y)$, $f^*(\dd z)$.
  \item Calculer $f^*(x\,\dd y \wedge \dd z)$.
  \item Montrer que $f^*(\dd x \wedge \dd y \wedge \dd z) = 0$.
    Interpréter géométriquement.
\end{enumerate}
\end{exercice}

\begin{exercice}[Formule de Cartan]\label{exo:formule-cartan}
Soit $X = x\,\partial_x + y\,\partial_y$ le champ radial sur $\R^2$ et
$\omega = x\,\dd y - y\,\dd x$.
\begin{enumerate}[label=\textup{(\alph*)}]
  \item Calculer $\iota_X \omega$.
  \item Calculer $\mathcal{L}_X \omega$ en utilisant la formule de Cartan.
  \item Vérifier en calculant $\mathcal{L}_X \omega$ directement via le
    flot de $X$.
\end{enumerate}
\end{exercice}

\begin{exercice}[Forme symplectique]\label{exo:forme-symplectique}%
\index{forme symplectique}
Sur $\R^{2n}$ avec coordonnées $(q^1, \ldots, q^n, p_1, \ldots, p_n)$,
on pose $\omega = \sum_{i=1}^n \dd q^i \wedge \dd p_i$.
\begin{enumerate}[label=\textup{(\alph*)}]
  \item Montrer que $\dd\omega = 0$.
  \item Montrer que $\omega^n = \underbrace{\omega \wedge \cdots
    \wedge \omega}_n = n!\, \dd q^1 \wedge \dd p_1 \wedge \cdots
    \wedge \dd q^n \wedge \dd p_n$.
  \item En déduire que $\omega^n \neq 0$ partout (forme volume).
  \item Montrer que $\omega$ n'est pas exacte sur $\R^{2n} \setminus
    \set{0}$ pour $n \geq 2$.
\end{enumerate}
\end{exercice}

\begin{exercice}[Orientation de $\RP^n$]\label{exo:orientation-RPn}
\begin{enumerate}[label=\textup{(\alph*)}]
  \item Montrer que $\RP^1 \cong S^1$ est orientable.
  \item Montrer que $\RP^2$ n'est pas orientable.
    \textit{Indication : considérer l'action de l'antipode sur les
    formes.}
  \item Plus généralement, montrer que $\RP^n$ est orientable si et
    seulement si $n$ est impair.
\end{enumerate}
\end{exercice}

\begin{exercice}[Intégration et degré]\label{exo:integration-degre}
Soit $f \colon S^1 \to S^1$ l'application $z \mapsto z^k$ (avec $k \in \Z$).
Soit $\omega = \dd\theta$ la $1$-forme angulaire standard.
\begin{enumerate}[label=\textup{(\alph*)}]
  \item Calculer $f^*\omega$.
  \item Calculer $\int_{S^1} f^*\omega$ et $\int_{S^1} \omega$.
  \item En déduire que $\deg(f) = k$.
\end{enumerate}
\end{exercice}

\begin{exercice}[$\dd^2 = 0$ et formes exactes]\label{exo:d2-exactes}
\begin{enumerate}[label=\textup{(\alph*)}]
  \item Soit $\omega = \frac{-y\,\dd x + x\,\dd y}{x^2 + y^2}$ sur
    $\R^2 \setminus \set{0}$. Montrer que $\dd\omega = 0$.
  \item Montrer que $\omega$ n'est pas exacte sur
    $\R^2 \setminus \set{0}$.
    \textit{Indication : calculer $\int_{S^1} \omega$.}
  \item Montrer que $\omega$ est exacte sur tout ouvert simplement
    connexe de $\R^2 \setminus \set{0}$.
\end{enumerate}
\end{exercice}

\begin{exercice}[Formes invariantes sur les groupes de Lie]%
\label{exo:formes-invariantes}\index{forme invariante}
Soit $G$ un groupe de Lie et $L_g \colon G \to G$, $h \mapsto gh$ la
translation à gauche.
\begin{enumerate}[label=\textup{(\alph*)}]
  \item Montrer que les formes invariantes à gauche (i.e.\
    $L_g^*\omega = \omega$ pour tout $g \in G$) forment une sous-algèbre
    de $(\Omega^*(G), \wedge)$ stable par $\dd$.
  \item Pour $G = \GL(n, \R)$, montrer que
    $\omega = \operatorname{tr}(A^{-1} \dd A)$ est une $1$-forme
    invariante à gauche.
  \item Calculer $\dd\omega$ pour $n = 1$ et $n = 2$.
\end{enumerate}
\end{exercice}

\begin{exercice}[Forme volume et métrique riemannienne]%
\label{exo:forme-volume}\index{forme volume}\index{métrique riemannienne}
Soit $(M^n, g)$ une variété riemannienne orientée.
\begin{enumerate}[label=\textup{(\alph*)}]
  \item Montrer qu'il existe une unique $n$-forme
    $\operatorname{vol}_g \in \Omega^n(M)$ telle que
    $\operatorname{vol}_g(e_1, \ldots, e_n) = 1$ pour toute base
    orthonormée directe $(e_1, \ldots, e_n)$.
  \item En coordonnées locales, montrer que
    $\operatorname{vol}_g = \sqrt{\det(g_{ij})}\, \dd x^1 \wedge
    \cdots \wedge \dd x^n$.
  \item Calculer $\operatorname{vol}_g$ pour $S^2$ munie de la métrique
    induite par $\R^3$.
\end{enumerate}
\end{exercice}

% === FIN DU CHUNK ch5-6 ===
% === DEBUT DU CHUNK ch7-8 ===

%%%%%%%%%%%%%%%%%%%%%%%%%%%%%%%%%%%%%%%%%%%%%%%%%%%%%%%%%%%%%%%%%%%%
%  CHAPITRE 7 : THÉORÈME DE STOKES
%%%%%%%%%%%%%%%%%%%%%%%%%%%%%%%%%%%%%%%%%%%%%%%%%%%%%%%%%%%%%%%%%%%%

\chapter{Théorème de Stokes}\label{ch:stokes}
\minitoc\bigskip

Le théorème de Stokes est l'un des résultats les plus profonds et les plus
élégants de la géométrie différentielle. Il unifie en une seule formule
les théorèmes classiques de Green, de la divergence et de Stokes en
$\R^3$, tout en ouvrant la voie à la cohomologie de de Rham. Ce chapitre
développe d'abord le formalisme des chaînes singulières lisses et de
l'intégration des formes différentielles, puis établit le théorème de Stokes
dans le cadre des variétés à bord compactes orientées. Les cas classiques en
analyse vectorielle apparaissent comme corollaires immédiats. La seconde moitié
du chapitre introduit la cohomologie de de Rham, le lemme de Poincaré, et
l'invariance homotopique, outils qui seront abondamment utilisés dans la suite
du cours.

% ============================================================
\section{Chaînes singulières lisses}\label{sec:chaines}
% ============================================================

\begin{definition}[Simplexe standard]\label{def:simplexe-standard}
\index{simplexe standard}
Le \emph{$k$-simplexe standard} est le sous-ensemble de $\R^{k+1}$ défini par
\[
  \Delta^k = \set{(t_0, t_1, \dots, t_k) \in \R^{k+1} \mid t_i \geq 0,\;
  \sum_{i=0}^{k} t_i = 1}.
\]
Ses \emph{sommets} sont les vecteurs $e_0, e_1, \dots, e_k$ de la base
canonique de $\R^{k+1}$.
\end{definition}

\begin{definition}[Applications de face]\label{def:face}
\index{application de face}
Pour $0 \leq i \leq k$, la \emph{$i$-ème application de face}
$\varepsilon_i^k \colon \Delta^{k-1} \to \Delta^k$ est l'unique
application affine envoyant $e_j$ sur $e_j$ si $j < i$ et sur $e_{j+1}$
si $j \geq i$. Explicitement,
\[
  \varepsilon_i^k(t_0, \dots, t_{k-1})
  = (t_0, \dots, t_{i-1}, 0, t_i, \dots, t_{k-1}).
\]
\end{definition}

\begin{definition}[Simplexe singulier lisse]\label{def:simplexe-singulier}
\index{simplexe singulier lisse}
Soit $M$ une variété lisse. Un \emph{$k$-simplexe singulier lisse} dans $M$
est une application lisse $\sigma \colon \Delta^k \to M$, où l'on entend
par «~lisse~» que $\sigma$ se prolonge en une application lisse sur un
voisinage ouvert de $\Delta^k$ dans le sous-espace affine engendré.
\end{definition}

\begin{definition}[Chaîne singulière lisse]\label{def:chaine-singuliere}
\index{chaîne singulière}
Le \emph{groupe des $k$-chaînes singulières lisses} de $M$, noté
$C_k^\infty(M)$, est le groupe abélien libre engendré par les $k$-simplexes
singuliers lisses dans $M$. Un élément $c \in C_k^\infty(M)$ s'écrit
\[
  c = \sum_{j=1}^{N} a_j \, \sigma_j, \qquad a_j \in \Z,
\]
où les $\sigma_j \colon \Delta^k \to M$ sont des $k$-simplexes singuliers
lisses.
\end{definition}

\begin{definition}[Opérateur de bord]\label{def:bord-chaines}
\index{opérateur de bord!chaînes singulières}
L'opérateur de bord $\partial_k \colon C_k^\infty(M) \to C_{k-1}^\infty(M)$
est défini sur les simplexes singuliers par
\[
  \partial_k \sigma = \sum_{i=0}^{k} (-1)^i \, \sigma \circ \varepsilon_i^k
\]
et étendu par linéarité à $C_k^\infty(M)$.
\end{definition}

\begin{proposition}\label{prop:d-carre-zero-chaines}
On a $\partial_{k-1} \circ \partial_k = 0$ pour tout $k \geq 1$.
\end{proposition}

\begin{proof}
Il suffit de vérifier sur un simplexe singulier $\sigma \colon \Delta^k \to M$.
On a
\[
  \partial_{k-1}(\partial_k \sigma)
  = \sum_{i=0}^{k} \sum_{j=0}^{k-1} (-1)^{i+j}
    \sigma \circ \varepsilon_i^k \circ \varepsilon_j^{k-1}.
\]
La relation $\varepsilon_i^k \circ \varepsilon_j^{k-1}
= \varepsilon_j^k \circ \varepsilon_{i-1}^{k-1}$ pour $j < i$ permet de
regrouper les termes par paires opposées, ce qui donne zéro.
\end{proof}

\begin{remarque}\label{rem:complexe-chaines}
On obtient ainsi un \emph{complexe de chaînes}
\[
  \cdots \xrightarrow{\partial_{k+1}} C_k^\infty(M)
  \xrightarrow{\partial_k} C_{k-1}^\infty(M)
  \xrightarrow{\partial_{k-1}} \cdots
  \xrightarrow{\partial_1} C_0^\infty(M) \to 0.
\]
L'homologie correspondante $H_k^\infty(M) = \ker \partial_k / \im \partial_{k+1}$
est isomorphe à l'homologie singulière usuelle $H_k(M;\Z)$ lorsque $M$ est lisse.
\end{remarque}

\begin{exemple}\label{ex:chaine-1}
\index{chaîne singulière!exemples}
Considérons $M = \R^2$ et le $1$-simplexe singulier $\sigma \colon \Delta^1 \to \R^2$
défini par $\sigma(t_0, t_1) = (1 - t_0) \cdot (0,0) + t_0 \cdot (1,1)$, avec
$t_0 + t_1 = 1$, soit en paramétrant par $t = t_0 \in [0,1]$, $\sigma(t) = (t,t)$.
Alors $\partial_1 \sigma = \sigma(e_1) - \sigma(e_0) = (1,1) - (0,0)$, ce qui
correspond à la différence des deux $0$-simplexes (points) aux extrémités du
segment.

Considérons maintenant le $2$-simplexe $\tau \colon \Delta^2 \to \R^2$ défini par
$\tau(t_0, t_1, t_2) = t_0 \cdot (0,0) + t_1 \cdot (1,0) + t_2 \cdot (0,1)$, qui
paramètre le triangle de sommets $(0,0)$, $(1,0)$, $(0,1)$. Son bord est
\[
  \partial_2 \tau = \tau \circ \varepsilon_0^2 - \tau \circ \varepsilon_1^2
  + \tau \circ \varepsilon_2^2,
\]
ce qui donne les trois arêtes du triangle avec les orientations correctes.
\end{exemple}

\begin{definition}[Cycle et bord]\label{def:cycle-bord}
\index{cycle}\index{bord!d'une chaîne}
Une $k$-chaîne $c \in C_k^\infty(M)$ est un \emph{$k$-cycle} si $\partial_k c = 0$.
Elle est un \emph{$k$-bord} s'il existe $b \in C_{k+1}^\infty(M)$ tel que
$c = \partial_{k+1} b$. L'ensemble des $k$-cycles est noté $Z_k^\infty(M)$ et
l'ensemble des $k$-bords $B_k^\infty(M)$.
\end{definition}

% ============================================================
\section{Intégration de formes sur des chaînes}\label{sec:integration-chaines}
% ============================================================

\begin{remarque}\label{rem:integration-orientation}
L'intégration des formes différentielles sur des chaînes dépend de l'orientation :
si l'on remplace $\sigma$ par $-\sigma$ (au sens de la chaîne), l'intégrale change
de signe. C'est l'analogue discret du fait que l'intégrale d'une $n$-forme sur une
variété orientée change de signe quand on renverse l'orientation.
\end{remarque}

\begin{definition}[Intégrale sur un simplexe singulier]\label{def:int-simplexe}
\index{intégrale!sur un simplexe singulier}
Soit $\omega$ une $k$-forme différentielle sur $M$ et
$\sigma \colon \Delta^k \to M$ un $k$-simplexe singulier lisse. On pose
\[
  \int_\sigma \omega = \int_{\Delta^k} \sigma^*\omega.
\]
L'intégrale de droite est l'intégrale usuelle de la $k$-forme $\sigma^*\omega$
sur le domaine compact $\Delta^k \subset \R^k$ (après identification du
sous-espace affine contenant $\Delta^k$ avec $\R^k$).
\end{definition}

\begin{definition}[Intégrale sur une chaîne]\label{def:int-chaine}
\index{intégrale!sur une chaîne}
Pour $c = \sum_j a_j \sigma_j \in C_k^\infty(M)$, on définit
\[
  \int_c \omega = \sum_j a_j \int_{\sigma_j} \omega.
\]
\end{definition}

\begin{proposition}[Stokes pour les chaînes]\label{prop:stokes-chaines}
\index{théorème de Stokes!pour les chaînes}
Soit $\omega \in \Omega^{k-1}(M)$ et $c \in C_k^\infty(M)$. Alors
\[
  \int_c \dd\omega = \int_{\partial c} \omega.
\]
\end{proposition}

\begin{proof}
Par linéarité, il suffit de traiter le cas $c = \sigma$, un $k$-simplexe
singulier. On a $\sigma^*(\dd\omega) = \dd(\sigma^*\omega)$. Par le
théorème de Stokes classique sur le domaine $\Delta^k$ (voir ci-dessous),
\[
  \int_{\Delta^k} \dd(\sigma^*\omega) = \int_{\partial\Delta^k} \sigma^*\omega
  = \sum_{i=0}^{k}(-1)^i \int_{\Delta^{k-1}}(\sigma\circ\varepsilon_i^k)^*\omega
  = \int_{\partial\sigma}\omega. \qedhere
\]
\end{proof}

% ============================================================
\section{Le théorème de Stokes}\label{sec:stokes-varietes}
% ============================================================

\begin{definition}[Variété à bord orientée]\label{def:variete-bord-orientee}
\index{variété à bord}
Soit $M$ une variété lisse de dimension $n$ à bord $\partial M$. On dit que
$M$ est \emph{orientée} si elle est munie d'une classe d'orientation, c'est-à-dire
d'un atlas maximal dont les changements de cartes ont un jacobien de déterminant
strictement positif. Le bord $\partial M$ hérite alors d'une \emph{orientation
induite}\index{orientation induite} : en un point $p \in \partial M$, si
$(v_1, \dots, v_{n-1})$ est un repère de $T_p(\partial M)$, on le déclare
positivement orienté si $(\nu, v_1, \dots, v_{n-1})$ est un repère positivement
orienté de $T_pM$, où $\nu$ est le vecteur normal sortant.
\end{definition}

\begin{theoreme}[Théorème de Stokes]\label{thm:stokes}
\index{théorème de Stokes}
Soit $M$ une variété lisse orientée de dimension $n$ à bord $\partial M$,
compacte. Pour toute $(n-1)$-forme différentielle $\omega \in \Omega^{n-1}(M)$,
on a
\[
  \boxed{\int_M \dd\omega = \int_{\partial M} \omega.}
\]
Si $\partial M = \varnothing$, le membre de droite est nul.
\end{theoreme}

\begin{proof}
La preuve se décompose en trois étapes.

\medskip\noindent\textbf{Étape 1 : Réduction locale via une partition de l'unité.}
\index{partition de l'unité!dans la preuve de Stokes}
Soit $\set{(U_\alpha, \varphi_\alpha)}_{\alpha \in A}$ un atlas fini (par
compacité) de cartes positives de $M$, compatible avec l'orientation, tel que
chaque $\varphi_\alpha(U_\alpha)$ soit soit un ouvert de $\R^n$, soit un ouvert
du demi-espace $\mathbb{H}^n = \set{x \in \R^n \mid x_n \geq 0}$. Soit
$\set{\rho_\alpha}$ une partition de l'unité subordonnée à $\set{U_\alpha}$.
On écrit
\[
  \omega = \sum_\alpha \rho_\alpha \, \omega,
\]
et par linéarité de l'intégrale et de $\dd$, il suffit de montrer
\[
  \int_M \dd(\rho_\alpha \omega) = \int_{\partial M} \rho_\alpha \omega
\]
pour chaque $\alpha$. Posons $\omega_\alpha = \rho_\alpha \omega$, qui est à
support compact dans $U_\alpha$.

\medskip\noindent\textbf{Étape 2 : Cas intérieur.}
Si $U_\alpha$ est une carte intérieure, i.e.\ $\varphi_\alpha(U_\alpha)$ est
un ouvert de $\R^n$, alors $\omega_\alpha$ est à support compact dans $\R^n$
et $\varphi_\alpha(U_\alpha) \cap \partial M = \varnothing$.

Écrivons, dans les coordonnées $(x_1, \dots, x_n)$,
\[
  (\varphi_\alpha^{-1})^*\omega_\alpha
  = \sum_{i=1}^{n} f_i \, \dd x_1 \wedge \cdots \wedge \widehat{\dd x_i}
    \wedge \cdots \wedge \dd x_n,
\]
où les $f_i$ sont à support compact. Alors
\[
  (\varphi_\alpha^{-1})^*\dd\omega_\alpha
  = \sum_{i=1}^{n} (-1)^{i-1} \frac{\partial f_i}{\partial x_i}
    \, \dd x_1 \wedge \cdots \wedge \dd x_n.
\]
Par le théorème de Fubini, chaque intégrale
$\int_{\R^n} \frac{\partial f_i}{\partial x_i}\, \dd x_1 \cdots \dd x_n$
se calcule en intégrant d'abord par rapport à $x_i$. Puisque $f_i$ est à
support compact, cette intégrale vaut zéro. Donc
$\int_M \dd\omega_\alpha = 0$, et $\int_{\partial M} \omega_\alpha = 0$
puisque $\operatorname{supp}(\omega_\alpha) \cap \partial M = \varnothing$.

\medskip\noindent\textbf{Étape 3 : Cas du bord.}
Si $U_\alpha$ est une carte de bord, $\varphi_\alpha(U_\alpha)$ est un ouvert
de $\mathbb{H}^n = \set{x \in \R^n \mid x_n \geq 0}$. On raisonne de manière
analogue. Pour $i < n$, les intégrales
$\int \frac{\partial f_i}{\partial x_i}$ s'annulent comme avant. Pour $i = n$,
on obtient
\[
  \int_{\mathbb{H}^n} \frac{\partial f_n}{\partial x_n}
  \, \dd x_1 \cdots \dd x_n
  = \int_{\R^{n-1}} \left[\int_0^\infty \frac{\partial f_n}{\partial x_n}
    \, \dd x_n \right] \dd x_1 \cdots \dd x_{n-1}.
\]
L'intégrale intérieure vaut $-f_n(x_1, \dots, x_{n-1}, 0)$ car $f_n$ est à
support compact et s'annule pour $x_n$ grand. En tenant compte des signes
$(-1)^{n-1}$ et de l'orientation induite sur $\partial M$ (correspondant à
$x_n = 0$ avec le vecteur normal sortant $-e_n$), on vérifie que
\[
  \int_M \dd\omega_\alpha
  = (-1)^{n-1} \cdot (-1) \int_{\R^{n-1}} f_n(x_1,\dots,x_{n-1},0)
    \, \dd x_1 \cdots \dd x_{n-1}
  = \int_{\partial M} \omega_\alpha.
\]

\medskip\noindent\textbf{Conclusion.}
En sommant sur $\alpha$, on obtient
\[
  \int_M \dd\omega = \sum_\alpha \int_M \dd\omega_\alpha
  = \sum_\alpha \int_{\partial M} \omega_\alpha
  = \int_{\partial M} \omega. \qedhere
\]
\end{proof}

\begin{remarque}\label{rem:stokes-sans-bord}
Lorsque $M$ est compacte sans bord, le théorème de Stokes donne
$\int_M \dd\omega = 0$ pour toute $(n-1)$-forme $\omega$. En particulier,
si $\int_M \eta \neq 0$ pour une $n$-forme $\eta$, alors $\eta$ n'est pas
exacte. Ceci fournit un outil fondamental pour détecter des classes de
cohomologie non triviales.
\end{remarque}

% ============================================================
\section{Cas classiques en $\R^2$ et $\R^3$}\label{sec:cas-classiques}
% ============================================================

\begin{corollaire}[Théorème de Green]\label{cor:green}
\index{théorème de Green}
Soit $D \subset \R^2$ un domaine compact à bord $\partial D$ lisse par morceaux,
orienté positivement (sens trigonométrique). Pour $P, Q \in C^\infty(\R^2)$,
\[
  \iint_D \left(\frac{\partial Q}{\partial x} - \frac{\partial P}{\partial y}\right)
  \dd x \, \dd y = \oint_{\partial D} P \, \dd x + Q \, \dd y.
\]
\end{corollaire}

\begin{proof}
On applique le théorème de Stokes à la $1$-forme $\omega = P \, \dd x + Q\, \dd y$
sur la variété à bord $M = D$. On a $\dd\omega = \left(\frac{\partial Q}{\partial x}
- \frac{\partial P}{\partial y}\right) \dd x \wedge \dd y$.
\end{proof}

\begin{corollaire}[Théorème de la divergence]\label{cor:divergence}
\index{théorème de la divergence}
Soit $\Omega \subset \R^3$ un domaine compact à bord lisse $\partial\Omega$,
orienté par la normale sortante. Pour un champ de vecteurs lisse
$F = (F_1, F_2, F_3)$,
\[
  \iiint_\Omega \operatorname{div}(F) \, \dd V
  = \iint_{\partial\Omega} F \cdot \nu \, \dd S,
\]
où $\nu$ est la normale unitaire sortante et $\dd S$ l'élément d'aire.
\end{corollaire}

\begin{proof}
On pose $\omega = F_1 \, \dd y \wedge \dd z + F_2 \, \dd z \wedge \dd x
+ F_3 \, \dd x \wedge \dd y$, qui est une $2$-forme sur $\R^3$. Alors
$\dd\omega = \operatorname{div}(F) \, \dd x \wedge \dd y \wedge \dd z$,
et l'intégrale de $\omega$ sur $\partial\Omega$ avec l'orientation induite
donne le flux $\iint_{\partial\Omega} F \cdot \nu \, \dd S$.
\end{proof}

\begin{corollaire}[Théorème de Stokes classique en $\R^3$]\label{cor:stokes-R3}
\index{théorème de Stokes!classique en $\R^3$}
Soit $S \subset \R^3$ une surface compacte orientée à bord $\partial S$
(courbe fermée orientée). Pour un champ de vecteurs lisse $F$,
\[
  \iint_S (\operatorname{rot} F) \cdot \nu \, \dd S
  = \oint_{\partial S} F \cdot \dd\ell.
\]
\end{corollaire}

\begin{proof}
On applique Stokes à la $1$-forme
$\omega = F_1\,\dd x + F_2\,\dd y + F_3\,\dd z$. Alors $\dd\omega$ est la
$2$-forme associée au rotationnel de $F$, et la formule en découle.
\end{proof}

\begin{exemple}\label{ex:green-aire}
\index{formule de l'aire}
Par le théorème de Green, l'aire d'un domaine $D \subset \R^2$ à bord lisse
$\gamma = \partial D$ s'écrit
\[
  \operatorname{Aire}(D) = \frac{1}{2} \oint_\gamma x\,\dd y - y\,\dd x.
\]
En effet, on prend $P = -y/2$ et $Q = x/2$, de sorte que
$\frac{\partial Q}{\partial x} - \frac{\partial P}{\partial y} = 1$. En
appliquant ceci au disque unité paramétré par $\gamma(t) = (\cos t, \sin t)$
pour $t \in [0,2\pi]$, on obtient
\[
  \operatorname{Aire}(D) = \frac{1}{2} \int_0^{2\pi}
  (\cos t \cdot \cos t + \sin t \cdot \sin t)\,\dd t
  = \frac{1}{2} \int_0^{2\pi} 1\,\dd t = \pi.
\]
\end{exemple}

\begin{exemple}\label{ex:flux-sphere}
\index{théorème de la divergence!exemple}
Considérons le champ $F(x,y,z) = (x,y,z)$ sur $\R^3$ et la boule unité
$B = \set{(x,y,z) \mid x^2 + y^2 + z^2 \leq 1}$. On a
$\operatorname{div}(F) = 3$, donc
\[
  \iiint_B 3\,\dd V = 3 \cdot \frac{4\pi}{3} = 4\pi.
\]
D'autre part, sur $S^2 = \partial B$, la normale sortante est $\nu = (x,y,z)$
et $F \cdot \nu = x^2 + y^2 + z^2 = 1$, d'où
\[
  \iint_{S^2} F \cdot \nu\,\dd S = \operatorname{Aire}(S^2) = 4\pi.
\]
Les deux résultats coïncident, conformément au théorème de la divergence.
\end{exemple}

\begin{remarque}\label{rem:stokes-physique}
\index{lois de conservation}
Les théorèmes de Green, de la divergence et de Stokes classique sont
fondamentaux en physique mathématique. Les équations de Maxwell, la
conservation de la masse en mécanique des fluides, et le théorème de Gauss
en électrostatique s'expriment naturellement en termes de ces théorèmes.
La formulation en termes de formes différentielles permet de les énoncer de
manière indépendante des coordonnées, ce qui est essentiel pour la
relativité générale et la théorie de jauge.
\end{remarque}

% ============================================================
\section{Cohomologie de de Rham}\label{sec:de-rham}
% ============================================================

L'une des applications les plus importantes du théorème de Stokes est la
construction de la cohomologie de de Rham\index{cohomologie de de Rham},
qui associe à toute variété lisse des invariants algébriques mesurant
l'obstruction à ce que les formes fermées soient exactes.

\begin{definition}[Complexe de de Rham]\label{def:complexe-de-rham}
\index{complexe de de Rham}
Le \emph{complexe de de Rham} de $M$ est la suite
\[
  0 \to \Omega^0(M) \xrightarrow{\dd} \Omega^1(M)
  \xrightarrow{\dd} \Omega^2(M) \xrightarrow{\dd} \cdots
  \xrightarrow{\dd} \Omega^n(M) \to 0.
\]
La relation $\dd \circ \dd = 0$ en fait un complexe de cochaînes.
\end{definition}

\begin{definition}[Groupes de cohomologie de de Rham]\label{def:cohomologie-de-rham}
\index{cohomologie de de Rham!groupes}
Le \emph{$k$-ème groupe de cohomologie de de Rham} de $M$ est le quotient
\[
  H^k_{\mathrm{dR}}(M)
  = \frac{\ker\bigl(\dd \colon \Omega^k(M) \to \Omega^{k+1}(M)\bigr)}
         {\im\bigl(\dd \colon \Omega^{k-1}(M) \to \Omega^k(M)\bigr)}
  = \frac{Z^k(M)}{B^k(M)},
\]
où $Z^k(M)$ désigne les \emph{$k$-formes fermées}\index{forme fermée}
($\dd\omega = 0$) et $B^k(M)$ les \emph{$k$-formes exactes}\index{forme exacte}
($\omega = \dd\eta$). La classe d'une forme fermée $\omega$ est notée $[\omega]$.
\end{definition}

\begin{proposition}\label{prop:H0}
Pour une variété lisse connexe $M$, on a $H^0_{\mathrm{dR}}(M) \cong \R$.
\end{proposition}

\begin{proof}
On a $\Omega^{-1}(M) = 0$, donc $B^0(M) = 0$, et $Z^0(M)$ est l'espace
des fonctions $f \in C^\infty(M)$ telles que $\dd f = 0$, c'est-à-dire les
fonctions localement constantes. Si $M$ est connexe, elles sont constantes,
d'où $H^0_{\mathrm{dR}}(M) = Z^0(M) \cong \R$.
\end{proof}

\begin{proposition}[Fonctorialité]\label{prop:fonctorialite-cohomologie}
\index{cohomologie de de Rham!fonctorialité}
Toute application lisse $f \colon M \to N$ induit un morphisme de $\R$-algèbres
graduées $f^* \colon H^*_{\mathrm{dR}}(N) \to H^*_{\mathrm{dR}}(M)$ par
$f^*[\omega] = [f^*\omega]$. De plus :
\begin{enumerate}[label=(\roman*)]
  \item $(\Id_M)^* = \Id$ ;
  \item $(g \circ f)^* = f^* \circ g^*$.
\end{enumerate}
\end{proposition}

\begin{proof}
La commutation $f^* \circ \dd = \dd \circ f^*$ assure que $f^*$ préserve les
formes fermées et les formes exactes, d'où un morphisme bien défini sur la
cohomologie. Les propriétés (i) et (ii) découlent des propriétés
correspondantes du pullback.
\end{proof}

\begin{exemple}\label{ex:H1-S1}
\index{cohomologie de de Rham!du cercle}
Calculons $H^1_{\mathrm{dR}}(S^1)$. Paramétrons $S^1$ par l'angle
$\theta \in [0, 2\pi]$. La $1$-forme $\dd\theta$ (définie globalement comme
forme angulaire) est fermée mais non exacte, car $\int_{S^1} \dd\theta = 2\pi
\neq 0$. Par le théorème de Stokes, une forme exacte $\dd f$ intégrée sur
$S^1$ donne zéro. Donc $[\dd\theta] \neq 0$ dans $H^1_{\mathrm{dR}}(S^1)$.
On peut montrer que $H^1_{\mathrm{dR}}(S^1) \cong \R$, engendré par
$[\dd\theta/(2\pi)]$.
\end{exemple}

\begin{proposition}[Produit en cohomologie]\label{prop:cup-product}
\index{produit cup!en cohomologie de de Rham}
Le produit extérieur induit un produit bien défini sur la cohomologie :
\[
  \smile \colon H^k_{\mathrm{dR}}(M) \times H^\ell_{\mathrm{dR}}(M)
  \to H^{k+\ell}_{\mathrm{dR}}(M), \qquad
  [\alpha] \smile [\beta] = [\alpha \wedge \beta].
\]
Ce produit est associatif, et graduellement commutatif :
$[\alpha] \smile [\beta] = (-1)^{k\ell} [\beta] \smile [\alpha]$.
Muni de ce produit, $H^*_{\mathrm{dR}}(M) = \bigoplus_{k=0}^n
H^k_{\mathrm{dR}}(M)$ est une $\R$-algèbre graduée commutative.
\end{proposition}

\begin{proof}
Le produit est bien défini car si $\alpha' = \alpha + \dd\eta$, alors
$\alpha' \wedge \beta = \alpha \wedge \beta + \dd(\eta \wedge \beta)$
(en utilisant $\dd\beta = 0$ et $\dd(\eta \wedge \beta)
= \dd\eta \wedge \beta + (-1)^{k-1}\eta \wedge \dd\beta
= \dd\eta \wedge \beta$). Les propriétés d'associativité et de commutativité
graduée se déduisent des propriétés correspondantes du produit extérieur.
\end{proof}

\begin{exemple}\label{ex:cohomologie-tore}
\index{cohomologie de de Rham!du tore}
Considérons le tore $T^2 = S^1 \times S^1$. Soient $\dd\theta_1$ et
$\dd\theta_2$ les formes angulaires sur les deux facteurs. Par la formule de
Künneth pour la cohomologie de de Rham, on a
\begin{align*}
  H^0_{\mathrm{dR}}(T^2) &\cong \R, \\
  H^1_{\mathrm{dR}}(T^2) &\cong \R^2,
    \quad \text{engendré par } [\dd\theta_1], [\dd\theta_2], \\
  H^2_{\mathrm{dR}}(T^2) &\cong \R,
    \quad \text{engendré par } [\dd\theta_1 \wedge \dd\theta_2].
\end{align*}
La forme volume $\dd\theta_1 \wedge \dd\theta_2$ n'est pas exacte car
$\int_{T^2} \dd\theta_1 \wedge \dd\theta_2 = (2\pi)^2 \neq 0$.
\end{exemple}

% ============================================================
\section{Lemme de Poincaré}\label{sec:poincare}
% ============================================================

\begin{definition}[Ouvert étoilé]\label{def:etoile}
\index{ouvert étoilé}
Un ouvert $U \subset \R^n$ est dit \emph{étoilé} par rapport au point
$p \in U$ si, pour tout $x \in U$, le segment $[p,x]$ est contenu dans $U$.
\end{definition}

\begin{theoreme}[Lemme de Poincaré]\label{thm:poincare}
\index{lemme de Poincaré}
Soit $U \subset \R^n$ un ouvert étoilé. Alors, pour tout $k \geq 1$,
\[
  H^k_{\mathrm{dR}}(U) = 0.
\]
Autrement dit, toute $k$-forme fermée sur $U$ est exacte.
\end{theoreme}

\begin{proof}
Quitte à translater, on peut supposer $U$ étoilé par rapport à l'origine.
On construit une \emph{homotopie de contraction}\index{homotopie de contraction}
$K \colon \Omega^k(U) \to \Omega^{k-1}(U)$ explicitement.

Pour $\omega \in \Omega^k(U)$, on écrit en coordonnées
\[
  \omega = \sum_{I} f_I(x) \, \dd x_{i_1} \wedge \cdots \wedge \dd x_{i_k},
\]
où la somme porte sur les multi-indices croissants $I = (i_1 < \cdots < i_k)$.
On définit l'opérateur $K$ par
\[
  (K\omega)(x)
  = \sum_{I} \sum_{j=1}^{k} (-1)^{j-1} x_{i_j}
    \left(\int_0^1 t^{k-1} f_I(tx) \, \dd t\right)
    \dd x_{i_1} \wedge \cdots \wedge \widehat{\dd x_{i_j}}
    \wedge \cdots \wedge \dd x_{i_k}.
\]

\noindent\textbf{Vérification de l'identité homotopique.}
Un calcul direct montre que
\begin{equation}\label{eq:homotopie-K}
  \dd(K\omega) + K(\dd\omega) = \omega
\end{equation}
pour tout $\omega \in \Omega^k(U)$ avec $k \geq 1$.

Vérifions-le. Considérons d'abord un monôme $\omega = f(x)\,\dd x_{i_1}
\wedge \cdots \wedge \dd x_{i_k}$. On a
\[
  K\omega = \sum_{j=1}^k (-1)^{j-1} x_{i_j}
  \left(\int_0^1 t^{k-1} f(tx)\,\dd t\right)
  \dd x_{i_1} \wedge \cdots \wedge \widehat{\dd x_{i_j}} \wedge \cdots
  \wedge \dd x_{i_k}.
\]
En calculant $\dd(K\omega)$, on fait apparaître deux types de termes :
ceux provenant de la dérivation de $x_{i_j} \int_0^1 t^{k-1}f(tx)\,\dd t$
par rapport à $x_{i_j}$ (qui contribuent par $\int_0^1 t^{k-1}f(tx)\,\dd t
+ x_{i_j}\int_0^1 t^{k-1}\frac{\partial f}{\partial x_{i_j}}(tx) \cdot t\,\dd t$)
et ceux provenant de la dérivation par rapport aux autres variables.

D'autre part, $K(\dd\omega)$ fait intervenir
$\dd\omega = \sum_\ell \frac{\partial f}{\partial x_\ell}\,\dd x_\ell
\wedge \dd x_{i_1} \wedge \cdots \wedge \dd x_{i_k}$.

En combinant et en utilisant la formule d'Euler pour les fonctions homogènes
(appliquée via le changement de variable $u = tx$),
\[
  \sum_{\ell=1}^n x_\ell \int_0^1 t^k \frac{\partial f}{\partial x_\ell}(tx)
  \,\dd t = \int_0^1 \frac{\dd}{\dd t}\bigl[t^k f(tx)\bigr]\,\dd t
  - k\int_0^1 t^{k-1}f(tx)\,\dd t,
\]
on obtient, après simplification,
\[
  \dd(K\omega) + K(\dd\omega) = f(x)\,\dd x_{i_1}\wedge\cdots\wedge\dd x_{i_k}
  = \omega.
\]

\noindent\textbf{Conclusion.}
Si $\omega$ est fermée ($\dd\omega = 0$), l'identité \eqref{eq:homotopie-K}
donne $\omega = \dd(K\omega)$, donc $\omega$ est exacte.
\end{proof}

\begin{corollaire}\label{cor:poincare-contractile}
Toute variété lisse difféomorphe à un ouvert étoilé de $\R^n$ a une
cohomologie de de Rham triviale en degrés $k \geq 1$ : $H^k_{\mathrm{dR}}(M) = 0$.
En particulier, $H^k_{\mathrm{dR}}(\R^n) = 0$ pour tout $k \geq 1$.
\end{corollaire}

\begin{remarque}\label{rem:poincare-geometrie}
\index{lemme de Poincaré!interprétation géométrique}
Le lemme de Poincaré admet l'interprétation géométrique suivante : sur un
ouvert étoilé, il n'y a pas de «~trous~» topologiques qui empêcheraient une
forme fermée d'être exacte. L'opérateur $K$ de la preuve réalise une
\emph{contraction} vers le centre de l'étoile, et la primitive $K\omega$
est obtenue en «~intégrant $\omega$ le long des rayons~» partant de l'origine.
\end{remarque}

\begin{exemple}\label{ex:poincare-calcul}
\index{lemme de Poincaré!calcul explicite}
Illustrons la construction sur un exemple concret. Soit
$\omega = 2xy\,\dd x + x^2\,\dd y$ sur $\R^2$ (étoilé par rapport à l'origine).
On vérifie que $\dd\omega = 0$ : en effet,
$\frac{\partial(x^2)}{\partial x} - \frac{\partial(2xy)}{\partial y} = 2x - 2x = 0$.

L'opérateur $K$ de la preuve donne, avec $f_1(x,y) = 2xy$ et
$f_2(x,y) = x^2$ :
\begin{align*}
  K\omega(x,y) &= x \int_0^1 f_1(tx,ty)\,\dd t
                 + y \int_0^1 f_2(tx,ty)\,\dd t \\
  &= x \int_0^1 2t^2 xy\,\dd t + y \int_0^1 t^2 x^2\,\dd t \\
  &= x \cdot \frac{2x y}{3} + y \cdot \frac{x^2}{3}
  = \frac{2x^2y}{3} + \frac{x^2 y}{3} = x^2 y.
\end{align*}
On vérifie : $\dd(x^2 y) = 2xy\,\dd x + x^2\,\dd y = \omega$.
\end{exemple}

% ============================================================
\section{Invariance homotopique}\label{sec:invariance-homotopique}
% ============================================================

\begin{definition}[Homotopie lisse]\label{def:homotopie-lisse}
\index{homotopie lisse}
Deux applications lisses $f_0, f_1 \colon M \to N$ sont \emph{lisses-homotopes}
s'il existe une application lisse $F \colon M \times [0,1] \to N$ telle que
$F(\cdot, 0) = f_0$ et $F(\cdot, 1) = f_1$.
\end{definition}

\begin{theoreme}[Invariance homotopique de la cohomologie]\label{thm:invariance-homotopique}
\index{cohomologie de de Rham!invariance homotopique}
Si $f_0, f_1 \colon M \to N$ sont lisses-homotopes, alors
$f_0^* = f_1^* \colon H^*_{\mathrm{dR}}(N) \to H^*_{\mathrm{dR}}(M)$.
\end{theoreme}

\begin{proof}
Soit $F \colon M \times [0,1] \to N$ l'homotopie. On construit un
\emph{opérateur d'homotopie de chaînes}\index{opérateur d'homotopie de chaînes}
$\mathcal{K} \colon \Omega^k(N) \to \Omega^{k-1}(M)$ par
\[
  \mathcal{K}\omega = \int_0^1 \iota_{\partial/\partial t}(F^*\omega)\,\dd t,
\]
où $\iota_{\partial/\partial t}$ désigne le produit intérieur par le champ
$\partial/\partial t$ sur $M \times [0,1]$, et l'intégrale est prise fibre
par fibre en $t$.

On vérifie par un calcul direct que cette construction satisfait
la \emph{formule magique de Cartan}\index{formule de Cartan} :
\[
  f_1^*\omega - f_0^*\omega = \dd(\mathcal{K}\omega) + \mathcal{K}(\dd\omega).
\]
En effet, en utilisant $\mathcal{L}_{\partial/\partial t} = \dd \circ
\iota_{\partial/\partial t} + \iota_{\partial/\partial t} \circ \dd$ (formule
de Cartan pour la dérivée de Lie), on obtient
\[
  F^*\omega\big|_{t=1} - F^*\omega\big|_{t=0}
  = \int_0^1 \mathcal{L}_{\partial/\partial t}(F^*\omega)\,\dd t
  = \int_0^1 \dd\bigl(\iota_{\partial/\partial t}F^*\omega\bigr)\,\dd t
  + \int_0^1 \iota_{\partial/\partial t}\bigl(\dd(F^*\omega)\bigr)\,\dd t.
\]
En restreignant à $M$ (identification de $M$ avec $M\times\{s\}$) et en
utilisant $F^*\circ\dd = \dd\circ F^*$ et la commutation de $\dd$ avec
l'intégration en $t$, on obtient la relation souhaitée.

Si $\omega$ est fermée, $\dd\omega = 0$ et
$f_1^*\omega - f_0^*\omega = \dd(\mathcal{K}\omega)$, donc
$[f_1^*\omega] = [f_0^*\omega]$.
\end{proof}

\begin{corollaire}\label{cor:equiv-homotopie-cohomologie}
\index{équivalence d'homotopie}
Si $f \colon M \to N$ est une équivalence d'homotopie lisse, alors
$f^* \colon H^*_{\mathrm{dR}}(N) \xrightarrow{\sim} H^*_{\mathrm{dR}}(M)$
est un isomorphisme.
\end{corollaire}

\begin{proof}
Soit $g \colon N \to M$ l'inverse homotopique de $f$. Alors $g \circ f
\simeq \Id_M$ et $f \circ g \simeq \Id_N$. Par le théorème,
$f^* \circ g^* = (g \circ f)^* = \Id_{H^*_{\mathrm{dR}}(M)}$ et
$g^* \circ f^* = \Id_{H^*_{\mathrm{dR}}(N)}$.
\end{proof}

\begin{corollaire}\label{cor:contractile-cohomologie}
\index{variété contractile}
Si $M$ est contractile (homotopiquement équivalente à un point), alors
$H^k_{\mathrm{dR}}(M) = 0$ pour tout $k \geq 1$.
\end{corollaire}

\begin{exemple}\label{ex:cohomologie-Rn}
Puisque $\R^n$ est contractile, on retrouve $H^k_{\mathrm{dR}}(\R^n) = 0$
pour $k \geq 1$, ce qui généralise le lemme de Poincaré.
\end{exemple}

\begin{exemple}\label{ex:cohomologie-Sn}
\index{cohomologie de de Rham!des sphères}
Calculons $H^*_{\mathrm{dR}}(S^n)$ pour $n \geq 1$ par la suite de
Mayer--Vietoris\index{Mayer--Vietoris}. Recouvrons $S^n$ par les ouverts
$U = S^n \setminus \{N\}$ (pôle nord retiré) et $V = S^n \setminus \{S\}$
(pôle sud retiré). Alors $U$ et $V$ sont difféomorphes à $\R^n$ (par
projection stéréographique) et $U \cap V$ est homotopiquement équivalent
à $S^{n-1}$. La suite exacte de Mayer--Vietoris et une récurrence sur $n$
donnent
\[
  H^k_{\mathrm{dR}}(S^n) \cong
  \begin{cases}
    \R & \text{si } k = 0 \text{ ou } k = n, \\
    0 & \text{sinon.}
  \end{cases}
\]
\end{exemple}

\begin{remarque}\label{rem:equivalence-homotopie-applications}
L'invariance homotopique de la cohomologie de de Rham a des conséquences
frappantes. Par exemple, puisque $\R^n \setminus \{0\}$ se rétracte par
déformation sur $S^{n-1}$, on a
$H^k_{\mathrm{dR}}(\R^n \setminus \{0\}) \cong H^k_{\mathrm{dR}}(S^{n-1})$.
De même, le cylindre $S^1 \times \R$ est homotope à $S^1$, donc
$H^1_{\mathrm{dR}}(S^1 \times \R) \cong \R$.
\end{remarque}

% ============================================================
\section{Degré d'une application via l'intégration}\label{sec:degre}
% ============================================================

\begin{definition}[Degré d'une application]\label{def:degre-integration}
\index{degré!via l'intégration}
Soit $f \colon M \to N$ une application lisse entre variétés compactes
orientées sans bord de même dimension $n$. Le \emph{degré de $f$} est
l'unique entier $\deg(f) \in \Z$ tel que, pour toute $n$-forme $\omega$
sur $N$,
\[
  \int_M f^*\omega = \deg(f) \int_N \omega.
\]
\end{definition}

\begin{proposition}\label{prop:degre-bien-defini}
Le degré est bien défini et ne dépend que de la classe d'homotopie de $f$.
\end{proposition}

\begin{proof}
L'application linéaire $f^* \colon H^n_{\mathrm{dR}}(N) \to
H^n_{\mathrm{dR}}(M)$ est un morphisme entre deux espaces de dimension $1$
(car $M$ et $N$ sont compactes, orientées, connexes et sans bord), donc
est la multiplication par un scalaire. L'intégrité de ce scalaire sera
montrée au chapitre~10 en utilisant le théorème de Sard. L'invariance
homotopique découle du 
ef{thm:invariance-homotopique}.
\end{proof}

\begin{remarque}\label{rem:degre-valeur-reguliere}
Si $q \in N$ est une valeur régulière de $f$, alors $\deg(f) = \sum_{p \in
f^{-1}(q)} \operatorname{sgn}\det(\dd f_p)$, somme des signes des
déterminants jacobiens aux préimages. Ce lien sera développé au chapitre~10.
\end{remarque}

\begin{exemple}\label{ex:degre-lacet}
\index{degré!exemples}
Considérons l'application $f_n \colon S^1 \to S^1$ définie par
$f_n(e^{i\theta}) = e^{in\theta}$ pour $n \in \Z$. La forme volume
$\omega = \dd\theta$ sur $S^1$ vérifie $\int_{S^1} \omega = 2\pi$, et
\[
  \int_{S^1} f_n^*\omega = \int_{S^1} n\,\dd\theta = 2\pi n.
\]
Donc $\deg(f_n) = n$, conformément à l'intuition que $f_n$ «~enroule $S^1$
autour de lui-même $n$ fois~».
\end{exemple}

\begin{proposition}[Propriétés du degré]\label{prop:proprietes-degre}
\index{degré!propriétés}
Le degré possède les propriétés suivantes :
\begin{enumerate}[label=(\roman*)]
  \item $\deg(\Id_M) = 1$ ;
  \item $\deg(g \circ f) = \deg(g) \cdot \deg(f)$ ;
  \item si $f$ est un difféomorphisme préservant l'orientation, $\deg(f) = 1$ ;
  si $f$ renverse l'orientation, $\deg(f) = -1$ ;
  \item si $f$ n'est pas surjective, $\deg(f) = 0$.
\end{enumerate}
\end{proposition}

\begin{proof}
(i) est immédiat car $\Id^*\omega = \omega$.
(ii) découle de $(g\circ f)^* = f^* \circ g^*$.
(iii) Un difféomorphisme bijectif a exactement une préimage en chaque point,
et le signe du jacobien donne $\pm 1$.
(iv) Si $q \notin f(M)$, on choisit $\omega$ à support dans un voisinage de
$q$ ne rencontrant pas $f(M)$, de sorte que $f^*\omega = 0$, d'où
$\deg(f) \int_N \omega = 0$ avec $\int_N \omega \neq 0$.
\end{proof}

% ============================================================
\section{Dualité de Poincaré}\label{sec:dualite-poincare}
% ============================================================

\begin{theoreme}[Dualité de Poincaré --- énoncé]\label{thm:dualite-poincare}
\index{dualité de Poincaré}
Soit $M$ une variété lisse compacte orientée sans bord de dimension $n$.
L'accouplement bilinéaire
\[
  H^k_{\mathrm{dR}}(M) \times H^{n-k}_{\mathrm{dR}}(M) \to \R,
  \qquad ([\alpha], [\beta]) \mapsto \int_M \alpha \wedge \beta
\]
est non dégénéré. Autrement dit, il induit un isomorphisme
\[
  H^k_{\mathrm{dR}}(M) \cong \bigl(H^{n-k}_{\mathrm{dR}}(M)\bigr)^*.
\]
\end{theoreme}

\begin{remarque}\label{rem:dualite-poincare}
La preuve complète de ce théorème requiert la suite exacte de
Mayer--Vietoris\index{Mayer--Vietoris} et un argument de récurrence sur
les recouvrements (voir, par exemple, Bott et Tu,
\textit{Differential Forms in Algebraic Topology}). L'idée essentielle est
que le produit extérieur $\alpha \wedge \beta$ donne une $n$-forme que l'on
intègre sur $M$, et le théorème de Stokes garantit que cet accouplement est
bien défini sur la cohomologie.
\end{remarque}

\begin{corollaire}\label{cor:betti-dualite}
\index{nombres de Betti}
Les nombres de Betti $b_k(M) = \dim H^k_{\mathrm{dR}}(M)$ satisfont
$b_k = b_{n-k}$ pour toute variété compacte orientée sans bord de dimension $n$.
\end{corollaire}

% ============================================================
\section{Exercices}\label{sec:exos-stokes}
% ============================================================

\begin{exercice}\label{exo:stokes-1}
Soit $\omega = x\,\dd y \wedge \dd z + y\,\dd z \wedge \dd x + z\,\dd x
\wedge \dd y$ la $2$-forme sur $\R^3$. Calculer $\int_{S^2} \omega$ de deux
manières :
\begin{enumerate}[label=(\alph*)]
  \item directement par paramétrisation ;
  \item en utilisant le théorème de la divergence (identifier le champ
  de vecteurs correspondant et son flux).
\end{enumerate}
\end{exercice}

\begin{exercice}\label{exo:stokes-2}
Vérifier le théorème de Stokes pour $\omega = xy\,\dd x + y^2\,\dd y$ sur le
triangle $T$ de sommets $(0,0)$, $(1,0)$ et $(0,1)$, en calculant séparément
$\int_T \dd\omega$ et $\int_{\partial T} \omega$.
\end{exercice}

\begin{exercice}\label{exo:stokes-3}
Soit $M$ une variété compacte orientée sans bord de dimension $n$. Montrer que
si $\alpha \in \Omega^k(M)$ et $\beta \in \Omega^{n-k-1}(M)$, alors
\[
  \int_M \dd\alpha \wedge \beta = (-1)^{k+1} \int_M \alpha \wedge \dd\beta.
\]
\end{exercice}

\begin{exercice}\label{exo:stokes-4}
\index{forme angulaire}
Soit $\omega = \frac{-y\,\dd x + x\,\dd y}{x^2 + y^2}$ la forme angulaire
sur $\R^2 \setminus \{0\}$.
\begin{enumerate}[label=(\alph*)]
  \item Vérifier que $\dd\omega = 0$.
  \item Calculer $\int_\gamma \omega$ où $\gamma$ est le cercle unité orienté
  dans le sens trigonométrique.
  \item En déduire que $[\omega] \neq 0$ dans
  $H^1_{\mathrm{dR}}(\R^2 \setminus \{0\})$.
\end{enumerate}
\end{exercice}

\begin{exercice}\label{exo:stokes-5}
Montrer que $H^1_{\mathrm{dR}}(\R^2 \setminus \{0\}) \cong \R$ en utilisant
l'homotopie avec $S^1$ et l'invariance homotopique de la cohomologie de
de Rham.
\end{exercice}

\begin{exercice}\label{exo:stokes-6}
\index{lemme de Poincaré!constructif}
Soit $\omega = e^{xyz}(yz\,\dd x + xz\,\dd y + xy\,\dd z)$ sur $\R^3$.
\begin{enumerate}[label=(\alph*)]
  \item Vérifier que $\omega$ est fermée.
  \item Utiliser la construction de la preuve du lemme de Poincaré pour
  trouver explicitement une fonction $f$ telle que $\omega = \dd f$.
\end{enumerate}
\end{exercice}

\begin{exercice}\label{exo:stokes-7}
Soient $M$ et $N$ deux variétés compactes connexes orientées sans bord de même
dimension. Montrer que si $f \colon M \to N$ n'est pas surjective, alors
$\deg(f) = 0$.
\textit{Indication : considérer le support d'une forme volume sur $N$ concentrée
au voisinage d'un point non atteint.}
\end{exercice}

\begin{exercice}\label{exo:stokes-8}
Soit $f \colon S^n \to S^n$ l'application antipodale $f(x) = -x$.
\begin{enumerate}[label=(\alph*)]
  \item Calculer $f^*\omega$ pour la forme volume standard $\omega$ de $S^n$.
  \item En déduire que $\deg(f) = (-1)^{n+1}$.
  \item Montrer que si $n$ est pair, $f$ n'est pas homotope à l'identité.
\end{enumerate}
\end{exercice}

\begin{exercice}\label{exo:stokes-9}
\index{suite exacte de Mayer--Vietoris}
Soit $M = U \cup V$ un recouvrement ouvert d'une variété lisse. Construire la
suite exacte longue de Mayer--Vietoris
\[
  \cdots \to H^{k-1}_{\mathrm{dR}}(U \cap V) \xrightarrow{\delta}
  H^k_{\mathrm{dR}}(M) \to H^k_{\mathrm{dR}}(U) \oplus H^k_{\mathrm{dR}}(V)
  \to H^k_{\mathrm{dR}}(U \cap V) \to \cdots
\]
et l'utiliser pour calculer $H^*_{\mathrm{dR}}(S^n)$ par récurrence sur $n$.
\end{exercice}

\begin{exercice}\label{exo:stokes-10}
Montrer que la caractéristique d'Euler--Poincaré $\chi(M) = \sum_{k=0}^n
(-1)^k b_k(M)$ est un invariant homotopique. Calculer $\chi(S^n)$,
$\chi(\R\mathrm{P}^2)$ et $\chi(T^2)$, où $T^2 = S^1 \times S^1$ est le tore.
\end{exercice}

\begin{exercice}\label{exo:stokes-11}
\index{formule de Künneth}
Soient $M$ et $N$ deux variétés compactes. Montrer que l'application
\[
  \bigoplus_{p+q=k} H^p_{\mathrm{dR}}(M) \otimes H^q_{\mathrm{dR}}(N)
  \to H^k_{\mathrm{dR}}(M \times N), \qquad
  [\alpha] \otimes [\beta] \mapsto [\pi_M^*\alpha \wedge \pi_N^*\beta]
\]
est un isomorphisme (formule de Künneth). En déduire
$H^*_{\mathrm{dR}}(T^n)$ pour le tore $T^n = (S^1)^n$.
\textit{Indication : utiliser le fait que $H^*_{\mathrm{dR}}(S^1) \cong
\R \oplus \R$ (en degrés $0$ et $1$).}
\end{exercice}

\begin{exercice}\label{exo:stokes-12}
\index{cohomologie de de Rham!à support compact}
Soit $\omega$ la $2$-forme sur $S^2$ définie comme la forme volume standard
(forme d'aire) normalisée à $\int_{S^2} \omega = 4\pi$.
\begin{enumerate}[label=(\alph*)]
  \item Montrer que $\omega$ n'est pas exacte.
  \item Montrer que $[\omega]$ engendre $H^2_{\mathrm{dR}}(S^2) \cong \R$.
  \item Soit $f \colon S^2 \to S^2$ lisse. Montrer que $\int_{S^2} f^*\omega
  = 4\pi \cdot \deg(f)$.
\end{enumerate}
\end{exercice}

%%%%%%%%%%%%%%%%%%%%%%%%%%%%%%%%%%%%%%%%%%%%%%%%%%%%%%%%%%%%%%%%%%%%
%  CHAPITRE 8 : THÉORÈME DE SARD
%%%%%%%%%%%%%%%%%%%%%%%%%%%%%%%%%%%%%%%%%%%%%%%%%%%%%%%%%%%%%%%%%%%%

\chapter{Théorème de Sard}\label{ch:sard}
\minitoc\bigskip

Le théorème de Sard affirme que les valeurs critiques d'une application lisse
sont «~rares~» au sens de la mesure de Lebesgue. Ce résultat, d'apparence
technique, est en réalité l'un des outils les plus puissants de la topologie
différentielle, comme en témoignent les applications spectaculaires qui
clôturent ce chapitre. La preuve repose sur une analyse fine utilisant le
développement de Taylor et un argument de comptage volumétrique. Les
applications incluent le théorème du point fixe de Brouwer, le théorème de
non-rétraction et l'existence d'applications transverses.

% ============================================================
\section{Valeurs critiques et valeurs régulières}\label{sec:critique-reguliere}
% ============================================================

\begin{definition}[Points et valeurs critiques]\label{def:critique}
\index{point critique}\index{valeur critique}\index{valeur régulière}
Soit $f \colon M^m \to N^n$ une application lisse entre variétés.
\begin{enumerate}[label=(\roman*)]
  \item Un point $p \in M$ est un \emph{point critique} de $f$ si l'application
  linéaire tangente $\dd f_p \colon T_pM \to T_{f(p)}N$ n'est pas surjective,
  c'est-à-dire $\rg(\dd f_p) < n$.
  \item L'ensemble des points critiques est noté $C_f \subset M$.
  \item Un point $q \in N$ est une \emph{valeur critique} si
  $q \in f(C_f)$.
  \item Un point $q \in N$ est une \emph{valeur régulière} si $q \notin f(C_f)$.
  En particulier, tout $q \notin f(M)$ est une valeur régulière.
\end{enumerate}
\end{definition}

\begin{remarque}\label{rem:reguliere-preimage}
Si $q$ est une valeur régulière de $f$ et $f^{-1}(q) \neq \varnothing$, alors
par le théorème des fonctions implicites (ou le théorème de la préimage
régulière), $f^{-1}(q)$ est une sous-variété lisse de $M$ de dimension $m - n$.
\end{remarque}

\begin{exemple}\label{ex:critique-projection}
\index{point critique!exemples}
Considérons la projection $\pi \colon S^2 \to \R$ définie par
$\pi(x,y,z) = z$ (la «~fonction hauteur~»). Ici $M = S^2$ (dimension $2$)
et $N = \R$ (dimension $1$). Les points critiques sont ceux où
$\dd\pi_p \colon T_pS^2 \to T_{\pi(p)}\R \cong \R$ est l'application nulle,
c'est-à-dire les points où le plan tangent à $S^2$ est horizontal. Ce sont
le pôle nord $(0,0,1)$ et le pôle sud $(0,0,-1)$, et les valeurs critiques
sont $\{-1, 1\}$. Pour tout $t \in (-1,1)$, la préimage $\pi^{-1}(t)$ est
le cercle $\{x^2 + y^2 = 1 - t^2,\, z = t\}$, qui est bien une sous-variété
de dimension $1$.
\end{exemple}

\begin{exemple}\label{ex:critique-determinant}
\index{point critique!application déterminant}
Considérons $f \colon M_{n}(\R) \to \R$ définie par $f(A) = \det(A)$. Ici
$M = M_n(\R) \cong \R^{n^2}$, $N = \R$. On montre que $\dd f_A(H) =
\det(A) \operatorname{tr}(A^{-1}H)$ lorsque $A$ est inversible. Donc les
matrices inversibles ne sont jamais des points critiques. Les points critiques
sont les matrices singulières, mais la seule valeur critique est $0$ (car
$f(C_f) = \{0\}$). Le théorème de Sard est vérifié de manière triviale ici :
$\{0\}$ est de mesure nulle dans $\R$.
\end{exemple}

% ============================================================
\section{Ensembles de mesure nulle}\label{sec:mesure-nulle}
% ============================================================

\begin{definition}[Mesure nulle dans $\R^n$]\label{def:mesure-nulle}
\index{mesure nulle}
Un sous-ensemble $A \subset \R^n$ est de \emph{mesure (de Lebesgue) nulle} si,
pour tout $\varepsilon > 0$, il existe une famille dénombrable de pavés
$\set{R_i}_{i \in \N}$ telle que $A \subset \bigcup_{i} R_i$ et
$\sum_{i} \operatorname{vol}(R_i) < \varepsilon$.
\end{definition}

\begin{lemme}\label{lem:mesure-nulle-denombrable}
\begin{enumerate}[label=(\roman*)]
  \item Une union dénombrable d'ensembles de mesure nulle est de mesure nulle.
  \item Un sous-ensemble de mesure nulle n'a pas de point intérieur (et en
  particulier n'est pas $\R^n$ entier).
  \item Si $A \subset \R^n$ est de mesure nulle et $f \colon \R^n \to \R^n$ est
  lipschitzienne sur un voisinage de $A$, alors $f(A)$ est de mesure nulle.
\end{enumerate}
\end{lemme}

\begin{proof}
(i) Si $A = \bigcup_j A_j$ avec $\mu(A_j) = 0$, pour chaque $j$ on recouvre
$A_j$ par des pavés de volume total $< \varepsilon/2^{j+1}$. Leur réunion
recouvre $A$ avec un volume total $< \varepsilon$.

(ii) Si $A$ contenait un ouvert $U$, alors $U$ aurait un volume positif,
contradiction puisque $U \subset A$ et $A$ est recouvert par des pavés de
volume total arbitrairement petit.

(iii) Si $f$ est $L$-lipschitzienne, elle envoie un pavé de côté $r$ dans un
pavé de côté $2Lr$. Donc les volumes sont multipliés par au plus $(2L)^n$, et
le résultat suit.
\end{proof}

\begin{definition}[Mesure nulle dans une variété]\label{def:mesure-nulle-variete}
\index{mesure nulle!dans une variété}
Un sous-ensemble $A \subset M$ d'une variété lisse $M$ est de \emph{mesure
nulle} si, pour toute carte $(U, \varphi)$ de $M$, l'ensemble $\varphi(A \cap U)$
est de mesure nulle dans $\R^n$. Par le 
ef{lem:mesure-nulle-denombrable}(iii),
cette définition ne dépend pas du choix d'atlas.
\end{definition}

\begin{lemme}[Lemme de Fubini pour les ensembles de mesure nulle]\label{lem:fubini-mesure}
\index{Fubini!lemme pour mesure nulle}
Soit $A \subset \R^n$ tel que, pour tout $t \in \R$, la \emph{tranche}
$A_t = \set{x' \in \R^{n-1} \mid (x', t) \in A}$ est de mesure nulle dans
$\R^{n-1}$. Alors $A$ est de mesure nulle dans $\R^n$.
\end{lemme}

\begin{proof}
Fixons $\varepsilon > 0$. Pour chaque rationnel $q_k \in \Q$ et chaque
$j \in \N$, recouvrons $A_{q_k}$ par des pavés de $\R^{n-1}$ de volume total
$< \varepsilon / 2^{k+j+2}$. En épaississant chaque pavé $(n{-}1)$-dimensionnel
en un pavé $n$-dimensionnel de hauteur $1/2^j$, et en prenant la réunion sur
tous les rationnels $q_k$ dans un intervalle $[q_k - 1/2^{j+1}, q_k + 1/2^{j+1}]$,
on construit un recouvrement de $A$ (car pour tout $(x', t) \in A$, la tranche
$A_t$ est de mesure nulle, et $t$ est approché par des rationnels) de volume
total fini arbitrairement petit.

Plus précisément, fixons $N \in \N$. Posons $I_k = [k/N, (k+1)/N]$ pour
$k \in \Z$. Pour $t \in I_k$, on a $A_t \subset \R^{n-1}$ de mesure nulle.
On recouvre $A_t$ par des $(n{-}1)$-pavés de volume total $< \varepsilon/(2^{\abs{k}+2} \cdot N)$.
En épaississant avec $I_k$ (de longueur $1/N$), on obtient des $n$-pavés
de volume total $< \varepsilon/2^{\abs{k}+2}$. En sommant sur $k \in \Z$,
le volume total est $< \varepsilon$.
\end{proof}

% ============================================================
\section{Théorème de Sard}\label{sec:sard}
% ============================================================

\begin{theoreme}[Sard, 1942]\label{thm:sard}
\index{théorème de Sard}
Soit $f \colon M^m \to N^n$ une application lisse entre variétés. L'ensemble
des valeurs critiques de $f$ est de mesure nulle dans $N$.
\end{theoreme}

La preuve se ramène au cas euclidien par des cartes.

\begin{lemme}\label{lem:sard-local}
Il suffit de démontrer le théorème de Sard pour $f \colon U \to \R^n$, où
$U \subset \R^m$ est ouvert.
\end{lemme}

\begin{proof}[Preuve du lemme]
Soit $\set{(U_\alpha, \varphi_\alpha)}$ un atlas dénombrable de $M$ et
$\set{(V_\beta, \psi_\beta)}$ un atlas dénombrable de $N$. L'ensemble des
valeurs critiques $f(C_f)$ s'écrit comme union dénombrable d'images
$\psi_\beta(f(C_f) \cap V_\beta)$. Pour chaque $(\alpha, \beta)$,
l'application $\psi_\beta \circ f \circ \varphi_\alpha^{-1}$ est lisse de
$\varphi_\alpha(U_\alpha \cap f^{-1}(V_\beta)) \subset \R^m$ vers $\R^n$, et
ses valeurs critiques contiennent $\psi_\beta(f(C_f \cap U_\alpha) \cap V_\beta)$.
Par union dénombrable, le résultat pour les ouverts de $\R^m$ implique le cas
général.
\end{proof}

\begin{proof}[Preuve du 
ef{thm:sard}]
Par le 
ef{lem:sard-local}, il suffit de traiter $f \colon U \to \R^n$
avec $U \subset \R^m$ ouvert. On procède par récurrence sur $m$.

\medskip\noindent\textbf{Cas $m = 0$.} Trivial : $U$ est discret, dénombrable,
et un ensemble dénombrable est de mesure nulle.

\medskip\noindent\textbf{Cas $m < n$.} Pour tout $p \in U$,
$\rg(\dd f_p) \leq m < n$, donc tout point est critique et $C_f = U$. On
doit montrer que $f(U)$ est de mesure nulle dans $\R^n$. Cela suit du fait
que $f(U)$ est contenu dans une union dénombrable d'images de cubes de
dimension $m < n$, et l'image d'un $m$-cube par une application lipschitzienne
dans $\R^n$ avec $m < n$ est de mesure nulle (l'image est contenue dans un
«~voisinage tubulaire~» de volume arbitrairement petit).

\medskip\noindent\textbf{Cas $m \geq n$ : décomposition de l'ensemble critique.}
Définissons les sous-ensembles suivants de $C_f$ :
\begin{align*}
  C_1 &= \set{p \in U \mid \dd f_p = 0}, \\
  C_k &= \set{p \in U \mid D^\alpha f(p) = 0
              \text{ pour tout multi-indice } \abs{\alpha} \leq k},
              \quad k \geq 2.
\end{align*}
On a la suite décroissante $C_f \supset C_1 \supset C_2 \supset \cdots$.

Nous montrons séparément :
\begin{enumerate}[label=(\Alph*)]
  \item $f(C_k)$ est de mesure nulle pour $k$ assez grand
  (plus précisément $k \geq m/n$) ;
  \item $f(C_{k} \setminus C_{k+1})$ est de mesure nulle pour tout $k \geq 1$ ;
  \item $f(C_f \setminus C_1)$ est de mesure nulle.
\end{enumerate}
Puisque $f(C_f) = f(C_f \setminus C_1) \cup f(C_1 \setminus C_2) \cup \cdots
\cup f(C_{k-1} \setminus C_k) \cup f(C_k)$, cela conclut.

\medskip\noindent\textbf{Preuve de (A) : $f(C_k)$ est de mesure nulle pour
$k \geq m/n$.}
Soit $p \in C_k$ et $Q$ un cube de côté $\delta$ centré en $p$, contenu dans
$U$. Par Taylor, pour $x \in Q$,
\[
  \abs{f(x) - f(p)} \leq C \abs{x - p}^{k+1} \leq C'(\delta)^{k+1}.
\]
Donc $f(Q)$ est contenu dans un cube de côté $\leq 2C'(\delta)^{k+1}$ dans
$\R^n$. Subdivisons $Q$ en $N^m$ sous-cubes de côté $\delta/N$. Chaque
sous-cube contenant un point de $C_k$ a son image contenue dans un cube de
$\R^n$ de côté $\leq 2C'(\delta/N)^{k+1}$. Le volume total des images est
donc majoré par
\[
  N^m \cdot \bigl(2C'(\delta/N)^{k+1}\bigr)^n
  = \text{const} \cdot N^{m - n(k+1)}.
\]
Si $k + 1 > m/n$, i.e.\ $k \geq \lceil m/n \rceil$, l'exposant $m - n(k+1) < 0$
et cette quantité tend vers $0$ quand $N \to \infty$. Donc $f(C_k \cap Q)$ est
de mesure nulle. Par un recouvrement dénombrable de $C_k$ par de tels cubes,
$f(C_k)$ est de mesure nulle.

\medskip\noindent\textbf{Preuve de (B) : $f(C_k \setminus C_{k+1})$ est de mesure
nulle.}
Soit $p \in C_k \setminus C_{k+1}$. Il existe un multi-indice $\abs{\alpha} = k+1$
avec $D^\alpha f_i(p) \neq 0$ pour une composante $f_i$. On peut supposer que
$\frac{\partial}{\partial x_1}(D^{\alpha'} f_i)(p) \neq 0$ pour un certain
$\abs{\alpha'} = k$. Posons $g = D^{\alpha'} f_i$; alors $g(p) = 0$ et
$\frac{\partial g}{\partial x_1}(p) \neq 0$.

Par le théorème des fonctions implicites, $\set{g = 0}$ est localement une
hypersurface lisse, difféomorphe à $\R^{m-1}$. L'application
$\Phi(x) = (g(x), x_2, \dots, x_m)$ est un difféomorphisme local. En composant
$f \circ \Phi^{-1}$, on se ramène à étudier les valeurs critiques d'une
application lisse sur $\set{0} \times \R^{m-1} \cong \R^{m-1}$, ce qui donne
le résultat par l'hypothèse de récurrence sur $m$.

\medskip\noindent\textbf{Preuve de (C) : $f(C_f \setminus C_1)$ est de mesure nulle.}
Soit $p \in C_f \setminus C_1$. Il existe $i,j$ tels que
$\frac{\partial f_i}{\partial x_j}(p) \neq 0$. Quitte à réordonner, supposons
$\frac{\partial f_1}{\partial x_1}(p) \neq 0$. L'application
\[
  \Phi(x) = (f_1(x), x_2, \dots, x_m)
\]
est un difféomorphisme local au voisinage de $p$. En posant
$\tilde{f} = f \circ \Phi^{-1}$, on obtient $\tilde{f}(y_1, y_2, \dots, y_m)
= (y_1, h_2(y), \dots, h_n(y))$.

Pour $y_1 = c$ fixé, considérons $\tilde{f}_c(y_2, \dots, y_m)
= (h_2(c, y'), \dots, h_n(c, y'))$. Un point critique de $\tilde{f}$
dans la tranche $\set{y_1 = c}$ est un point critique de $\tilde{f}_c$.
Par récurrence sur $m$ (appliquée à $\tilde{f}_c \colon \R^{m-1} \to \R^{n-1}$),
l'ensemble des valeurs critiques de $\tilde{f}_c$ est de mesure nulle dans
$\R^{n-1}$ pour chaque $c$. Par le lemme de Fubini
(
ef{lem:fubini-mesure}), l'ensemble des valeurs critiques de $\tilde{f}$
(dans $\R^n$) est de mesure nulle, d'où le résultat.
\end{proof}

\begin{corollaire}\label{cor:valeurs-regulieres-denses}
\index{valeurs régulières!densité}
L'ensemble des valeurs régulières d'une application lisse $f \colon M \to N$
est dense dans $N$. Plus précisément, c'est un résiduel (intersection
dénombrable d'ouverts denses) dans $N$.
\end{corollaire}

\begin{proof}
Un ensemble de mesure nulle a un complémentaire dense (car un ouvert non vide
de $\R^n$ n'est pas de mesure nulle). Le complémentaire de l'ensemble des
valeurs critiques est l'ensemble des valeurs régulières. La structure de
résiduel résulte du fait qu'un ensemble de mesure nulle est contenu dans une
intersection dénombrable d'ouverts de mesure arbitrairement petite.
\end{proof}

% ============================================================
\section{Applications du théorème de Sard}\label{sec:sard-applications}
% ============================================================

\subsection{Existence de valeurs régulières}\label{ssec:existence-val-reg}

\begin{proposition}\label{prop:existence-val-reg}
Soit $f \colon M \to N$ une application lisse entre variétés. Si $\dim M < \dim N$,
alors $f(M)$ est de mesure nulle dans $N$; en particulier, $f$ n'est pas
surjective.
\end{proposition}

\begin{proof}
Si $\dim M < \dim N$, tout point de $M$ est critique. Par Sard, $f(C_f) = f(M)$
est de mesure nulle.
\end{proof}

\begin{corollaire}\label{cor:pas-immersion-dim}
Il n'existe pas de submersion $f \colon M^m \to N^n$ si $m < n$.
\end{corollaire}

\subsection{Théorème de non-rétraction}\label{ssec:non-retraction}

\begin{theoreme}[Non-rétraction]\label{thm:non-retraction}
\index{non-rétraction}
Il n'existe pas de rétraction lisse $r \colon D^n \to S^{n-1}$, c'est-à-dire
d'application lisse $r \colon D^n \to S^{n-1}$ telle que
$r|_{S^{n-1}} = \Id_{S^{n-1}}$.
\end{theoreme}

\begin{proof}
Supposons par l'absurde qu'une telle rétraction lisse $r$ existe.

\medskip\noindent\textbf{Étape 1.}
Par le théorème de Sard, $r$ admet une valeur régulière
$q \in S^{n-1}$. Puisque $r|_{S^{n-1}} = \Id$, on a $q \in r(D^n)$,
donc $r^{-1}(q) \neq \varnothing$.

\medskip\noindent\textbf{Étape 2.}
Puisque $q$ est une valeur régulière, $r^{-1}(q)$ est une sous-variété lisse
de $D^n$ de dimension $n - (n-1) = 1$. C'est donc une courbe lisse (compacte,
car $D^n$ est compact) dans $D^n$.

\medskip\noindent\textbf{Étape 3.}
Le bord de $r^{-1}(q)$ dans $D^n$ est $r^{-1}(q) \cap S^{n-1} = \{q\}$
(puisque $r|_{S^{n-1}} = \Id$ implique que le seul point de $S^{n-1}$ envoyé
sur $q$ est $q$ lui-même). Ainsi, $r^{-1}(q)$ est une $1$-variété compacte
à bord consistant en un seul point.

\medskip\noindent\textbf{Étape 4.}
Or, toute variété compacte de dimension $1$ à bord est une union disjointe de
copies de $[0,1]$ et de $S^1$. Le nombre de points de bord est donc pair.
Avoir un seul point de bord est impossible. Contradiction.
\end{proof}

\subsection{Théorème du point fixe de Brouwer}\label{ssec:brouwer}

\begin{theoreme}[Brouwer, version lisse]\label{thm:brouwer}
\index{théorème de Brouwer}
Toute application lisse $f \colon D^n \to D^n$ admet un point fixe.
\end{theoreme}

\begin{proof}
Supposons que $f$ n'ait pas de point fixe. Pour chaque $x \in D^n$,
considérons la demi-droite issue de $f(x)$ passant par $x$ et prolongée
au-delà de $x$. Soit $r(x) \in S^{n-1}$ le point d'intersection de cette
demi-droite avec $S^{n-1}$. Plus précisément, $r(x) = x + t(x)(x - f(x))$
où $t(x) \geq 0$ est l'unique valeur telle que $\abs{r(x)} = 1$.

L'application $r \colon D^n \to S^{n-1}$ est lisse (car $t(x)$ dépend
lissement de $x$ par le théorème des fonctions implicites, puisque la
condition $\abs{x + t(x-f(x))}^2 = 1$ est résolue par une équation du
second degré en $t$ dont le discriminant est strictement positif).

De plus, si $x \in S^{n-1}$, la demi-droite issue de $f(x)$ passant par $x$
atteint $S^{n-1}$ en $x$ (car $\abs{x} = 1$ et $t = 0$ convient), donc
$r|_{S^{n-1}} = \Id_{S^{n-1}}$.

Ceci contredit le 
ef{thm:non-retraction}.
\end{proof}

\begin{remarque}\label{rem:brouwer-continu}
Le théorème de Brouwer reste vrai pour les applications continues. On se
ramène au cas lisse par approximation : toute application continue
$f \colon D^n \to D^n$ est uniformément approchée par des applications lisses
$f_\varepsilon$ (par convolution). Si $f$ n'a pas de point fixe, alors
$\abs{f(x) - x} \geq \delta > 0$ pour tout $x$, et pour $\varepsilon$
assez petit, $f_\varepsilon$ n'a pas non plus de point fixe, contredisant
le cas lisse.
\end{remarque}

\begin{remarque}\label{rem:brouwer-dimension-1}
\index{théorème de Brouwer!dimension 1}
En dimension $1$, le théorème de Brouwer est le théorème des valeurs
intermédiaires : si $f \colon [0,1] \to [0,1]$ est continue, la fonction
$g(x) = f(x) - x$ satisfait $g(0) = f(0) \geq 0$ et $g(1) = f(1) - 1 \leq 0$,
donc $g$ s'annule en un point, qui est un point fixe de $f$.
\end{remarque}

\begin{corollaire}\label{cor:borsuk-ulam-faible}
\index{Borsuk--Ulam}
Il n'existe pas d'application continue $f \colon D^n \to \partial D^n = S^{n-1}$
sans point fixe (au sens où $f(x) \neq x$ pour tout $x \in \partial D^n$) qui
soit l'identité sur le bord.
\end{corollaire}

\subsection{Théorème de Brown}\label{ssec:brown}

\begin{theoreme}[Brown]\label{thm:brown}
\index{théorème de Brown}
Soit $h \colon \R^n \to \R^n$ un homéomorphisme. Alors $h$ envoie les
ensembles de mesure nulle sur des ensembles de mesure nulle.
\end{theoreme}

\begin{remarque}\label{rem:brown-apercu}
La preuve de ce théorème utilise le théorème de Sard de manière indirecte, via
l'invariance du domaine et des propriétés topologiques des homéomorphismes de
$\R^n$. Nous ne donnons pas la preuve complète ici, mais signalons que
l'ingrédient clé est le fait (conséquence de Sard) que l'image d'un compact
de mesure nulle par une application lipschitzienne est de mesure nulle, combiné
à l'approximation d'un homéomorphisme par des applications lipschitziennes
par morceaux.
\end{remarque}

\subsection{Applications transverses}\label{ssec:transverses}

\begin{definition}[Transversalité]\label{def:transversalite-rappel}
\index{transversalité}
Soit $f \colon M \to N$ lisse et $W \subset N$ une sous-variété. On dit que
$f$ est \emph{transverse} à $W$, noté $f \pitchfork W$, si pour tout
$p \in f^{-1}(W)$,
\[
  \im(\dd f_p) + T_{f(p)}W = T_{f(p)}N.
\]
\end{definition}

\begin{theoreme}[Théorème de transversalité]\label{thm:transversalite-sard}
\index{théorème de transversalité}
Soit $F \colon M \times S \to N$ une application lisse telle que
$F \pitchfork W$, où $W \subset N$ est une sous-variété fermée. Alors,
pour presque tout $s \in S$ (i.e.\ pour tout $s$ sauf un ensemble de
mesure nulle), l'application $f_s = F(\cdot, s) \colon M \to N$ est
transverse à $W$.
\end{theoreme}

\begin{proof}
Posons $Z = F^{-1}(W)$. Puisque $F \pitchfork W$, $Z$ est une sous-variété
de $M \times S$. Considérons la projection $\pi \colon Z \to S$,
$\pi(p, s) = s$. Par Sard, l'ensemble des valeurs critiques de $\pi$ est de
mesure nulle dans $S$. Un calcul montre que $s$ est une valeur régulière de
$\pi$ si et seulement si $f_s \pitchfork W$.
\end{proof}

\begin{corollaire}\label{cor:approx-transverse}
\index{approximation transverse}
Soit $f \colon M \to N$ lisse et $W \subset N$ une sous-variété fermée.
Alors $f$ peut être approchée arbitrairement (en topologie $C^\infty$) par
des applications transverses à $W$.
\end{corollaire}

\begin{proof}[Esquisse de preuve]
On plonge $N$ dans $\R^k$ (par Whitney) et on considère un voisinage tubulaire
$\nu \colon E \to N$ de $N$ dans $\R^k$. On définit
$F \colon M \times B_\varepsilon(0) \to N$ par $F(p,s) = \nu^{-1}(f(p) + s)$
(pour $\varepsilon$ assez petit). On vérifie que $F \pitchfork W$ en utilisant
la surjectivité des translations. Par le 
ef{thm:transversalite-sard},
$f_s = F(\cdot, s)$ est transverse à $W$ pour presque tout $s \in B_\varepsilon(0)$,
et $f_s \to f$ quand $s \to 0$.
\end{proof}

\begin{remarque}\label{rem:transversalite-generique}
\index{transversalité!généricité}
Le 
ef{cor:approx-transverse} est souvent résumé par le slogan :
«~la transversalité est une propriété générique~». Plus précisément,
l'ensemble des applications $f \colon M \to N$ transverses à une sous-variété
fermée $W \subset N$ est un résiduel (et en particulier dense) dans l'espace
$C^\infty(M, N)$ muni de la topologie de Whitney. Ce résultat fondamental,
dû à Thom, est à la base de la théorie de la transversalité développée au
chapitre~5.
\end{remarque}

% ============================================================
\section{Compléments}\label{sec:sard-complements}
% ============================================================

\begin{proposition}[Sard paramétrique]\label{prop:sard-parametrique}
\index{Sard!version paramétrique}
Soit $f \colon M \to N$ une application lisse de classe $C^r$ avec
$r \geq \max(m - n + 1, 1)$. Alors les valeurs critiques de $f$ sont de
mesure nulle dans $N$.
\end{proposition}

\begin{remarque}\label{rem:sard-Cr}
L'hypothèse de régularité $C^r$ avec $r \geq \max(m-n+1, 1)$ est optimale :
Whitney a construit des exemples d'applications $C^{r-1}$ dont l'ensemble des
valeurs critiques a une mesure de Lebesgue positive. Plus précisément, pour
$m \geq n \geq 1$ et $r < \max(m - n + 1, 1)$, il existe $f \in C^r(\R^m, \R^n)$
telle que l'ensemble des valeurs critiques de $f$ contient un ouvert non vide
de $\R^n$.
\end{remarque}

\begin{exemple}\label{ex:sard-Morse}
\index{fonctions de Morse!et Sard}
Soit $M$ une variété compacte et $f \colon M \to \R$ une fonction lisse.
Les valeurs critiques de $f$ forment un compact de mesure nulle dans $\R$.
Les valeurs régulières forment donc un ouvert dense. Si de plus $f$ est
une \emph{fonction de Morse} (tous les points critiques sont non dégénérés),
l'ensemble des valeurs critiques est fini. Par Sard, les fonctions de Morse
sont denses dans $C^\infty(M, \R)$ : étant donné $f \colon M \to \R$, la
famille $f_a(x) = f(x) + \langle a, \varphi(x) \rangle$ (où $\varphi$ est
un plongement de $M$ dans $\R^k$) est une fonction de Morse pour presque
tout $a \in \R^k$.
\end{exemple}

\begin{proposition}\label{prop:preimage-reguliere-sard}
Soit $f \colon M^m \to N^n$ lisse avec $m \geq n$. Alors pour presque toute
valeur $q \in N$, la préimage $f^{-1}(q)$ est soit vide, soit une sous-variété
lisse de $M$ de dimension $m - n$.
\end{proposition}

\begin{proof}
Par Sard, presque tout $q$ est une valeur régulière. Pour une valeur régulière,
le théorème de la préimage régulière donne le résultat.
\end{proof}

% ============================================================
\section{Exercices}\label{sec:exos-sard}
% ============================================================

\begin{exercice}\label{exo:sard-1}
Soit $f \colon \R \to \R$ définie par $f(x) = x^3$. Déterminer l'ensemble des
points critiques et des valeurs critiques de $f$. Le théorème de Sard
est-il vérifié?
\end{exercice}

\begin{exercice}\label{exo:sard-2}
\index{mesure nulle!exemples}
Montrer que l'ensemble triadique de Cantor $\mathcal{C} \subset [0,1]$ est de
mesure nulle. Plus généralement, montrer qu'un compact totalement discontinu
de $\R$ est de mesure nulle.
\end{exercice}

\begin{exercice}\label{exo:sard-3}
Soit $\gamma \colon \R \to \R^2$ une courbe lisse. Montrer que l'image
$\gamma(\R)$ est de mesure nulle dans $\R^2$.
\textit{Indication : utiliser la proposition \ref{prop:existence-val-reg}.}
\end{exercice}

\begin{exercice}\label{exo:sard-4}
Soit $f \colon S^2 \to \R$ une application lisse. Montrer que presque toute
valeur $t \in \R$ est régulière, et que pour une telle valeur $t$,
$f^{-1}(t)$ est une union finie de courbes fermées lisses plongées dans $S^2$.
\end{exercice}

\begin{exercice}\label{exo:sard-5}
\index{théorème de Brouwer!constructif}
Soit $n \geq 1$.
\begin{enumerate}[label=(\alph*)]
  \item Reprendre la preuve du 
ef{thm:brouwer} et vérifier tous les
  détails de la construction de la rétraction $r$.
  \item Montrer que $r$ est bien lisse en détaillant le calcul du discriminant.
\end{enumerate}
\end{exercice}

\begin{exercice}\label{exo:sard-6}
Soit $f \colon \R^n \to \R^n$ une application lisse propre telle que
$\det(\dd f_p) > 0$ pour tout $p$. Montrer que $f$ est surjective.
\textit{Indication : montrer que $f(\R^n)$ est ouvert et fermé dans $\R^n$.}
\end{exercice}

\begin{exercice}\label{exo:sard-7}
\index{transversalité!exercice}
Soient $M = S^1 \times S^1$ (le tore) et $f \colon M \to \R^3$ un plongement
lisse. Montrer que pour presque tout plan affine $P \subset \R^3$ de dimension
$2$, l'intersection $f(M) \cap P$ est soit vide, soit une union finie de courbes
lisses.
\textit{Indication : paramétrer les plans par la grassmannienne et un vecteur
de translation, puis appliquer le théorème de transversalité.}
\end{exercice}

\begin{exercice}\label{exo:sard-8}
Soit $f \colon M \to N$ une application lisse entre variétés compactes de même
dimension, avec $M$ connexe et $N$ connexe. Montrer que si $f$ est une
submersion (pas de points critiques), alors $f$ est un revêtement lisse.
\textit{Indication : montrer que $f$ est propre et un difféomorphisme local, puis
que les fibres sont finies.}
\end{exercice}

\begin{exercice}\label{exo:sard-9}
Montrer qu'il n'existe pas d'application lisse $f \colon S^2 \to S^1$ qui
soit une submersion.
\textit{Indication : raisonner par les préimages de valeurs régulières.}
\end{exercice}

\begin{exercice}\label{exo:sard-10}
\index{théorème du point fixe!conséquences}
Déduire du théorème de Brouwer les résultats suivants :
\begin{enumerate}[label=(\alph*)]
  \item Toute application continue $f \colon [0,1]^n \to [0,1]^n$ admet un
  point fixe.
  \item Il n'existe pas de champ de vecteurs continu non nul sur $S^{2n}$
  (théorème de la boule chevelue pour les sphères paires).
  \textit{Indication pour (b) : si $v$ est un tel champ, construire une
  rétraction $D^{2n+1} \to S^{2n}$.}
\end{enumerate}
\end{exercice}

% === FIN DU CHUNK ch7-8 ===
% === DEBUT DU CHUNK ch9-11+end ===

% ============================================================
%  CHAPITRE 9 : THÉORIE DE MORSE --- INTRODUCTION
% ============================================================
\chapter{Théorie de Morse --- Introduction}\label{chap:morse}
\minitoc

\vspace{1cm}

La théorie de Morse\index{théorie de Morse} constitue l'un des ponts les plus
remarquables entre l'analyse et la topologie. L'idée fondatrice, due à Marston
Morse dans les années 1920--1930, est d'étudier la topologie d'une variété $M$
à travers le comportement des points critiques d'une fonction lisse
$f : M \to \R$. Plus précisément, on montre que les changements de topologie
des ensembles de sous-niveaux $f^{-1}((-\infty, a])$ sont entièrement contrôlés
par les points critiques de $f$ et par leur \emph{indice}.

% ============================================================
\section{Fonctions de Morse}\label{sec:morse-fonctions}
% ============================================================

\begin{definition}[Point critique non dégénéré]\label{def:point-critique-non-degenere}
\index{point critique!non dégénéré}
Soit $M$ une variété lisse de dimension $n$ et $f : M \to \R$ une application
lisse. Un point $p \in M$ est un \textbf{point critique}\index{point critique}
de $f$ si $\dd f_p = 0$, c'est-à-dire si la différentielle de $f$ en $p$
s'annule. Un point critique $p$ est dit \textbf{non dégénéré} si la
\textbf{matrice hessienne}\index{matrice hessienne}
\[
  H_f(p) = \left(\frac{\partial^2 f}{\partial x_i \partial x_j}(p)\right)_{1 \leqslant i,j \leqslant n}
\]
est inversible, pour un choix (et donc pour tout choix) de coordonnées locales
$(x_1, \ldots, x_n)$ au voisinage de $p$.
\end{definition}

\begin{remarque}
La non-dégénérescence de la hessienne ne dépend pas du choix de coordonnées
locales. En effet, si $\varphi$ désigne un changement de coordonnées, la matrice
hessienne se transforme par la règle
\[
  H_{f \circ \varphi}(q) = J_\varphi(q)^T \, H_f(\varphi(q)) \, J_\varphi(q)
  + \sum_k \frac{\partial f}{\partial y_k}(\varphi(q)) \, \frac{\partial^2 \varphi_k}{\partial x_i \partial x_j}(q).
\]
Si $p$ est un point critique, le second terme disparaît et l'inversibilité est
préservée par congruence.
\end{remarque}

\begin{definition}[Fonction de Morse]\label{def:fonction-morse}
\index{fonction de Morse}
Une application lisse $f : M \to \R$ est une \textbf{fonction de Morse} si tous
ses points critiques sont non dégénérés.
\end{definition}

\begin{definition}[Indice d'un point critique]\label{def:indice-morse}
\index{indice!d'un point critique}
Soit $p$ un point critique non dégénéré de $f : M \to \R$. L'\textbf{indice} de
$f$ en $p$, noté $\lambda(p)$ ou $\operatorname{ind}(f, p)$, est le nombre de
valeurs propres strictement négatives de la hessienne $H_f(p)$. Puisque
$H_f(p)$ est une matrice symétrique réelle inversible, on a
$0 \leqslant \lambda(p) \leqslant n$.
\end{definition}

\begin{remarque}
Par la loi d'inertie de Sylvester\index{loi d'inertie de Sylvester}, l'indice
est un invariant bien défini, indépendant du choix de coordonnées locales.
\end{remarque}

% ============================================================
\section{Le lemme de Morse}\label{sec:lemme-morse}
% ============================================================

Le lemme suivant montre qu'au voisinage d'un point critique non dégénéré, une
fonction de Morse admet une forme normale particulièrement simple.

\begin{lemme}[Lemme de Morse]\label{lem:morse}
\index{lemme de Morse}
Soit $f : M \to \R$ une application lisse et $p \in M$ un point critique non
dégénéré d'indice $\lambda$. Alors il existe un système de coordonnées locales
$(y_1, \ldots, y_n)$ centré en $p$ tel que, dans ces coordonnées,
\[
  f(y) = f(p) - y_1^2 - \cdots - y_\lambda^2 + y_{\lambda+1}^2 + \cdots + y_n^2.
\]
\end{lemme}

\begin{proof}
On peut supposer, quitte à travailler en coordonnées, que $M = \R^n$, $p = 0$
et $f(0) = 0$. On procède en deux étapes.

\medskip
\noindent\textbf{Étape 1 : Décomposition quadratique.}
Par le lemme de Hadamard\index{lemme de Hadamard}, puisque $f(0) = 0$, il existe
des fonctions lisses $g_i$ avec $f(x) = \sum_{i=1}^n x_i \, g_i(x)$ et
$g_i(0) = \frac{\partial f}{\partial x_i}(0) = 0$ (car $0$ est un point
critique). Appliquant le lemme de Hadamard à chaque $g_i$, on obtient des
fonctions lisses $h_{ij}$ telles que
\[
  f(x) = \sum_{i,j=1}^n x_i \, x_j \, h_{ij}(x).
\]
On peut supposer $h_{ij} = h_{ji}$ en posant $h_{ij}' = \frac{1}{2}(h_{ij} + h_{ji})$.
De plus, $(h_{ij}(0)) = \frac{1}{2} H_f(0)$ est inversible.

\medskip
\noindent\textbf{Étape 2 : Diagonalisation par récurrence.}
On procède par récurrence sur le nombre de variables déjà mises sous forme
canonique. Supposons que, dans des coordonnées $(u_1, \ldots, u_n)$, on a
\[
  f = \pm u_1^2 \pm \cdots \pm u_{r-1}^2 + \sum_{i,j \geqslant r} u_i \, u_j \, h_{ij}(u)
\]
avec $(h_{ij}(0))_{i,j \geqslant r}$ inversible. En particulier, quitte à
effectuer un changement linéaire dans les variables $u_r, \ldots, u_n$, on peut
supposer $h_{rr}(0) \neq 0$. Posons $\varepsilon = \operatorname{sgn}(h_{rr}(0))$
et, au voisinage de $0$,
\[
  v_r = \sqrt{\abs{h_{rr}(u)}} \left(u_r + \sum_{j > r} \frac{h_{rj}(u)}{h_{rr}(u)} \, u_j \right),
  \quad v_i = u_i \text{ pour } i \neq r.
\]
La fonction $\abs{h_{rr}}^{1/2}$ est lisse et non nulle au voisinage de $0$.
Le jacobien de la transformation $(u_1, \ldots, u_n) \mapsto (v_1, \ldots, v_n)$
est non nul en $0$, donc c'est un difféomorphisme local. Un calcul direct montre
que dans les nouvelles coordonnées :
\[
  f = \pm v_1^2 \pm \cdots \pm v_{r-1}^2 + \varepsilon \, v_r^2
    + \sum_{i,j > r} v_i \, v_j \, h'_{ij}(v),
\]
ce qui achève la récurrence. Au bout de $n$ étapes, on obtient la forme
quadratique souhaitée, et le nombre de signes $-$ est exactement l'indice
$\lambda$ par la loi d'inertie de Sylvester.
\end{proof}

\begin{corollaire}\label{cor:morse-isoles}
Les points critiques non dégénérés sont isolés. En particulier, si $M$ est
compacte, une fonction de Morse $f : M \to \R$ n'a qu'un nombre fini de points
critiques.
\end{corollaire}

\begin{proof}
Au voisinage d'un point critique non dégénéré $p$, le lemme de Morse donne
$\dd f = -2y_1 \, \dd y_1 - \cdots - 2y_\lambda \, \dd y_\lambda + 2y_{\lambda+1} \, \dd y_{\lambda+1} + \cdots + 2y_n \, \dd y_n$,
qui ne s'annule qu'en $p$.
\end{proof}

\begin{corollaire}\label{cor:morse-valeurs-distinctes}
Soit $f : M \to \R$ une fonction de Morse sur une variété compacte. Alors, par
une perturbation arbitrairement petite, on peut supposer que les valeurs critiques
$f(p_1), \ldots, f(p_k)$ sont toutes distinctes.
\end{corollaire}

\begin{proof}
Les points critiques sont isolés (et en nombre fini sur une variété compacte).
Si $f(p_i) = f(p_j)$ pour $i \neq j$, on modifie $f$ en ajoutant une petite
fonction à support compact au voisinage de $p_i$ pour séparer les valeurs
critiques, sans créer de nouveaux points critiques.
\end{proof}

\begin{proposition}[Unicité des coordonnées de Morse]\label{prop:morse-unicite}
\index{lemme de Morse!unicité}
Les coordonnées de Morse ne sont pas uniques, mais l'indice $\lambda$ est
univoquement déterminé. De plus, si $(y_1, \ldots, y_n)$ et $(z_1, \ldots, z_n)$
sont deux systèmes de coordonnées de Morse en un même point critique $p$, alors
la transformation $y \mapsto z$ préserve les sous-espaces
$\set{y_1 = \cdots = y_\lambda = 0}$ et $\set{y_{\lambda+1} = \cdots = y_n = 0}$
au premier ordre.
\end{proposition}

% ============================================================
\section{Densité des fonctions de Morse}\label{sec:morse-densite}
% ============================================================

\begin{theoreme}[Densité des fonctions de Morse]\label{thm:morse-dense}
\index{fonction de Morse!densité}
Soit $M$ une variété lisse compacte. L'ensemble des fonctions de Morse
$f : M \to \R$ est un ouvert dense dans $\mathcal{C}^\infty(M, \R)$ muni de la
topologie $\mathcal{C}^2$.
\end{theoreme}

\begin{proof}[Idée de la preuve]
La preuve utilise le théorème de Sard\index{théorème de Sard} et la
transversalité. On plonge $M$ dans $\R^N$ (par Whitney) et on considère les
fonctions hauteur $f_a(x) = \langle a, x \rangle$ pour $a \in S^{N-1}$.
L'application
\[
  F : M \times S^{N-1} \to T^*M, \quad (x, a) \mapsto (x, \dd f_a(x))
\]
est transverse à la section nulle de $T^*M$. Par le théorème de transversalité
de Thom\index{théorème de transversalité de Thom}, pour presque tout $a$
(au sens de Sard), la fonction $f_a$ est de Morse. Ceci montre la densité.

Pour l'ouverture, on observe que la condition <<~la hessienne est inversible en
tout point critique~>> est ouverte pour la topologie $\mathcal{C}^2$, car le
déterminant de la hessienne dépend continûment des dérivées secondes, et les
points critiques dépendent continûment de la fonction dans la topologie
$\mathcal{C}^1$.
\end{proof}

\begin{remarque}
On peut aussi démontrer la densité en perturbant directement une fonction
lisse : si $p$ est un point critique dégénéré de $f$, on peut ajouter une petite
forme linéaire pour lever la dégénérescence.
\end{remarque}

% ============================================================
\section{Sous-niveaux et attachement de cellules}\label{sec:morse-cellules}
% ============================================================

Notons $M^a = f^{-1}((-\infty, a])$ l'ensemble de sous-niveau\index{ensemble de sous-niveau}
de $f$ en $a$.

\begin{theoreme}[Passage sans point critique]\label{thm:morse-sans-critique}
\index{sous-niveau!passage sans point critique}
Soit $f : M \to \R$ une fonction lisse propre et $a < b$ deux valeurs régulières
de $f$ telles que $f^{-1}([a,b])$ ne contienne aucun point critique. Alors $M^a$
est difféomorphe à $M^b$, et l'inclusion $M^a \hookrightarrow M^b$ est une
équivalence d'homotopie.
\end{theoreme}

\begin{proof}[Idée de la preuve]
On utilise le champ de vecteurs gradient $\nabla f / \abs{\nabla f}^2$ (pour une
métrique riemannienne fixée) et son flot. L'absence de point critique assure que
ce champ est bien défini sur $f^{-1}([a,b])$. Le flot de ce champ fournit le
difféomorphisme désiré entre $M^a$ et $M^b$.
\end{proof}

\begin{theoreme}[Attachement d'une cellule]\label{thm:morse-attachement}
\index{attachement de cellule}
Soit $f : M \to \R$ une fonction de Morse propre et $p$ un point critique
d'indice $\lambda$ avec $f(p) = c$. Supposons que $f^{-1}([c-\varepsilon, c+\varepsilon])$
ne contienne aucun autre point critique (pour $\varepsilon > 0$ assez petit).
Alors $M^{c+\varepsilon}$ a le type d'homotopie de $M^{c-\varepsilon}$ avec une
$\lambda$-cellule attachée :
\[
  M^{c+\varepsilon} \simeq M^{c-\varepsilon} \cup_\varphi e^\lambda,
\]
où $\varphi : \partial D^\lambda \to \partial M^{c-\varepsilon}$ est une
application d'attachement convenable.
\end{theoreme}

\begin{proof}[Idée de la preuve]
Dans les coordonnées de Morse $(y_1, \ldots, y_n)$ au voisinage de $p$, on a
\[
  f(y) = c - y_1^2 - \cdots - y_\lambda^2 + y_{\lambda+1}^2 + \cdots + y_n^2.
\]
Le sous-ensemble $M^{c+\varepsilon} \setminus \operatorname{int}(M^{c-\varepsilon})$
est un <<~col~>> qui se déforme par rétraction sur le disque
$D^\lambda = \set{y_1^2 + \cdots + y_\lambda^2 \leqslant \varepsilon, \; y_{\lambda+1} = \cdots = y_n = 0}$.
Ce disque est la $\lambda$-cellule qui s'attache à $M^{c-\varepsilon}$. Le bord
$\partial D^\lambda = S^{\lambda-1}$ s'envoie dans le niveau
$f^{-1}(c-\varepsilon) = \partial M^{c-\varepsilon}$ par l'application
d'attachement $\varphi$.

La vérification que $M^{c+\varepsilon}$ a le type d'homotopie de
$M^{c-\varepsilon} \cup_\varphi e^\lambda$ utilise le flot du gradient modifié
et un argument de déformation par rétraction. On renvoie à
\cite{MilnorMorse} pour les détails techniques.
\end{proof}

\begin{figure}[ht]
\centering
\begin{tikzpicture}[scale=1.0]
  % Sous-niveau avant
  \fill[blue!10] (-3,-1.5) rectangle (3,0);
  \draw[thick, blue!60!black] (-3,0) -- (3,0);
  \node[below, blue!60!black] at (0,-1.2) {$M^{c-\varepsilon}$};

  % Cellule attachée
  \fill[red!20] (0,0) ellipse (1.2 and 1);
  \draw[thick, red!60!black] (0,0) ellipse (1.2 and 1);
  \node[red!60!black] at (0, 0.3) {$e^\lambda$};

  % Sous-niveau après
  \draw[thick, dashed, green!50!black] (-3,1.3) -- (3,1.3);
  \node[above, green!50!black] at (0,1.3) {$M^{c+\varepsilon}$};

  % Annotations
  \draw[->, thick] (4, -0.7) -- (4, 0.7);
  \node[right] at (4, 0) {$\nabla f$};
\end{tikzpicture}
\caption{Attachement d'une cellule au passage d'un point critique.}
\label{fig:morse-attachement}
\end{figure}

% ============================================================
\section{Décomposition en anses}\label{sec:morse-anses}
% ============================================================

\begin{definition}[Anse]\label{def:anse}
\index{anse (handle)}\index{décomposition en anses}
Soit $M$ une variété lisse compacte de dimension $n$. Une \textbf{anse d'indice
$\lambda$} (ou $\lambda$-anse) est un exemplaire de $D^\lambda \times D^{n-\lambda}$
attaché à la frontière de la variété existante par un plongement
$\varphi : \partial D^\lambda \times D^{n-\lambda} \to \partial M$. La variété
résultante est
\[
  M' = M \cup_\varphi (D^\lambda \times D^{n-\lambda}).
\]
\end{definition}

\begin{theoreme}[Décomposition en anses]\label{thm:decomp-anses}
\index{décomposition en anses}
Toute variété lisse compacte $M$ admet une \textbf{décomposition en anses} :
partant de $D^n$ (une $0$-anse), on construit $M$ en attachant successivement
des anses d'indices croissants.
\end{theoreme}

\begin{proof}[Esquisse]
On choisit une fonction de Morse $f : M \to \R$ (qui existe par le théorème de
densité~\ref{thm:morse-dense}) avec des valeurs critiques distinctes
$c_1 < c_2 < \cdots < c_k$. D'après le
théorème~\ref{thm:morse-sans-critique}, la topologie ne change qu'aux niveaux
critiques, et le théorème~\ref{thm:morse-attachement} montre que chaque passage
correspond à l'attachement d'une anse. Le minimum global de $f$ fournit la
$0$-anse initiale.
\end{proof}

\begin{corollaire}[CW-structure]\label{cor:morse-cw}
\index{CW-complexe!induit par Morse}
Toute variété lisse compacte admet la structure d'un CW-complexe fini, avec
exactement une cellule de dimension $\lambda$ pour chaque point critique d'indice
$\lambda$ de toute fonction de Morse.
\end{corollaire}

\begin{exemple}[Décomposition en anses du tore $T^2$]\label{ex:anses-tore}
\index{décomposition en anses!du tore}
En utilisant la fonction hauteur sur le tore $T^2$, la décomposition en anses
est la suivante :
\begin{enumerate}
  \item une $0$-anse $D^0 \times D^2 \cong D^2$ (le minimum) ;
  \item une première $1$-anse $D^1 \times D^1$ (le premier col) ;
  \item une deuxième $1$-anse $D^1 \times D^1$ (le deuxième col) ;
  \item une $2$-anse $D^2 \times D^0 \cong D^2$ (le maximum).
\end{enumerate}

\begin{center}
\begin{tikzpicture}[scale=0.8]
  % 0-anse
  \fill[blue!20] (-1.5, 0) rectangle (1.5, 1);
  \draw[thick, blue!60!black] (-1.5, 0) rectangle (1.5, 1);
  \node at (0, 0.5) {$0$-anse};

  % Flèche
  \draw[->, thick] (2, 0.5) -- (3, 0.5);

  % 0-anse + 1-anse
  \fill[blue!20] (3.5, 0) rectangle (6.5, 1);
  \fill[red!20] (4.2, 1) rectangle (5.8, 1.8);
  \draw[thick, blue!60!black] (3.5, 0) rectangle (6.5, 1);
  \draw[thick, red!60!black] (4.2, 1) rectangle (5.8, 1.8);
  \node at (5, 0.5) {$0$};
  \node at (5, 1.4) {$1$};

  % Flèche
  \draw[->, thick] (7, 0.5) -- (8, 0.5);

  % + 2eme 1-anse
  \fill[blue!20] (8.5, 0) rectangle (11.5, 1);
  \fill[red!20] (9.2, 1) rectangle (10.8, 1.8);
  \fill[orange!20] (8.5, 0) -- (8.5, -0.8) -- (11.5, -0.8) -- (11.5, 0) -- cycle;
  \draw[thick, blue!60!black] (8.5, 0) rectangle (11.5, 1);
  \draw[thick, red!60!black] (9.2, 1) rectangle (10.8, 1.8);
  \draw[thick, orange!60!black] (8.5, 0) -- (8.5, -0.8) -- (11.5, -0.8) -- (11.5, 0);
  \node at (10, 0.5) {$0$};
  \node at (10, 1.4) {$1$};
  \node at (10, -0.4) {$1$};
\end{tikzpicture}
\captionof{figure}{Construction progressive du tore par attachement d'anses.}
\label{fig:anses-tore}
\end{center}
\end{exemple}

\begin{proposition}[Réarrangement des anses]\label{prop:rearrangement}
\index{réarrangement des anses}
Soit $M$ une variété compacte munie d'une décomposition en anses. On peut
toujours réarranger l'ordre d'attachement de sorte que les anses d'indice
$\lambda$ soient attachées avant celles d'indice $\lambda + 1$.
\end{proposition}

\begin{proof}[Idée]
On modifie la fonction de Morse par une isotopie du flot gradient de façon à
séparer les valeurs critiques d'indices différents. Si deux anses d'indices
$\lambda < \mu$ sont attachées dans le <<~mauvais~>> ordre, on peut permuter les
valeurs critiques correspondantes sans changer le type de difféomorphisme de $M$,
à condition que $\lambda < \mu$.
\end{proof}

% ============================================================
\section{Inégalités de Morse}\label{sec:inegalites-morse}
% ============================================================

Soit $f : M \to \R$ une fonction de Morse sur une variété compacte $M$. Notons
$c_\lambda$ le nombre de points critiques d'indice $\lambda$ et $b_\lambda = \dim H_\lambda(M; \Q)$
le $\lambda$-ième nombre de Betti\index{nombre de Betti}.

\begin{theoreme}[Inégalités de Morse faibles]\label{thm:morse-ineg-faibles}
\index{inégalités de Morse!faibles}
Pour tout $\lambda \in \set{0, 1, \ldots, n}$,
\[
  c_\lambda \geqslant b_\lambda.
\]
\end{theoreme}

\begin{proof}
Le corollaire~\ref{cor:morse-cw} fournit un CW-complexe avec $c_\lambda$
cellules de dimension $\lambda$. Le complexe cellulaire associé a un groupe de
chaînes $C_\lambda \cong \Z^{c_\lambda}$. Par le théorème des coefficients
universels,
\[
  b_\lambda = \dim_\Q H_\lambda(M;\Q) \leqslant \dim_\Q (C_\lambda \otimes \Q) = c_\lambda.
  \qedhere
\]
\end{proof}

\begin{theoreme}[Inégalités de Morse fortes]\label{thm:morse-ineg-fortes}
\index{inégalités de Morse!fortes}
Pour tout $k \in \set{0, 1, \ldots, n}$,
\[
  \sum_{\lambda=0}^k (-1)^{k-\lambda} c_\lambda \geqslant
  \sum_{\lambda=0}^k (-1)^{k-\lambda} b_\lambda.
\]
En particulier, pour $k = n$ on retrouve l'égalité d'Euler--Poincaré :
\[
  \sum_{\lambda=0}^n (-1)^\lambda c_\lambda = \sum_{\lambda=0}^n (-1)^\lambda b_\lambda = \chi(M).
\]
\end{theoreme}

\begin{proof}
Notons $d_\lambda : C_\lambda \to C_{\lambda-1}$ les différentielles du complexe
cellulaire, et posons $z_\lambda = \rg(\ker d_\lambda)$,
$\beta_\lambda = \rg(\im d_{\lambda+1})$. Alors $b_\lambda = z_\lambda - \beta_\lambda$
et $c_\lambda = z_\lambda + \beta_{\lambda-1}$. On a donc
\[
  c_\lambda - b_\lambda = 2\beta_{\lambda-1} + (z_\lambda - z_\lambda) + (\beta_\lambda - \beta_\lambda) + \beta_\lambda - \beta_{\lambda-1} + \beta_{\lambda-1}.
\]
Plus précisément, un calcul par somme alternée donne
\[
  \sum_{\lambda=0}^k (-1)^{k-\lambda}(c_\lambda - b_\lambda) = \beta_k \geqslant 0,
\]
ce qui est exactement l'inégalité forte annoncée.
\end{proof}

% ============================================================
\section{Exemples fondamentaux}\label{sec:morse-exemples}
% ============================================================

\begin{exemple}[Hauteur sur la sphère $S^n$]\label{ex:morse-sphere}
\index{fonction de Morse!sur la sphère}
Considérons $S^n = \set{x \in \R^{n+1} \mid \abs{x}^2 = 1}$ et la fonction
hauteur $f(x_0, \ldots, x_n) = x_n$. Les points critiques sont :
\begin{itemize}
  \item le pôle sud $S = (0, \ldots, 0, -1)$, minimum, d'indice $\lambda = 0$ ;
  \item le pôle nord $N = (0, \ldots, 0, 1)$, maximum, d'indice $\lambda = n$.
\end{itemize}
La fonction $f$ est de Morse avec exactement $2$ points critiques. Les
inégalités de Morse faibles donnent $c_0 = 1 \geqslant b_0 = 1$ et
$c_n = 1 \geqslant b_n = 1$, avec $c_\lambda = 0 \geqslant b_\lambda = 0$
pour $0 < \lambda < n$, ce qui est cohérent avec $H_k(S^n) = \Z$ pour $k = 0, n$
et $0$ sinon.
\end{exemple}

\begin{exemple}[Hauteur sur le tore $T^2$]\label{ex:morse-tore}
\index{fonction de Morse!sur le tore}
Plongeons le tore $T^2$ dans $\R^3$ de la manière standard et considérons la
fonction hauteur $f : T^2 \to \R$ définie par $f(x,y,z) = z$.

\begin{center}
\begin{tikzpicture}[scale=0.9]
  % Tore stylisé (profil)
  \draw[thick] (0,0) ellipse (2.5 and 3.5);
  \draw[thick] (0,0.8) ellipse (0.8 and 0.5);

  % Lignes de niveau
  \draw[dashed, blue!60] (-2.5, -3.5) -- (3.5, -3.5);
  \draw[dashed, blue!60] (-2.5, -1.2) -- (3.5, -1.2);
  \draw[dashed, blue!60] (-2.5, 1.2) -- (3.5, 1.2);
  \draw[dashed, blue!60] (-2.5, 3.5) -- (3.5, 3.5);

  % Points critiques
  \fill[red] (0, -3.5) circle (3pt) node[right=5pt] {$p_0$ : min, $\lambda = 0$};
  \fill[red] (0, -1.2) circle (3pt);
  \node[right=5pt] at (0.1, -1.2) {$p_1$ : col, $\lambda = 1$};
  \fill[red] (0, 1.2) circle (3pt);
  \node[right=5pt] at (0.1, 1.2) {$p_2$ : col, $\lambda = 1$};
  \fill[red] (0, 3.5) circle (3pt) node[right=5pt] {$p_3$ : max, $\lambda = 2$};

  % Axe f
  \draw[->, thick] (4, -4) -- (4, 4.5) node[above] {$f = z$};
\end{tikzpicture}
\captionof{figure}{Fonction hauteur sur le tore $T^2$ : quatre points critiques.}
\label{fig:morse-tore}
\end{center}

La fonction $f$ possède exactement $4$ points critiques :
\begin{enumerate}
  \item un minimum ($\lambda = 0$) ;
  \item deux cols ($\lambda = 1$) ;
  \item un maximum ($\lambda = 2$).
\end{enumerate}
On a donc $c_0 = 1$, $c_1 = 2$, $c_2 = 1$. Les nombres de Betti du tore sont
$b_0 = 1$, $b_1 = 2$, $b_2 = 1$. Les inégalités de Morse sont des égalités :
la fonction hauteur est une fonction de Morse \textbf{parfaite}\index{fonction de Morse!parfaite}.

La caractéristique d'Euler vaut $\chi(T^2) = 1 - 2 + 1 = 0$.

L'évolution des ensembles de sous-niveaux est la suivante :
\begin{itemize}
  \item Pour $a$ juste au-dessus du minimum : $M^a \cong D^2$ (un disque).
  \item Après le premier col : $M^a$ est un cylindre (homotope à $S^1$).
  \item Après le deuxième col : $M^a$ est homéomorphe à $T^2$ privé d'un disque.
  \item Après le maximum : $M^a = T^2$.
\end{itemize}

\begin{center}
\begin{tikzpicture}[scale=0.75]
  % Disque
  \begin{scope}[shift={(-5,0)}]
    \fill[blue!15] (0,0) circle (1);
    \draw[thick, blue!50!black] (0,0) circle (1);
    \node[below] at (0,-1.3) {$a < c_1$};
    \node at (0,0) {$D^2$};
  \end{scope}

  % Cylindre
  \begin{scope}[shift={(-1.5,0)}]
    \draw[thick, blue!50!black] (-0.8,-1) -- (-0.8,1);
    \draw[thick, blue!50!black] (0.8,-1) -- (0.8,1);
    \draw[thick, blue!50!black] (-0.8,1) arc (180:360:0.8 and 0.3);
    \draw[thick, dashed, blue!50!black] (-0.8,1) arc (180:0:0.8 and 0.3);
    \draw[thick, blue!50!black] (-0.8,-1) arc (180:360:0.8 and 0.3);
    \fill[blue!15] (-0.8,-1) -- (-0.8,1) arc (180:360:0.8 and 0.3) -- (0.8,-1) arc (360:180:0.8 and 0.3) -- cycle;
    \node[below] at (0,-1.6) {$c_1 < a < c_2$};
  \end{scope}

  % T^2 privé d'un disque
  \begin{scope}[shift={(2.5,0)}]
    \draw[thick, blue!50!black] (0,0) ellipse (1.2 and 1);
    \draw[thick, blue!50!black, dashed] (0,0.3) ellipse (0.4 and 0.2);
    \node[below] at (0,-1.3) {$c_2 < a < c_3$};
  \end{scope}

  % Tore complet
  \begin{scope}[shift={(6,0)}]
    \draw[thick, blue!50!black] (0,0) ellipse (1.2 and 1);
    \draw[thick, blue!50!black] (0,0.3) ellipse (0.4 and 0.2);
    \node[below] at (0,-1.3) {$a > c_3$};
    \node at (0,-0.4) {$T^2$};
  \end{scope}
\end{tikzpicture}
\captionof{figure}{Évolution des sous-niveaux de la fonction hauteur sur $T^2$.}
\label{fig:morse-sous-niveaux-tore}
\end{center}
\end{exemple}

\begin{exemple}[Fonctions de Morse sur $\RP^2$]\label{ex:morse-rp2}
\index{fonction de Morse!sur $\RP^2$}
On réalise $\RP^2$ comme le quotient de $S^2$ par l'identification antipodale.
La fonction $f([x_0, x_1, x_2]) = \frac{x_0^2 - x_2^2}{x_0^2 + x_1^2 + x_2^2}$
est bien définie sur $\RP^2$ et constitue une fonction de Morse avec trois
points critiques : $c_0 = 1$, $c_1 = 1$, $c_2 = 1$. Les nombres de Betti
rationnels de $\RP^2$ sont $b_0 = 1$, $b_1 = 0$, $b_2 = 0$. Les inégalités de
Morse sont strictes pour $\lambda = 1, 2$, ce qui reflète la torsion dans
$H_1(\RP^2; \Z) \cong \Z/2\Z$.
\end{exemple}

% ============================================================
\section{Complexe de Morse et connexion avec la topologie algébrique}%
\label{sec:complexe-morse}
% ============================================================

\begin{definition}[Complexe de Morse--Smale--Witten]\label{def:complexe-morse}
\index{complexe de Morse}\index{complexe de Morse--Smale--Witten}
Soit $f : M \to \R$ une fonction de Morse sur une variété compacte orientée $M$,
et $g$ une métrique riemannienne telle que le couple $(f, g)$ satisfait la
condition de Morse--Smale (les variétés stables et instables se coupent
transversalement). Le \textbf{complexe de Morse} est le complexe de chaînes
$(C_*^{\mathrm{Morse}}, \partial)$ défini par :
\begin{itemize}
  \item $C_\lambda^{\mathrm{Morse}} = \bigoplus_{\operatorname{ind}(p) = \lambda} \Z \cdot p$,
    le groupe abélien libre engendré par les points critiques d'indice $\lambda$ ;
  \item la différentielle $\partial : C_\lambda^{\mathrm{Morse}} \to C_{\lambda-1}^{\mathrm{Morse}}$
    est définie par $\partial p = \sum_{\operatorname{ind}(q) = \lambda - 1} n(p, q) \cdot q$,
    où $n(p, q)$ est le nombre algébrique de lignes de flot du gradient
    $-\nabla f$ allant de $p$ à $q$.
\end{itemize}
\end{definition}

\begin{theoreme}[Homologie de Morse]\label{thm:homologie-morse}
\index{homologie de Morse}
Le complexe de Morse vérifie $\partial^2 = 0$, et son homologie est
canoniquement isomorphe à l'homologie singulière :
\[
  H_\lambda^{\mathrm{Morse}}(M, f, g) \cong H_\lambda(M; \Z).
\]
\end{theoreme}

\begin{remarque}
Ce résultat remarquable permet de calculer l'homologie d'une variété compacte à
l'aide d'un complexe de chaînes de rang fini (engendré par les points critiques
d'une fonction de Morse), ce qui est bien plus économique que le complexe
singulier.
\end{remarque}

% ============================================================
\section{Exercices}\label{sec:morse-exos}
% ============================================================

\begin{exercice}[Points critiques sur des surfaces]\label{exo:morse-1}
Soit $\Sigma_g$ la surface orientable compacte de genre $g$. Montrer que toute
fonction de Morse sur $\Sigma_g$ possède au moins $2g + 2$ points critiques.
\emph{Indication :} utiliser les inégalités de Morse faibles et les nombres de
Betti $b_0 = 1$, $b_1 = 2g$, $b_2 = 1$.
\end{exercice}

\begin{exercice}[Minimum du nombre de points critiques]\label{exo:morse-2}
Montrer qu'il existe une fonction de Morse sur $T^2$ ayant exactement $4$ points
critiques (la fonction hauteur), et qu'aucune fonction de Morse sur $T^2$ ne
peut en avoir moins de $4$.
\end{exercice}

\begin{exercice}[Fonction de Morse sur $S^1 \times S^2$]\label{exo:morse-3}
Construire une fonction de Morse sur $S^1 \times S^2$ et déterminer ses points
critiques et leurs indices. Vérifier les inégalités de Morse.
\end{exercice}

\begin{exercice}[Morse et caractéristique d'Euler]\label{exo:morse-4}
Soit $f : M \to \R$ une fonction de Morse sur une variété compacte de dimension
$n$. Montrer que $\chi(M) = \sum_{\lambda=0}^n (-1)^\lambda c_\lambda$, où
$c_\lambda$ est le nombre de points critiques d'indice $\lambda$. En déduire que
$\chi(S^n) = 1 + (-1)^n$.
\end{exercice}

\begin{exercice}[Morse sur les espaces projectifs]\label{exo:morse-5}
Considérer la fonction $f : \CP^n \to \R$ définie par
\[
  f([z_0 : \cdots : z_n]) = \frac{\sum_{k=0}^n k \, \abs{z_k}^2}{\sum_{k=0}^n \abs{z_k}^2}.
\]
Montrer que $f$ est une fonction de Morse avec exactement $n+1$ points critiques,
d'indices $0, 2, 4, \ldots, 2n$. En déduire que $\chi(\CP^n) = n + 1$.
\end{exercice}

\begin{exercice}[Lemme de Morse en dimension 1]\label{exo:morse-6}
Vérifier directement le lemme de Morse pour une fonction
$f : \R \to \R$ ayant un point critique non dégénéré en $0$ : trouver
explicitement le changement de variable $y = y(x)$.
\end{exercice}

\begin{exercice}[Indice et signature]\label{exo:morse-7}
Soit $f : \R^n \to \R$ définie par $f(x) = x^T A x$ où $A$ est une matrice
symétrique inversible. Montrer que $0$ est un point critique non dégénéré et que
l'indice est le nombre de valeurs propres négatives de $A$.
\end{exercice}

\begin{exercice}[Fonctions de Morse et homotopie]\label{exo:morse-8}
Soit $M$ une variété compacte connexe admettant une fonction de Morse
$f : M \to \R$ avec exactement deux points critiques (un minimum et un maximum).
\begin{enumerate}[label=(\alph*)]
  \item Montrer que $M$ est homéomorphe à $S^n$.
  \item L'application est-elle nécessairement un difféomorphisme ?
    \emph{(Considérer les sphères exotiques de Milnor.)}
\end{enumerate}
\end{exercice}

\begin{exercice}[Perturbation en fonction de Morse]\label{exo:morse-9}
Soit $f : \R^2 \to \R$ définie par $f(x,y) = x^3 - 3xy^2$ (singularité de type
$D_4^-$). Montrer que $0$ est un point critique dégénéré. Trouver une
perturbation $f_\varepsilon$ de $f$ qui est une fonction de Morse au voisinage
de l'origine, et déterminer les points critiques et indices de $f_\varepsilon$.
\end{exercice}

\begin{exercice}[Lacet polynomial de Morse]\label{exo:morse-10}
Soit $M$ une variété compacte orientée de dimension $n$ et $f : M \to \R$ une
fonction de Morse. On définit le \textbf{polynôme de Morse}\index{polynôme de Morse} par
\[
  \mathcal{M}_f(t) = \sum_{\lambda = 0}^n c_\lambda \, t^\lambda,
\]
où $c_\lambda$ est le nombre de points critiques d'indice $\lambda$. De même, le
\textbf{polynôme de Poincaré} est $\mathcal{P}_M(t) = \sum_\lambda b_\lambda \, t^\lambda$.
\begin{enumerate}[label=(\alph*)]
  \item Montrer que $\mathcal{M}_f(t) - \mathcal{P}_M(t) = (1+t)\,Q(t)$ pour un
    certain polynôme $Q(t)$ à coefficients entiers non négatifs.
  \item Vérifier cette relation pour la fonction hauteur sur $T^2$, $S^2$, $\RP^2$.
  \item En déduire les inégalités de Morse fortes.
\end{enumerate}
\end{exercice}

\begin{exercice}[Fonction de Morse sur les groupes de Lie]\label{exo:morse-11}
\index{fonction de Morse!sur un groupe de Lie}
Soit $G = \SO(3)$. En identifiant $\SO(3) \cong \RP^3$, construire une fonction
de Morse sur $G$ et en déduire les nombres de Betti rationnels de $\SO(3)$.
\end{exercice}


% ============================================================
%  CHAPITRE 10 : DEGRÉ D'UNE APPLICATION LISSE ET APPLICATIONS
% ============================================================
\chapter{Degré d'une application lisse et applications}\label{chap:degre}
\minitoc

\vspace{1cm}

Le \emph{degré} d'une application lisse entre variétés compactes est un invariant
entier (ou modulo $2$ dans le cas non orientable) qui mesure combien de fois le
domaine <<~enveloppe~>> le but. Cet invariant, malgré sa simplicité apparente,
possède des conséquences topologiques profondes : il permet de démontrer le
théorème de Brouwer sur le point fixe, le théorème de la sphère chevelue, et
bien d'autres résultats fondamentaux.

% ============================================================
\section{Degré modulo 2}\label{sec:degre-mod2}
% ============================================================

\begin{definition}[Degré mod 2]\label{def:degre-mod2}
\index{degré!modulo 2}
Soient $M$ et $N$ deux variétés lisses compactes sans bord de même dimension $n$,
et $f : M \to N$ une application lisse. Soit $q \in N$ une valeur
régulière\index{valeur régulière} de $f$. Le \textbf{degré modulo 2} de $f$ est
\[
  \deg_2(f) = \#f^{-1}(q) \pmod{2} \in \Z/2\Z.
\]
\end{definition}

\begin{theoreme}[Bonne définition du degré mod 2]\label{thm:degre-mod2-bdef}
Le degré $\deg_2(f)$ ne dépend pas du choix de la valeur régulière $q$.
\end{theoreme}

\begin{proof}
Soient $q_0, q_1 \in N$ deux valeurs régulières de $f$. On relie $q_0$ à $q_1$
par un chemin lisse $\gamma : [0,1] \to N$. Par le théorème de Sard, on peut
choisir $\gamma$ transverse à $f$, de sorte que $W = \set{(t, p) \in [0,1] \times M \mid f(p) = \gamma(t)}$
est une sous-variété de dimension $1$ à bord de $[0,1] \times M$.

Le bord de $W$ est
\[
  \partial W = (\set{0} \times f^{-1}(q_0)) \cup (\set{1} \times f^{-1}(q_1)).
\]
Puisque toute variété compacte de dimension $1$ à bord a un nombre pair de points
de bord, on obtient $\#f^{-1}(q_0) + \#f^{-1}(q_1) \equiv 0 \pmod{2}$.
\end{proof}

\begin{theoreme}[Invariance homotopique du degré mod 2]\label{thm:degre-mod2-htpie}
\index{degré!invariance homotopique}
Si $f_0, f_1 : M \to N$ sont homotopes, alors $\deg_2(f_0) = \deg_2(f_1)$.
\end{theoreme}

\begin{proof}
Soit $F : M \times [0,1] \to N$ une homotopie lisse entre $f_0$ et $f_1$ (on
peut toujours lisser une homotopie continue). Soit $q \in N$ une valeur
régulière de $F$ (qui est aussi valeur régulière de $f_0$ et de $f_1$ pour un
choix générique). Alors $W = F^{-1}(q)$ est une variété compacte de dimension
$1$ à bord, avec $\partial W = f_0^{-1}(q) \sqcup f_1^{-1}(q)$. On conclut
comme précédemment par parité.
\end{proof}

% ============================================================
\section{Degré orienté}\label{sec:degre-oriente}
% ============================================================

\begin{definition}[Degré orienté]\label{def:degre-oriente}
\index{degré!orienté}
Soient $M$ et $N$ deux variétés lisses compactes sans bord, connexes, orientées,
de même dimension $n$, et $f : M \to N$ une application lisse. Pour une valeur
régulière $q \in N$ et $p \in f^{-1}(q)$, on définit le \textbf{signe local} de
$f$ en $p$ par
\[
  \varepsilon(p) = \operatorname{sgn}(\det \dd f_p) =
  \begin{cases}
    +1 & \text{si } \dd f_p : T_pM \to T_qN \text{ préserve l'orientation}, \\
    -1 & \text{si } \dd f_p \text{ renverse l'orientation}.
  \end{cases}
\]
Le \textbf{degré} de $f$ est l'entier
\[
  \deg(f) = \sum_{p \in f^{-1}(q)} \varepsilon(p) \in \Z.
\]
\end{definition}

\begin{theoreme}[Bonne définition et invariance homotopique]\label{thm:degre-oriente-bdef}
\index{degré!bonne définition}
Soient $M, N$ comme ci-dessus.
\begin{enumerate}[label=(\roman*)]
  \item Le degré $\deg(f)$ ne dépend pas du choix de la valeur régulière $q$.
  \item Si $f_0 \simeq f_1$ (homotopes), alors $\deg(f_0) = \deg(f_1)$.
  \item $\deg(\Id_M) = 1$.
  \item $\deg(g \circ f) = \deg(g) \cdot \deg(f)$.
\end{enumerate}
\end{theoreme}

\begin{proof}
\textit{(i)} et \textit{(ii)} : La preuve est analogue au cas mod $2$, mais en
comptant les signes. Si $F : M \times [0,1] \to N$ est une homotopie lisse et
$q$ une valeur régulière de $F$, alors $W = F^{-1}(q)$ est une variété orientée
de dimension $1$ à bord. Les composantes de $W$ sont soit des cercles (qui ne
contribuent pas au bord), soit des arcs reliant un point de $f_0^{-1}(q)$ à un
point de $f_1^{-1}(q)$. Un examen des orientations montre que les signes des
deux extrémités de chaque arc sont opposés (en tenant compte de l'orientation du
bord de $W$), d'où $\deg(f_0) = \deg(f_1)$.

\textit{(iii)} est immédiat. \textit{(iv)} découle du fait que, pour une valeur
régulière $q$ de $g \circ f$, on a
$\varepsilon_{g \circ f}(p) = \varepsilon_g(f(p)) \cdot \varepsilon_f(p)$, et
l'on somme en regroupant par fibre.
\end{proof}

\begin{proposition}[Degré d'un difféomorphisme]\label{prop:degre-diffeo}
Un difféomorphisme $f : M \to N$ entre variétés compactes orientées connexes a
pour degré $\deg(f) = +1$ s'il préserve l'orientation, et $\deg(f) = -1$ s'il
la renverse.
\end{proposition}

\begin{proposition}[Surjectivité et degré]\label{prop:degre-surjectif}
Soit $f : M \to N$ une application lisse entre variétés compactes connexes
orientées de même dimension. Si $\deg(f) \neq 0$, alors $f$ est surjective.
\end{proposition}

\begin{proof}
Si $f$ n'est pas surjective, soit $q \in N \setminus f(M)$. Puisque $f(M)$ est
compact (donc fermé) dans $N$, $q$ est une valeur régulière de $f$ avec
$f^{-1}(q) = \varnothing$, donc $\deg(f) = 0$.
\end{proof}

\begin{proposition}[Degré et applications constantes]\label{prop:degre-constante}
Toute application constante $f : M \to N$ a pour degré $\deg(f) = 0$ (avec la
convention $\deg_2(f) = 0$ dans le cas non orienté).
\end{proposition}

\begin{remarque}
La réciproque est fausse : $\deg(f) = 0$ n'implique pas que $f$ est homotope à
une application constante. Par exemple, l'application de Hopf
$h : S^3 \to S^2$ n'est pas de degré zéro au sens usuel (elle est entre variétés
de dimensions différentes), mais toute application $S^n \to S^n$ de degré zéro
qui se factorise par un espace de dimension inférieure est homotopiquement
triviale.
\end{remarque}

% ============================================================
\section{Calcul du degré via les formes différentielles}\label{sec:degre-formes}
% ============================================================

\begin{theoreme}[Formule intégrale du degré]\label{thm:degre-integral}
\index{degré!formule intégrale}
Soient $M, N$ des variétés compactes orientées sans bord de dimension $n$, et
$f : M \to N$ lisse. Pour toute $n$-forme $\omega$ sur $N$ telle que
$\int_N \omega \neq 0$,
\[
  \deg(f) = \frac{\int_M f^*\omega}{\int_N \omega}.
\]
\end{theoreme}

\begin{proof}
Il suffit de vérifier la formule pour une forme $\omega$ à support dans un petit
voisinage d'une valeur régulière $q$. Au voisinage de chaque $p_i \in f^{-1}(q)$,
$f$ est un difféomorphisme local, et $f^*\omega$ a pour intégrale
$\varepsilon(p_i) \int_N \omega$. En sommant sur les préimages, on obtient
$\int_M f^*\omega = \deg(f) \int_N \omega$.
\end{proof}

\begin{exemple}[Degré de $z \mapsto z^k$ sur $S^1$]\label{ex:degre-cercle}
\index{degré!de $z^k$ sur $S^1$}
Soit $f_k : S^1 \to S^1$ définie par $f_k(z) = z^k$ (en voyant $S^1 \subset \C$).
La forme $\omega = \dd\theta$ vérifie $\int_{S^1} \omega = 2\pi$, et
$f_k^*(\dd\theta) = k \, \dd\theta$, donc
\[
  \deg(f_k) = \frac{\int_{S^1} k \, \dd\theta}{\int_{S^1} \dd\theta} = k.
\]
On peut aussi le voir directement : pour $q = 1$, on a
$f_k^{-1}(1) = \set{e^{2i\pi j/k} \mid 0 \leqslant j < k}$, et le jacobien est
positif en chaque préimage (pour $k > 0$).
\end{exemple}

% ============================================================
\section{Applications du degré}\label{sec:degre-applications}
% ============================================================

\subsection{Théorème de Brouwer via le degré}

\begin{theoreme}[Théorème du point fixe de Brouwer]\label{thm:brouwer-degre}
\index{théorème!de Brouwer (point fixe)}
Toute application continue $g : D^n \to D^n$ du disque fermé dans lui-même
possède un point fixe.
\end{theoreme}

\begin{proof}
Supposons par l'absurde que $g$ n'a pas de point fixe. On construit alors une
rétraction $r : D^n \to S^{n-1}$ en envoyant chaque $x \in D^n$ sur le point
d'intersection de la demi-droite partant de $g(x)$ et passant par $x$ avec
$S^{n-1}$.

L'application $r$ est continue et vérifie $r|_{S^{n-1}} = \Id_{S^{n-1}}$.
Considérons l'inclusion $\iota : S^{n-1} \hookrightarrow D^n$. Alors
$r \circ \iota = \Id_{S^{n-1}}$, donc
$\deg(r \circ \iota) = \deg(\Id_{S^{n-1}}) = 1$. Mais $\iota$ se prolonge en
$\iota' : D^n \to D^n$ qui est contractile (puisque $D^n$ est contractile),
donc $r \circ \iota$ est homotope à une application constante, qui a degré $0$.
Contradiction.
\end{proof}

\subsection{Théorème de la sphère chevelue}

\begin{theoreme}[Théorème de la sphère chevelue]\label{thm:sphere-chevelue}
\index{théorème!de la sphère chevelue}
Il n'existe pas de champ de vecteurs continu ne s'annulant en aucun point sur
$S^{2n}$ (sphère de dimension paire).
\end{theoreme}

\begin{proof}
Supposons qu'il existe un champ de vecteurs $V : S^{2n} \to \R^{2n+1}$ continu,
tangent et non nul en chaque point. On normalise : $v(x) = V(x)/\abs{V(x)}$, de
sorte que $\abs{v(x)} = 1$ et $\langle v(x), x \rangle = 0$ pour tout
$x \in S^{2n}$. Considérons l'homotopie
\[
  H_t(x) = \cos(\pi t) \, x + \sin(\pi t) \, v(x), \quad t \in [0,1].
\]
On vérifie que $\abs{H_t(x)} = 1$ pour tout $t$ (par orthogonalité de $x$ et
$v(x)$), donc $H_t : S^{2n} \to S^{2n}$. De plus, $H_0 = \Id$ et
$H_1 = -\Id$ (l'application antipodale). Par invariance homotopique,
$\deg(\Id) = \deg(-\Id)$, soit $1 = (-1)^{2n+1} = -1$, ce qui est absurde.
\end{proof}

\begin{remarque}
En revanche, $S^{2n+1}$ admet toujours un champ de vecteurs non nul.
Par exemple, sur $S^{2n+1} \subset \R^{2n+2}$, le champ
$V(x_1, x_2, \ldots, x_{2n+1}, x_{2n+2}) = (-x_2, x_1, \ldots, -x_{2n+2}, x_{2n+1})$
convient.
\end{remarque}

\subsection{Théorème de Borsuk--Ulam}

\begin{theoreme}[Borsuk--Ulam]\label{thm:borsuk-ulam}
\index{théorème!de Borsuk--Ulam}
Pour toute application continue $f : S^n \to \R^n$, il existe $x \in S^n$ tel
que $f(x) = f(-x)$.
\end{theoreme}

\begin{proof}
Supposons par l'absurde que $f(x) \neq f(-x)$ pour tout $x$. On pose
\[
  g(x) = \frac{f(x) - f(-x)}{\abs{f(x) - f(-x)}},
\]
qui est une application continue $g : S^n \to S^{n-1}$ vérifiant $g(-x) = -g(x)$
(application \textbf{impaire}\index{application impaire}).

On montre que toute application impaire $S^n \to S^{n-1}$ a un degré mod $2$
non nul (par récurrence sur $n$). Pour $n = 1$, une application impaire
$g : S^1 \to S^0 = \set{\pm 1}$ envoie des points antipodaux sur des points
distincts, ce qui est impossible par connexité de $S^1$.

Pour le pas de récurrence, on restreint $g$ à l'équateur
$S^{n-1} \subset S^n$. La restriction $g|_{S^{n-1}} : S^{n-1} \to S^{n-1}$ est
encore impaire, donc de degré impair par hypothèse de récurrence. Mais
$g|_{S^{n-1}}$ se prolonge à l'hémisphère $D^n_+$, donc serait d'extension
possible à $D^n$, ce qui impliquerait degré nul. Contradiction.
\end{proof}

\subsection{Théorème du point fixe de Lefschetz}

\begin{theoreme}[Lefschetz, esquisse]\label{thm:lefschetz}
\index{théorème!de Lefschetz (point fixe)}
Soit $M$ une variété compacte et $f : M \to M$ une application lisse. On
définit le \textbf{nombre de Lefschetz}\index{nombre de Lefschetz}
\[
  L(f) = \sum_{k=0}^n (-1)^k \operatorname{tr}(f_* : H_k(M;\Q) \to H_k(M;\Q)).
\]
Si $L(f) \neq 0$, alors $f$ admet un point fixe.
\end{theoreme}

\begin{remarque}
Le nombre de Lefschetz de l'identité est $L(\Id) = \chi(M)$. On retrouve ainsi
le fait qu'une variété compacte de caractéristique d'Euler non nulle vérifie la
propriété du point fixe pour toute application homotope à l'identité.
\end{remarque}

% ============================================================
\section{Nombre d'intersection comme degré}\label{sec:intersection-degre}
% ============================================================

\begin{definition}[Nombre d'intersection]\label{def:nombre-intersection}
\index{nombre d'intersection}
Soient $M$ une variété compacte orientée de dimension $n$, et $A$, $B$ des
sous-variétés compactes orientées de dimensions complémentaires
$\dim A + \dim B = n$, se coupant transversalement. Le \textbf{nombre
d'intersection} de $A$ et $B$ est
\[
  I(A, B) = \sum_{p \in A \cap B} \varepsilon(p),
\]
où $\varepsilon(p) = +1$ si l'orientation de $T_pA \oplus T_pB$ coïncide avec
celle de $T_pM$, et $\varepsilon(p) = -1$ sinon.
\end{definition}

\begin{proposition}\label{prop:intersection-degre}
Le nombre d'intersection peut s'exprimer comme un degré. Si $A$ et $B$ sont des
sous-variétés compactes orientées sans bord de $M$ de dimensions complémentaires,
et si $\tau : M \to M \times M$ est l'inclusion diagonale tandis que
$A \times B \hookrightarrow M \times M$, alors $I(A, B)$ coïncide avec le
nombre d'intersection de $\tau(M)$ et $A \times B$ dans $M \times M$.
\end{proposition}

% ============================================================
\section{Exercices}\label{sec:degre-exos}
% ============================================================

\begin{exercice}[Degré d'applications entre sphères]\label{exo:degre-1}
\begin{enumerate}[label=(\alph*)]
  \item Calculer le degré de la réflexion $r : S^n \to S^n$ définie par
    $r(x_0, \ldots, x_n) = (-x_0, x_1, \ldots, x_n)$.
  \item En déduire le degré de l'application antipodale
    $a : S^n \to S^n$, $a(x) = -x$.
  \item Montrer que si $n$ est pair, l'application antipodale n'est pas
    homotope à l'identité.
\end{enumerate}
\end{exercice}

\begin{exercice}[Degré et revêtement]\label{exo:degre-2}
Soit $p : M \to N$ un revêtement lisse à $k$ feuillets entre variétés compactes
orientées connexes. Montrer que $\deg(p) = k$ (si $p$ préserve l'orientation) ou
$\deg(p) = -k$ (sinon).
\end{exercice}

\begin{exercice}[Application sans point fixe]\label{exo:degre-3}
Soit $f : S^n \to S^n$ une application lisse sans point fixe. Montrer que
$\deg(f) = (-1)^{n+1}$. \emph{Indication :} montrer que $f$ est homotope à
l'application antipodale.
\end{exercice}

\begin{exercice}[Théorème fondamental de l'algèbre via le degré]\label{exo:degre-4}
\index{théorème!fondamental de l'algèbre}
Soit $P(z) = z^n + a_{n-1}z^{n-1} + \cdots + a_0$ un polynôme de degré $n \geqslant 1$.
\begin{enumerate}[label=(\alph*)]
  \item Pour $R > 0$ assez grand, montrer que l'application
    $f_R : S^1 \to S^1$ définie par $f_R(z) = P(Rz)/\abs{P(Rz)}$ est bien
    définie et a pour degré $n$.
  \item En déduire que $P$ a au moins une racine dans $\C$.
\end{enumerate}
\end{exercice}

\begin{exercice}[Degré et forme volume]\label{exo:degre-5}
Soit $f : S^2 \to S^2$ une application lisse. En utilisant la forme volume
$\omega = \sin\varphi \, \dd\varphi \wedge \dd\theta$ de $S^2$, calculer
$\deg(f)$ par la formule intégrale pour
\begin{enumerate}[label=(\alph*)]
  \item $f$ = projection stéréographique composée avec l'inversion,
  \item $f(x,y,z) = (x,y,-z)$.
\end{enumerate}
\end{exercice}

\begin{exercice}[Degré de l'application de Hopf]\label{exo:degre-6}
\index{application de Hopf}
L'application de Hopf $h : S^3 \to S^2$ est définie par
$h(z_1, z_2) = [z_1 : z_2] \in \CP^1 \cong S^2$, où $(z_1, z_2) \in S^3 \subset \C^2$.
Expliquer pourquoi la notion de degré ne s'applique pas directement (dimensions
différentes). Quel invariant peut-on associer à $h$ ?
\end{exercice}

\begin{exercice}[Borsuk--Ulam et coloriage]\label{exo:degre-7}
En utilisant le théorème de Borsuk--Ulam, montrer que pour toute application
continue $f : S^2 \to \R^2$, il existe un grand cercle de $S^2$ sur lequel $f$
est constante en au moins un point et son antipodal.
\end{exercice}

\begin{exercice}[Nombre d'intersection]\label{exo:degre-8}
Soit $T^2 = S^1 \times S^1$ le tore, et $A = S^1 \times \set{1}$,
$B = \set{1} \times S^1$ les deux cercles générateurs. Calculer $I(A,B)$ en
munissant $T^2$, $A$ et $B$ de leurs orientations standard.
\end{exercice}

\begin{exercice}[Degré et polynômes trigonométriques]\label{exo:degre-9}
Soit $f : S^1 \to S^1$ définie par $f(e^{i\theta}) = e^{iP(\theta)}$ où
$P(\theta) = n\theta + \sum_{k=1}^{N} a_k \sin(k\theta)$ avec des $a_k$ assez
petits. Montrer que $\deg(f) = n$.
\end{exercice}

\begin{exercice}[Degré et indice d'un champ de vecteurs]\label{exo:degre-10}
\index{indice!d'un champ de vecteurs}
Soit $V$ un champ de vecteurs lisse sur $\R^2$ ayant un zéro isolé en $p$.
On définit l'\textbf{indice} de $V$ en $p$ par
\[
  \operatorname{ind}_p(V) = \deg\!\left(\frac{V}{\abs{V}} : S^1_\varepsilon(p) \to S^1\right)
\]
pour $\varepsilon > 0$ assez petit (où $S^1_\varepsilon(p)$ est le cercle de rayon
$\varepsilon$ centré en $p$).
\begin{enumerate}[label=(\alph*)]
  \item Calculer l'indice de $V(x,y) = (x, y)$ en l'origine.
  \item Calculer l'indice de $V(x,y) = (-y, x)$ en l'origine.
  \item Calculer l'indice de $V(x,y) = (x^2 - y^2, 2xy)$ en l'origine.
  \item Énoncer le théorème de Poincaré--Hopf\index{théorème!de Poincaré--Hopf}
    reliant la somme des indices à la caractéristique d'Euler.
\end{enumerate}
\end{exercice}

\begin{exercice}[Degré mod 2 et Borsuk--Ulam]\label{exo:degre-11}
Montrer que toute application impaire $f : S^n \to S^n$ a un degré mod $2$ égal
à $1$, c'est-à-dire $\deg_2(f) = 1$. En déduire une preuve alternative du
théorème de Borsuk--Ulam.
\end{exercice}


% ============================================================
%  CHAPITRE 11 : THÉORÈME DE WHITNEY ET PLONGEMENTS
% ============================================================
\chapter{Théorème de Whitney et plongements}\label{chap:whitney}
\minitoc

\vspace{1cm}

Une question naturelle en topologie différentielle est de savoir si toute variété
abstraite (définie par un atlas) peut être <<~réalisée~>> comme une sous-variété
d'un espace euclidien. Les théorèmes de Whitney\index{théorème!de Whitney}
répondent par l'affirmative et fournissent des bornes optimales sur la dimension
de l'espace ambiant nécessaire.

% ============================================================
\section{Rappels : plongements et immersions}\label{sec:whitney-rappels}
% ============================================================

\begin{definition}[Rappel]\label{def:plongement-immersion-rappel}
Soit $f : M \to N$ une application lisse entre variétés.
\begin{enumerate}[label=(\roman*)]
  \item $f$ est une \textbf{immersion}\index{immersion} si $\dd f_p$ est
    injective pour tout $p \in M$.
  \item $f$ est un \textbf{plongement}\index{plongement} (embedding) si $f$ est
    une immersion injective qui est un homéomorphisme sur son image $f(M)$
    (munie de la topologie induite).
\end{enumerate}
\end{definition}

\begin{remarque}
Si $M$ est compacte, toute immersion injective $f : M \to N$ est automatiquement
un plongement (car une bijection continue d'un compact vers un espace séparé est
un homéomorphisme).
\end{remarque}

% ============================================================
\section{Théorème de Whitney : version faible}\label{sec:whitney-faible}
% ============================================================

\begin{theoreme}[Plongement de Whitney, version faible]\label{thm:whitney-faible}
\index{théorème!de Whitney (version faible)}
Toute variété lisse compacte de dimension $n$ admet un plongement dans
$\R^{2n+1}$.
\end{theoreme}

\begin{proof}
La preuve se décompose en plusieurs étapes.

\medskip
\noindent\textbf{Étape 1 : Plongement dans $\R^N$ pour $N$ grand.}
On recouvre $M$ par un nombre fini de cartes $(U_i, \varphi_i)_{1 \leqslant i \leqslant k}$
et on choisit une partition de l'unité $(\rho_i)$ subordonnée à ce recouvrement.
On définit
\[
  F : M \to \R^{k(n+1)}, \quad
  F(p) = \bigl(\rho_1(p)\varphi_1(p), \ldots, \rho_k(p)\varphi_k(p),
    \rho_1(p), \ldots, \rho_k(p)\bigr).
\]
L'application $F$ est une immersion injective (l'injectivité utilise les $\rho_i$,
et le fait que $\dd f$ est injective se vérifie dans chaque ouvert $U_i$ où
$\rho_i \neq 0$). Comme $M$ est compacte, $F$ est un plongement. Posons
$N = k(n+1)$.

\medskip
\noindent\textbf{Étape 2 : Réduction de la dimension.}
On montre que si $N > 2n + 1$, on peut projeter $F(M) \subset \R^N$ sur un
hyperplan $\R^{N-1}$ tout en conservant le caractère de plongement. Il suffit de
trouver une direction $v \in S^{N-1}$ telle que la projection
$\pi_v : \R^N \to v^\perp \cong \R^{N-1}$ restreinte à $F(M)$ reste un
plongement.

La projection $\pi_v \circ F$ échoue à être injective si $v$ est de la forme
\[
  v = \frac{F(p) - F(q)}{\abs{F(p) - F(q)}} \quad \text{pour } p \neq q.
\]
Définissons $\alpha : M \times M \setminus \Delta \to S^{N-1}$ par
$\alpha(p,q) = \frac{F(p)-F(q)}{\abs{F(p)-F(q)}}$. L'ensemble source est de
dimension $2n$, et le but $S^{N-1}$ de dimension $N - 1 > 2n$ (car $N > 2n+1$).
Par le théorème de Sard, $\im(\alpha)$ est de mesure nulle dans $S^{N-1}$.

De même, $\pi_v \circ F$ échoue à être une immersion si $v$ appartient à l'image
de l'application de Gauss $\beta : TM \setminus \set{0} \to S^{N-1}$ définie par
$\beta(p, w) = \dd F_p(w)/\abs{\dd F_p(w)}$. Le domaine est de dimension
$2n - 1$, donc $\im(\beta)$ est aussi de mesure nulle pour $N - 1 > 2n - 1$,
c'est-à-dire $N > 2n$.

On choisit $v \notin \im(\alpha) \cup \im(\beta)$, et la projection $\pi_v \circ F$
est un plongement dans $\R^{N-1}$. On itère jusqu'à atteindre $\R^{2n+1}$.
\end{proof}

% ============================================================
\section{Théorème de Whitney : version forte}\label{sec:whitney-fort}
% ============================================================

\begin{theoreme}[Plongement de Whitney, version forte]\label{thm:whitney-fort}
\index{théorème!de Whitney (version forte)}
Toute variété lisse compacte de dimension $n \geqslant 1$ admet un plongement
dans $\R^{2n}$.
\end{theoreme}

\begin{remarque}
La preuve de la version forte est considérablement plus délicate que celle de la
version faible. Elle nécessite la technique du <<~truc de Whitney~>>
\index{truc de Whitney} (Whitney trick) qui permet d'éliminer des points doubles
par paires via des disques de Whitney. Nous renvoyons à \cite{Whitney1944} et
\cite{AdachiBook} pour les détails.
\end{remarque}

\begin{remarque}
La borne $2n$ est optimale en général. Par exemple, la bouteille de
Klein\index{bouteille de Klein} (variété de dimension $2$) ne se plonge pas dans
$\R^3$ mais se plonge dans $\R^4$.
\end{remarque}

\begin{theoreme}[Immersion de Whitney]\label{thm:whitney-immersion}
\index{théorème!de Whitney (immersion)}
Toute variété lisse compacte de dimension $n \geqslant 1$ admet une immersion
dans $\R^{2n-1}$.
\end{theoreme}

\begin{remarque}
Pour les immersions, la borne $2n - 1$ est elle aussi optimale en général. Un
résultat plus fin de Hirsch et Smale montre que l'existence d'une immersion
$M^n \looparrowright \R^k$ est contrôlée par l'existence d'un monomorphisme de
fibrés vectoriels $TM \to \varepsilon^k$.
\end{remarque}

% ============================================================
\section{Théorème d'approximation de Whitney}\label{sec:whitney-approx}
% ============================================================

\begin{theoreme}[Approximation de Whitney]\label{thm:whitney-approx}
\index{théorème!d'approximation de Whitney}
Soient $M$ et $N$ des variétés lisses, et $f : M \to N$ une application
continue. Pour tout $\varepsilon > 0$, il existe une application lisse
$g : M \to N$ telle que $g$ est $\varepsilon$-proche de $f$ (pour une métrique
fixée sur $N$) et $g$ est homotope à $f$.
\end{theoreme}

\begin{proof}[Idée de la preuve]
On plonge $N$ dans $\R^K$ par Whitney. L'application $f : M \to N \subset \R^K$
est approchée par une application lisse $\tilde{g} : M \to \R^K$ par convolution
avec un noyau régularisant. On projette ensuite $\tilde{g}$ sur $N$ via la
rétraction par voisinage tubulaire $r : \mathcal{U} \to N$, ce qui donne
$g = r \circ \tilde{g}$. L'homotopie entre $f$ et $g$ est fournie par le
segment $[f, \tilde{g}]$ composé avec $r$.
\end{proof}

\begin{corollaire}\label{cor:approx-homotopie}
Deux applications lisses $f, g : M \to N$ qui sont homotopes par une homotopie
continue sont aussi homotopes par une homotopie lisse.
\end{corollaire}

\begin{remarque}
Le théorème d'approximation de Whitney est l'un des outils fondamentaux
permettant de transférer les résultats de topologie algébrique (qui opèrent avec
des applications continues) au cadre de la topologie différentielle. Par exemple,
c'est grâce à ce théorème que l'invariance homotopique du degré, démontrée pour
des homotopies lisses, s'étend automatiquement aux homotopies continues.
\end{remarque}

\begin{proposition}[Approximation relative]\label{prop:approx-relative}
\index{approximation!relative}
Soient $M, N$ des variétés lisses, $A \subset M$ un sous-ensemble fermé, et
$f : M \to N$ une application continue qui est déjà lisse sur un voisinage de
$A$. Alors pour tout $\varepsilon > 0$, il existe une application lisse
$g : M \to N$ telle que $g|_A = f|_A$, $g$ est $\varepsilon$-proche de $f$, et
$g$ est homotope à $f$ relativement à $A$.
\end{proposition}

\begin{proof}[Idée]
On utilise une partition de l'unité pour construire l'approximation seulement
loin de $A$, en interpolant avec $f$ au voisinage de $A$ où $f$ est déjà lisse.
\end{proof}

% ============================================================
\section{Isotopie et plongements}\label{sec:isotopie}
% ============================================================

\begin{definition}[Isotopie]\label{def:isotopie}
\index{isotopie}
Deux plongements $f_0, f_1 : M \hookrightarrow N$ sont \textbf{isotopes} s'il
existe une application lisse $F : M \times [0,1] \to N$ telle que :
\begin{enumerate}[label=(\roman*)]
  \item $F(-, 0) = f_0$ et $F(-, 1) = f_1$ ;
  \item pour tout $t \in [0,1]$, $F(-, t) : M \hookrightarrow N$ est un plongement.
\end{enumerate}
\end{definition}

\begin{definition}[Isotopie ambiante]\label{def:isotopie-ambiante}
\index{isotopie!ambiante}
Deux plongements $f_0, f_1 : M \hookrightarrow N$ sont \textbf{ambiamment
isotopes} s'il existe un difféomorphisme $\Phi : N \to N$ isotope à l'identité
tel que $\Phi \circ f_0 = f_1$.
\end{definition}

\begin{theoreme}[Théorème d'extension d'isotopie]\label{thm:extension-isotopie}
\index{théorème!d'extension d'isotopie}
Soit $M$ une variété compacte et $N$ une variété sans bord. Toute isotopie entre
deux plongements $f_0, f_1 : M \hookrightarrow N$ s'étend en une isotopie
ambiante de $N$.
\end{theoreme}

\begin{proof}[Idée de la preuve]
L'isotopie $F : M \times [0,1] \to N$ engendre un champ de vecteurs dépendant
du temps sur $N$, défini le long de l'image $F(M, t)$ pour chaque $t$. On
étend ce champ à $N$ tout entier à l'aide d'une partition de l'unité et d'un
voisinage tubulaire de $F(M, t)$. Le flot intégré de ce champ étendu fournit
le difféomorphisme $\Phi_t : N \to N$ avec $\Phi_0 = \Id_N$ et
$\Phi_1 \circ f_0 = f_1$.
\end{proof}

\begin{remarque}
Le théorème d'extension d'isotopie est faux si $M$ n'est pas compacte. Un
contre-exemple est donné par une isotopie qui <<~pousse~>> un plongement de $\R$
dans $\R^2$ vers l'infini.
\end{remarque}

\begin{proposition}[Unicité du plongement de la boule]\label{prop:plongement-boule}
\index{plongement!de la boule}
Deux plongements lisses $f_0, f_1 : D^n \hookrightarrow M^n$ qui préservent
l'orientation sont ambiamment isotopes (théorème du disque de Palais).
\end{proposition}

% ============================================================
\section{Applications et exemples}\label{sec:whitney-exemples}
% ============================================================

\begin{exemple}[Plongement du tore dans $\R^3$]\label{ex:tore-R3}
\index{tore!plongement dans $\R^3$}
Le tore $T^2 = S^1 \times S^1$ admet le plongement classique dans $\R^3$ :
\[
  \iota(\theta, \varphi) = \bigl((R + r\cos\varphi)\cos\theta,
    \; (R + r\cos\varphi)\sin\theta, \; r\sin\varphi\bigr)
\]
pour $0 < r < R$. Ce plongement est bien une immersion (la différentielle est
de rang $2$ partout) et l'injectivité est assurée par le choix $r < R$.
\end{exemple}

\begin{exemple}[Plongement de $\RP^2$ dans $\R^4$ et non-plongement dans $\R^3$]%
\label{ex:rp2-plongement}
\index{espace projectif réel!plongement dans $\R^4$}
Le plan projectif réel $\RP^2$ admet un plongement dans $\R^4$ mais ne se
plonge pas dans $\R^3$. Le non-plongement résulte du fait que $\RP^2$ est une
surface compacte non orientable, et toute surface compacte sans bord plongée
dans $\R^3$ est orientable (par le théorème de Jordan--Brouwer\index{théorème!de Jordan--Brouwer}).

Un plongement explicite de $\RP^2$ dans $\R^4$ est donné par
\[
  F([x:y:z]) = \bigl(xy, \, xz, \, y^2 - z^2, \, 2yz\bigr) \cdot \frac{1}{x^2+y^2+z^2}.
\]
On vérifie que cette application est bien définie, lisse, injective et de rang $2$
partout.
\end{exemple}

\begin{exemple}[Immersion de Boy]\label{ex:boy}
\index{immersion de Boy}\index{surface de Boy}
L'immersion de Boy est une immersion de $\RP^2$ dans $\R^3$ (mais pas un
plongement, car elle possède des auto-intersections). Elle a été découverte par
Werner Boy en 1901 et constitue un exemple remarquable d'immersion d'une surface
non orientable dans $\R^3$. L'immersion de Boy possède une symétrie d'ordre $3$
et un unique point triple.
\end{exemple}

\begin{figure}[ht]
\centering
\begin{tikzpicture}[scale=1.2]
  % Schéma simplifié du "truc de Whitney"
  \draw[thick] (-3, 0) to[out=30, in=150] (-0.5, 0.5);
  \draw[thick] (-0.5, 0.5) to[out=-30, in=210] (0.5, 0.5);
  \draw[thick] (0.5, 0.5) to[out=30, in=150] (3, 0);

  \draw[thick] (-3, -0.2) to[out=-30, in=210] (-0.5, -0.5);
  \draw[thick] (-0.5, -0.5) to[out=30, in=150] (0.5, -0.5);
  \draw[thick] (0.5, -0.5) to[out=-30, in=210] (3, -0.2);

  % Points d'intersection
  \fill[red] (-0.5, 0) circle (2pt) node[above left] {$p$};
  \fill[red] (0.5, 0) circle (2pt) node[above right] {$q$};

  % Disque de Whitney
  \draw[dashed, blue!60!black, thick] (-0.5, 0) to[out=-60, in=240] (0.5, 0);
  \draw[dashed, blue!60!black, thick] (-0.5, 0) to[out=60, in=120] (0.5, 0);
  \node[blue!60!black, below] at (0, -0.6) {disque de Whitney};

  \draw[->, thick] (3.5, -0.1) -- (5, -0.1) node[midway, above] {Whitney trick};

  % Après le trick
  \draw[thick] (5.5, 0.3) to[out=0, in=180] (8.5, 0.3);
  \draw[thick] (5.5, -0.3) to[out=0, in=180] (8.5, -0.3);
  \node at (7, -1) {pas d'intersection};
\end{tikzpicture}
\caption{Le <<~truc de Whitney~>> : élimination de deux points d'intersection
  de signes opposés via un disque de Whitney.}
\label{fig:whitney-trick}
\end{figure}

% ============================================================
\section{Classification des surfaces lisses compactes}\label{sec:classif-surfaces}
% ============================================================

\begin{theoreme}[Classification des surfaces compactes]\label{thm:classif-surfaces}
\index{classification des surfaces}
Toute surface lisse compacte connexe sans bord est difféomorphe à l'une des
surfaces suivantes :
\begin{enumerate}[label=(\roman*)]
  \item la sphère $S^2$ ;
  \item la somme connexe de $g$ tores : $\Sigma_g = T^2 \# \cdots \# T^2$
    ($g$ fois), pour $g \geqslant 1$ ;
  \item la somme connexe de $k$ plans projectifs : $\RP^2 \# \cdots \# \RP^2$
    ($k$ fois), pour $k \geqslant 1$.
\end{enumerate}
Les surfaces \textup{(i)} et \textup{(ii)} sont orientables ; les surfaces
\textup{(iii)} ne le sont pas.
\end{theoreme}

\begin{remarque}
Ce théorème, combiné avec les résultats de plongement de Whitney, montre que
toute surface lisse compacte orientable se plonge dans $\R^3$, tandis que les
surfaces non orientables nécessitent au moins $\R^4$.
\end{remarque}

\begin{remarque}
La classification est complètement déterminée par deux invariants :
\begin{itemize}
  \item l'\textbf{orientabilité} (oui ou non) ;
  \item la \textbf{caractéristique d'Euler} $\chi(\Sigma)$, ou de manière
    équivalente, le genre.
\end{itemize}
Pour les surfaces orientables, $\chi(\Sigma_g) = 2 - 2g$. Pour les surfaces non
orientables à $k$ plans projectifs, $\chi = 2 - k$. Notons que
$\RP^2 \# \RP^2$ est la bouteille de Klein, et $\RP^2 \# \RP^2 \# \RP^2$ est
homéomorphe à $T^2 \# \RP^2$.
\end{remarque}

\begin{proposition}[Plongement des surfaces orientables]\label{prop:plongement-surfaces}
\index{surfaces!plongement}
Toute surface orientable compacte $\Sigma_g$ se plonge dans $\R^3$. En
particulier :
\begin{enumerate}[label=(\roman*)]
  \item $S^2$ est le bord de la boule $D^3$ ;
  \item $\Sigma_g$ pour $g \geqslant 1$ se réalise comme un <<~bretzel~>> à $g$
    trous dans $\R^3$.
\end{enumerate}
\end{proposition}

% ============================================================
\section{Aperçu : variétés exotiques}\label{sec:varietes-exotiques}
% ============================================================

Les théorèmes de Whitney montrent que toute variété lisse compacte se plonge dans
un espace euclidien, mais la question de l'\emph{unicité} de la structure lisse
est bien plus subtile.

\begin{theoreme}[Milnor, 1956]\label{thm:milnor-exotiques}
\index{sphère exotique}\index{Milnor}
Il existe des variétés lisses homéomorphes à $S^7$ mais non difféomorphes à
$S^7$. Plus précisément, le groupe $\Theta_7$ des classes de difféomorphisme de
structures lisses sur $S^7$ est isomorphe à $\Z/28\Z$.
\end{theoreme}

\begin{remarque}
La construction originale de Milnor utilise des fibrés en sphères $S^3$ sur
$S^4$. Ces fibrés sont classifiés par $\pi_3(\SO(4)) \cong \Z \oplus \Z$, et
pour certains choix, l'espace total est homéomorphe mais non difféomorphe à $S^7$.
Milnor détecte la non-trivialité de la structure lisse en utilisant un invariant
basé sur la signature et les classes de Pontryaguine\index{classes de Pontryaguine}.
\end{remarque}

\begin{remarque}
Ce résultat fondamental de Milnor a ouvert un champ immense de recherche. On
sait aujourd'hui que :
\begin{itemize}
  \item En dimension $n \leqslant 6$, la sphère $S^n$ admet une unique
    structure lisse (à difféomorphisme près).
  \item En dimension $4$, la situation est encore largement ouverte : on ne sait
    pas s'il existe des $S^4$ exotiques (c'est le problème lisse de la
    conjecture de Poincaré en dimension $4$\index{conjecture de Poincaré}).
  \item L'espace $\R^4$ admet une infinité non dénombrable de structures lisses
    non difféomorphes entre elles (résultat remarquable de Donaldson et
    Freedman\index{Donaldson}\index{Freedman}).
  \item Le groupe des sphères exotiques $\Theta_n$ est connu pour de nombreuses
    valeurs de $n$. Par exemple, $\abs{\Theta_{11}} = 992$ et $\abs{\Theta_{15}} = 16256$.
\end{itemize}
\end{remarque}

\begin{remarque}[Sphères exotiques et théorie de Morse]
La connexion avec la théorie de Morse est la suivante : une sphère exotique
$\Sigma^n$ admet une fonction de Morse avec exactement deux points critiques
(un minimum et un maximum), puisqu'elle est homéomorphe à $S^n$. La décomposition
en anses correspondante consiste en une $0$-anse et une $n$-anse, mais
l'application de recollement entre ces deux anses définit un difféomorphisme
$\varphi : S^{n-1} \to S^{n-1}$ qui n'est pas isotope à l'identité.
\end{remarque}

% ============================================================
\section{Exercices}\label{sec:whitney-exos}
% ============================================================

\begin{exercice}[Plongement de graphe]\label{exo:whitney-1}
Soit $f : M \to N$ une application lisse. Montrer que l'application
$\Gamma_f : M \to M \times N$ définie par $\Gamma_f(p) = (p, f(p))$ est un
plongement. En déduire que toute variété compacte de dimension $n$ se plonge
dans $\R^{2n+1}$ à l'aide d'une application $f : M \to \R^{n+1}$.
\end{exercice}

\begin{exercice}[Immersions de $S^1$ dans $\R^2$]\label{exo:whitney-2}
\index{immersion!de $S^1$}
\begin{enumerate}[label=(\alph*)]
  \item Montrer que le plongement standard $S^1 \hookrightarrow \R^2$ et
    l'immersion en huit $\gamma(t) = (\sin(2t), \sin(t))$ ne sont pas
    régulièrement homotopes.
  \item Énoncer le théorème de classification de Whitney--Graustein pour les
    courbes fermées régulières dans le plan.
\end{enumerate}
\end{exercice}

\begin{exercice}[Non-plongement de la bouteille de Klein]\label{exo:whitney-3}
\index{bouteille de Klein!non-plongement}
Montrer que la bouteille de Klein ne se plonge pas dans $\R^3$.
\emph{Indication :} utiliser l'argument d'orientabilité (théorème de
Jordan--Brouwer).
\end{exercice}

\begin{exercice}[Plongements de $\RP^n$]\label{exo:whitney-4}
\begin{enumerate}[label=(\alph*)]
  \item Construire un plongement de $\RP^1 \cong S^1$ dans $\R^2$.
  \item Montrer que $\RP^2$ ne se plonge pas dans $\R^3$.
  \item Montrer que $\RP^n$ se plonge dans $\R^{n(n+2)/2}$ via l'application
    $[x] \mapsto \frac{x x^T}{\abs{x}^2}$ (matrices symétriques).
\end{enumerate}
\end{exercice}

\begin{exercice}[Immersion et fibré normal]\label{exo:whitney-5}
Soit $f : M^n \hookrightarrow \R^{n+k}$ une immersion. On définit le
\textbf{fibré normal}\index{fibré normal} $\nu_f$ de $f$ comme le fibré
vectoriel sur $M$ dont la fibre en $p$ est $(\dd f_p(T_pM))^\perp$ dans
$\R^{n+k}$.
\begin{enumerate}[label=(\alph*)]
  \item Montrer que $TM \oplus \nu_f \cong \varepsilon^{n+k}$ (fibré trivial).
  \item En déduire que $w(TM) \cdot w(\nu_f) = 1$ dans
    $H^*(M; \Z/2\Z)$, où $w$ désigne la classe de Stiefel--Whitney totale.
\end{enumerate}
\end{exercice}

\begin{exercice}[Whitney et Sard]\label{exo:whitney-6}
Reprendre la preuve du théorème de Whitney faible (théorème~\ref{thm:whitney-faible})
et vérifier en détail que l'argument de dimension fonctionne : si $N > 2n + 1$,
les <<~mauvaises directions~>> forment un ensemble de mesure nulle dans
$S^{N-1}$.
\end{exercice}

\begin{exercice}[Approximation par des plongements]\label{exo:whitney-7}
Soit $M$ une variété compacte de dimension $n$ et $N > 2n$. Montrer que toute
application lisse $f : M \to \R^N$ peut être approchée arbitrairement près (en
topologie $\mathcal{C}^1$) par un plongement.
\end{exercice}

\begin{exercice}[Isotopie de noeuds]\label{exo:whitney-8}
\index{noeud}
Un \textbf{n\oe ud}\index{noeud} est un plongement $\gamma : S^1 \hookrightarrow \R^3$.
\begin{enumerate}[label=(\alph*)]
  \item Montrer que tout n\oe ud $\gamma : S^1 \hookrightarrow \R^n$ avec
    $n \geqslant 4$ est isotope au n\oe ud trivial (le plongement standard de
    $S^1$ dans un plan).
  \item Pourquoi cet argument échoue-t-il pour $n = 3$ ?
\end{enumerate}
\end{exercice}

\begin{exercice}[Surfaces exotiques]\label{exo:whitney-9}
\begin{enumerate}[label=(\alph*)]
  \item Montrer que toute variété de dimension $\leqslant 3$ admet une unique
    structure lisse (à difféomorphisme près) compatible avec sa structure
    topologique.
    \emph{(Admis : citer les théorèmes de Radó, Moise et Munkres.)}
  \item Pourquoi la dimension $4$ est-elle exceptionnelle du point de vue des
    structures lisses ?
\end{enumerate}
\end{exercice}

\begin{exercice}[Immersions en codimension 1]\label{exo:whitney-10}
\index{immersion!en codimension 1}
\begin{enumerate}[label=(\alph*)]
  \item Montrer que toute variété compacte orientable de dimension $n$ admet une
    immersion dans $\R^{n+1}$ si et seulement si son fibré tangent est stablement
    trivial.
  \item La sphère $S^n$ admet-elle toujours une immersion dans $\R^{n+1}$ ?
    Justifier.
  \item Montrer que $\RP^n$ admet une immersion dans $\R^{n+1}$ si et seulement
    si $n$ est de la forme $2^k - 1$.
    \emph{(Admis : on pourra utiliser les classes de Stiefel--Whitney.)}
\end{enumerate}
\end{exercice}

\begin{exercice}[Plongement de variétés produits]\label{exo:whitney-11}
\begin{enumerate}[label=(\alph*)]
  \item Si $M^m$ se plonge dans $\R^p$ et $N^n$ se plonge dans $\R^q$, montrer
    que $M \times N$ se plonge dans $\R^{p+q}$.
  \item Donner un plongement explicite de $T^n = (S^1)^n$ dans $\R^{2n}$.
  \item Peut-on toujours faire mieux que la borne $p + q$ ? Donner un exemple où
    $M \times N$ se plonge dans $\R^k$ avec $k < p + q$.
\end{enumerate}
\end{exercice}


% ============================================================
% ============================================================
%  APPENDICE
% ============================================================
% ============================================================
\appendix

\chapter{Rappels d'analyse}\label{app:analyse}
\minitoc

\vspace{0.5cm}

Nous rassemblons ici quelques résultats classiques d'analyse utilisés tout au
long de ce cours.

% ============================================================
\section{Théorème d'inversion locale}\label{app:inversion-locale}
% ============================================================

\begin{theoreme}[Inversion locale]\label{thm:app-inversion}
\index{théorème!d'inversion locale}
Soit $U \subset \R^n$ un ouvert, $f : U \to \R^n$ une application de classe
$\mathcal{C}^k$ ($k \geqslant 1$) et $a \in U$ tel que $\det(\dd f_a) \neq 0$.
Alors il existe des ouverts $V \ni a$ et $W \ni f(a)$ tels que
$f|_V : V \to W$ est un $\mathcal{C}^k$-difféomorphisme.
\end{theoreme}

\begin{corollaire}[Théorème des fonctions implicites]\label{cor:app-implicite}
\index{théorème!des fonctions implicites}
Soit $F : U \times V \to \R^m$ de classe $\mathcal{C}^k$, où
$U \subset \R^n$, $V \subset \R^m$ sont ouverts, et $(a,b) \in U \times V$ tel
que $F(a,b) = 0$ et $\frac{\partial F}{\partial y}(a,b)$ est inversible. Alors
il existe un voisinage $U' \ni a$ et une unique application
$g : U' \to V$ de classe $\mathcal{C}^k$ telle que $g(a) = b$ et
$F(x, g(x)) = 0$ pour tout $x \in U'$.
\end{corollaire}

\begin{corollaire}[Théorème du rang constant]\label{cor:app-rang}
\index{théorème!du rang constant}
Soit $f : U \to \R^m$ de classe $\mathcal{C}^k$ ($k \geqslant 1$) de rang
constant $r$ sur $U \subset \R^n$. Alors pour tout $a \in U$, il existe des
$\mathcal{C}^k$-difféomorphismes $\varphi$ et $\psi$ au voisinage de $a$ et
$f(a)$ respectivement, tels que
\[
  \psi \circ f \circ \varphi^{-1}(x_1, \ldots, x_n) = (x_1, \ldots, x_r, 0, \ldots, 0).
\]
\end{corollaire}

% ============================================================
\section{Théorème de Cauchy--Lipschitz}\label{app:cauchy-lipschitz}
% ============================================================

\begin{theoreme}[Cauchy--Lipschitz / Picard--Lindelöf]\label{thm:app-cauchy-lip}
\index{théorème!de Cauchy--Lipschitz}
Soit $U \subset \R^n$ un ouvert, $I \subset \R$ un intervalle ouvert et
$X : I \times U \to \R^n$ une application de classe $\mathcal{C}^k$
($k \geqslant 1$). Pour tout $(t_0, x_0) \in I \times U$, le problème de
Cauchy
\[
  \begin{cases}
    \dot{x}(t) = X(t, x(t)), \\
    x(t_0) = x_0,
  \end{cases}
\]
admet une unique solution maximale $x : J \to U$ de classe $\mathcal{C}^{k+1}$,
définie sur un intervalle ouvert maximal $J \subset I$ contenant $t_0$.
\end{theoreme}

\begin{theoreme}[Dépendance régulière]\label{thm:app-dependance}
\index{dépendance régulière en les conditions initiales}
Si $X$ est de classe $\mathcal{C}^k$, le flot
$\Phi : \Omega \to U$, $\Phi(t, x_0) = x(t; t_0, x_0)$ est de classe
$\mathcal{C}^k$ sur son domaine de définition $\Omega \subset \R \times U$.
\end{theoreme}

\begin{remarque}
Pour les champs de vecteurs lisses ($\mathcal{C}^\infty$) sur une variété
compacte, le flot est défini pour tout temps : $\Omega = \R \times M$.
\end{remarque}

% ============================================================
\section{Mesure de Lebesgue et ensembles de mesure nulle}\label{app:lebesgue}
% ============================================================

\begin{definition}[Ensemble de mesure nulle]\label{def:app-mesure-nulle}
\index{ensemble de mesure nulle}
Un sous-ensemble $A \subset \R^n$ est de \textbf{mesure de Lebesgue nulle} si
pour tout $\varepsilon > 0$, il existe une famille dénombrable de pavés
$(R_k)_{k \in \N}$ telle que $A \subset \bigcup_{k \in \N} R_k$ et
$\sum_{k \in \N} \operatorname{vol}(R_k) < \varepsilon$.
\end{definition}

\begin{proposition}[Propriétés]\label{prop:app-mesure-nulle}
\begin{enumerate}[label=(\roman*)]
  \item Toute réunion dénombrable d'ensembles de mesure nulle est de mesure nulle.
  \item Tout sous-ensemble d'un ensemble de mesure nulle est de mesure nulle.
  \item Si $f : U \to \R^n$ est localement lipschitzienne et $A \subset U$ est de
    mesure nulle, alors $f(A)$ est de mesure nulle.
  \item Si $k < n$, tout sous-ensemble $A \subset \R^k \times \set{0}^{n-k} \subset \R^n$
    est de mesure nulle dans $\R^n$.
\end{enumerate}
\end{proposition}

\begin{theoreme}[Sard]\label{thm:app-sard}
\index{théorème de Sard}
Soit $f : U \to \R^m$ une application de classe $\mathcal{C}^k$ avec
$U \subset \R^n$ ouvert et $k \geqslant \max(1, n - m + 1)$. Alors l'ensemble
des valeurs critiques de $f$ est de mesure de Lebesgue nulle dans $\R^m$.
\end{theoreme}

% ============================================================
\section{Partition de l'unité}\label{app:partition}
% ============================================================

\begin{theoreme}[Existence de partitions de l'unité]\label{thm:app-partition}
\index{partition de l'unité}
Soit $M$ une variété lisse paracompacte et $(U_\alpha)_{\alpha \in A}$ un
recouvrement ouvert de $M$. Il existe une partition de l'unité
$(\rho_\alpha)_{\alpha \in A}$ subordonnée à $(U_\alpha)$, c'est-à-dire :
\begin{enumerate}[label=(\roman*)]
  \item chaque $\rho_\alpha : M \to [0,1]$ est lisse ;
  \item $\operatorname{supp}(\rho_\alpha) \subset U_\alpha$ pour tout $\alpha$ ;
  \item la famille $(\operatorname{supp}(\rho_\alpha))$ est localement finie ;
  \item $\sum_\alpha \rho_\alpha(p) = 1$ pour tout $p \in M$.
\end{enumerate}
\end{theoreme}

\begin{remarque}
L'existence des partitions de l'unité est un outil fondamental en topologie
différentielle. Elle permet de <<~recoller~>> des constructions locales en
constructions globales, et joue un rôle essentiel dans les preuves des théorèmes
de plongement de Whitney, de l'existence de métriques riemanniennes, de connexions,
et de bien d'autres résultats.
\end{remarque}

% ============================================================
\section{Rappels de topologie différentielle de base}\label{app:rappels-diff}
% ============================================================

\begin{theoreme}[Existence de métriques riemanniennes]\label{thm:app-metrique}
\index{métrique riemannienne}
Toute variété lisse paracompacte admet une métrique riemannienne.
\end{theoreme}

\begin{proof}[Esquisse]
On choisit un recouvrement par des cartes $(U_i, \varphi_i)$ et une partition de
l'unité $(\rho_i)$ subordonnée. Sur chaque $U_i$, la métrique euclidienne
tirée en arrière par $\varphi_i$ donne une métrique locale $g_i$. On pose
$g = \sum_i \rho_i \, g_i$, qui est une métrique riemannienne globale car une
combinaison convexe de formes définies positives est définie positive.
\end{proof}

\begin{theoreme}[Voisinage tubulaire]\label{thm:app-tubulaire}
\index{voisinage tubulaire}
Soit $S \subset M$ une sous-variété compacte d'une variété lisse $M$. Alors
il existe un voisinage ouvert $U$ de $S$ dans $M$ et un difféomorphisme
$\Phi : \nu(S) \supset V \to U$ d'un voisinage ouvert $V$ du plongement par la
section nulle du fibré normal $\nu(S) \to S$ sur $U$, tel que
$\Phi|_S = \Id_S$.
\end{theoreme}

\begin{proof}[Esquisse]
On choisit une métrique riemannienne $g$ sur $M$. L'application exponentielle
restreinte au fibré normal $\exp^\perp : \nu(S) \to M$ est un difféomorphisme
local au voisinage de la section nulle (par le théorème d'inversion locale). Pour
$\varepsilon > 0$ assez petit, l'application $\exp^\perp$ est un difféomorphisme
de $\set{v \in \nu(S) \mid \abs{v}_g < \varepsilon}$ sur un voisinage tubulaire
ouvert $U$ de $S$.
\end{proof}

\begin{remarque}
Le voisinage tubulaire est unique à isotopie près : deux voisinages tubulaires
de $S$ dans $M$ sont ambiamment isotopes (c'est un corollaire du théorème
d'extension d'isotopie et de l'unicité du fibré normal).
\end{remarque}

% ============================================================
\section{Lemme de Hadamard}\label{app:hadamard}
% ============================================================

\begin{lemme}[Hadamard]\label{lem:app-hadamard}
\index{lemme de Hadamard}
Soit $U \subset \R^n$ un ouvert convexe contenant l'origine, et
$f : U \to \R$ une application lisse avec $f(0) = 0$. Alors il existe des
fonctions lisses $g_1, \ldots, g_n : U \to \R$ telles que
\[
  f(x) = \sum_{i=1}^n x_i \, g_i(x) \quad \text{et} \quad
  g_i(0) = \frac{\partial f}{\partial x_i}(0).
\]
\end{lemme}

\begin{proof}
On écrit
\[
  f(x) = f(x) - f(0) = \int_0^1 \frac{\dd}{\dd t} f(tx) \, \dd t
  = \int_0^1 \sum_{i=1}^n x_i \frac{\partial f}{\partial x_i}(tx) \, \dd t
  = \sum_{i=1}^n x_i \underbrace{\int_0^1 \frac{\partial f}{\partial x_i}(tx) \, \dd t}_{= g_i(x)}.
\]
La fonction $g_i$ est lisse par le théorème de différentiation sous le signe
intégral, et $g_i(0) = \frac{\partial f}{\partial x_i}(0)$.
\end{proof}


% ============================================================
%  BIBLIOGRAPHIE
% ============================================================
\begin{thebibliography}{99}

\bibitem{Milnor1965}
J.\,W. Milnor,
\emph{Topology from the Differentiable Viewpoint},
The University Press of Virginia, 1965.
\index{Milnor}

\bibitem{MilnorMorse}
J.\,W. Milnor,
\emph{Morse Theory},
Annals of Mathematics Studies, No.~51, Princeton University Press, 1963.

\bibitem{GP2010}
V. Guillemin et A. Pollack,
\emph{Differential Topology},
AMS Chelsea Publishing, 2010 (réédition).
\index{Guillemin}\index{Pollack}

\bibitem{Hirsch1976}
M.\,W. Hirsch,
\emph{Differential Topology},
Graduate Texts in Mathematics~33, Springer-Verlag, 1976.
\index{Hirsch}

\bibitem{Lee2013}
J.\,M. Lee,
\emph{Introduction to Smooth Manifolds},
2e édition, Graduate Texts in Mathematics~218, Springer, 2013.
\index{Lee}

\bibitem{BrockerJanich}
T. Bröcker et K. Jänich,
\emph{Introduction to Differential Topology},
Cambridge University Press, 1982.
\index{Bröcker}\index{Jänich}

\bibitem{BottTu1982}
R. Bott et L.\,W. Tu,
\emph{Differential Forms in Algebraic Topology},
Graduate Texts in Mathematics~82, Springer-Verlag, 1982.
\index{Bott}\index{Tu}

\bibitem{Warner1983}
F.\,W. Warner,
\emph{Foundations of Differentiable Manifolds and Lie Groups},
Graduate Texts in Mathematics~94, Springer-Verlag, 1983.
\index{Warner}

\bibitem{Whitney1944}
H. Whitney,
\emph{The self-intersections of a smooth $n$-manifold in $2n$-space},
Annals of Mathematics 45 (1944), 220--246.
\index{Whitney}

\bibitem{AdachiBook}
M. Adachi,
\emph{Embeddings and Immersions},
Translations of Mathematical Monographs~124, AMS, 1993.

\bibitem{Kosinski1993}
A.\,A. Kosinski,
\emph{Differential Manifolds},
Academic Press, 1993.

\bibitem{MilnorStasheff}
J.\,W. Milnor et J.\,D. Stasheff,
\emph{Characteristic Classes},
Annals of Mathematics Studies, No.~76, Princeton University Press, 1974.

\bibitem{milnor1965topology}
J.\,W. Milnor,
\emph{Topology from the Differentiable Viewpoint},
The University Press of Virginia, 1965.

\end{thebibliography}


% ============================================================
%  INDEX
% ============================================================
\printindex

\end{document}
% === FIN DU CHUNK ch9-11+end ===
